% ============================================================
% refer/papers 全量内容总结报告
% 生成器：docs/build_refer_papers_report.py
% 编译：XeLaTeX 两遍；本文件不依赖 refer/papers 的图片或外部数据
% ============================================================
\documentclass[UTF8,aspectratio=169,11pt]{ctexbeamer}
\usepackage{amsmath,amssymb,mathtools}
\usepackage{booktabs}
\usepackage{tabularx}
\usepackage{array}
\usepackage{tikz}
\usepackage{xcolor}
\usepackage{microtype}
\usepackage{hyperref}
\hypersetup{hidelinks}
\usetikzlibrary{arrows.meta,positioning,calc,fit,shapes.geometric}

\definecolor{primary}{HTML}{17365D}
\definecolor{accent}{HTML}{1F77B4}
\definecolor{evidence}{HTML}{1E8449}
\definecolor{warning}{HTML}{C55A11}
\definecolor{muted}{HTML}{5B6573}
\definecolor{panel}{HTML}{F4F7FA}
\setbeamersize{text margin left=0.55cm,text margin right=0.55cm}
\setlength{\emergencystretch}{3em}
\setbeamertemplate{navigation symbols}{}
\setbeamertemplate{blocks}[rounded][shadow=false]
\setbeamercolor{normal text}{fg=black!86,bg=white}
\setbeamercolor{structure}{fg=primary}
\setbeamercolor{frametitle}{fg=primary,bg=white}
\setbeamercolor{block title}{fg=primary,bg=panel}
\setbeamercolor{block body}{fg=black!86,bg=panel}
\setbeamerfont{frametitle}{size=\large,series=\bfseries}
\setbeamertemplate{frametitle}{\vskip0.12cm\insertframetitle\par\vskip0.08cm}
\setbeamertemplate{footline}{%
  \hbox{\begin{beamercolorbox}[wd=\paperwidth,ht=2.3ex,dp=0.9ex,
    leftskip=0.55cm,rightskip=0.55cm]{author in head/foot}%
    \scriptsize\color{muted}\insertshortauthor\hfill
    \insertshorttitle\hfill\insertframenumber/\inserttotalframenumber
    \end{beamercolorbox}}}
\newcommand{\source}[1]{\vspace{0.08em}\par{\raggedright\scriptsize\color{muted}来源：#1\par}}
\newcommand{\verified}{\textcolor{evidence}{确证}}
\newcommand{\inferred}{\textcolor{accent}{推断}}
\newcommand{\unverified}{\textcolor{warning}{未验证}}
\newcolumntype{Y}{>{\raggedright\arraybackslash}X}
\title[refer/papers 全量总结]{refer/papers 全量论文内容总结}
\author[Codex]{Codex}
\institute{格点 QCD、部分子物理与数值方法文献档案}
\date{2026 年 8 月 28 日}

\begin{document}
%% frame_manifest|frame_id=G-01|paper_id=GLOBAL|claim_id=G1|visual_type=title|width_budget=0.98linewidth|height_budget=0.86textheight|min_font_pt=10|split_allowed=no|expected_page=1|risks=title-table-source-footline
\begin{frame}[plain]
  \titlepage
\end{frame}

%% frame_manifest|frame_id=G-02|paper_id=GLOBAL|claim_id=G2|visual_type=scope-table|width_budget=0.98linewidth|height_budget=0.86textheight|min_font_pt=10|split_allowed=no|expected_page=2|risks=title-table-source-footline
\begin{frame}[t]{这份报告总结了什么，以及哪些内容没有被越界推断}
  \small
  % 主张：报告对象是 50 篇索引论文，而不是过程日志或重复选集。
  \begingroup
\renewcommand{\arraystretch}{0.84}
\begin{tabularx}{0.98\linewidth}{@{}p{2.55cm}>{\raggedright\arraybackslash}X@{}}
\toprule
论文集合 & \textbf{50 篇}；按 \detokenize{refer/papers/INDEX.md} 的英中目录对去重。 \\
直接证据 & 中文/英文 \detokenize{main.tex}、\detokenize{chapters/*.tex}、转换/翻译指南与索引。 \\
不纳入科学内容 & 已有 Essential Papers 选集不重复展开；\detokenize{.agent.*} 日志和构建缓存只作过程记录。 \\
当前限制 & 目录中未发现论文 PDF；PDF 原文逐页核验、原始数据复现和独立数值检查均标为\unverified。 \\
\bottomrule
\end{tabularx}
\endgroup
  \source{证据：\detokenize{refer/papers/INDEX.md}、\detokenize{refer/papers/AGENTS.md}；报告边界为本次任务定义。}
\end{frame}

%% frame_manifest|frame_id=G-03|paper_id=GLOBAL|claim_id=G3|visual_type=corpus-statistics|width_budget=0.98linewidth|height_budget=0.86textheight|min_font_pt=10|split_allowed=no|expected_page=3|risks=title-table-source-footline
\begin{frame}[t]{语料规模显示：原文转排库足以支撑逐论文档案，但不是 PDF 复核}
  \small
  \begingroup
\renewcommand{\arraystretch}{0.84}
\begin{tabularx}{0.98\linewidth}{@{}p{2.55cm}>{\raggedright\arraybackslash}X@{}}
\toprule
目录结构 & 50 个英/中目录对，另有 2 个已有精华选集目录；论文正文只按 50 篇计数。 \\
索引页数 & 英文转排页数合计 1285 页；中文译本页数合计 1175 页；库内索引合计 2460 页。 \\
报告规模 & 每篇论文 4 个固定 frame，加总览、主题综合、交叉结论与来源附录；预期页数在生成时实测。 \\
物理主线 & 格点规范底座 → 流/重正化 → LaMET/赝 PDF → 胶子、强子 PDF/TMD → 采样算法。 \\
\bottomrule
\end{tabularx}
\endgroup
  \source{页数证据：\detokenize{refer/papers/INDEX.md}；报告页数以新 PDF 的 \detokenize{pdfinfo} 实测为准。}
\end{frame}

%% frame_manifest|frame_id=G-04|paper_id=GLOBAL|claim_id=G4|visual_type=evidence-table|width_budget=0.98linewidth|height_budget=0.86textheight|min_font_pt=10|split_allowed=no|expected_page=4|risks=title-table-source-footline
\begin{frame}[t]{每一个判断都分成“确证、推断、未验证”三层}
  \small
  \begingroup
\renewcommand{\arraystretch}{0.84}
\begin{tabularx}{0.98\linewidth}{@{}p{2.55cm}>{\raggedright\arraybackslash}X@{}}
\toprule
确证 & 源 LaTeX 摘要、章节、公式、图表或结论中可以直接定位的内容；报告保留文件和行号。 \\
推断 & 由一个或多个确证内容得到的物理关系或跨论文接口；文字明确写出“综合判断/推断”。 \\
未验证 & 缺少 PDF、原始数据、独立复现或源文件证据的数值/外推；只写缺口与核验动作。 \\
物理检查 & 优先检查规范对称性、量纲、连续/大动量/短距离极限和守恒关系；检查不替代实测。 \\
\bottomrule
\end{tabularx}
\endgroup
  \source{证据规则：本报告规格与每篇论文的中文/英文源文件；颜色不是唯一状态编码。}
\end{frame}

%% frame_manifest|frame_id=G-05|paper_id=GLOBAL|claim_id=G5|visual_type=workflow-table|width_budget=0.98linewidth|height_budget=0.86textheight|min_font_pt=10|split_allowed=no|expected_page=5|risks=title-table-source-footline
\begin{frame}[t]{全库的共同计算链条是“欧氏可计算量 → 匹配/演化 → 光锥观测量”}
  \small
  \begingroup
\renewcommand{\arraystretch}{0.84}
\begin{tabularx}{0.98\linewidth}{@{}p{2.55cm}>{\raggedright\arraybackslash}X@{}}
\toprule
输入 & 格点规范场系综、强子关联函数、空间 Wilson 线或流场；对象必须保持规范不变/协变。 \\
中间量 & 矩阵元、Ioffe 时间分布、准/赝分布、流时间复合算符、非微扰重正化因子。 \\
理论桥梁 & OPE、LaMET 因子化、重正化方案转换、微扰匹配、DGLAP 或 Collins--Soper 演化。 \\
输出 & PDF/GPD/TMD/强子结构或经物理可观测量检验的采样系综；有限动量、有限间距和系统误差不能省略。 \\
\bottomrule
\end{tabularx}
\endgroup
  \source{综合依据：50 篇论文档案的章节路线；此页是跨论文结构示意，不是单篇论文的数值结果。}
\end{frame}

%% frame_manifest|frame_id=T-006|paper_id=GLOBAL|claim_id=T-格点规范基础|visual_type=topic-map-formula|width_budget=0.98linewidth|height_budget=0.86textheight|min_font_pt=10|split_allowed=no|expected_page=6|risks=title-table-source-footline
\begin{frame}[t]{格点规范理论与禁闭：3 篇论文围绕同一个中间量链条展开}
  \small
\begingroup
\renewcommand{\arraystretch}{0.84}
\begin{tabularx}{0.98\linewidth}{@{}p{2.55cm}>{\raggedright\arraybackslash}X@{}}
\toprule
共同问题 & 从规范不变离散化、Wilson 圈和强耦合结构出发，建立欧氏格点上的规范场、夸克传播与禁闭图像。 \\
代表结构式 & $\langle W(C)\rangle\sim\exp[-\sigma A(C)]$\\[-0.2em]\scriptsize 结构式/示意：用于说明变量和极限，不冒充单篇原文公式。 \\
物理校验 & 量纲检查：$\sigma$ 具有面积倒数维度；大圈极限把面积律与线性势联系起来。 \\
本组论文 & 夸克禁闭、SU3链接变量解析涂抹、高动量强子新夸克涂抹 \\
\bottomrule
\end{tabularx}
\endgroup
  \source{综合来源：本组论文的 \detokenize{main.tex} 与章节源；结构式为本报告的主题示意。}
\end{frame}

%% frame_manifest|frame_id=T-007|paper_id=GLOBAL|claim_id=T-梯度流与能动量张量|visual_type=topic-map-formula|width_budget=0.98linewidth|height_budget=0.86textheight|min_font_pt=10|split_allowed=no|expected_page=7|risks=title-table-source-footline
\begin{frame}[t]{梯度流、涂抹与能动量张量：12 篇论文围绕同一个中间量链条展开}
  \small
\begingroup
\renewcommand{\arraystretch}{0.84}
\begin{tabularx}{0.98\linewidth}{@{}p{2.55cm}>{\raggedright\arraybackslash}X@{}}
\toprule
共同问题 & 用流时间平滑短距离涨落，在保持规范协变性的同时构造可控的复合算符、耦合、质量和能动量张量。 \\
代表结构式 & $\partial_t B_\mu(t,x)=D_\nu G_{\nu\mu}(t,x),\qquad t\sim L_{\rm sm}^{2}$\\[-0.2em]\scriptsize 结构式/示意：用于说明变量和极限，不冒充单篇原文公式。 \\
物理校验 & 流时间与平滑长度满足平方关系；t→0 回到原始场，有限 t 则抑制紫外模式。 \\
本组论文 & 平凡化映射威尔逊流与HMC、威尔逊流的性质与应用、梯度流的微扰分析、手征对称性与杨米尔斯梯度流、杨米尔斯梯度流提取能动量张量、准部分子分布与梯度流、含费米子场的格点能动量张量、梯度流形式的双圈能动量张量、梯度流高阶计算的技术与结果、局域涂抹算符乘积展开、梯度流中的非光锥威尔逊线算符、任意阶部分子分布矩 \\
\bottomrule
\end{tabularx}
\endgroup
  \source{综合来源：本组论文的 \detokenize{main.tex} 与章节源；结构式为本报告的主题示意。}
\end{frame}

%% frame_manifest|frame_id=T-008|paper_id=GLOBAL|claim_id=T-LaMET、准 PDF 与赝 PDF|visual_type=topic-map-formula|width_budget=0.98linewidth|height_budget=0.86textheight|min_font_pt=10|split_allowed=no|expected_page=8|risks=title-table-source-footline
\begin{frame}[t]{LaMET、准分布与赝分布：4 篇论文围绕同一个中间量链条展开}
  \small
\begingroup
\renewcommand{\arraystretch}{0.84}
\begin{tabularx}{0.98\linewidth}{@{}p{2.55cm}>{\raggedright\arraybackslash}X@{}}
\toprule
共同问题 & 以大动量强子和等时空间关联替代直接光锥算符，通过有效场论匹配与演化提取部分子信息。 \\
代表结构式 & $\widetilde q(x,P_z,\mu)=\int dy\,C(x,y,\mu/P_z)q(y,\mu)+\mathcal O(\Lambda_{\rm QCD}^{2}/P_z^{2})$\\[-0.2em]\scriptsize 结构式/示意：用于说明变量和极限，不冒充单篇原文公式。 \\
物理校验 & 大动量极限 P\_z→∞ 压低幂修正；有限动量、有限间距和核子激发态是必须单独控制的误差源。 \\
本组论文 & 欧氏格点上的部分子物理、大动量有效理论与部分子物理、部分子伪分布的格点探索、伪部分子分布的单圈演化 \\
\bottomrule
\end{tabularx}
\endgroup
  \source{综合来源：本组论文的 \detokenize{main.tex} 与章节源；结构式为本报告的主题示意。}
\end{frame}

%% frame_manifest|frame_id=T-009|paper_id=GLOBAL|claim_id=T-重正化与匹配|visual_type=topic-map-formula|width_budget=0.98linewidth|height_budget=0.86textheight|min_font_pt=10|split_allowed=no|expected_page=9|risks=title-table-source-footline
\begin{frame}[t]{非微扰重正化、OPE 与匹配：19 篇论文围绕同一个中间量链条展开}
  \small
\begingroup
\renewcommand{\arraystretch}{0.84}
\begin{tabularx}{0.98\linewidth}{@{}p{2.55cm}>{\raggedright\arraybackslash}X@{}}
\toprule
共同问题 & 处理 Wilson 线自能、线性发散、算符混合和方案转换，建立格点裸矩阵元与连续方案物理分布之间的桥梁。 \\
代表结构式 & $O^{R}(z,\mu)=Z(z,\mu,a)\,O^{\rm bare}(z,a)$\\[-0.2em]\scriptsize 结构式/示意：用于说明变量和极限，不冒充单篇原文公式。 \\
物理校验 & 重正化因子抵消截止依赖；a→0 和短距离 z→0 极限必须与匹配阶数、方案和尺度保持一致。 \\
本组论文 & 夸克场产生算符构造、连续威尔逊圈的无穷N相变、准部分子分布单圈匹配、大动量有效理论再论、欧氏与光锥分布因子化定理、格点算符非微扰重正化一般方法、准部分子分布的可重正性、威尔逊线重正化改进准分布、大动量有效理论的重正化、非定域夸克双线性的非微扰重正化、准PDF完整非微扰重正化方案、准部分子算符乘法可重正性、大动量有效理论获取胶子分布、胶子伪分布的短距行为、准光前关联函数的自重正化、准光前关联的混合重正化方案、首个核子胶子部分子分布、准分布的幂修正与renormalon、同位旋矢量分布的重正化匹配系统分析 \\
\bottomrule
\end{tabularx}
\endgroup
  \source{综合来源：本组论文的 \detokenize{main.tex} 与章节源；结构式为本报告的主题示意。}
\end{frame}

%% frame_manifest|frame_id=T-010|paper_id=GLOBAL|claim_id=T-胶子、强子 PDF 与 TMD|visual_type=topic-map-formula|width_budget=0.98linewidth|height_budget=0.86textheight|min_font_pt=10|split_allowed=no|expected_page=10|risks=title-table-source-footline
\begin{frame}[t]{胶子、强子部分子分布与 TMD：9 篇论文围绕同一个中间量链条展开}
  \small
\begingroup
\renewcommand{\arraystretch}{0.84}
\begin{tabularx}{0.98\linewidth}{@{}p{2.55cm}>{\raggedright\arraybackslash}X@{}}
\toprule
共同问题 & 把上述格点工具应用到胶子、核子、$\pi$/K 介子和横向动量依赖结构，并与全局拟合或 Collins–Soper 演化衔接。 \\
代表结构式 & $\mathcal O_{\rm Euclid}\xrightarrow[\rm matching\,/\,evolution]{}f_{g/h}(x,\mu)\;{\rm or}\;F_{\rm TMD}(x,b_T,\mu,\zeta)$\\[-0.2em]\scriptsize 结构式/示意：用于说明变量和极限，不冒充单篇原文公式。 \\
物理校验 & 纵向分数、横向距离和重正化尺度是不同变量；不能把有限动量或有限 b\_T 的结果直接当作光锥分布。 \\
本组论文 & 准分布动量分布与伪分布、大动量有效理论综述、CT18全局分析、NNPDF17高精度数据部分子分布、微扰论中的涂抹准分布、$\pi$介子胶子部分子分布、K介子胶子分布初探、从准TMD波函数确定CS核、LaMET计算TMD软函数 \\
\bottomrule
\end{tabularx}
\endgroup
  \source{综合来源：本组论文的 \detokenize{main.tex} 与章节源；结构式为本报告的主题示意。}
\end{frame}

%% frame_manifest|frame_id=T-011|paper_id=GLOBAL|claim_id=T-机器学习采样与随机量化|visual_type=topic-map-formula|width_budget=0.98linewidth|height_budget=0.86textheight|min_font_pt=10|split_allowed=no|expected_page=11|risks=title-table-source-footline
\begin{frame}[t]{机器学习采样与随机量化：3 篇论文围绕同一个中间量链条展开}
  \small
\begingroup
\renewcommand{\arraystretch}{0.84}
\begin{tabularx}{0.98\linewidth}{@{}p{2.55cm}>{\raggedright\arraybackslash}X@{}}
\toprule
共同问题 & 以流模型、规范等变结构、扩散过程或随机量化改造格点场配置生成与采样，关注可逆性、规范对称性和分布一致性。 \\
代表结构式 & $z\sim p(z),\qquad U=f_\theta(z),\qquad p_\theta(U)\;\text{与目标系综比较}$\\[-0.2em]\scriptsize 结构式/示意：用于说明变量和极限，不冒充单篇原文公式。 \\
物理校验 & 规范不变性/等变性是结构约束；采样质量须由可观测量、接受率、自相关或分布距离验证，不能只看训练损失。 \\
本组论文 & 基于流的格点MCMC生成模型、规范等变的流采样、扩散模型即随机量化 \\
\bottomrule
\end{tabularx}
\endgroup
  \source{综合来源：本组论文的 \detokenize{main.tex} 与章节源；结构式为本报告的主题示意。}
\end{frame}

%% frame_manifest|frame_id=G-12|paper_id=GLOBAL|claim_id=G12|visual_type=risk-matrix|width_budget=0.98linewidth|height_budget=0.86textheight|min_font_pt=10|split_allowed=no|expected_page=12|risks=title-table-source-footline
\begin{frame}[t]{最需要警惕的不是页数，而是把不同极限下的量混为一谈}
  \small
  \begingroup
\renewcommand{\arraystretch}{0.84}
\begin{tabularx}{0.98\linewidth}{@{}p{2.55cm}>{\raggedright\arraybackslash}X@{}}
\toprule
紫外/线性发散 & 裸 Wilson 线和非局域算符必须先重正化；未经方案转换的矩阵元不能直接比较。 \\
大动量幂修正 & LaMET/准分布的 \(P_z\) 有限；外推到无限动量需要多个动量、匹配和稳定性证据。 \\
短距离/流时间 & 赝分布、梯度流和 OPE 的适用区间受 \(z\)、\(t\)、\(a\) 共同约束，不能只看单一变量。 \\
采样正确性 & 生成模型的损失下降不等于系综正确；必须回到规范不变可观测量、自相关和分布检验。 \\
\bottomrule
\end{tabularx}
\endgroup
  \source{综合判断：各主题论文的误差/讨论章节；具体数值与置信度仍以各篇源文件为准。}
\end{frame}

%% frame_manifest|frame_id=G-13|paper_id=GLOBAL|claim_id=G13|visual_type=reading-guide|width_budget=0.98linewidth|height_budget=0.86textheight|min_font_pt=10|split_allowed=no|expected_page=13|risks=title-table-source-footline
\begin{frame}[t]{阅读方式：先看每篇四页档案，再沿主题接口回到原文}
  \small
  \begingroup
\renewcommand{\arraystretch}{0.84}
\begin{tabularx}{0.98\linewidth}{@{}p{2.55cm}>{\raggedright\arraybackslash}X@{}}
\toprule
第 1 页 & 问题、对象、摘要摘录和论文定位；回答“论文试图测量/证明什么”。 \\
第 2 页 & 输入、算符/状态、方法链和结构式；回答“物理量如何从格点对象产生”。 \\
第 3 页 & 结果、验证和物理含义；直接引用源文件，未提供的数据明确降级为未验证。 \\
第 4 页 & 局限、跨论文接口、量纲/极限检查和完整来源；回答“下一步应核验什么”。 \\
\bottomrule
\end{tabularx}
\endgroup
  \source{报告结构：\detokenize{docs/superpowers/specs/2026-08-28-refer-papers-report-design.md}。}
\end{frame}

%% frame_manifest|frame_id=P01-01|paper_id=P01|claim_id=P01-定位|visual_type=paper-card|width_budget=0.98linewidth|height_budget=0.86textheight|min_font_pt=10|split_allowed=no|expected_page=14|risks=title-table-source-footline
\section{夸克禁闭}
\begin{frame}[t]{夸克禁闭：研究问题与论文定位}
  \fontsize{10}{10.8}\selectfont
\begingroup
\renewcommand{\arraystretch}{0.84}
\begin{tabularx}{0.98\linewidth}{@{}p{2.55cm}>{\raggedright\arraybackslash}X@{}}
\toprule
论文编号 & P01；主题：格点规范理论与禁闭 \\
书目信息 & Confinement of Quarks；books/复用(Wilson74)；英/中转排页数 30/24 \\
研究对象 & 0.86 本文定义了一种与施温格（Schwinger）类似的使夸克完全禁闭的机制， 该机制要求存在阿贝尔或非阿贝尔规范场。文中展示了如何在欧几里得时空的离散 格点上对规范场理论进行量子化，既保持精确的规范不变性，又把规范场视为角变……有夸克路径以及连接这些路径的所有（格点上的）曲面求和。 这一结构与强子的相对论性弦模型十分相似 \\
章节证据 & 绪论；夸克束缚机制；规范场的格点量子化；格点上的经典作用量；量子化；强耦合近似；弱耦合近似；相变 \\
证据状态 & \verified：摘要和章节来自现有 LaTeX；\unverified：没有论文 PDF/原始数据独立复核。 \\
\bottomrule
\end{tabularx}
\endgroup
  \source{索引：\detokenize{refer/papers/INDEX.md:1}；正文：\detokenize{refer/papers/夸克禁闭_latex/main.tex}；\detokenize{refer/papers/Confinement_of_quarks_latex/main.tex}}
\end{frame}

%% frame_manifest|frame_id=P01-02|paper_id=P01|claim_id=P01-方法|visual_type=method-flow-formula|width_budget=0.98linewidth|height_budget=0.86textheight|min_font_pt=10|split_allowed=no|expected_page=15|risks=title-table-source-footline
\begin{frame}[t]{夸克禁闭：理论结构与方法链}
  \fontsize{10}{10.8}\selectfont
\begingroup
\renewcommand{\arraystretch}{0.84}
\begin{tabularx}{0.98\linewidth}{@{}p{2.55cm}>{\raggedright\arraybackslash}X@{}}
\toprule
输入与状态 & 格点/解析输入 → 格点规范理论与禁闭中的算符、矩阵元或场配置；具体输入见源章节。 \\
章节路线 & 绪论；夸克束缚机制；规范场的格点量子化；格点上的经典作用量；量子化；强耦合近似；弱耦合近似；相变 \\
方法摘录 & sec:intro 夸克组元图像（quark-constituent picture）在共振态和深度非弹性电子、中微子过程 两方面所取得的成功，使我们很难相信夸克并不存在。问题在于夸克从未被观测到。 这表明夸克出于某种原因不能作为独立粒子出现在……点规范理论的强耦合展开。第五节对弱耦合作粗略讨论。 第六节简要讨论洛伦兹不变性问题以及与弦模型的关系 \\
结构式/示意 & $\langle W(C)\rangle\sim\exp[-\sigma A(C)]$\\[-0.2em]\scriptsize 该式只表达主题变量关系；论文特定推导以源文件为准。 \\
物理校验 & 量纲检查：$\sigma$ 具有面积倒数维度；大圈极限把面积律与线性势联系起来。 \\
\bottomrule
\end{tabularx}
\endgroup
  \source{方法证据：\detokenize{refer/papers/夸克禁闭_latex/chapters/section01.tex:4-87}；主题结构式为示意。}
\end{frame}

%% frame_manifest|frame_id=P01-03|paper_id=P01|claim_id=P01-结果|visual_type=result-evidence-table|width_budget=0.98linewidth|height_budget=0.86textheight|min_font_pt=10|split_allowed=no|expected_page=16|risks=title-table-source-footline
\begin{frame}[t]{夸克禁闭：结果、验证与物理含义}
  \fontsize{10}{10.8}\selectfont
\begingroup
\renewcommand{\arraystretch}{0.84}
\begin{tabularx}{0.98\linewidth}{@{}p{2.55cm}>{\raggedright\arraybackslash}X@{}}
\toprule
源文结果 & ……取值（公式）；令相互作用为（公式或图表位置）其中（公式）与自旋-自旋耦合有关，（公式）正比于外场。于是（公式或图表位置）其中（公式）表示对所有自旋的所有组态求和，而（公式或图表位置）在平均场近似中，假定与（公式）耦合的自旋（公式）可以用它们的平均值（公式）代替。结果，（公式）的公式化简为只对（公式）求和，即（公式或图表位置）其中（公式）是维数（通常为 3），而（公式或图表位……的哈密顿量形式由 Kogut 和 Susskind 给出。（公式）格点理论中夸克禁闭 的清晰综述见文献 20。强耦合规范理论与弦模型之间联系的另一种表述见文献 21 \\
验证状态 & 确证：源章节有结果/讨论摘录。 \\
库内可见证据 & 中文源约 1200 行；公式命中 64，图命中 20，表/表格命中 0。 \\
物理含义 & 在本主题共同链条中，本论文把上述输入推进到可比较的中间量或观测量；具体强度、误差和外推范围不超出源文档。 \\
未验证项 & 当前工作区没有论文 PDF、原始数据和独立复现脚本；因此数值复核、图像再现和统计显著性不作额外断言。 \\
\bottomrule
\end{tabularx}
\endgroup
  \source{结果证据：\detokenize{refer/papers/夸克禁闭_latex/chapters/section06.tex:4-139}；图表/公式计数由本报告生成器对中文源的静态统计得到。}
\end{frame}

%% frame_manifest|frame_id=P01-04|paper_id=P01|claim_id=P01-局限|visual_type=risk-relation-table|width_budget=0.98linewidth|height_budget=0.86textheight|min_font_pt=10|split_allowed=no|expected_page=17|risks=title-table-source-footline
\begin{frame}[t]{夸克禁闭：局限、跨论文接口与可复查入口}
  \fontsize{10}{10.8}\selectfont
\begingroup
\renewcommand{\arraystretch}{0.84}
\begin{tabularx}{0.98\linewidth}{@{}p{2.55cm}>{\raggedright\arraybackslash}X@{}}
\toprule
源文局限 & 需要回到结果/讨论/展望章节核对适用区间；本档案不把缺失的独立数据复核补成已证实结论。 \\
跨论文接口 & 为后续梯度流、重正化和部分子算符提供规范不变的格点底座。 \\
极限检查 & 量纲检查：$\sigma$ 具有面积倒数维度；大圈极限把面积律与线性势联系起来。 \\
下一步核验 & 补齐 PDF 或原始数据；按源文参数复现关键图表；比较不同动量、流时间、距离、格距或方案转换后的稳定性。 \\
完整来源 & \detokenize{refer/papers/夸克禁闭_latex/main.tex}；\detokenize{refer/papers/Confinement_of_quarks_latex/main.tex}；章节源：\detokenize{refer/papers/夸克禁闭_latex/chapters/*.tex}；英中目录：\detokenize{refer/papers/Confinement_of_quarks_latex}. \\
\bottomrule
\end{tabularx}
\endgroup
  \source{状态：\inferred 为主题接口推断；\unverified 为当前材料边界；完整文件清单见 manifest。}
\end{frame}

%% frame_manifest|frame_id=P02-01|paper_id=P02|claim_id=P02-定位|visual_type=paper-card|width_budget=0.98linewidth|height_budget=0.86textheight|min_font_pt=10|split_allowed=no|expected_page=18|risks=title-table-source-footline
\section{SU3链接变量解析涂抹}
\begin{frame}[t]{SU3链接变量解析涂抹：研究问题与论文定位}
  \fontsize{10}{10.8}\selectfont
\begingroup
\renewcommand{\arraystretch}{0.84}
\begin{tabularx}{0.98\linewidth}{@{}p{2.55cm}>{\raggedright\arraybackslash}X@{}}
\toprule
论文编号 & P02；主题：格点规范理论与禁闭 \\
书目信息 & Analytic Smearing of（公式）Link Variables in Lattice QCD；arXiv:hep-lat/0311018；英/中转排页数 15/12 \\
研究对象 & 本文提出并检验了一种格点 QCD 中链接变量的解析涂抹方法。该涂抹方案对链接 变量的可微性，使得对用这种涂抹链接构造的规范场作用量，可以采用基于分子 动力学演化的现代蒙特卡罗更新方法。在考察涂抹平均方格值与静态夸克-反夸克势 时，观察到与现行链接涂抹方法相比效果并无退化，只是发现对涂抹参数的敏感性 有所增大 \\
章节证据 & 引言；链接变量的解析涂抹；涂抹的实现；分子动力学演化；数值检验；结论 \\
证据状态 & \verified：摘要和章节来自现有 LaTeX；\unverified：没有论文 PDF/原始数据独立复核。 \\
\bottomrule
\end{tabularx}
\endgroup
  \source{索引：\detokenize{refer/papers/INDEX.md:2}；正文：\detokenize{refer/papers/SU3链接变量解析涂抹_latex/main.tex}；\detokenize{refer/papers/Analytic_Smearing_of_SU3_Link_Variables_latex/main.tex}}
\end{frame}

%% frame_manifest|frame_id=P02-02|paper_id=P02|claim_id=P02-方法|visual_type=method-flow-formula|width_budget=0.98linewidth|height_budget=0.86textheight|min_font_pt=10|split_allowed=no|expected_page=19|risks=title-table-source-footline
\begin{frame}[t]{SU3链接变量解析涂抹：理论结构与方法链}
  \fontsize{10}{10.8}\selectfont
\begingroup
\renewcommand{\arraystretch}{0.84}
\begin{tabularx}{0.98\linewidth}{@{}p{2.55cm}>{\raggedright\arraybackslash}X@{}}
\toprule
输入与状态 & 格点/解析输入 → 格点规范理论与禁闭中的算符、矩阵元或场配置；具体输入见源章节。 \\
章节路线 & 引言；链接变量的解析涂抹；涂抹的实现；分子动力学演化；数值检验；结论 \\
方法摘录 & 格点 QCD 中从欧几里得空间关联函数的蒙特卡罗估计提取强子质量与矩阵元时， 如果所用算符对目标态耦合更强、对较高能级的污染态耦合更弱，则提取会更加 可靠和精确。对于包含胶子的态，在格点 QCD 中构造此类算符的一个关键要素是 链接变量涂抹（s……变量构造的格点作用量与算符受辐射修正的影响可能要小得多，因为通常占主导 地位的大蝌蚪图贡献被急剧削减 \\
结构式/示意 & $\langle W(C)\rangle\sim\exp[-\sigma A(C)]$\\[-0.2em]\scriptsize 该式只表达主题变量关系；论文特定推导以源文件为准。 \\
物理校验 & 量纲检查：$\sigma$ 具有面积倒数维度；大圈极限把面积律与线性势联系起来。 \\
\bottomrule
\end{tabularx}
\endgroup
  \source{方法证据：\detokenize{refer/papers/SU3链接变量解析涂抹_latex/chapters/section01.tex:1-48}；主题结构式为示意。}
\end{frame}

%% frame_manifest|frame_id=P02-03|paper_id=P02|claim_id=P02-结果|visual_type=result-evidence-table|width_budget=0.98linewidth|height_budget=0.86textheight|min_font_pt=10|split_allowed=no|expected_page=20|risks=title-table-source-footline
\begin{frame}[t]{SU3链接变量解析涂抹：结果、验证与物理含义}
  \fontsize{10}{10.8}\selectfont
\begingroup
\renewcommand{\arraystretch}{0.84}
\begin{tabularx}{0.98\linewidth}{@{}p{2.55cm}>{\raggedright\arraybackslash}X@{}}
\toprule
源文结果 & sec:conclude 链接变量光滑化是构造胶子算符的关键要素，这类算符与理论高频模式的混合 被大幅压低。链接涂抹在改进格点作用量的构造中也正发挥着日益重要的作用。 标准涂抹方案相对底层链接变量缺乏可微性，阻碍了基于分子动力学演化的高效 蒙特卡罗更新方法的使用。 本文提出并检验了一种规避这些问题的链接涂抹方法。该链接涂抹方法在整个 有限复平面上处处解析，并以保持在（公式）内的……件设计自然而高效。利用涂抹平均方格值以及与静态夸克-反夸克势相联系的 有效能量，我们表明：与现行的链接涂抹方法相比效果并无退化，只是发现对涂抹 参数的敏感性有所增大 \\
验证状态 & 确证：源章节有结果/讨论摘录。 \\
库内可见证据 & 中文源约 907 行；公式命中 50，图命中 8，表/表格命中 0。 \\
物理含义 & 在本主题共同链条中，本论文把上述输入推进到可比较的中间量或观测量；具体强度、误差和外推范围不超出源文档。 \\
未验证项 & 当前工作区没有论文 PDF、原始数据和独立复现脚本；因此数值复核、图像再现和统计显著性不作额外断言。 \\
\bottomrule
\end{tabularx}
\endgroup
  \source{结果证据：\detokenize{refer/papers/SU3链接变量解析涂抹_latex/chapters/section06.tex:1-17}；图表/公式计数由本报告生成器对中文源的静态统计得到。}
\end{frame}

%% frame_manifest|frame_id=P02-04|paper_id=P02|claim_id=P02-局限|visual_type=risk-relation-table|width_budget=0.98linewidth|height_budget=0.86textheight|min_font_pt=10|split_allowed=no|expected_page=21|risks=title-table-source-footline
\begin{frame}[t]{SU3链接变量解析涂抹：局限、跨论文接口与可复查入口}
  \fontsize{10}{10.8}\selectfont
\begingroup
\renewcommand{\arraystretch}{0.84}
\begin{tabularx}{0.98\linewidth}{@{}p{2.55cm}>{\raggedright\arraybackslash}X@{}}
\toprule
源文局限 & 需要回到结果/讨论/展望章节核对适用区间；本档案不把缺失的独立数据复核补成已证实结论。 \\
跨论文接口 & 为后续梯度流、重正化和部分子算符提供规范不变的格点底座。 索引备注：文选卡片ID有误(0307022)已核实纠正 \\
极限检查 & 量纲检查：$\sigma$ 具有面积倒数维度；大圈极限把面积律与线性势联系起来。 \\
下一步核验 & 补齐 PDF 或原始数据；按源文参数复现关键图表；比较不同动量、流时间、距离、格距或方案转换后的稳定性。 \\
完整来源 & \detokenize{refer/papers/SU3链接变量解析涂抹_latex/main.tex}；\detokenize{refer/papers/Analytic_Smearing_of_SU3_Link_Variables_latex/main.tex}；章节源：\detokenize{refer/papers/SU3链接变量解析涂抹_latex/chapters/*.tex}；英中目录：\detokenize{refer/papers/Analytic_Smearing_of_SU3_Link_Variables_latex}. \\
\bottomrule
\end{tabularx}
\endgroup
  \source{状态：\inferred 为主题接口推断；\unverified 为当前材料边界；完整文件清单见 manifest。}
\end{frame}

%% frame_manifest|frame_id=P03-01|paper_id=P03|claim_id=P03-定位|visual_type=paper-card|width_budget=0.98linewidth|height_budget=0.86textheight|min_font_pt=10|split_allowed=no|expected_page=22|risks=title-table-source-footline
\section{夸克场产生算符构造}
\begin{frame}[t]{夸克场产生算符构造：研究问题与论文定位}
  \fontsize{10}{10.8}\selectfont
\begingroup
\renewcommand{\arraystretch}{0.84}
\begin{tabularx}{0.98\linewidth}{@{}p{2.55cm}>{\raggedright\arraybackslash}X@{}}
\toprule
论文编号 & P03；主题：非微扰重正化、OPE 与匹配 \\
书目信息 & A novel quark-field creation operator construction for hadronic physics in lattice QCD；arXiv:0905.2160；英/中转排页数 17/14 \\
研究对象 & 本文定义了一种新的夸克场涂抹（smearing）算法，它使大范围强子关联函数的高效计算成为可能。 该技术施加一个低秩算符来定义光滑的场，供强子产生算符使用。 所得光滑场空间足够小，因而能够在合理的计算代价下 精确计算出约化夸克传播……i） 计算，而无需新的格点狄拉克算符求逆。 该方法在带有轻动力学夸克的现实格点尺寸上进行了检验 \\
章节证据 & 引言；算符构造 sec:theory；介子两点关联函数；重子两点关联函数；介子三点关联函数；重子三点关联函数；检验结果 sec:results；蒸馏算符的轮廓分布；介子关联函数；重子关联函数；多介子关联函数；讨论与未来方向 sec:discussion；总结 \\
证据状态 & \verified：摘要和章节来自现有 LaTeX；\unverified：没有论文 PDF/原始数据独立复核。 \\
\bottomrule
\end{tabularx}
\endgroup
  \source{索引：\detokenize{refer/papers/INDEX.md:3}；正文：\detokenize{refer/papers/夸克场产生算符构造_latex/main.tex}；\detokenize{refer/papers/Novel_Quark_Field_Creation_Operator_latex/main.tex}}
\end{frame}

%% frame_manifest|frame_id=P03-02|paper_id=P03|claim_id=P03-方法|visual_type=method-flow-formula|width_budget=0.98linewidth|height_budget=0.86textheight|min_font_pt=10|split_allowed=no|expected_page=23|risks=title-table-source-footline
\begin{frame}[t]{夸克场产生算符构造：理论结构与方法链}
  \fontsize{10}{10.8}\selectfont
\begingroup
\renewcommand{\arraystretch}{0.84}
\begin{tabularx}{0.98\linewidth}{@{}p{2.55cm}>{\raggedright\arraybackslash}X@{}}
\toprule
输入与状态 & 格点/解析输入 → 非微扰重正化、OPE 与匹配中的算符、矩阵元或场配置；具体输入见源章节。 \\
章节路线 & 引言；算符构造 sec:theory；介子两点关联函数；重子两点关联函数；介子三点关联函数；重子三点关联函数；检验结果 sec:results；蒸馏算符的轮廓分布；介子关联函数；重子关联函数；多介子关联函数；讨论与未来方向 sec:discussion；总结 \\
方法摘录 & 格点计算的目标之一是仅从 QCD 拉格朗日量出发， 预言禁闭夸克与胶子的低能强子谱。这一谱学进路必须有测量具有所研究 量子数的场算符两点关联函数的手段。对 QCD 谱的完整理解 还必须包括介子与重子的激发态以及奇特态。 把散射实验中看到的共振态……给出该方法有效性的初步检验结果； 第（交叉引用）节讨论该方法引发的实际问题， 并概述若干未来研究方向 \\
结构式/示意 & $O^{R}(z,\mu)=Z(z,\mu,a)\,O^{\rm bare}(z,a)$\\[-0.2em]\scriptsize 该式只表达主题变量关系；论文特定推导以源文件为准。 \\
物理校验 & 重正化因子抵消截止依赖；a→0 和短距离 z→0 极限必须与匹配阶数、方案和尺度保持一致。 \\
\bottomrule
\end{tabularx}
\endgroup
  \source{方法证据：\detokenize{refer/papers/夸克场产生算符构造_latex/chapters/section01.tex:1-56}；主题结构式为示意。}
\end{frame}

%% frame_manifest|frame_id=P03-03|paper_id=P03|claim_id=P03-结果|visual_type=result-evidence-table|width_budget=0.98linewidth|height_budget=0.86textheight|min_font_pt=10|split_allowed=no|expected_page=24|risks=title-table-source-footline
\begin{frame}[t]{夸克场产生算符构造：结果、验证与物理含义}
  \fontsize{10}{10.8}\selectfont
\begingroup
\renewcommand{\arraystretch}{0.84}
\begin{tabularx}{0.98\linewidth}{@{}p{2.55cm}>{\raggedright\arraybackslash}X@{}}
\toprule
源文结果 & 本文提出了一种简单的新夸克场涂抹方法，并用带（公式）味 动力学夸克的 QCD 规范场做了数值检验。 强子产生算符的构造存在很大的自由度。 本工作利用这一点定义了一种算法： 在每个时间片上把一个低秩投影算符作用于路径积分的夸克场， 为后续算符构造定义光滑模式。 这使得测量蒸馏强子关联函数所需的关于夸克传播的全部信息 的计算变得可以负担。一旦传播矩阵已经求出， 任意产生与湮灭算符选……，本文概述了若干研究方向。 合作组发展工作的首要焦点是以最优方式把关联函数的随机评估 纳入进来，结果将很快发表。合作组将在近期的一系列谱学测定中 使用本文概述的技术 \\
验证状态 & 确证：源章节有结果/讨论摘录。 \\
库内可见证据 & 中文源约 942 行；公式命中 30，图命中 19，表/表格命中 2。 \\
物理含义 & 在本主题共同链条中，本论文把上述输入推进到可比较的中间量或观测量；具体强度、误差和外推范围不超出源文档。 \\
未验证项 & 当前工作区没有论文 PDF、原始数据和独立复现脚本；因此数值复核、图像再现和统计显著性不作额外断言。 \\
\bottomrule
\end{tabularx}
\endgroup
  \source{结果证据：\detokenize{refer/papers/夸克场产生算符构造_latex/chapters/section05.tex:1-27}；图表/公式计数由本报告生成器对中文源的静态统计得到。}
\end{frame}

%% frame_manifest|frame_id=P03-04|paper_id=P03|claim_id=P03-局限|visual_type=risk-relation-table|width_budget=0.98linewidth|height_budget=0.86textheight|min_font_pt=10|split_allowed=no|expected_page=25|risks=title-table-source-footline
\begin{frame}[t]{夸克场产生算符构造：局限、跨论文接口与可复查入口}
  \fontsize{10}{10.8}\selectfont
\begingroup
\renewcommand{\arraystretch}{0.84}
\begin{tabularx}{0.98\linewidth}{@{}p{2.55cm}>{\raggedright\arraybackslash}X@{}}
\toprule
源文局限 & 需要回到结果/讨论/展望章节核对适用区间；本档案不把缺失的独立数据复核补成已证实结论。 \\
跨论文接口 & 是准/赝分布从格点数据走向 MS-bar 或其他连续方案的必要中间层。 \\
极限检查 & 重正化因子抵消截止依赖；a→0 和短距离 z→0 极限必须与匹配阶数、方案和尺度保持一致。 \\
下一步核验 & 补齐 PDF 或原始数据；按源文参数复现关键图表；比较不同动量、流时间、距离、格距或方案转换后的稳定性。 \\
完整来源 & \detokenize{refer/papers/夸克场产生算符构造_latex/main.tex}；\detokenize{refer/papers/Novel_Quark_Field_Creation_Operator_latex/main.tex}；章节源：\detokenize{refer/papers/夸克场产生算符构造_latex/chapters/*.tex}；英中目录：\detokenize{refer/papers/Novel_Quark_Field_Creation_Operator_latex}. \\
\bottomrule
\end{tabularx}
\endgroup
  \source{状态：\inferred 为主题接口推断；\unverified 为当前材料边界；完整文件清单见 manifest。}
\end{frame}

%% frame_manifest|frame_id=P04-01|paper_id=P04|claim_id=P04-定位|visual_type=paper-card|width_budget=0.98linewidth|height_budget=0.86textheight|min_font_pt=10|split_allowed=no|expected_page=26|risks=title-table-source-footline
\section{高动量强子新夸克涂抹}
\begin{frame}[t]{高动量强子新夸克涂抹：研究问题与论文定位}
  \fontsize{10}{10.8}\selectfont
\begingroup
\renewcommand{\arraystretch}{0.84}
\begin{tabularx}{0.98\linewidth}{@{}p{2.55cm}>{\raggedright\arraybackslash}X@{}}
\toprule
论文编号 & P04；主题：格点规范理论与禁闭 \\
书目信息 & Novel Quark Smearing High Momenta；arXiv:1602.05525；英/中转排页数 27/23 \\
研究对象 & —————- quote \\
章节证据 & 引言；动量涂抹：基本思想；动量涂抹与强子两点函数；Wuppertal 涂抹；两点函数的构造；动量（Wuppertal）涂抹；自由场研究；非迭代（动量）涂抹；格点系综与色散关系；结果；涂抹的实……验；助推（动量）涂抹的检验；与色散关系的比较；结论；附录：助推（动量）涂抹 \\
证据状态 & \verified：摘要和章节来自现有 LaTeX；\unverified：没有论文 PDF/原始数据独立复核。 \\
\bottomrule
\end{tabularx}
\endgroup
  \source{索引：\detokenize{refer/papers/INDEX.md:4}；正文：\detokenize{refer/papers/高动量强子新夸克涂抹_latex/main.tex}；\detokenize{refer/papers/Novel_Quark_Smearing_High_Momenta_latex/main.tex}}
\end{frame}

%% frame_manifest|frame_id=P04-02|paper_id=P04|claim_id=P04-方法|visual_type=method-flow-formula|width_budget=0.98linewidth|height_budget=0.86textheight|min_font_pt=10|split_allowed=no|expected_page=27|risks=title-table-source-footline
\begin{frame}[t]{高动量强子新夸克涂抹：理论结构与方法链}
  \fontsize{10}{10.8}\selectfont
\begingroup
\renewcommand{\arraystretch}{0.84}
\begin{tabularx}{0.98\linewidth}{@{}p{2.55cm}>{\raggedright\arraybackslash}X@{}}
\toprule
输入与状态 & 格点/解析输入 → 格点规范理论与禁闭中的算符、矩阵元或场配置；具体输入见源章节。 \\
章节路线 & 引言；动量涂抹：基本思想；动量涂抹与强子两点函数；Wuppertal 涂抹；两点函数的构造；动量（Wuppertal）涂抹；自由场研究；非迭代（动量）涂抹；格点系综与色散关系；结果；涂抹的实……验；助推（动量）涂抹的检验；与色散关系的比较；结论；附录：助推（动量）涂抹 \\
方法摘录 & 格点 QCD 模拟所预言的、与粒子和强子物理唯象学相关的可观测量正日益增多。 这些结果通常从大欧氏时间间隔处（公式）点函数的期望值中提取。 由于这些函数随时间衰减，当时间间隔取得很大时， 统计噪声与信号之比呈指数式增大（零动量赝标介子是一个显著……可行性， 优化涂抹参数并对新旧方法进行比较。最后， 我们研究（公式）介子和核子的色散关系，然后作结论 \\
结构式/示意 & $\langle W(C)\rangle\sim\exp[-\sigma A(C)]$\\[-0.2em]\scriptsize 该式只表达主题变量关系；论文特定推导以源文件为准。 \\
物理校验 & 量纲检查：$\sigma$ 具有面积倒数维度；大圈极限把面积律与线性势联系起来。 \\
\bottomrule
\end{tabularx}
\endgroup
  \source{方法证据：\detokenize{refer/papers/高动量强子新夸克涂抹_latex/chapters/section01.tex:1-105}；主题结构式为示意。}
\end{frame}

%% frame_manifest|frame_id=P04-03|paper_id=P04|claim_id=P04-结果|visual_type=result-evidence-table|width_budget=0.98linewidth|height_budget=0.86textheight|min_font_pt=10|split_allowed=no|expected_page=28|risks=title-table-source-footline
\begin{frame}[t]{高动量强子新夸克涂抹：结果、验证与物理含义}
  \fontsize{10}{10.8}\selectfont
\begingroup
\renewcommand{\arraystretch}{0.84}
\begin{tabularx}{0.98\linewidth}{@{}p{2.55cm}>{\raggedright\arraybackslash}X@{}}
\toprule
源文结果 & 在许多格点规范理论应用中需要携带高动量的强子。由于（公式）点函数的相对误差随欧氏时间距离指数增长以及基态采样的减弱，高动量以往非常困难 甚至无法实现。在第（交叉引用）节 我们引入了一类新的夸克涂抹方法用于构造强子内插算符， 它们应对并显著缓解了这些问题。这些方法的一个具体实现——实现平凡且计算开销极小——是动量 Wuppertal 涂抹，定义见式（交叉引用）。 我们在第（交叉引……的格点模拟。对三点函数而言，这些方法 甚至可能允许虚度（公式）：把源动量（公式）切换为汇动量（公式）。 用新涂抹计算这些量的工作视可观测量而定正在进行或已在计划之中 \\
验证状态 & 确证：源章节有结果/讨论摘录。 \\
库内可见证据 & 中文源约 2244 行；公式命中 57，图命中 22，表/表格命中 0。 \\
物理含义 & 在本主题共同链条中，本论文把上述输入推进到可比较的中间量或观测量；具体强度、误差和外推范围不超出源文档。 \\
未验证项 & 当前工作区没有论文 PDF、原始数据和独立复现脚本；因此数值复核、图像再现和统计显著性不作额外断言。 \\
\bottomrule
\end{tabularx}
\endgroup
  \source{结果证据：\detokenize{refer/papers/高动量强子新夸克涂抹_latex/chapters/section06.tex:1-40}；图表/公式计数由本报告生成器对中文源的静态统计得到。}
\end{frame}

%% frame_manifest|frame_id=P04-04|paper_id=P04|claim_id=P04-局限|visual_type=risk-relation-table|width_budget=0.98linewidth|height_budget=0.86textheight|min_font_pt=10|split_allowed=no|expected_page=29|risks=title-table-source-footline
\begin{frame}[t]{高动量强子新夸克涂抹：局限、跨论文接口与可复查入口}
  \fontsize{10}{10.8}\selectfont
\begingroup
\renewcommand{\arraystretch}{0.84}
\begin{tabularx}{0.98\linewidth}{@{}p{2.55cm}>{\raggedright\arraybackslash}X@{}}
\toprule
源文局限 & 需要回到结果/讨论/展望章节核对适用区间；本档案不把缺失的独立数据复核补成已证实结论。 \\
跨论文接口 & 为后续梯度流、重正化和部分子算符提供规范不变的格点底座。 \\
极限检查 & 量纲检查：$\sigma$ 具有面积倒数维度；大圈极限把面积律与线性势联系起来。 \\
下一步核验 & 补齐 PDF 或原始数据；按源文参数复现关键图表；比较不同动量、流时间、距离、格距或方案转换后的稳定性。 \\
完整来源 & \detokenize{refer/papers/高动量强子新夸克涂抹_latex/main.tex}；\detokenize{refer/papers/Novel_Quark_Smearing_High_Momenta_latex/main.tex}；章节源：\detokenize{refer/papers/高动量强子新夸克涂抹_latex/chapters/*.tex}；英中目录：\detokenize{refer/papers/Novel_Quark_Smearing_High_Momenta_latex}. \\
\bottomrule
\end{tabularx}
\endgroup
  \source{状态：\inferred 为主题接口推断；\unverified 为当前材料边界；完整文件清单见 manifest。}
\end{frame}

%% frame_manifest|frame_id=P05-01|paper_id=P05|claim_id=P05-定位|visual_type=paper-card|width_budget=0.98linewidth|height_budget=0.86textheight|min_font_pt=10|split_allowed=no|expected_page=30|risks=title-table-source-footline
\section{平凡化映射威尔逊流与HMC}
\begin{frame}[t]{平凡化映射威尔逊流与HMC：研究问题与论文定位}
  \fontsize{10}{10.8}\selectfont
\begingroup
\renewcommand{\arraystretch}{0.84}
\begin{tabularx}{0.98\linewidth}{@{}p{2.55cm}>{\raggedright\arraybackslash}X@{}}
\toprule
论文编号 & P05；主题：梯度流、涂抹与能动量张量 \\
书目信息 & Trivializing Maps Wilson Flow HMC；arXiv:0907.5491；英/中转排页数 20/18 \\
研究对象 & 在格点规范理论中，存在一类场变换，能把该理论映射到基本场变量完全彼此退耦的平凡理论。 这类映射可以通过对场空间中某些流方程作积分而系统地构造出来。 本文较为细致地给出了这一构造，并建议将威尔逊流 （它为 Wilson 规范作用量生成近似的平凡化映射） 与 HMC 模拟算法结合起来， 以提高格点 QCD 模拟的效率 \\
章节证据 & 引言；场变换；场空间；右不变微分算符；雅可比矩阵；HMC 算法的变换行为；由流方程生成的变换；场空间中的流；积分后的变换；平凡化映射；平凡化流；平凡化流的存在性；（公式）的幂级数展开；Wil……微性；场流形上的 Sobolev 空间；（公式）之逆的性质；（公式）的展开 \\
证据状态 & \verified：摘要和章节来自现有 LaTeX；\unverified：没有论文 PDF/原始数据独立复核。 \\
\bottomrule
\end{tabularx}
\endgroup
  \source{索引：\detokenize{refer/papers/INDEX.md:5}；正文：\detokenize{refer/papers/平凡化映射威尔逊流与HMC_latex/main.tex}；\detokenize{refer/papers/Trivializing_Maps_Wilson_Flow_HMC_latex/main.tex}}
\end{frame}

%% frame_manifest|frame_id=P05-02|paper_id=P05|claim_id=P05-方法|visual_type=method-flow-formula|width_budget=0.98linewidth|height_budget=0.86textheight|min_font_pt=10|split_allowed=no|expected_page=31|risks=title-table-source-footline
\begin{frame}[t]{平凡化映射威尔逊流与HMC：理论结构与方法链}
  \fontsize{10}{10.8}\selectfont
\begingroup
\renewcommand{\arraystretch}{0.84}
\begin{tabularx}{0.98\linewidth}{@{}p{2.55cm}>{\raggedright\arraybackslash}X@{}}
\toprule
输入与状态 & 格点/解析输入 → 梯度流、涂抹与能动量张量中的算符、矩阵元或场配置；具体输入见源章节。 \\
章节路线 & 引言；场变换；场空间；右不变微分算符；雅可比矩阵；HMC 算法的变换行为；由流方程生成的变换；场空间中的流；积分后的变换；平凡化映射；平凡化流；平凡化流的存在性；（公式）的幂级数展开；Wil……微性；场流形上的 Sobolev 空间；（公式）之逆的性质；（公式）的展开 \\
方法摘录 & 作为说明，现在对方块作用量 [引文]（公式或图表位置）（其中（公式）为逆规范耦合） 显式算出级数 (4.11) 的前两项。 利用完备性关系 (A.5) 作一简短计算， 可知领头项为（公式或图表位置）精确到这一阶，并且在时间参数（公式）的重新标度……格点上的圈数所补偿 （对于（公式）、（公式）、（公式）， 圈数分别等于（公式）、（公式）和（公式）） \\
结构式/示意 & $\partial_t B_\mu(t,x)=D_\nu G_{\nu\mu}(t,x),\qquad t\sim L_{\rm sm}^{2}$\\[-0.2em]\scriptsize 该式只表达主题变量关系；论文特定推导以源文件为准。 \\
物理校验 & 流时间与平滑长度满足平方关系；t→0 回到原始场，有限 t 则抑制紫外模式。 \\
\bottomrule
\end{tabularx}
\endgroup
  \source{方法证据：\detokenize{refer/papers/平凡化映射威尔逊流与HMC_latex/chapters/section04.tex:136-232}；主题结构式为示意。}
\end{frame}

%% frame_manifest|frame_id=P05-03|paper_id=P05|claim_id=P05-结果|visual_type=result-evidence-table|width_budget=0.98linewidth|height_budget=0.86textheight|min_font_pt=10|split_allowed=no|expected_page=32|risks=title-table-source-footline
\begin{frame}[t]{平凡化映射威尔逊流与HMC：结果、验证与物理含义}
  \fontsize{10}{10.8}\selectfont
\begingroup
\renewcommand{\arraystretch}{0.84}
\begin{tabularx}{0.98\linewidth}{@{}p{2.55cm}>{\raggedright\arraybackslash}X@{}}
\toprule
源文结果 & 本文以一个具体应用为目标讨论了 场空间中的平凡化映射与流。 然而，其底层概念相当普遍， 例如在严格的构造性工作、数值微扰论、 或与可重整化的光滑化技术相结合等方面， 可能还有其他用途。 所提议的威尔逊流与 HMC 算法的组合， 是否确实比迄今使用的模拟算法 更高效地对格点 QCD 的拓扑扇区采样，仍有待确定。 当格距越取越小时， 夸克场最终可能不得不被纳入流中， 或许还要包括第 4 节讨论的次领头阶修正。 无论如何，存在物质场时的平凡化映射 本身就是一个值得研究的课题 \\
验证状态 & 确证：源章节有结果/讨论摘录。 \\
库内可见证据 & 中文源约 1493 行；公式命中 111，图命中 6，表/表格命中 0。 \\
物理含义 & 在本主题共同链条中，本论文把上述输入推进到可比较的中间量或观测量；具体强度、误差和外推范围不超出源文档。 \\
未验证项 & 当前工作区没有论文 PDF、原始数据和独立复现脚本；因此数值复核、图像再现和统计显著性不作额外断言。 \\
\bottomrule
\end{tabularx}
\endgroup
  \source{结果证据：\detokenize{refer/papers/平凡化映射威尔逊流与HMC_latex/chapters/section07.tex:1-18}；图表/公式计数由本报告生成器对中文源的静态统计得到。}
\end{frame}

%% frame_manifest|frame_id=P05-04|paper_id=P05|claim_id=P05-局限|visual_type=risk-relation-table|width_budget=0.98linewidth|height_budget=0.86textheight|min_font_pt=10|split_allowed=no|expected_page=33|risks=title-table-source-footline
\begin{frame}[t]{平凡化映射威尔逊流与HMC：局限、跨论文接口与可复查入口}
  \fontsize{10}{10.8}\selectfont
\begingroup
\renewcommand{\arraystretch}{0.84}
\begin{tabularx}{0.98\linewidth}{@{}p{2.55cm}>{\raggedright\arraybackslash}X@{}}
\toprule
源文局限 & 需要回到结果/讨论/展望章节核对适用区间；本档案不把缺失的独立数据复核补成已证实结论。 \\
跨论文接口 & 把连续极限、微扰有限性和能动量张量重构连接到可计算的格点观测量。 \\
极限检查 & 流时间与平滑长度满足平方关系；t→0 回到原始场，有限 t 则抑制紫外模式。 \\
下一步核验 & 补齐 PDF 或原始数据；按源文参数复现关键图表；比较不同动量、流时间、距离、格距或方案转换后的稳定性。 \\
完整来源 & \detokenize{refer/papers/平凡化映射威尔逊流与HMC_latex/main.tex}；\detokenize{refer/papers/Trivializing_Maps_Wilson_Flow_HMC_latex/main.tex}；章节源：\detokenize{refer/papers/平凡化映射威尔逊流与HMC_latex/chapters/*.tex}；英中目录：\detokenize{refer/papers/Trivializing_Maps_Wilson_Flow_HMC_latex}. \\
\bottomrule
\end{tabularx}
\endgroup
  \source{状态：\inferred 为主题接口推断；\unverified 为当前材料边界；完整文件清单见 manifest。}
\end{frame}

%% frame_manifest|frame_id=P06-01|paper_id=P06|claim_id=P06-定位|visual_type=paper-card|width_budget=0.98linewidth|height_budget=0.86textheight|min_font_pt=10|split_allowed=no|expected_page=34|risks=title-table-source-footline
\section{连续威尔逊圈的无穷N相变}
\begin{frame}[t]{连续威尔逊圈的无穷N相变：研究问题与论文定位}
  \fontsize{10}{10.8}\selectfont
\begingroup
\renewcommand{\arraystretch}{0.84}
\begin{tabularx}{0.98\linewidth}{@{}p{2.55cm}>{\raggedright\arraybackslash}X@{}}
\toprule
论文编号 & P06；主题：非微扰重正化、OPE 与匹配 \\
书目信息 & Infinite（公式）phase transitions in continuum Wilson loop operators Published in JHEP 03 (2006) 064; arXiv:hep-th/0601210; DOI: 10.1088/1126-6708/2006/03/064.；arXiv:hep-th/0601210；英/中转排页数 30/25 \\
研究对象 & 我们在四维格点上定义了光滑化（涂抹）的威尔逊圈算符，并数值地验证它们具有有限而非平庸的连续统极限。这些连续统算符保持其作为幺正矩阵的品格，并在无穷大（公式）处经历一次相变；当定义性的光滑圈从小尺寸被胀大到大尺寸时，本征值分布在其谱……相似性表明，我们找到的这些相变只影响本征值分布，而威尔逊圈算符有限次幂的迹在标度变换下保持光滑 \\
章节证据 & 引言；格点研究概述；涂抹的格点威尔逊圈；格点大（公式）相变在连续统极限中存活；涂抹消除紫外发散；格点研究的数值细节；（公式）的确定；标度变换下相变两侧的研究；（公式）的一般特征；连续统的观点……普适类的一个猜想；总结与讨论；与其他工作的联系；本文主要的新思想；未来计划 \\
证据状态 & \verified：摘要和章节来自现有 LaTeX；\unverified：没有论文 PDF/原始数据独立复核。 \\
\bottomrule
\end{tabularx}
\endgroup
  \source{索引：\detokenize{refer/papers/INDEX.md:6}；正文：\detokenize{refer/papers/连续威尔逊圈的无穷N相变_latex/main.tex}；\detokenize{refer/papers/Infinite_N_Phase_Transitions_Wilson_Loops_latex/main.tex}}
\end{frame}

%% frame_manifest|frame_id=P06-02|paper_id=P06|claim_id=P06-方法|visual_type=method-flow-formula|width_budget=0.98linewidth|height_budget=0.86textheight|min_font_pt=10|split_allowed=no|expected_page=35|risks=title-table-source-footline
\begin{frame}[t]{连续威尔逊圈的无穷N相变：理论结构与方法链}
  \fontsize{10}{10.8}\selectfont
\begingroup
\renewcommand{\arraystretch}{0.84}
\begin{tabularx}{0.98\linewidth}{@{}p{2.55cm}>{\raggedright\arraybackslash}X@{}}
\toprule
输入与状态 & 格点/解析输入 → 非微扰重正化、OPE 与匹配中的算符、矩阵元或场配置；具体输入见源章节。 \\
章节路线 & 引言；格点研究概述；涂抹的格点威尔逊圈；格点大（公式）相变在连续统极限中存活；涂抹消除紫外发散；格点研究的数值细节；（公式）的确定；标度变换下相变两侧的研究；（公式）的一般特征；连续统的观点……普适类的一个猜想；总结与讨论；与其他工作的联系；本文主要的新思想；未来计划 \\
方法摘录 & 本文的大多数具体结果是用格点场论方法得到的。这些结果构成一个更大、相当雄心勃勃且颇为思辨性的纲领的第一部分；该纲领是在连续统中定义的，不依赖于任何正规化方案。本节的主要内容就是勾画这一纲领，其目的在于为后面各节给出的格点结果提供恰当的视角。在整……计算本身将完全是解析的。不言而喻，任何解析步骤都需要对照数值工作加以检验，因为这是一片未经勘测的领域 \\
结构式/示意 & $O^{R}(z,\mu)=Z(z,\mu,a)\,O^{\rm bare}(z,a)$\\[-0.2em]\scriptsize 该式只表达主题变量关系；论文特定推导以源文件为准。 \\
物理校验 & 重正化因子抵消截止依赖；a→0 和短距离 z→0 极限必须与匹配阶数、方案和尺度保持一致。 \\
\bottomrule
\end{tabularx}
\endgroup
  \source{方法证据：\detokenize{refer/papers/连续威尔逊圈的无穷N相变_latex/chapters/section01.tex:1-45}；主题结构式为示意。}
\end{frame}

%% frame_manifest|frame_id=P06-03|paper_id=P06|claim_id=P06-结果|visual_type=result-evidence-table|width_budget=0.98linewidth|height_budget=0.86textheight|min_font_pt=10|split_allowed=no|expected_page=36|risks=title-table-source-footline
\begin{frame}[t]{连续威尔逊圈的无穷N相变：结果、验证与物理含义}
  \fontsize{10}{10.8}\selectfont
\begingroup
\renewcommand{\arraystretch}{0.84}
\begin{tabularx}{0.98\linewidth}{@{}p{2.55cm}>{\raggedright\arraybackslash}X@{}}
\toprule
源文结果 & 我们首先概述各种影响了我们、却尚未被我们直接提及的想法。随后我们进行总结，说明本文所作的实质性新贡献是什么。最后简短描述计划中的未来工作 \\
验证状态 & 确证：源章节有结果/讨论摘录。 \\
库内可见证据 & 中文源约 931 行；公式命中 9，图命中 18，表/表格命中 0。 \\
物理含义 & 在本主题共同链条中，本论文把上述输入推进到可比较的中间量或观测量；具体强度、误差和外推范围不超出源文档。 \\
未验证项 & 当前工作区没有论文 PDF、原始数据和独立复现脚本；因此数值复核、图像再现和统计显著性不作额外断言。 \\
\bottomrule
\end{tabularx}
\endgroup
  \source{结果证据：\detokenize{refer/papers/连续威尔逊圈的无穷N相变_latex/chapters/section05.tex:1-5}；图表/公式计数由本报告生成器对中文源的静态统计得到。}
\end{frame}

%% frame_manifest|frame_id=P06-04|paper_id=P06|claim_id=P06-局限|visual_type=risk-relation-table|width_budget=0.98linewidth|height_budget=0.86textheight|min_font_pt=10|split_allowed=no|expected_page=37|risks=title-table-source-footline
\begin{frame}[t]{连续威尔逊圈的无穷N相变：局限、跨论文接口与可复查入口}
  \fontsize{10}{10.8}\selectfont
\begingroup
\renewcommand{\arraystretch}{0.84}
\begin{tabularx}{0.98\linewidth}{@{}p{2.55cm}>{\raggedright\arraybackslash}X@{}}
\toprule
源文局限 & 需要回到结果/讨论/展望章节核对适用区间；本档案不把缺失的独立数据复核补成已证实结论。 \\
跨论文接口 & 是准/赝分布从格点数据走向 MS-bar 或其他连续方案的必要中间层。 \\
极限检查 & 重正化因子抵消截止依赖；a→0 和短距离 z→0 极限必须与匹配阶数、方案和尺度保持一致。 \\
下一步核验 & 补齐 PDF 或原始数据；按源文参数复现关键图表；比较不同动量、流时间、距离、格距或方案转换后的稳定性。 \\
完整来源 & \detokenize{refer/papers/连续威尔逊圈的无穷N相变_latex/main.tex}；\detokenize{refer/papers/Infinite_N_Phase_Transitions_Wilson_Loops_latex/main.tex}；章节源：\detokenize{refer/papers/连续威尔逊圈的无穷N相变_latex/chapters/*.tex}；英中目录：\detokenize{refer/papers/Infinite_N_Phase_Transitions_Wilson_Loops_latex}. \\
\bottomrule
\end{tabularx}
\endgroup
  \source{状态：\inferred 为主题接口推断；\unverified 为当前材料边界；完整文件清单见 manifest。}
\end{frame}

%% frame_manifest|frame_id=P07-01|paper_id=P07|claim_id=P07-定位|visual_type=paper-card|width_budget=0.98linewidth|height_budget=0.86textheight|min_font_pt=10|split_allowed=no|expected_page=38|risks=title-table-source-footline
\section{威尔逊流的性质与应用}
\begin{frame}[t]{威尔逊流的性质与应用：研究问题与论文定位}
  \fontsize{10}{10.8}\selectfont
\begingroup
\renewcommand{\arraystretch}{0.84}
\begin{tabularx}{0.98\linewidth}{@{}p{2.55cm}>{\raggedright\arraybackslash}X@{}}
\toprule
论文编号 & P07；主题：梯度流、涂抹与能动量张量 \\
书目信息 & Properties and uses of the Wilson flow in lattice QCD CERN-PH-TH/2010-143. Journal publication: JHEP 08 (2010) 071; preprint arXiv:1006.4518 [hep-lat].；arXiv:1006.4518；英/中转排页数 15/13 \\
研究对象 & 对格点 QCD 中威尔逊流的理论与数值研究表明，在流时间（公式）处得到的规范场是一 个光滑的重正化场。因此，该场中局域规范不变表达式的期望值是良定义的物理量，它 们探测的是长度尺度（公式）量级上的理论。此外，通过把 QCD 泛函积……出这是由一种动力学效应引起的：当晶格间距从（公式）fm 减小到更小值时，各扇区被 迅速分离开来 \\
章节证据 & 引言；小耦合下威尔逊流的性质；规范固定；修改后流方程的解；（公式）的展开；（公式）的计算；（公式）到（公式）阶的计算；重正化；威尔逊流的格点研究；模拟参数；观测量；（公式）的时间依赖性；晶格……区的动力学分离；拓扑磁化率；结束语；记号约定；一圈积分；威尔逊流的数值积分 \\
证据状态 & \verified：摘要和章节来自现有 LaTeX；\unverified：没有论文 PDF/原始数据独立复核。 \\
\bottomrule
\end{tabularx}
\endgroup
  \source{索引：\detokenize{refer/papers/INDEX.md:7}；正文：\detokenize{refer/papers/威尔逊流的性质与应用_latex/main.tex}；\detokenize{refer/papers/Properties_Uses_Wilson_Flow_latex/main.tex}}
\end{frame}

%% frame_manifest|frame_id=P07-02|paper_id=P07|claim_id=P07-方法|visual_type=method-flow-formula|width_budget=0.98linewidth|height_budget=0.86textheight|min_font_pt=10|split_allowed=no|expected_page=39|risks=title-table-source-footline
\begin{frame}[t]{威尔逊流的性质与应用：理论结构与方法链}
  \fontsize{10}{10.8}\selectfont
\begingroup
\renewcommand{\arraystretch}{0.84}
\begin{tabularx}{0.98\linewidth}{@{}p{2.55cm}>{\raggedright\arraybackslash}X@{}}
\toprule
输入与状态 & 格点/解析输入 → 梯度流、涂抹与能动量张量中的算符、矩阵元或场配置；具体输入见源章节。 \\
章节路线 & 引言；小耦合下威尔逊流的性质；规范固定；修改后流方程的解；（公式）的展开；（公式）的计算；（公式）到（公式）阶的计算；重正化；威尔逊流的格点研究；模拟参数；观测量；（公式）的时间依赖性；晶格……区的动力学分离；拓扑磁化率；结束语；记号约定；一圈积分；威尔逊流的数值积分 \\
方法摘录 & 在三个格点上，采用 Wilson 规范作用量和几种知名链更新算法的组合，生成了代表性规范 场系综。表（交叉引用）所引晶格间距值由 Guagnelli 等人（引文）发表的 Sommer 参考标度（公式）fm（引文）的格点单位结果导出。在所选 耦合…… 迅速（引文）。特别地，早在（公式）fm 处， 所有其他常用感兴趣的量在模拟时间上的关联就已远远更弱 \\
结构式/示意 & $\partial_t B_\mu(t,x)=D_\nu G_{\nu\mu}(t,x),\qquad t\sim L_{\rm sm}^{2}$\\[-0.2em]\scriptsize 该式只表达主题变量关系；论文特定推导以源文件为准。 \\
物理校验 & 流时间与平滑长度满足平方关系；t→0 回到原始场，有限 t 则抑制紫外模式。 \\
\bottomrule
\end{tabularx}
\endgroup
  \source{方法证据：\detokenize{refer/papers/威尔逊流的性质与应用_latex/chapters/section03.tex:26-41}；主题结构式为示意。}
\end{frame}

%% frame_manifest|frame_id=P07-03|paper_id=P07|claim_id=P07-结果|visual_type=result-evidence-table|width_budget=0.98linewidth|height_budget=0.86textheight|min_font_pt=10|split_allowed=no|expected_page=40|risks=title-table-source-footline
\begin{frame}[t]{威尔逊流的性质与应用：结果、验证与物理含义}
  \fontsize{10}{10.8}\selectfont
\begingroup
\renewcommand{\arraystretch}{0.84}
\begin{tabularx}{0.98\linewidth}{@{}p{2.55cm}>{\raggedright\arraybackslash}X@{}}
\toprule
源文结果 & 本文报告的结果为非阿贝尔规范理论的本质以及格点 QCD 的连续极限提供了若干新的认识。 它们显然在多个方面仍不完整，并引出一些有趣的问题。特别是，要对其重正化性质获得结 构性的理解，很可能需要对威尔逊流作全阶的微扰分析。 第 2 节中的一圈计算提示，无论夸克味数多少，在正流时间得到的场都是重正化场。迄今 这个问题只在纯规范理论中被数值地研究过；但令人鼓舞的是可以指出：作为规范 ……逆的，这个积分仍然编码了理论所描述的全部物理。不过， 理论的某些性质（例如拓扑扇区的分割）在这种表述下变得特别透明，而其在高能下的行为 则用基本规范场来讨论更为合适 \\
验证状态 & 确证：源章节有结果/讨论摘录。 \\
库内可见证据 & 中文源约 951 行；公式命中 49，图命中 8，表/表格命中 2。 \\
物理含义 & 在本主题共同链条中，本论文把上述输入推进到可比较的中间量或观测量；具体强度、误差和外推范围不超出源文档。 \\
未验证项 & 当前工作区没有论文 PDF、原始数据和独立复现脚本；因此数值复核、图像再现和统计显著性不作额外断言。 \\
\bottomrule
\end{tabularx}
\endgroup
  \source{结果证据：\detokenize{refer/papers/威尔逊流的性质与应用_latex/chapters/section05.tex:1-15}；图表/公式计数由本报告生成器对中文源的静态统计得到。}
\end{frame}

%% frame_manifest|frame_id=P07-04|paper_id=P07|claim_id=P07-局限|visual_type=risk-relation-table|width_budget=0.98linewidth|height_budget=0.86textheight|min_font_pt=10|split_allowed=no|expected_page=41|risks=title-table-source-footline
\begin{frame}[t]{威尔逊流的性质与应用：局限、跨论文接口与可复查入口}
  \fontsize{10}{10.8}\selectfont
\begingroup
\renewcommand{\arraystretch}{0.84}
\begin{tabularx}{0.98\linewidth}{@{}p{2.55cm}>{\raggedright\arraybackslash}X@{}}
\toprule
源文局限 & 需要回到结果/讨论/展望章节核对适用区间；本档案不把缺失的独立数据复核补成已证实结论。 \\
跨论文接口 & 把连续极限、微扰有限性和能动量张量重构连接到可计算的格点观测量。 \\
极限检查 & 流时间与平滑长度满足平方关系；t→0 回到原始场，有限 t 则抑制紫外模式。 \\
下一步核验 & 补齐 PDF 或原始数据；按源文参数复现关键图表；比较不同动量、流时间、距离、格距或方案转换后的稳定性。 \\
完整来源 & \detokenize{refer/papers/威尔逊流的性质与应用_latex/main.tex}；\detokenize{refer/papers/Properties_Uses_Wilson_Flow_latex/main.tex}；章节源：\detokenize{refer/papers/威尔逊流的性质与应用_latex/chapters/*.tex}；英中目录：\detokenize{refer/papers/Properties_Uses_Wilson_Flow_latex}. \\
\bottomrule
\end{tabularx}
\endgroup
  \source{状态：\inferred 为主题接口推断；\unverified 为当前材料边界；完整文件清单见 manifest。}
\end{frame}

%% frame_manifest|frame_id=P08-01|paper_id=P08|claim_id=P08-定位|visual_type=paper-card|width_budget=0.98linewidth|height_budget=0.86textheight|min_font_pt=10|split_allowed=no|expected_page=42|risks=title-table-source-footline
\section{梯度流的微扰分析}
\begin{frame}[t]{梯度流的微扰分析：研究问题与论文定位}
  \fontsize{10}{10.8}\selectfont
\begingroup
\renewcommand{\arraystretch}{0.84}
\begin{tabularx}{0.98\linewidth}{@{}p{2.55cm}>{\raggedright\arraybackslash}X@{}}
\toprule
论文编号 & P08；主题：梯度流、涂抹与能动量张量 \\
书目信息 & Perturbative Analysis Gradient Flow；arXiv:1101.0963；英/中转排页数 19/19 \\
研究对象 & 非阿贝尔规范理论中的梯度流由一个定域扩散方程定义，它以规范协变的方式 使规范场随流时间演化。与 Langevin 方程的情形类似， 随时间变化的场的关联函数可以按微扰论展开， 其费曼规则就是（公式）上一个可重正化场论的费曼规则。 ……式重正化， 它们便不需要附加的重正化。 因此，该流把规范场映射为一族带单个参数的光滑的重正化场 \\
章节证据 & 引言；流方程的迭代解；梯度流的定义；按基本规范场幂次的展开；流线图；微扰论；规范固定；规范场传播子；（公式）维中的费曼规则；（公式）维中的场论；作用量；传播子；顶点；流线与流线圈；单圈阶示例……质场的理论；格点正规化；结论；记号约定；流顶点；图（交叉引用）中各图的计算 \\
证据状态 & \verified：摘要和章节来自现有 LaTeX；\unverified：没有论文 PDF/原始数据独立复核。 \\
\bottomrule
\end{tabularx}
\endgroup
  \source{索引：\detokenize{refer/papers/INDEX.md:8}；正文：\detokenize{refer/papers/梯度流的微扰分析_latex/main.tex}；\detokenize{refer/papers/Perturbative_Analysis_Gradient_Flow_latex/main.tex}}
\end{frame}

%% frame_manifest|frame_id=P08-02|paper_id=P08|claim_id=P08-方法|visual_type=method-flow-formula|width_budget=0.98linewidth|height_budget=0.86textheight|min_font_pt=10|split_allowed=no|expected_page=43|risks=title-table-source-footline
\begin{frame}[t]{梯度流的微扰分析：理论结构与方法链}
  \fontsize{10}{10.8}\selectfont
\begingroup
\renewcommand{\arraystretch}{0.84}
\begin{tabularx}{0.98\linewidth}{@{}p{2.55cm}>{\raggedright\arraybackslash}X@{}}
\toprule
输入与状态 & 格点/解析输入 → 梯度流、涂抹与能动量张量中的算符、矩阵元或场配置；具体输入见源章节。 \\
章节路线 & 引言；流方程的迭代解；梯度流的定义；按基本规范场幂次的展开；流线图；微扰论；规范固定；规范场传播子；（公式）维中的费曼规则；（公式）维中的场论；作用量；传播子；顶点；流线与流线圈；单圈阶示例……质场的理论；格点正规化；结论；记号约定；流顶点；图（交叉引用）中各图的计算 \\
方法摘录 & ssec:defflow 梯度流使规范场作为一个参数（公式）（称为流时间）的函数而演化。 从基本规范场出发，（公式或图表位置）随时间变化的场（公式）由微分方程（公式或图表位置）确定。``梯度流''这一名称源于如下事实：式（交叉引用）右端第一项 …… 因为在不同（公式）值下得到的方程（交叉引用）的解 彼此通过一个（依赖时间的）规范变换相联系（引文） \\
结构式/示意 & $\partial_t B_\mu(t,x)=D_\nu G_{\nu\mu}(t,x),\qquad t\sim L_{\rm sm}^{2}$\\[-0.2em]\scriptsize 该式只表达主题变量关系；论文特定推导以源文件为准。 \\
物理校验 & 流时间与平滑长度满足平方关系；t→0 回到原始场，有限 t 则抑制紫外模式。 \\
\bottomrule
\end{tabularx}
\endgroup
  \source{方法证据：\detokenize{refer/papers/梯度流的微扰分析_latex/chapters/section02.tex:14-42}；主题结构式为示意。}
\end{frame}

%% frame_manifest|frame_id=P08-03|paper_id=P08|claim_id=P08-结果|visual_type=result-evidence-table|width_budget=0.98linewidth|height_budget=0.86textheight|min_font_pt=10|split_allowed=no|expected_page=44|risks=title-table-source-footline
\begin{frame}[t]{梯度流的微扰分析：结果、验证与物理含义}
  \fontsize{10}{10.8}\selectfont
\begingroup
\renewcommand{\arraystretch}{0.84}
\begin{tabularx}{0.98\linewidth}{@{}p{2.55cm}>{\raggedright\arraybackslash}X@{}}
\toprule
源文结果 & sec:conclusions 梯度流生成的规范场不需要重正化这一事实， 是流方程的定域性、对称性与抛物型性质共同导致的结果。 特别是在与之相联系的（公式）维场论中， 体抵消项之所以被排除，仅仅是因为体场的演化在大时间处是推迟的， 从而由树图描述。 然而，除流时间为零处理论重正化所需的那些抵消项之外 其他边界抵消项的不存在却是非平凡的， 只能借助幂次计数与 BRS 对称性来证明。 存在物质场时，只要只有规范场随时间演化，情形便基本不变。 因此场在正流时间处关联函数的有限性仍然得到保证。 不过，把全部或部分物质场纳入流中是一个有待探索的有趣选择 \\
验证状态 & 确证：源章节有结果/讨论摘录。 \\
库内可见证据 & 中文源约 1304 行；公式命中 79，图命中 18，表/表格命中 0。 \\
物理含义 & 在本主题共同链条中，本论文把上述输入推进到可比较的中间量或观测量；具体强度、误差和外推范围不超出源文档。 \\
未验证项 & 当前工作区没有论文 PDF、原始数据和独立复现脚本；因此数值复核、图像再现和统计显著性不作额外断言。 \\
\bottomrule
\end{tabularx}
\endgroup
  \source{结果证据：\detokenize{refer/papers/梯度流的微扰分析_latex/chapters/section09.tex:1-15}；图表/公式计数由本报告生成器对中文源的静态统计得到。}
\end{frame}

%% frame_manifest|frame_id=P08-04|paper_id=P08|claim_id=P08-局限|visual_type=risk-relation-table|width_budget=0.98linewidth|height_budget=0.86textheight|min_font_pt=10|split_allowed=no|expected_page=45|risks=title-table-source-footline
\begin{frame}[t]{梯度流的微扰分析：局限、跨论文接口与可复查入口}
  \fontsize{10}{10.8}\selectfont
\begingroup
\renewcommand{\arraystretch}{0.84}
\begin{tabularx}{0.98\linewidth}{@{}p{2.55cm}>{\raggedright\arraybackslash}X@{}}
\toprule
源文局限 & 需要回到结果/讨论/展望章节核对适用区间；本档案不把缺失的独立数据复核补成已证实结论。 \\
跨论文接口 & 把连续极限、微扰有限性和能动量张量重构连接到可计算的格点观测量。 \\
极限检查 & 流时间与平滑长度满足平方关系；t→0 回到原始场，有限 t 则抑制紫外模式。 \\
下一步核验 & 补齐 PDF 或原始数据；按源文参数复现关键图表；比较不同动量、流时间、距离、格距或方案转换后的稳定性。 \\
完整来源 & \detokenize{refer/papers/梯度流的微扰分析_latex/main.tex}；\detokenize{refer/papers/Perturbative_Analysis_Gradient_Flow_latex/main.tex}；章节源：\detokenize{refer/papers/梯度流的微扰分析_latex/chapters/*.tex}；英中目录：\detokenize{refer/papers/Perturbative_Analysis_Gradient_Flow_latex}. \\
\bottomrule
\end{tabularx}
\endgroup
  \source{状态：\inferred 为主题接口推断；\unverified 为当前材料边界；完整文件清单见 manifest。}
\end{frame}

%% frame_manifest|frame_id=P09-01|paper_id=P09|claim_id=P09-定位|visual_type=paper-card|width_budget=0.98linewidth|height_budget=0.86textheight|min_font_pt=10|split_allowed=no|expected_page=46|risks=title-table-source-footline
\section{手征对称性与杨米尔斯梯度流}
\begin{frame}[t]{手征对称性与杨米尔斯梯度流：研究问题与论文定位}
  \fontsize{10}{10.8}\selectfont
\begingroup
\renewcommand{\arraystretch}{0.84}
\begin{tabularx}{0.98\linewidth}{@{}p{2.55cm}>{\raggedright\arraybackslash}X@{}}
\toprule
论文编号 & P09；主题：梯度流、涂抹与能动量张量 \\
书目信息 & Chiral symmetry and the Yang–Mills gradient flow；arXiv:1302.5246；英/中转排页数 34/34 \\
研究对象 & 近几年来，杨–米尔斯梯度流已被证明是非阿贝尔规范理论非微扰研究的一种有吸引力的工具。 本文考虑将该流向 QCD 夸克场的一个简单推广。 与纯规范梯度流的情形一样， 涉及正流时间局域场的关联函数的可重正化性， 可以借助一个（公式）维……括： 非微扰重正化与 O(（公式）) 改善， 以及在手征极限下手征凝聚和赝标衰变常数的精确计算 \\
章节证据 & 引言；QCD 中的梯度流；流方程；局域复合场；微扰论；非零流时间下的夸克凝聚；（公式）+1 维场论；场与作用量；关联函数；费米子积分；重正化；手征对称性关系；夸克场方程与 PCAC 关系；重…… (7.10) 的积分；伴随龙格–库塔积分；稳定性问题；改进方案；分级方案 \\
证据状态 & \verified：摘要和章节来自现有 LaTeX；\unverified：没有论文 PDF/原始数据独立复核。 \\
\bottomrule
\end{tabularx}
\endgroup
  \source{索引：\detokenize{refer/papers/INDEX.md:9}；正文：\detokenize{refer/papers/手征对称性与杨米尔斯梯度流_latex/main.tex}；\detokenize{refer/papers/Chiral_Symmetry_YM_Gradient_Flow_latex/main.tex}}
\end{frame}

%% frame_manifest|frame_id=P09-02|paper_id=P09|claim_id=P09-方法|visual_type=method-flow-formula|width_budget=0.98linewidth|height_budget=0.86textheight|min_font_pt=10|split_allowed=no|expected_page=47|risks=title-table-source-footline
\begin{frame}[t]{手征对称性与杨米尔斯梯度流：理论结构与方法链}
  \fontsize{10}{10.8}\selectfont
\begingroup
\renewcommand{\arraystretch}{0.84}
\begin{tabularx}{0.98\linewidth}{@{}p{2.55cm}>{\raggedright\arraybackslash}X@{}}
\toprule
输入与状态 & 格点/解析输入 → 梯度流、涂抹与能动量张量中的算符、矩阵元或场配置；具体输入见源章节。 \\
章节路线 & 引言；QCD 中的梯度流；流方程；局域复合场；微扰论；非零流时间下的夸克凝聚；（公式）+1 维场论；场与作用量；关联函数；费米子积分；重正化；手征对称性关系；夸克场方程与 PCAC 关系；重…… (7.10) 的积分；伴随龙格–库塔积分；稳定性问题；改进方案；分级方案 \\
方法摘录 & 现在假定夸克多重态包含两个轻夸克， 称为（公式）、（公式）夸克，二者具有相同的 重正化质量（公式）。 若再无太多其余夸克，（公式）道的手征对称性预期会发生自发破缺， 于是存在一个赝标介子三重态——``（公式）介子''， 其质量随（公式）而趋于零……）消失了， 因此例如衰变常数（公式）可以直接由裸的随时间演化凝聚 与赝标真空到（公式）介子矩阵元给出 \\
结构式/示意 & $\partial_t B_\mu(t,x)=D_\nu G_{\nu\mu}(t,x),\qquad t\sim L_{\rm sm}^{2}$\\[-0.2em]\scriptsize 该式只表达主题变量关系；论文特定推导以源文件为准。 \\
物理校验 & 流时间与平滑长度满足平方关系；t→0 回到原始场，有限 t 则抑制紫外模式。 \\
\bottomrule
\end{tabularx}
\endgroup
  \source{方法证据：\detokenize{refer/papers/手征对称性与杨米尔斯梯度流_latex/chapters/zh04.tex:206-300}；主题结构式为示意。}
\end{frame}

%% frame_manifest|frame_id=P09-03|paper_id=P09|claim_id=P09-结果|visual_type=result-evidence-table|width_budget=0.98linewidth|height_budget=0.86textheight|min_font_pt=10|split_allowed=no|expected_page=48|risks=title-table-source-footline
\begin{frame}[t]{手征对称性与杨米尔斯梯度流：结果、验证与物理含义}
  \fontsize{10}{10.8}\selectfont
\begingroup
\renewcommand{\arraystretch}{0.84}
\begin{tabularx}{0.98\linewidth}{@{}p{2.55cm}>{\raggedright\arraybackslash}X@{}}
\toprule
源文结果 & 通过向夸克场的推广， 梯度流成为研究 QCD 中手征对称性自发破缺 以及轻赝标介子物理的有力工具。 本文的重点是理论框架及其在广泛使用的格点 QCD 表述之一中的实现。 鉴于流的重正化性质， 可以预期：只要格点理论的局域性与规范不变性得到保证， 格点正规化的选择在连续极限下无关紧要。 显然，这一理论分析同样适用于 具有其他规范群与费米子多重态的场论。 在迄今具体完成的例子中， ……尚未被考虑， 而随时间演化的凝聚很可能是非零温度 QCD 中 一个易于获取的手征序参量。 此外，把流应用于共形窗附近的类 QCD 理论， 或许可以带来有趣的定性洞见 \\
验证状态 & 确证：源章节有结果/讨论摘录。 \\
库内可见证据 & 中文源约 3170 行；公式命中 0，图命中 1，表/表格命中 2。 \\
物理含义 & 在本主题共同链条中，本论文把上述输入推进到可比较的中间量或观测量；具体强度、误差和外推范围不超出源文档。 \\
未验证项 & 当前工作区没有论文 PDF、原始数据和独立复现脚本；因此数值复核、图像再现和统计显著性不作额外断言。 \\
\bottomrule
\end{tabularx}
\endgroup
  \source{结果证据：\detokenize{refer/papers/手征对称性与杨米尔斯梯度流_latex/chapters/zh10.tex:2-24}；图表/公式计数由本报告生成器对中文源的静态统计得到。}
\end{frame}

%% frame_manifest|frame_id=P09-04|paper_id=P09|claim_id=P09-局限|visual_type=risk-relation-table|width_budget=0.98linewidth|height_budget=0.86textheight|min_font_pt=10|split_allowed=no|expected_page=49|risks=title-table-source-footline
\begin{frame}[t]{手征对称性与杨米尔斯梯度流：局限、跨论文接口与可复查入口}
  \fontsize{10}{10.8}\selectfont
\begingroup
\renewcommand{\arraystretch}{0.84}
\begin{tabularx}{0.98\linewidth}{@{}p{2.55cm}>{\raggedright\arraybackslash}X@{}}
\toprule
源文局限 & 需要回到结果/讨论/展望章节核对适用区间；本档案不把缺失的独立数据复核补成已证实结论。 \\
跨论文接口 & 把连续极限、微扰有限性和能动量张量重构连接到可计算的格点观测量。 \\
极限检查 & 流时间与平滑长度满足平方关系；t→0 回到原始场，有限 t 则抑制紫外模式。 \\
下一步核验 & 补齐 PDF 或原始数据；按源文参数复现关键图表；比较不同动量、流时间、距离、格距或方案转换后的稳定性。 \\
完整来源 & \detokenize{refer/papers/手征对称性与杨米尔斯梯度流_latex/main.tex}；\detokenize{refer/papers/Chiral_Symmetry_YM_Gradient_Flow_latex/main.tex}；章节源：\detokenize{refer/papers/手征对称性与杨米尔斯梯度流_latex/chapters/*.tex}；英中目录：\detokenize{refer/papers/Chiral_Symmetry_YM_Gradient_Flow_latex}. \\
\bottomrule
\end{tabularx}
\endgroup
  \source{状态：\inferred 为主题接口推断；\unverified 为当前材料边界；完整文件清单见 manifest。}
\end{frame}

%% frame_manifest|frame_id=P10-01|paper_id=P10|claim_id=P10-定位|visual_type=paper-card|width_budget=0.98linewidth|height_budget=0.86textheight|min_font_pt=10|split_allowed=no|expected_page=50|risks=title-table-source-footline
\section{杨米尔斯梯度流提取能动量张量}
\begin{frame}[t]{杨米尔斯梯度流提取能动量张量：研究问题与论文定位}
  \fontsize{10}{10.8}\selectfont
\begingroup
\renewcommand{\arraystretch}{0.84}
\begin{tabularx}{0.98\linewidth}{@{}p{2.55cm}>{\raggedright\arraybackslash}X@{}}
\toprule
论文编号 & P10；主题：梯度流、涂抹与能动量张量 \\
书目信息 & Energy–momentum tensor from the Yang–Mills gradient flow Prog. Theor. Exp. Phys. 2013, 083B03; arXiv:1304.0533 [hep-lat]. Report number: RIKEN-QHP-82.；arXiv:1304.0533；英/中转排页数 17/16 \\
研究对象 & 对于正的流时间，由杨–米尔斯梯度流产生的规范场的乘积不表现出重合点奇异性， 因而这样的局域乘积与正规化方案无关。而且，这样的局域乘积可以用零流时间处 的重正化局域算符展开，展开系数有限并受重正化群方程支配。利用这些事实，我 们导出……供了一种可能的方 法，通过格点正规化与蒙特卡罗模拟来计算一个良定义的能量–动量张量的关联函 数 \\
章节证据 & 引言；杨–米尔斯理论与能量–动量张量；维数正规化下的能量–动量张量；迹反常的含义；杨–米尔斯梯度流与小流时间展开；重正化群方程与渐近公式；系数的重正化群方程；最低阶近似与渐近公式；结论；系数函数的单圈计算 \\
证据状态 & \verified：摘要和章节来自现有 LaTeX；\unverified：没有论文 PDF/原始数据独立复核。 \\
\bottomrule
\end{tabularx}
\endgroup
  \source{索引：\detokenize{refer/papers/INDEX.md:10}；正文：\detokenize{refer/papers/杨米尔斯梯度流提取能动量张量_latex/main.tex}；\detokenize{refer/papers/EMT_From_YM_Gradient_Flow_latex/main.tex}}
\end{frame}

%% frame_manifest|frame_id=P10-02|paper_id=P10|claim_id=P10-方法|visual_type=method-flow-formula|width_budget=0.98linewidth|height_budget=0.86textheight|min_font_pt=10|split_allowed=no|expected_page=51|risks=title-table-source-footline
\begin{frame}[t]{杨米尔斯梯度流提取能动量张量：理论结构与方法链}
  \fontsize{10}{10.8}\selectfont
\begingroup
\renewcommand{\arraystretch}{0.84}
\begin{tabularx}{0.98\linewidth}{@{}p{2.55cm}>{\raggedright\arraybackslash}X@{}}
\toprule
输入与状态 & 格点/解析输入 → 梯度流、涂抹与能动量张量中的算符、矩阵元或场配置；具体输入见源章节。 \\
章节路线 & 引言；杨–米尔斯理论与能量–动量张量；维数正规化下的能量–动量张量；迹反常的含义；杨–米尔斯梯度流与小流时间展开；重正化群方程与渐近公式；系数的重正化群方程；最低阶近似与渐近公式；结论；系数函数的单圈计算 \\
方法摘录 & sec:1 尽管格点正规化为场论提供了非常强有力的非微扰表述，但遗憾的是，它与基本的 整体对称性常常不相容。最著名的例子是手征对称性（引文）；超对称性是另一个臭名昭著的例子（引文），不用 说，平移不变性也是如此。当一种正规化在某个对称性下不是不……模拟来计算正确归一化的守恒能量–动量张 量的关联函数。 除非另有说明，本文沿用文献（引文）的记号约定 \\
结构式/示意 & $\partial_t B_\mu(t,x)=D_\nu G_{\nu\mu}(t,x),\qquad t\sim L_{\rm sm}^{2}$\\[-0.2em]\scriptsize 该式只表达主题变量关系；论文特定推导以源文件为准。 \\
物理校验 & 流时间与平滑长度满足平方关系；t→0 回到原始场，有限 t 则抑制紫外模式。 \\
\bottomrule
\end{tabularx}
\endgroup
  \source{方法证据：\detokenize{refer/papers/杨米尔斯梯度流提取能动量张量_latex/chapters/section01.tex:1-79}；主题结构式为示意。}
\end{frame}

%% frame_manifest|frame_id=P10-03|paper_id=P10|claim_id=P10-结果|visual_type=result-evidence-table|width_budget=0.98linewidth|height_budget=0.86textheight|min_font_pt=10|split_allowed=no|expected_page=52|risks=title-table-source-footline
\begin{frame}[t]{杨米尔斯梯度流提取能动量张量：结果、验证与物理含义}
  \fontsize{10}{10.8}\selectfont
\begingroup
\renewcommand{\arraystretch}{0.84}
\begin{tabularx}{0.98\linewidth}{@{}p{2.55cm}>{\raggedright\arraybackslash}X@{}}
\toprule
源文结果 & 本文导出了一个公式，它把杨–米尔斯梯度流产生的某些规范不变局域乘积的短流 时间行为，与杨–米尔斯理论中正确归一化的守恒能量–动量张量联系起来。我们 的主要结果是式（交叉引用）。式（交叉引用）的右端可以在格点规 范理论中用 Wilson 流、配合算符（交叉引用）与（交叉引用）的适当 离散化来计算（例如见文献（引文））。这里必须 首先取连续极限（公式），然后才取（公式）极限；否则我……中的 算符通过梯度流联系起来——并不限于能量–动量张量。例如，或许可以从局域乘 积的小流时间极限出发，在格点上构造理想的手征流或理想超流。追求这一想法将 是有意义的 \\
验证状态 & 确证：源章节有结果/讨论摘录。 \\
库内可见证据 & 中文源约 1482 行；公式命中 80，图命中 24，表/表格命中 2。 \\
物理含义 & 在本主题共同链条中，本论文把上述输入推进到可比较的中间量或观测量；具体强度、误差和外推范围不超出源文档。 \\
未验证项 & 当前工作区没有论文 PDF、原始数据和独立复现脚本；因此数值复核、图像再现和统计显著性不作额外断言。 \\
\bottomrule
\end{tabularx}
\endgroup
  \source{结果证据：\detokenize{refer/papers/杨米尔斯梯度流提取能动量张量_latex/chapters/section05.tex:1-27}；图表/公式计数由本报告生成器对中文源的静态统计得到。}
\end{frame}

%% frame_manifest|frame_id=P10-04|paper_id=P10|claim_id=P10-局限|visual_type=risk-relation-table|width_budget=0.98linewidth|height_budget=0.86textheight|min_font_pt=10|split_allowed=no|expected_page=53|risks=title-table-source-footline
\begin{frame}[t]{杨米尔斯梯度流提取能动量张量：局限、跨论文接口与可复查入口}
  \fontsize{10}{10.8}\selectfont
\begingroup
\renewcommand{\arraystretch}{0.84}
\begin{tabularx}{0.98\linewidth}{@{}p{2.55cm}>{\raggedright\arraybackslash}X@{}}
\toprule
源文局限 & 需要回到结果/讨论/展望章节核对适用区间；本档案不把缺失的独立数据复核补成已证实结论。 \\
跨论文接口 & 把连续极限、微扰有限性和能动量张量重构连接到可计算的格点观测量。 \\
极限检查 & 流时间与平滑长度满足平方关系；t→0 回到原始场，有限 t 则抑制紫外模式。 \\
下一步核验 & 补齐 PDF 或原始数据；按源文参数复现关键图表；比较不同动量、流时间、距离、格距或方案转换后的稳定性。 \\
完整来源 & \detokenize{refer/papers/杨米尔斯梯度流提取能动量张量_latex/main.tex}；\detokenize{refer/papers/EMT_From_YM_Gradient_Flow_latex/main.tex}；章节源：\detokenize{refer/papers/杨米尔斯梯度流提取能动量张量_latex/chapters/*.tex}；英中目录：\detokenize{refer/papers/EMT_From_YM_Gradient_Flow_latex}. \\
\bottomrule
\end{tabularx}
\endgroup
  \source{状态：\inferred 为主题接口推断；\unverified 为当前材料边界；完整文件清单见 manifest。}
\end{frame}

%% frame_manifest|frame_id=P11-01|paper_id=P11|claim_id=P11-定位|visual_type=paper-card|width_budget=0.98linewidth|height_budget=0.86textheight|min_font_pt=10|split_allowed=no|expected_page=54|risks=title-table-source-footline
\section{欧氏格点上的部分子物理}
\begin{frame}[t]{欧氏格点上的部分子物理：研究问题与论文定位}
  \fontsize{10}{10.8}\selectfont
\begingroup
\renewcommand{\arraystretch}{0.84}
\begin{tabularx}{0.98\linewidth}{@{}p{2.55cm}>{\raggedright\arraybackslash}X@{}}
\toprule
论文编号 & P11；主题：LaMET、准分布与赝分布 \\
书目信息 & Parton Physics on Euclidean Lattice Published as Phys. Rev. Lett. 110, 262002 (2013). DOI: 10.1103/PhysRevLett.110.262002; arXiv:1305.1539 [hep-ph].；arXiv:1305.1539；英/中转排页数 7/6 \\
研究对象 & 我证明，与光锥上夸克、胶子关联相联系的部分子物理，可以在大动量极限下通过依赖参考系的 等时关联函数的矩阵元素加以研究。这一观察使在欧氏格点上实际计算部分子性质成为可能。 作为示例，我演示了如何通过把一个等时关联函数助推到大动量来恢复领头扭度的夸克分布 \\
章节证据 & 未验证：源文件未提供可用的散文摘录。 \\
证据状态 & \verified：摘要和章节来自现有 LaTeX；\unverified：没有论文 PDF/原始数据独立复核。 \\
\bottomrule
\end{tabularx}
\endgroup
  \source{索引：\detokenize{refer/papers/INDEX.md:11}；正文：\detokenize{refer/papers/欧氏格点上的部分子物理_latex/main.tex}；\detokenize{refer/papers/Parton_Physics_on_Euclidean_Lattice_latex/main.tex}}
\end{frame}

%% frame_manifest|frame_id=P11-02|paper_id=P11|claim_id=P11-方法|visual_type=method-flow-formula|width_budget=0.98linewidth|height_budget=0.86textheight|min_font_pt=10|split_allowed=no|expected_page=55|risks=title-table-source-footline
\begin{frame}[t]{欧氏格点上的部分子物理：理论结构与方法链}
  \fontsize{10}{10.8}\selectfont
\begingroup
\renewcommand{\arraystretch}{0.84}
\begin{tabularx}{0.98\linewidth}{@{}p{2.55cm}>{\raggedright\arraybackslash}X@{}}
\toprule
输入与状态 & 格点/解析输入 → LaMET、准分布与赝分布中的算符、矩阵元或场配置；具体输入见源章节。 \\
章节路线 & 未验证：源文件未提供可用的散文摘录。 \\
方法摘录 & 自 20 世纪 60 年代末斯坦福直线加速器中心的深度非弹性散射实验以来，质子的结构已在 各种硬散射过程中得到探测（引文）。其结果是夸克与胶子（部分子）的动量分布 和关联，以及分布振幅等。部分子物理本质上涉及光锥关联：所有夸克与胶子场沿光锥分开……诚然，在格点上研究大动量的强子在计算上仍具挑战性，但随着计算能力的持续提升，这至少是 可以实现的目标 \\
结构式/示意 & $\widetilde q(x,P_z,\mu)=\int dy\,C(x,y,\mu/P_z)q(y,\mu)+\mathcal O(\Lambda_{\rm QCD}^{2}/P_z^{2})$\\[-0.2em]\scriptsize 该式只表达主题变量关系；论文特定推导以源文件为准。 \\
物理校验 & 大动量极限 P\_z→∞ 压低幂修正；有限动量、有限间距和核子激发态是必须单独控制的误差源。 \\
\bottomrule
\end{tabularx}
\endgroup
  \source{方法证据：\detokenize{refer/papers/欧氏格点上的部分子物理_latex/chapters/section01.tex:1-211}；主题结构式为示意。}
\end{frame}

%% frame_manifest|frame_id=P11-03|paper_id=P11|claim_id=P11-结果|visual_type=result-evidence-table|width_budget=0.98linewidth|height_budget=0.86textheight|min_font_pt=10|split_allowed=no|expected_page=56|risks=title-table-source-footline
\begin{frame}[t]{欧氏格点上的部分子物理：结果、验证与物理含义}
  \fontsize{10}{10.8}\selectfont
\begingroup
\renewcommand{\arraystretch}{0.84}
\begin{tabularx}{0.98\linewidth}{@{}p{2.55cm}>{\raggedright\arraybackslash}X@{}}
\toprule
源文结果 & ……（引文），它是 GPD 与 TMD 两者的生成函数。其定义来自夹在非前向核子态之间的 TMD 算符。于是可以在大的（公式）极限下计算它们，例如（公式或图表位置）其中（公式）；同样，为保证有限格点上的规范不变性还需要一条横向 规范链。我们还注意到，在大的（公式）极限下，（公式）依赖并不消失。 光锥振幅，即强子态与 QCD 真空之间、夸克与胶子插值场光锥关联的矩阵元素。我们现在……过去从第一性原理处理非常困难——的格点 QCD 计算尤为有用。 诚然，在格点上研究大动量的强子在计算上仍具挑战性，但随着计算能力的持续提升，这至少是 可以实现的目标 \\
验证状态 & 确证：源章节有结果/讨论摘录。 \\
库内可见证据 & 中文源约 459 行；公式命中 16，图命中 2，表/表格命中 0。 \\
物理含义 & 在本主题共同链条中，本论文把上述输入推进到可比较的中间量或观测量；具体强度、误差和外推范围不超出源文档。 \\
未验证项 & 当前工作区没有论文 PDF、原始数据和独立复现脚本；因此数值复核、图像再现和统计显著性不作额外断言。 \\
\bottomrule
\end{tabularx}
\endgroup
  \source{结果证据：\detokenize{refer/papers/欧氏格点上的部分子物理_latex/chapters/section01.tex:1-211}；图表/公式计数由本报告生成器对中文源的静态统计得到。}
\end{frame}

%% frame_manifest|frame_id=P11-04|paper_id=P11|claim_id=P11-局限|visual_type=risk-relation-table|width_budget=0.98linewidth|height_budget=0.86textheight|min_font_pt=10|split_allowed=no|expected_page=57|risks=title-table-source-footline
\begin{frame}[t]{欧氏格点上的部分子物理：局限、跨论文接口与可复查入口}
  \fontsize{10}{10.8}\selectfont
\begingroup
\renewcommand{\arraystretch}{0.84}
\begin{tabularx}{0.98\linewidth}{@{}p{2.55cm}>{\raggedright\arraybackslash}X@{}}
\toprule
源文局限 & 需要回到结果/讨论/展望章节核对适用区间；本档案不把缺失的独立数据复核补成已证实结论。 \\
跨论文接口 & 承接格点矩阵元，向重正化、匹配、演化和最终 PDF/TMD 现象学输出。 \\
极限检查 & 大动量极限 P\_z→∞ 压低幂修正；有限动量、有限间距和核子激发态是必须单独控制的误差源。 \\
下一步核验 & 补齐 PDF 或原始数据；按源文参数复现关键图表；比较不同动量、流时间、距离、格距或方案转换后的稳定性。 \\
完整来源 & \detokenize{refer/papers/欧氏格点上的部分子物理_latex/main.tex}；\detokenize{refer/papers/Parton_Physics_on_Euclidean_Lattice_latex/main.tex}；章节源：\detokenize{refer/papers/欧氏格点上的部分子物理_latex/chapters/*.tex}；英中目录：\detokenize{refer/papers/Parton_Physics_on_Euclidean_Lattice_latex}. \\
\bottomrule
\end{tabularx}
\endgroup
  \source{状态：\inferred 为主题接口推断；\unverified 为当前材料边界；完整文件清单见 manifest。}
\end{frame}

%% frame_manifest|frame_id=P12-01|paper_id=P12|claim_id=P12-定位|visual_type=paper-card|width_budget=0.98linewidth|height_budget=0.86textheight|min_font_pt=10|split_allowed=no|expected_page=58|risks=title-table-source-footline
\section{大动量有效理论与部分子物理}
\begin{frame}[t]{大动量有效理论与部分子物理：研究问题与论文定位}
  \fontsize{10}{10.8}\selectfont
\begingroup
\renewcommand{\arraystretch}{0.84}
\begin{tabularx}{0.98\linewidth}{@{}p{2.55cm}>{\raggedright\arraybackslash}X@{}}
\toprule
论文编号 & P12；主题：LaMET、准分布与赝分布 \\
书目信息 & Parton Physics from Large-Momentum Effective Field Theory；arXiv:1404.6680；英/中转排页数 9/7 \\
研究对象 & 当把部分子物理表述为光前关联时，尽管光前量子化给出了希望，它仍然难以作非微扰研究。最近，通过重新审视强子以渐近大动量运动的 Feynman 原始图像，人们提出了一条研究部分子的替代途径。本文用格点 QCD 中强子大动量（公式）的有……式）的系统展开提取出来，正如部分子分布可以从几个 GeV 动量标度下的硬散射数据中提取出来一样 \\
章节证据 & 未验证：源文件未提供可用的散文摘录。 \\
证据状态 & \verified：摘要和章节来自现有 LaTeX；\unverified：没有论文 PDF/原始数据独立复核。 \\
\bottomrule
\end{tabularx}
\endgroup
  \source{索引：\detokenize{refer/papers/INDEX.md:12}；正文：\detokenize{refer/papers/大动量有效理论与部分子物理_latex/main.tex}；\detokenize{refer/papers/Parton_Physics_From_LaMET_latex/main.tex}}
\end{frame}

%% frame_manifest|frame_id=P12-02|paper_id=P12|claim_id=P12-方法|visual_type=method-flow-formula|width_budget=0.98linewidth|height_budget=0.86textheight|min_font_pt=10|split_allowed=no|expected_page=59|risks=title-table-source-footline
\begin{frame}[t]{大动量有效理论与部分子物理：理论结构与方法链}
  \fontsize{10}{10.8}\selectfont
\begingroup
\renewcommand{\arraystretch}{0.84}
\begin{tabularx}{0.98\linewidth}{@{}p{2.55cm}>{\raggedright\arraybackslash}X@{}}
\toprule
输入与状态 & 格点/解析输入 → LaMET、准分布与赝分布中的算符、矩阵元或场配置；具体输入见源章节。 \\
章节路线 & 未验证：源文件未提供可用的散文摘录。 \\
方法摘录 & 描述高能强子散射物理（例如在大型强子对撞机 LHC 上产生希格斯玻色子）的最大简化之一，是 R. Feynman 引入的部分子模型（引文）。按照这一图像，一个快速运动的强子（例如质子）可以被看作一束无相互作用的夸克和胶子（部分子），由其动量密度……微扰可算的匹配系数和高阶幂次修正。这使得格点数据在研究 QCD 束缚态性质方面与真实实验数据同样有用 \\
结构式/示意 & $\widetilde q(x,P_z,\mu)=\int dy\,C(x,y,\mu/P_z)q(y,\mu)+\mathcal O(\Lambda_{\rm QCD}^{2}/P_z^{2})$\\[-0.2em]\scriptsize 该式只表达主题变量关系；论文特定推导以源文件为准。 \\
物理校验 & 大动量极限 P\_z→∞ 压低幂修正；有限动量、有限间距和核子激发态是必须单独控制的误差源。 \\
\bottomrule
\end{tabularx}
\endgroup
  \source{方法证据：\detokenize{refer/papers/大动量有效理论与部分子物理_latex/chapters/section01.tex:1-129}；主题结构式为示意。}
\end{frame}

%% frame_manifest|frame_id=P12-03|paper_id=P12|claim_id=P12-结果|visual_type=result-evidence-table|width_budget=0.98linewidth|height_budget=0.86textheight|min_font_pt=10|split_allowed=no|expected_page=60|risks=title-table-source-footline
\begin{frame}[t]{大动量有效理论与部分子物理：结果、验证与物理含义}
  \fontsize{10}{10.8}\selectfont
\begingroup
\renewcommand{\arraystretch}{0.84}
\begin{tabularx}{0.98\linewidth}{@{}p{2.55cm}>{\raggedright\arraybackslash}X@{}}
\toprule
源文结果 & ……中质心运动与内部运动之间便没有干净的分离，除非强子非常重。当考虑以一定速度运动的强子的空间图像时会出现更多问题，此时人们常画一个洛伦兹收缩的圆。强子态中长度收缩的物理意义往往是不清楚的。（公式或图表位置）可以转而考虑平面波强子态中的场关联函数。一般地说，当强子静止（质心动量为零）时，场关联函数沿所有空间方向的关联长度约为 1fm（公式）。超过这个范围，关联应逐渐降到零。现在……为几个 GeV 量级的实际计算中提取光前部分子物理。相关展开包含微扰可算的匹配系数和高阶幂次修正。这使得格点数据在研究 QCD 束缚态性质方面与真实实验数据同样有用 \\
验证状态 & 确证：源章节有结果/讨论摘录。 \\
库内可见证据 & 中文源约 513 行；公式命中 15，图命中 2，表/表格命中 0。 \\
物理含义 & 在本主题共同链条中，本论文把上述输入推进到可比较的中间量或观测量；具体强度、误差和外推范围不超出源文档。 \\
未验证项 & 当前工作区没有论文 PDF、原始数据和独立复现脚本；因此数值复核、图像再现和统计显著性不作额外断言。 \\
\bottomrule
\end{tabularx}
\endgroup
  \source{结果证据：\detokenize{refer/papers/大动量有效理论与部分子物理_latex/chapters/section01.tex:1-129}；图表/公式计数由本报告生成器对中文源的静态统计得到。}
\end{frame}

%% frame_manifest|frame_id=P12-04|paper_id=P12|claim_id=P12-局限|visual_type=risk-relation-table|width_budget=0.98linewidth|height_budget=0.86textheight|min_font_pt=10|split_allowed=no|expected_page=61|risks=title-table-source-footline
\begin{frame}[t]{大动量有效理论与部分子物理：局限、跨论文接口与可复查入口}
  \fontsize{10}{10.8}\selectfont
\begingroup
\renewcommand{\arraystretch}{0.84}
\begin{tabularx}{0.98\linewidth}{@{}p{2.55cm}>{\raggedright\arraybackslash}X@{}}
\toprule
源文局限 & 需要回到结果/讨论/展望章节核对适用区间；本档案不把缺失的独立数据复核补成已证实结论。 \\
跨论文接口 & 承接格点矩阵元，向重正化、匹配、演化和最终 PDF/TMD 现象学输出。 \\
极限检查 & 大动量极限 P\_z→∞ 压低幂修正；有限动量、有限间距和核子激发态是必须单独控制的误差源。 \\
下一步核验 & 补齐 PDF 或原始数据；按源文参数复现关键图表；比较不同动量、流时间、距离、格距或方案转换后的稳定性。 \\
完整来源 & \detokenize{refer/papers/大动量有效理论与部分子物理_latex/main.tex}；\detokenize{refer/papers/Parton_Physics_From_LaMET_latex/main.tex}；章节源：\detokenize{refer/papers/大动量有效理论与部分子物理_latex/chapters/*.tex}；英中目录：\detokenize{refer/papers/Parton_Physics_From_LaMET_latex}. \\
\bottomrule
\end{tabularx}
\endgroup
  \source{状态：\inferred 为主题接口推断；\unverified 为当前材料边界；完整文件清单见 manifest。}
\end{frame}

%% frame_manifest|frame_id=P13-01|paper_id=P13|claim_id=P13-定位|visual_type=paper-card|width_budget=0.98linewidth|height_budget=0.86textheight|min_font_pt=10|split_allowed=no|expected_page=62|risks=title-table-source-footline
\section{准部分子分布单圈匹配}
\begin{frame}[t]{准部分子分布单圈匹配：研究问题与论文定位}
  \fontsize{10}{10.8}\selectfont
\begingroup
\renewcommand{\arraystretch}{0.84}
\begin{tabularx}{0.98\linewidth}{@{}p{2.55cm}>{\raggedright\arraybackslash}X@{}}
\toprule
论文编号 & P13；主题：非微扰重正化、OPE 与匹配 \\
书目信息 & One-Loop Matching for Parton Distributions: Non-Singlet Case Published as Phys. Rev. D 90 , 014051 (2014); D……hysRevD.90.014051; arXiv:1310.7471 [hep-ph].；arXiv:1310.7471；英/中转排页数 9/9 \\
研究对象 & 我们在横向动量截断方案下推导了非单态夸克分布的单圈匹配条件，其中包括非极化、 螺旋度与横率（transversity）分布。该匹配连接的两个对象分别是：由有限核子动量下的 静态关联定义的准分布（quasi-distribution），以及实验中可测量的光锥分布。这一结果 有助于在格点 QCD 计算中从前者提取后者 \\
章节证据 & 引言；非极化准夸克分布的单圈结果；单圈因子化；螺旋度与横率分布；结论 \\
证据状态 & \verified：摘要和章节来自现有 LaTeX；\unverified：没有论文 PDF/原始数据独立复核。 \\
\bottomrule
\end{tabularx}
\endgroup
  \source{索引：\detokenize{refer/papers/INDEX.md:13}；正文：\detokenize{refer/papers/准部分子分布单圈匹配_latex/main.tex}；\detokenize{refer/papers/One_Loop_Matching_Nonsinglet_latex/main.tex}}
\end{frame}

%% frame_manifest|frame_id=P13-02|paper_id=P13|claim_id=P13-方法|visual_type=method-flow-formula|width_budget=0.98linewidth|height_budget=0.86textheight|min_font_pt=10|split_allowed=no|expected_page=63|risks=title-table-source-footline
\begin{frame}[t]{准部分子分布单圈匹配：理论结构与方法链}
  \fontsize{10}{10.8}\selectfont
\begingroup
\renewcommand{\arraystretch}{0.84}
\begin{tabularx}{0.98\linewidth}{@{}p{2.55cm}>{\raggedright\arraybackslash}X@{}}
\toprule
输入与状态 & 格点/解析输入 → 非微扰重正化、OPE 与匹配中的算符、矩阵元或场配置；具体输入见源章节。 \\
章节路线 & 引言；非极化准夸克分布的单圈结果；单圈因子化；螺旋度与横率分布；结论 \\
方法摘录 & 部分子分布在高能散射过程中以部分子的语言刻画核子的结构。它们是给出强子-强子或轻子-强子碰撞实验物理预言的一个基本要素。尽管人们已经投入大量努力从各种实验数据中提取部分子分布（引文），但由于部分子分布的非微扰本质，从强相互作用的基本理论——量子……提取准分布与光锥分布之间的匹配因子；第 4 节给出夸克螺旋度与横率分布的结果；最后在第 5 节作结论 \\
结构式/示意 & $O^{R}(z,\mu)=Z(z,\mu,a)\,O^{\rm bare}(z,a)$\\[-0.2em]\scriptsize 该式只表达主题变量关系；论文特定推导以源文件为准。 \\
物理校验 & 重正化因子抵消截止依赖；a→0 和短距离 z→0 极限必须与匹配阶数、方案和尺度保持一致。 \\
\bottomrule
\end{tabularx}
\endgroup
  \source{方法证据：\detokenize{refer/papers/准部分子分布单圈匹配_latex/chapters/section01.tex:1-19}；主题结构式为示意。}
\end{frame}

%% frame_manifest|frame_id=P13-03|paper_id=P13|claim_id=P13-结果|visual_type=result-evidence-table|width_budget=0.98linewidth|height_budget=0.86textheight|min_font_pt=10|split_allowed=no|expected_page=64|risks=title-table-source-footline
\begin{frame}[t]{准部分子分布单圈匹配：结果、验证与物理含义}
  \fontsize{10}{10.8}\selectfont
\begingroup
\renewcommand{\arraystretch}{0.84}
\begin{tabularx}{0.98\linewidth}{@{}p{2.55cm}>{\raggedright\arraybackslash}X@{}}
\toprule
源文结果 & 我们推导了非单态夸克分布（包括非极化、螺旋度与横率分布）的单圈匹配条件。该匹配条件也可以被视为一种因子化：它把准夸克分布与实验中可测量的光锥分布联系起来，从而使人们能够从前者提取后者；前者定义为有限核子动量下不含时间的关联函数，因此可以在格点上计算。为了开展准夸克分布的格点模拟，需要采用横向动量截断作为紫外（UV）发散的调节器。这一紫外调节器的选择会在轴向规范（公式）下产生一个……40762、上海市科学技术委员会基金（No. 11DZ2260700）以及中国国家自然科学基金 （No. 11175114）的资助。 acknowledgments \\
验证状态 & 确证：源章节有结果/讨论摘录。 \\
库内可见证据 & 中文源约 482 行；公式命中 31，图命中 3，表/表格命中 0。 \\
物理含义 & 在本主题共同链条中，本论文把上述输入推进到可比较的中间量或观测量；具体强度、误差和外推范围不超出源文档。 \\
未验证项 & 当前工作区没有论文 PDF、原始数据和独立复现脚本；因此数值复核、图像再现和统计显著性不作额外断言。 \\
\bottomrule
\end{tabularx}
\endgroup
  \source{结果证据：\detokenize{refer/papers/准部分子分布单圈匹配_latex/chapters/section05.tex:1-11}；图表/公式计数由本报告生成器对中文源的静态统计得到。}
\end{frame}

%% frame_manifest|frame_id=P13-04|paper_id=P13|claim_id=P13-局限|visual_type=risk-relation-table|width_budget=0.98linewidth|height_budget=0.86textheight|min_font_pt=10|split_allowed=no|expected_page=65|risks=title-table-source-footline
\begin{frame}[t]{准部分子分布单圈匹配：局限、跨论文接口与可复查入口}
  \fontsize{10}{10.8}\selectfont
\begingroup
\renewcommand{\arraystretch}{0.84}
\begin{tabularx}{0.98\linewidth}{@{}p{2.55cm}>{\raggedright\arraybackslash}X@{}}
\toprule
源文局限 & 需要回到结果/讨论/展望章节核对适用区间；本档案不把缺失的独立数据复核补成已证实结论。 \\
跨论文接口 & 是准/赝分布从格点数据走向 MS-bar 或其他连续方案的必要中间层。 \\
极限检查 & 重正化因子抵消截止依赖；a→0 和短距离 z→0 极限必须与匹配阶数、方案和尺度保持一致。 \\
下一步核验 & 补齐 PDF 或原始数据；按源文参数复现关键图表；比较不同动量、流时间、距离、格距或方案转换后的稳定性。 \\
完整来源 & \detokenize{refer/papers/准部分子分布单圈匹配_latex/main.tex}；\detokenize{refer/papers/One_Loop_Matching_Nonsinglet_latex/main.tex}；章节源：\detokenize{refer/papers/准部分子分布单圈匹配_latex/chapters/*.tex}；英中目录：\detokenize{refer/papers/One_Loop_Matching_Nonsinglet_latex}. \\
\bottomrule
\end{tabularx}
\endgroup
  \source{状态：\inferred 为主题接口推断；\unverified 为当前材料边界；完整文件清单见 manifest。}
\end{frame}

%% frame_manifest|frame_id=P14-01|paper_id=P14|claim_id=P14-定位|visual_type=paper-card|width_budget=0.98linewidth|height_budget=0.86textheight|min_font_pt=10|split_allowed=no|expected_page=66|risks=title-table-source-footline
\section{大动量有效理论再论}
\begin{frame}[t]{大动量有效理论再论：研究问题与论文定位}
  \fontsize{10}{10.8}\selectfont
\begingroup
\renewcommand{\arraystretch}{0.84}
\begin{tabularx}{0.98\linewidth}{@{}p{2.55cm}>{\raggedright\arraybackslash}X@{}}
\toprule
论文编号 & P14；主题：非微扰重正化、OPE 与匹配 \\
书目信息 & More On Large-Momentum Effective Theory Approach to Parton Physics Published in Nucl. Phys. B 924, 366–376 (……er MIT-CTP/4914; arXiv:1706.07416 [hep-lat].；arXiv:1706.07416；英/中转排页数 10/9 \\
研究对象 & 本文作者所倡导的大动量有效理论（Large-Momentum Effective Theory，简称 LaMET）为在欧氏格点 QCD 理论中模拟部分子物理提供了一条直接途径。近来，文献中对这一理论表现出浓厚兴趣，其中不乏 对其正…… 相同的大动量（或短距离） 极限。在对误差作恰当量化之后，两种提取方法应与同一批格点数据相比较 \\
章节证据 & 引言；闵氏空间中的微扰匹配与欧氏空间中的非微扰重正化；幂次发散与格点 QCD 中准分布的重正化；准分布与赝分布：异同之比较；结论 \\
证据状态 & \verified：摘要和章节来自现有 LaTeX；\unverified：没有论文 PDF/原始数据独立复核。 \\
\bottomrule
\end{tabularx}
\endgroup
  \source{索引：\detokenize{refer/papers/INDEX.md:14}；正文：\detokenize{refer/papers/大动量有效理论再论_latex/main.tex}；\detokenize{refer/papers/More_On_LaMET_Parton_Physics_latex/main.tex}}
\end{frame}

%% frame_manifest|frame_id=P14-02|paper_id=P14|claim_id=P14-方法|visual_type=method-flow-formula|width_budget=0.98linewidth|height_budget=0.86textheight|min_font_pt=10|split_allowed=no|expected_page=67|risks=title-table-source-footline
\begin{frame}[t]{大动量有效理论再论：理论结构与方法链}
  \fontsize{10}{10.8}\selectfont
\begingroup
\renewcommand{\arraystretch}{0.84}
\begin{tabularx}{0.98\linewidth}{@{}p{2.55cm}>{\raggedright\arraybackslash}X@{}}
\toprule
输入与状态 & 格点/解析输入 → 非微扰重正化、OPE 与匹配中的算符、矩阵元或场配置；具体输入见源章节。 \\
章节路线 & 引言；闵氏空间中的微扰匹配与欧氏空间中的非微扰重正化；幂次发散与格点 QCD 中准分布的重正化；准分布与赝分布：异同之比较；结论 \\
方法摘录 & intro 费曼发明的部分子模型一直是描述高能强相互作用物理最有用的语言之一（引文）。 部分子是强子以光速运动时才存在的理想化对象。因此用现代语言来说，它们是某个有效理论中的有效 自由度，例如软共线有效理论（引文）。强子的部分子性质，如部分子分……因而与 LaMET 方法 完全等价。此外，在该方法中观察到的表观标度行为对收敛问题是否有帮助尚无定论 \\
结构式/示意 & $O^{R}(z,\mu)=Z(z,\mu,a)\,O^{\rm bare}(z,a)$\\[-0.2em]\scriptsize 该式只表达主题变量关系；论文特定推导以源文件为准。 \\
物理校验 & 重正化因子抵消截止依赖；a→0 和短距离 z→0 极限必须与匹配阶数、方案和尺度保持一致。 \\
\bottomrule
\end{tabularx}
\endgroup
  \source{方法证据：\detokenize{refer/papers/大动量有效理论再论_latex/chapters/section01.tex:11-76}；主题结构式为示意。}
\end{frame}

%% frame_manifest|frame_id=P14-03|paper_id=P14|claim_id=P14-结果|visual_type=result-evidence-table|width_budget=0.98linewidth|height_budget=0.86textheight|min_font_pt=10|split_allowed=no|expected_page=68|risks=title-table-source-footline
\begin{frame}[t]{大动量有效理论再论：结果、验证与物理含义}
  \fontsize{10}{10.8}\selectfont
\begingroup
\renewcommand{\arraystretch}{0.84}
\begin{tabularx}{0.98\linewidth}{@{}p{2.55cm}>{\raggedright\arraybackslash}X@{}}
\toprule
源文结果 & 我们给出了旨在进一步阐明 LaMET 方法的若干讨论。我们澄清了该方法中因子化的有效性，并强调了 解析延拓在从欧氏积分恢复闵氏空间积分正确红外行为中的重要性。我们还指出，困扰传统矩方法的幂次 发散在采用非局部表述的准分布方法中不构成问题。最后，我们证明了从 LaMET 所用的同一组格点观测量 中提取部分子分布的 Ioffe 时间分布方法需要完全相同的物理条件。在对误差作恰当量化之后比较两种方法 的结果将是很有意思的 \\
验证状态 & 确证：源章节有结果/讨论摘录。 \\
库内可见证据 & 中文源约 792 行；公式命中 17，图命中 0，表/表格命中 0。 \\
物理含义 & 在本主题共同链条中，本论文把上述输入推进到可比较的中间量或观测量；具体强度、误差和外推范围不超出源文档。 \\
未验证项 & 当前工作区没有论文 PDF、原始数据和独立复现脚本；因此数值复核、图像再现和统计显著性不作额外断言。 \\
\bottomrule
\end{tabularx}
\endgroup
  \source{结果证据：\detokenize{refer/papers/大动量有效理论再论_latex/chapters/section05.tex:1-8}；图表/公式计数由本报告生成器对中文源的静态统计得到。}
\end{frame}

%% frame_manifest|frame_id=P14-04|paper_id=P14|claim_id=P14-局限|visual_type=risk-relation-table|width_budget=0.98linewidth|height_budget=0.86textheight|min_font_pt=10|split_allowed=no|expected_page=69|risks=title-table-source-footline
\begin{frame}[t]{大动量有效理论再论：局限、跨论文接口与可复查入口}
  \fontsize{10}{10.8}\selectfont
\begingroup
\renewcommand{\arraystretch}{0.84}
\begin{tabularx}{0.98\linewidth}{@{}p{2.55cm}>{\raggedright\arraybackslash}X@{}}
\toprule
源文局限 & 需要回到结果/讨论/展望章节核对适用区间；本档案不把缺失的独立数据复核补成已证实结论。 \\
跨论文接口 & 是准/赝分布从格点数据走向 MS-bar 或其他连续方案的必要中间层。 \\
极限检查 & 重正化因子抵消截止依赖；a→0 和短距离 z→0 极限必须与匹配阶数、方案和尺度保持一致。 \\
下一步核验 & 补齐 PDF 或原始数据；按源文参数复现关键图表；比较不同动量、流时间、距离、格距或方案转换后的稳定性。 \\
完整来源 & \detokenize{refer/papers/大动量有效理论再论_latex/main.tex}；\detokenize{refer/papers/More_On_LaMET_Parton_Physics_latex/main.tex}；章节源：\detokenize{refer/papers/大动量有效理论再论_latex/chapters/*.tex}；英中目录：\detokenize{refer/papers/More_On_LaMET_Parton_Physics_latex}. \\
\bottomrule
\end{tabularx}
\endgroup
  \source{状态：\inferred 为主题接口推断；\unverified 为当前材料边界；完整文件清单见 manifest。}
\end{frame}

%% frame_manifest|frame_id=P15-01|paper_id=P15|claim_id=P15-定位|visual_type=paper-card|width_budget=0.98linewidth|height_budget=0.86textheight|min_font_pt=10|split_allowed=no|expected_page=70|risks=title-table-source-footline
\section{准分布动量分布与伪分布}
\begin{frame}[t]{准分布动量分布与伪分布：研究问题与论文定位}
  \fontsize{10}{10.8}\selectfont
\begingroup
\renewcommand{\arraystretch}{0.84}
\begin{tabularx}{0.98\linewidth}{@{}p{2.55cm}>{\raggedright\arraybackslash}X@{}}
\toprule
论文编号 & P15；主题：胶子、强子部分子分布与 TMD \\
书目信息 & Quasi-parton distribution functions, momentum distributions, and pseudo-parton distribution functions Phys. ……heory), 12.38.Gc (Lattice QCD calculations).；arXiv:1705.01488；英/中转排页数 10/9 \\
研究对象 & 我们证明，准部分子分布（quasi-PDFs）可以被看作部分子分布函数（PDFs） 与强子静止系中部分子原始动量分布的混合体。 由此导致 quasi-PDFs 具有复杂的卷积结构， 因而必须使用（公式）GeV 量级的较大强子动量 ……范链接重正化中 棘手因子的（公式）依赖性约去。（公式）依赖性保持不变，并决定 PDFs 的形状 \\
章节证据 & 引言；部分子分布；一般矩阵元素与 Lorentz 不变性；共线 PDFs；横动量依赖分布；准分布（Quasi-Distributions）；定义及与 TMDs 的关系；量子色动力学（QCD）情形；动量分布；因子化模型；伪部分子分布（Pseudo-PDFs）；格点实现；总结 \\
证据状态 & \verified：摘要和章节来自现有 LaTeX；\unverified：没有论文 PDF/原始数据独立复核。 \\
\bottomrule
\end{tabularx}
\endgroup
  \source{索引：\detokenize{refer/papers/INDEX.md:15}；正文：\detokenize{refer/papers/准分布动量分布与伪分布_latex/main.tex}；\detokenize{refer/papers/Quasi_Pseudo_PDF_Radyushkin_latex/main.tex}}
\end{frame}

%% frame_manifest|frame_id=P15-02|paper_id=P15|claim_id=P15-方法|visual_type=method-flow-formula|width_budget=0.98linewidth|height_budget=0.86textheight|min_font_pt=10|split_allowed=no|expected_page=71|risks=title-table-source-footline
\begin{frame}[t]{准分布动量分布与伪分布：理论结构与方法链}
  \fontsize{10}{10.8}\selectfont
\begingroup
\renewcommand{\arraystretch}{0.84}
\begin{tabularx}{0.98\linewidth}{@{}p{2.55cm}>{\raggedright\arraybackslash}X@{}}
\toprule
输入与状态 & 格点/解析输入 → 胶子、强子部分子分布与 TMD中的算符、矩阵元或场配置；具体输入见源章节。 \\
章节路线 & 引言；部分子分布；一般矩阵元素与 Lorentz 不变性；共线 PDFs；横动量依赖分布；准分布（Quasi-Distributions）；定义及与 TMDs 的关系；量子色动力学（QCD）情形；动量分布；因子化模型；伪部分子分布（Pseudo-PDFs）；格点实现；总结 \\
方法摘录 & 历史上（引文），PDFs 的引入 是为了描述 自旋-1/2 夸克。 由于与自旋相关的复杂性 并不影响 部分子分布这一概念本身， 我们从一个简单的标量理论例子 出发。在那种情形下，关于靶的信息 累积在一般矩阵元素（公式）中。 由 Lorentz ……意，式 (（交叉引用）) 给出了（公式）的一个协变定义。 无需假设（公式）或（公式）即可定义（公式） \\
结构式/示意 & $\mathcal O_{\rm Euclid}\xrightarrow[\rm matching\,/\,evolution]{}f_{g/h}(x,\mu)\;{\rm or}\;F_{\rm TMD}(x,b_T,\mu,\zeta)$\\[-0.2em]\scriptsize 该式只表达主题变量关系；论文特定推导以源文件为准。 \\
物理校验 & 纵向分数、横向距离和重正化尺度是不同变量；不能把有限动量或有限 b\_T 的结果直接当作光锥分布。 \\
\bottomrule
\end{tabularx}
\endgroup
  \source{方法证据：\detokenize{refer/papers/准分布动量分布与伪分布_latex/chapters/section02.tex:3-35}；主题结构式为示意。}
\end{frame}

%% frame_manifest|frame_id=P15-03|paper_id=P15|claim_id=P15-结果|visual_type=result-evidence-table|width_budget=0.98linewidth|height_budget=0.86textheight|min_font_pt=10|split_allowed=no|expected_page=72|risks=title-table-source-footline
\begin{frame}[t]{准分布动量分布与伪分布：结果、验证与物理含义}
  \fontsize{10}{10.8}\selectfont
\begingroup
\renewcommand{\arraystretch}{0.84}
\begin{tabularx}{0.98\linewidth}{@{}p{2.55cm}>{\raggedright\arraybackslash}X@{}}
\toprule
源文结果 & 在本文中，我们证明了 quasi-PDFs 可以看作 PDFs 与部分子在强子静止系中 原始动量分布的混合体。 在此图像下，部分子的（公式）动量 既来自强子整体的 运动，也来自静止系中的原始 动量分布。quasi-PDFs 复杂的卷积本质 要求使用（公式）3 GeV 的动量才能消除 原始动量分布的效应并 足够接近 PDF 极限。 作为替代方案，我们建议使用 pseudo-PDF……TMD（公式）软部分的（公式）依赖性与（公式）依赖性的因子化。 研究还证明，在小（公式）处残余的（公式）依赖性 可以用微扰演化解释，相应的（公式）值 对应于（公式） \\
验证状态 & 确证：源章节有结果/讨论摘录。 \\
库内可见证据 & 中文源约 832 行；公式命中 23，图命中 6，表/表格命中 0。 \\
物理含义 & 在本主题共同链条中，本论文把上述输入推进到可比较的中间量或观测量；具体强度、误差和外推范围不超出源文档。 \\
未验证项 & 当前工作区没有论文 PDF、原始数据和独立复现脚本；因此数值复核、图像再现和统计显著性不作额外断言。 \\
\bottomrule
\end{tabularx}
\endgroup
  \source{结果证据：\detokenize{refer/papers/准分布动量分布与伪分布_latex/chapters/section04.tex:1-53}；图表/公式计数由本报告生成器对中文源的静态统计得到。}
\end{frame}

%% frame_manifest|frame_id=P15-04|paper_id=P15|claim_id=P15-局限|visual_type=risk-relation-table|width_budget=0.98linewidth|height_budget=0.86textheight|min_font_pt=10|split_allowed=no|expected_page=73|risks=title-table-source-footline
\begin{frame}[t]{准分布动量分布与伪分布：局限、跨论文接口与可复查入口}
  \fontsize{10}{10.8}\selectfont
\begingroup
\renewcommand{\arraystretch}{0.84}
\begin{tabularx}{0.98\linewidth}{@{}p{2.55cm}>{\raggedright\arraybackslash}X@{}}
\toprule
源文局限 & 需要回到结果/讨论/展望章节核对适用区间；本档案不把缺失的独立数据复核补成已证实结论。 \\
跨论文接口 & 检验 LaMET/赝 PDF/准 TMD 方法是否能够覆盖真实强子结构与实验拟合所需的观测量。 \\
极限检查 & 纵向分数、横向距离和重正化尺度是不同变量；不能把有限动量或有限 b\_T 的结果直接当作光锥分布。 \\
下一步核验 & 补齐 PDF 或原始数据；按源文参数复现关键图表；比较不同动量、流时间、距离、格距或方案转换后的稳定性。 \\
完整来源 & \detokenize{refer/papers/准分布动量分布与伪分布_latex/main.tex}；\detokenize{refer/papers/Quasi_Pseudo_PDF_Radyushkin_latex/main.tex}；章节源：\detokenize{refer/papers/准分布动量分布与伪分布_latex/chapters/*.tex}；英中目录：\detokenize{refer/papers/Quasi_Pseudo_PDF_Radyushkin_latex}. \\
\bottomrule
\end{tabularx}
\endgroup
  \source{状态：\inferred 为主题接口推断；\unverified 为当前材料边界；完整文件清单见 manifest。}
\end{frame}

%% frame_manifest|frame_id=P16-01|paper_id=P16|claim_id=P16-定位|visual_type=paper-card|width_budget=0.98linewidth|height_budget=0.86textheight|min_font_pt=10|split_allowed=no|expected_page=74|risks=title-table-source-footline
\section{部分子伪分布的格点探索}
\begin{frame}[t]{部分子伪分布的格点探索：研究问题与论文定位}
  \fontsize{10}{10.8}\selectfont
\begingroup
\renewcommand{\arraystretch}{0.84}
\begin{tabularx}{0.98\linewidth}{@{}p{2.55cm}>{\raggedright\arraybackslash}X@{}}
\toprule
论文编号 & P16；主题：LaMET、准分布与赝分布 \\
书目信息 & Lattice QCD exploration of parton pseudo-distribution functions Phys. Rev. D 96 , 094503 (2017). DOI: 10.110……4-1 two-column layout has not been retained.；arXiv:1706.05373；英/中转排页数 21/18 \\
研究对象 & 我们演示了一种从格点计算中抽取部分子分布的新方法。 出发点是把一般性的等时矩阵元（公式）视为 Ioffe 时间（公式）与距离（公式）的函数。 下一步是用静止系密度（公式）去除（公式）。 我们的格点计算表明，静止系函数中存在线性指数……与（公式）依赖的因子化。 对于小的（公式），残余的（公式）依赖可以由（公式）的微扰演化加以解释 \\
章节证据 & 引言；部分子分布；一般矩阵元与部分子分布；横动量依赖分布与准分布；量子色动力学（QCD）情形；因子化模型；伪部分子分布；数值研究；结果讨论；静止系密度与（公式）因子；约化 Ioffe 时间分布；夸克—反夸克分解；演化；总结 \\
证据状态 & \verified：摘要和章节来自现有 LaTeX；\unverified：没有论文 PDF/原始数据独立复核。 \\
\bottomrule
\end{tabularx}
\endgroup
  \source{索引：\detokenize{refer/papers/INDEX.md:16}；正文：\detokenize{refer/papers/部分子伪分布的格点探索_latex/main.tex}；\detokenize{refer/papers/Pseudo_PDF_Lattice_Exploration_latex/main.tex}}
\end{frame}

%% frame_manifest|frame_id=P16-02|paper_id=P16|claim_id=P16-方法|visual_type=method-flow-formula|width_budget=0.98linewidth|height_budget=0.86textheight|min_font_pt=10|split_allowed=no|expected_page=75|risks=title-table-source-footline
\begin{frame}[t]{部分子伪分布的格点探索：理论结构与方法链}
  \fontsize{10}{10.8}\selectfont
\begingroup
\renewcommand{\arraystretch}{0.84}
\begin{tabularx}{0.98\linewidth}{@{}p{2.55cm}>{\raggedright\arraybackslash}X@{}}
\toprule
输入与状态 & 格点/解析输入 → LaMET、准分布与赝分布中的算符、矩阵元或场配置；具体输入见源章节。 \\
章节路线 & 引言；部分子分布；一般矩阵元与部分子分布；横动量依赖分布与准分布；量子色动力学（QCD）情形；因子化模型；伪部分子分布；数值研究；结果讨论；静止系密度与（公式）因子；约化 Ioffe 时间分布；夸克—反夸克分解；演化；总结 \\
方法摘录 & 定义部分子分布的基本对象是一个双定域算符的矩阵元， （略去其自旋结构中无关紧要的细节后）它可以一般地写成（公式）。 由于 Lorentz 变换下的不变性，它是两个标量（公式）（记作（公式））与（公式）（或写成（公式），以便对类空的（公式）取正值……位置）其中（公式）是 Euler 常数。（公式）与（公式）之间的重标度因子非常接近 1，因为（公式） \\
结构式/示意 & $\widetilde q(x,P_z,\mu)=\int dy\,C(x,y,\mu/P_z)q(y,\mu)+\mathcal O(\Lambda_{\rm QCD}^{2}/P_z^{2})$\\[-0.2em]\scriptsize 该式只表达主题变量关系；论文特定推导以源文件为准。 \\
物理校验 & 大动量极限 P\_z→∞ 压低幂修正；有限动量、有限间距和核子激发态是必须单独控制的误差源。 \\
\bottomrule
\end{tabularx}
\endgroup
  \source{方法证据：\detokenize{refer/papers/部分子伪分布的格点探索_latex/chapters/section02.tex:3-80}；主题结构式为示意。}
\end{frame}

%% frame_manifest|frame_id=P16-03|paper_id=P16|claim_id=P16-结果|visual_type=result-evidence-table|width_budget=0.98linewidth|height_budget=0.86textheight|min_font_pt=10|split_allowed=no|expected_page=76|risks=title-table-source-footline
\begin{frame}[t]{部分子伪分布的格点探索：结果、验证与物理含义}
  \fontsize{10}{10.8}\selectfont
\begingroup
\renewcommand{\arraystretch}{0.84}
\begin{tabularx}{0.98\linewidth}{@{}p{2.55cm}>{\raggedright\arraybackslash}X@{}}
\toprule
源文结果 & 本文演示了一种从格点计算抽取部分子分布的新方法。 它建立在文献（引文）所阐述的思想之上。 首先，我们把一般性等时矩阵元处理为 Ioffe 时间（公式）与距离（公式）的函数（公式）。 下一步的想法是构造 Ioffe 时间分布（公式）与静止系密度（公式）之比（公式）。 我们的格点计算清楚地表明，静止系函数的（公式）依赖中存在线性成分， 它可以归因于规范链接所生成的预期行为（公式）。……然需要更小的格距以及更大的核子动量范围。 此外，我们还需研究有限体积效应， 并纳入 $\pi$ 介子质量更接近物理点的动力学费米子。 我们计划在未来的工作中处理所有这些问题 \\
验证状态 & 确证：源章节有结果/讨论摘录。 \\
库内可见证据 & 中文源约 1442 行；公式命中 43，图命中 37，表/表格命中 0。 \\
物理含义 & 在本主题共同链条中，本论文把上述输入推进到可比较的中间量或观测量；具体强度、误差和外推范围不超出源文档。 \\
未验证项 & 当前工作区没有论文 PDF、原始数据和独立复现脚本；因此数值复核、图像再现和统计显著性不作额外断言。 \\
\bottomrule
\end{tabularx}
\endgroup
  \source{结果证据：\detokenize{refer/papers/部分子伪分布的格点探索_latex/chapters/section06.tex:1-92}；图表/公式计数由本报告生成器对中文源的静态统计得到。}
\end{frame}

%% frame_manifest|frame_id=P16-04|paper_id=P16|claim_id=P16-局限|visual_type=risk-relation-table|width_budget=0.98linewidth|height_budget=0.86textheight|min_font_pt=10|split_allowed=no|expected_page=77|risks=title-table-source-footline
\begin{frame}[t]{部分子伪分布的格点探索：局限、跨论文接口与可复查入口}
  \fontsize{10}{10.8}\selectfont
\begingroup
\renewcommand{\arraystretch}{0.84}
\begin{tabularx}{0.98\linewidth}{@{}p{2.55cm}>{\raggedright\arraybackslash}X@{}}
\toprule
源文局限 & 需要回到结果/讨论/展望章节核对适用区间；本档案不把缺失的独立数据复核补成已证实结论。 \\
跨论文接口 & 承接格点矩阵元，向重正化、匹配、演化和最终 PDF/TMD 现象学输出。 \\
极限检查 & 大动量极限 P\_z→∞ 压低幂修正；有限动量、有限间距和核子激发态是必须单独控制的误差源。 \\
下一步核验 & 补齐 PDF 或原始数据；按源文参数复现关键图表；比较不同动量、流时间、距离、格距或方案转换后的稳定性。 \\
完整来源 & \detokenize{refer/papers/部分子伪分布的格点探索_latex/main.tex}；\detokenize{refer/papers/Pseudo_PDF_Lattice_Exploration_latex/main.tex}；章节源：\detokenize{refer/papers/部分子伪分布的格点探索_latex/chapters/*.tex}；英中目录：\detokenize{refer/papers/Pseudo_PDF_Lattice_Exploration_latex}. \\
\bottomrule
\end{tabularx}
\endgroup
  \source{状态：\inferred 为主题接口推断；\unverified 为当前材料边界；完整文件清单见 manifest。}
\end{frame}

%% frame_manifest|frame_id=P17-01|paper_id=P17|claim_id=P17-定位|visual_type=paper-card|width_budget=0.98linewidth|height_budget=0.86textheight|min_font_pt=10|split_allowed=no|expected_page=78|risks=title-table-source-footline
\section{欧氏与光锥分布因子化定理}
\begin{frame}[t]{欧氏与光锥分布因子化定理：研究问题与论文定位}
  \fontsize{10}{10.8}\selectfont
\begingroup
\renewcommand{\arraystretch}{0.84}
\begin{tabularx}{0.98\linewidth}{@{}p{2.55cm}>{\raggedright\arraybackslash}X@{}}
\toprule
论文编号 & P17；主题：非微扰重正化、OPE 与匹配 \\
书目信息 & Factorization Euclidean LightCone；arXiv:1801.03917；英/中转排页数 31/28 \\
研究对象 & 在大动量核子态中，可以从格点 QCD 计算得到的规范不变欧氏 Wilson 线算符的矩阵元，通过大动量有效理论（LaMET）展开与标准的光锥部分子分布函数联系起来。这一关系由一个带非平凡匹配系数的因子化公式给出。我们利用算符乘积展……的单圈显式结果并作了比较。我们还证明，（公式）方案中的匹配系数依赖于大的部分子动量而非核子动量 \\
章节证据 & 引言；由算符乘积展开得到的因子化；空间关联函数的 OPE 与因子化；准 PDF 的因子化；赝 PDF 的因子化；单圈阶的等价性；其他重正化方案；数值结果；对格点计算的影响；结论；傅里叶变换；准 PDF 的计算；（公式）展开与 plus 函数 \\
证据状态 & \verified：摘要和章节来自现有 LaTeX；\unverified：没有论文 PDF/原始数据独立复核。 \\
\bottomrule
\end{tabularx}
\endgroup
  \source{索引：\detokenize{refer/papers/INDEX.md:17}；正文：\detokenize{refer/papers/欧氏与光锥分布因子化定理_latex/main.tex}；\detokenize{refer/papers/Factorization_Euclidean_LightCone_latex/main.tex}}
\end{frame}

%% frame_manifest|frame_id=P17-02|paper_id=P17|claim_id=P17-方法|visual_type=method-flow-formula|width_budget=0.98linewidth|height_budget=0.86textheight|min_font_pt=10|split_allowed=no|expected_page=79|risks=title-table-source-footline
\begin{frame}[t]{欧氏与光锥分布因子化定理：理论结构与方法链}
  \fontsize{10}{10.8}\selectfont
\begingroup
\renewcommand{\arraystretch}{0.84}
\begin{tabularx}{0.98\linewidth}{@{}p{2.55cm}>{\raggedright\arraybackslash}X@{}}
\toprule
输入与状态 & 格点/解析输入 → 非微扰重正化、OPE 与匹配中的算符、矩阵元或场配置；具体输入见源章节。 \\
章节路线 & 引言；由算符乘积展开得到的因子化；空间关联函数的 OPE 与因子化；准 PDF 的因子化；赝 PDF 的因子化；单圈阶的等价性；其他重正化方案；数值结果；对格点计算的影响；结论；傅里叶变换；准 PDF 的计算；（公式）展开与 plus 函数 \\
方法摘录 & 部分子分布函数（PDF）是理解强子结构以及对高能散射实验的截面作出预言的关键量。在硬散射过程的 QCD 因子化定理中（引文），相关的 PDF 定义为光锥关联算符在核子态之间的矩阵元。例如，在维数正规化（公式）下，裸的非极化夸克 PDF 为（公式……对准 PDF 与赝 PDF 两种方法中 PDF 格点计算的影响。最后，我们在第（交叉引用）节给出结论 \\
结构式/示意 & $O^{R}(z,\mu)=Z(z,\mu,a)\,O^{\rm bare}(z,a)$\\[-0.2em]\scriptsize 该式只表达主题变量关系；论文特定推导以源文件为准。 \\
物理校验 & 重正化因子抵消截止依赖；a→0 和短距离 z→0 极限必须与匹配阶数、方案和尺度保持一致。 \\
\bottomrule
\end{tabularx}
\endgroup
  \source{方法证据：\detokenize{refer/papers/欧氏与光锥分布因子化定理_latex/chapters/section01.tex:1-98}；主题结构式为示意。}
\end{frame}

%% frame_manifest|frame_id=P17-03|paper_id=P17|claim_id=P17-结果|visual_type=result-evidence-table|width_budget=0.98linewidth|height_budget=0.86textheight|min_font_pt=10|split_allowed=no|expected_page=80|risks=title-table-source-footline
\begin{frame}[t]{欧氏与光锥分布因子化定理：结果、验证与物理含义}
  \fontsize{10}{10.8}\selectfont
\begingroup
\renewcommand{\arraystretch}{0.84}
\begin{tabularx}{0.98\linewidth}{@{}p{2.55cm}>{\raggedright\arraybackslash}X@{}}
\toprule
源文结果 & sec:summary 从 QCD 中规范不变算符乘积的欧氏算符乘积展开出发，我们推导了准 PDF、 与赝 PDF 的因子化公式。这三个欧氏分布函数是相互关联的可观测量，都源自同一基本因子化。对 而言，这一推导意味着式 (（交叉引用）)中的比值确实按（公式）标度，但需要式 (（交叉引用）)的小（公式）因子化公式来提取 PDF。 我们的因子化公式推导适用于重正化 在任何方案下的定……量来减小高扭度修正。这限制了格点计算中有用数据点的数目，如 lattice 所述。要在不作模型假设的情况下实现精确计算，强烈期望走向更细的格距以增加有效数据点的数目 \\
验证状态 & 确证：源章节有结果/讨论摘录。 \\
库内可见证据 & 中文源约 1937 行；公式命中 84，图命中 12，表/表格命中 2。 \\
物理含义 & 在本主题共同链条中，本论文把上述输入推进到可比较的中间量或观测量；具体强度、误差和外推范围不超出源文档。 \\
未验证项 & 当前工作区没有论文 PDF、原始数据和独立复现脚本；因此数值复核、图像再现和统计显著性不作额外断言。 \\
\bottomrule
\end{tabularx}
\endgroup
  \source{结果证据：\detokenize{refer/papers/欧氏与光锥分布因子化定理_latex/chapters/section07.tex:1-16}；图表/公式计数由本报告生成器对中文源的静态统计得到。}
\end{frame}

%% frame_manifest|frame_id=P17-04|paper_id=P17|claim_id=P17-局限|visual_type=risk-relation-table|width_budget=0.98linewidth|height_budget=0.86textheight|min_font_pt=10|split_allowed=no|expected_page=81|risks=title-table-source-footline
\begin{frame}[t]{欧氏与光锥分布因子化定理：局限、跨论文接口与可复查入口}
  \fontsize{10}{10.8}\selectfont
\begingroup
\renewcommand{\arraystretch}{0.84}
\begin{tabularx}{0.98\linewidth}{@{}p{2.55cm}>{\raggedright\arraybackslash}X@{}}
\toprule
源文局限 & 需要回到结果/讨论/展望章节核对适用区间；本档案不把缺失的独立数据复核补成已证实结论。 \\
跨论文接口 & 是准/赝分布从格点数据走向 MS-bar 或其他连续方案的必要中间层。 \\
极限检查 & 重正化因子抵消截止依赖；a→0 和短距离 z→0 极限必须与匹配阶数、方案和尺度保持一致。 \\
下一步核验 & 补齐 PDF 或原始数据；按源文参数复现关键图表；比较不同动量、流时间、距离、格距或方案转换后的稳定性。 \\
完整来源 & \detokenize{refer/papers/欧氏与光锥分布因子化定理_latex/main.tex}；\detokenize{refer/papers/Factorization_Euclidean_LightCone_latex/main.tex}；章节源：\detokenize{refer/papers/欧氏与光锥分布因子化定理_latex/chapters/*.tex}；英中目录：\detokenize{refer/papers/Factorization_Euclidean_LightCone_latex}. \\
\bottomrule
\end{tabularx}
\endgroup
  \source{状态：\inferred 为主题接口推断；\unverified 为当前材料边界；完整文件清单见 manifest。}
\end{frame}

%% frame_manifest|frame_id=P18-01|paper_id=P18|claim_id=P18-定位|visual_type=paper-card|width_budget=0.98linewidth|height_budget=0.86textheight|min_font_pt=10|split_allowed=no|expected_page=82|risks=title-table-source-footline
\section{大动量有效理论综述}
\begin{frame}[t]{大动量有效理论综述：研究问题与论文定位}
  \fontsize{10}{10.8}\selectfont
\begingroup
\renewcommand{\arraystretch}{0.84}
\begin{tabularx}{0.98\linewidth}{@{}p{2.55cm}>{\raggedright\arraybackslash}X@{}}
\toprule
论文编号 & P18；主题：胶子、强子部分子分布与 TMD \\
书目信息 & Large-Momentum Effective Theory；arXiv:2004.03543；英/中转排页数 121/105 \\
研究对象 & 自 Feynman 五十多年前提出部分子模型以来，我们通过大量高能实验数据和专门的全局拟合，对质子的部分子结构有了很多了解。然而，从强相互作用基本理论 QCD 出发计算部分子可观测量（如部分子分布函数 PDF）的进展有限。近年来，……文综述了近几年 LaMET 形式体系及其应用的发展，特别是其有效性的初步格点 QCD 模拟验证 \\
章节证据 & 引言；经由无穷大动量态理解部分子；经由光前关联函数理解部分子；部分子结构的其他途径；大动量有效理论；有限动量下的质子结构；动量重正化群；到 PDF 的有效场论匹配；LaMET 处理部分子物理……分子性质；TMD；结论与展望；附录 A：缩略语、简称与术语；附录 B：约定 \\
证据状态 & \verified：摘要和章节来自现有 LaTeX；\unverified：没有论文 PDF/原始数据独立复核。 \\
\bottomrule
\end{tabularx}
\endgroup
  \source{索引：\detokenize{refer/papers/INDEX.md:18}；正文：\detokenize{refer/papers/大动量有效理论综述_latex/main.tex}；\detokenize{refer/papers/LaMET_RMP_Review_latex/main.tex}}
\end{frame}

%% frame_manifest|frame_id=P18-02|paper_id=P18|claim_id=P18-方法|visual_type=method-flow-formula|width_budget=0.98linewidth|height_budget=0.86textheight|min_font_pt=10|split_allowed=no|expected_page=83|risks=title-table-source-footline
\begin{frame}[t]{大动量有效理论综述：理论结构与方法链}
  \fontsize{10}{10.8}\selectfont
\begingroup
\renewcommand{\arraystretch}{0.84}
\begin{tabularx}{0.98\linewidth}{@{}p{2.55cm}>{\raggedright\arraybackslash}X@{}}
\toprule
输入与状态 & 格点/解析输入 → 胶子、强子部分子分布与 TMD中的算符、矩阵元或场配置；具体输入见源章节。 \\
章节路线 & 引言；经由无穷大动量态理解部分子；经由光前关联函数理解部分子；部分子结构的其他途径；大动量有效理论；有限动量下的质子结构；动量重正化群；到 PDF 的有效场论匹配；LaMET 处理部分子物理……分子性质；TMD；结论与展望；附录 A：缩略语、简称与术语；附录 B：约定 \\
方法摘录 & sec:lamet 正如 lc 所解释的，Feynman 部分子的动机来自描述以大动量（公式）运动的束缚态结构；另一方面，在 QCD 因子化中，它们是忽略 UV 发散的无穷动量极限下出现的有效自由度。调和这两种图像便得到关于强子部分子结构的大动…… LFQ 可能为质子提供诱人的物理图像，但欧氏等时表述更便于开展计算，LaMET 正是沟通二者的桥梁 \\
结构式/示意 & $\mathcal O_{\rm Euclid}\xrightarrow[\rm matching\,/\,evolution]{}f_{g/h}(x,\mu)\;{\rm or}\;F_{\rm TMD}(x,b_T,\mu,\zeta)$\\[-0.2em]\scriptsize 该式只表达主题变量关系；论文特定推导以源文件为准。 \\
物理校验 & 纵向分数、横向距离和重正化尺度是不同变量；不能把有限动量或有限 b\_T 的结果直接当作光锥分布。 \\
\bottomrule
\end{tabularx}
\endgroup
  \source{方法证据：\detokenize{refer/papers/大动量有效理论综述_latex/chapters/section02.tex:1-13}；主题结构式为示意。}
\end{frame}

%% frame_manifest|frame_id=P18-03|paper_id=P18|claim_id=P18-结果|visual_type=result-evidence-table|width_budget=0.98linewidth|height_budget=0.86textheight|min_font_pt=10|split_allowed=no|expected_page=84|risks=title-table-source-footline
\begin{frame}[t]{大动量有效理论综述：结果、验证与物理含义}
  \fontsize{10}{10.8}\selectfont
\begingroup
\renewcommand{\arraystretch}{0.84}
\begin{tabularx}{0.98\linewidth}{@{}p{2.55cm}>{\raggedright\arraybackslash}X@{}}
\toprule
源文结果 & sec:conclusion 自 Feynman 五十多年前提出部分子模型以来，我们对质子部分子结构的理解已大为推进。一方面，SLAC、DESY、CERN、Fermi 实验室、JLab、BNL 等世界各地装置上开展的大量高能实验，使我们得以在不同能量与极化下探测强子结构的方方面面。另一方面，许多部分子观测量被并行提出，为质子结构提供多维描述，包括共线 PDF、TMDPDF、GP……及新技术与新算法的发展，我们预期上述系统误差未来会逐步得到控制。这对确立 LaMET 作为计算部分子物理的系统方法、并让格点计算在 EIC 时代扮演关键角色十分重要 \\
验证状态 & 确证：源章节有结果/讨论摘录。 \\
库内可见证据 & 中文源约 9094 行；公式命中 246，图命中 54，表/表格命中 0。 \\
物理含义 & 在本主题共同链条中，本论文把上述输入推进到可比较的中间量或观测量；具体强度、误差和外推范围不超出源文档。 \\
未验证项 & 当前工作区没有论文 PDF、原始数据和独立复现脚本；因此数值复核、图像再现和统计显著性不作额外断言。 \\
\bottomrule
\end{tabularx}
\endgroup
  \source{结果证据：\detokenize{refer/papers/大动量有效理论综述_latex/chapters/section07.tex:1-33}；图表/公式计数由本报告生成器对中文源的静态统计得到。}
\end{frame}

%% frame_manifest|frame_id=P18-04|paper_id=P18|claim_id=P18-局限|visual_type=risk-relation-table|width_budget=0.98linewidth|height_budget=0.86textheight|min_font_pt=10|split_allowed=no|expected_page=85|risks=title-table-source-footline
\begin{frame}[t]{大动量有效理论综述：局限、跨论文接口与可复查入口}
  \fontsize{10}{10.8}\selectfont
\begingroup
\renewcommand{\arraystretch}{0.84}
\begin{tabularx}{0.98\linewidth}{@{}p{2.55cm}>{\raggedright\arraybackslash}X@{}}
\toprule
源文局限 & 需要回到结果/讨论/展望章节核对适用区间；本档案不把缺失的独立数据复核补成已证实结论。 \\
跨论文接口 & 检验 LaMET/赝 PDF/准 TMD 方法是否能够覆盖真实强子结构与实验拟合所需的观测量。 \\
极限检查 & 纵向分数、横向距离和重正化尺度是不同变量；不能把有限动量或有限 b\_T 的结果直接当作光锥分布。 \\
下一步核验 & 补齐 PDF 或原始数据；按源文参数复现关键图表；比较不同动量、流时间、距离、格距或方案转换后的稳定性。 \\
完整来源 & \detokenize{refer/papers/大动量有效理论综述_latex/main.tex}；\detokenize{refer/papers/LaMET_RMP_Review_latex/main.tex}；章节源：\detokenize{refer/papers/大动量有效理论综述_latex/chapters/*.tex}；英中目录：\detokenize{refer/papers/LaMET_RMP_Review_latex}. \\
\bottomrule
\end{tabularx}
\endgroup
  \source{状态：\inferred 为主题接口推断；\unverified 为当前材料边界；完整文件清单见 manifest。}
\end{frame}

%% frame_manifest|frame_id=P19-01|paper_id=P19|claim_id=P19-定位|visual_type=paper-card|width_budget=0.98linewidth|height_budget=0.86textheight|min_font_pt=10|split_allowed=no|expected_page=86|risks=title-table-source-footline
\section{格点算符非微扰重正化一般方法}
\begin{frame}[t]{格点算符非微扰重正化一般方法：研究问题与论文定位}
  \fontsize{10}{10.8}\selectfont
\begingroup
\renewcommand{\arraystretch}{0.84}
\begin{tabularx}{0.98\linewidth}{@{}p{2.55cm}>{\raggedright\arraybackslash}X@{}}
\toprule
论文编号 & P19；主题：非微扰重正化、OPE 与匹配 \\
书目信息 & General Method Nonpert Renorm；arXiv:hep-lat/9411010；英/中转排页数 27/26 \\
研究对象 & 我们提出一种非微扰方法，用于计算一般复合算符的重正化常数。该方法旨在减小某些系统误差， 这些误差在人们试图从格点算符的矩阵元得到物理预言时是存在的。我们还给出了若干两费米子 算符重正化常数的计算结果，这些结果是用我们的方法、通过（……式）的（公式）格点上得到的。模拟结果令人鼓舞，并建议进一步应用于 四费米子算符和重夸克有效理论 \\
章节证据 & 引言；非微扰重正化条件；方法对一般复合算符的应用；有限算符；对数发散算符；幂次发散算符；方法在数值模拟中的实现；格点微扰论；非微扰结果；结论；附录 \\
证据状态 & \verified：摘要和章节来自现有 LaTeX；\unverified：没有论文 PDF/原始数据独立复核。 \\
\bottomrule
\end{tabularx}
\endgroup
  \source{索引：\detokenize{refer/papers/INDEX.md:19}；正文：\detokenize{refer/papers/格点算符非微扰重正化一般方法_latex/main.tex}；\detokenize{refer/papers/General_Method_Nonpert_Renorm_latex/main.tex}}
\end{frame}

%% frame_manifest|frame_id=P19-02|paper_id=P19|claim_id=P19-方法|visual_type=method-flow-formula|width_budget=0.98linewidth|height_budget=0.86textheight|min_font_pt=10|split_allowed=no|expected_page=87|risks=title-table-source-footline
\begin{frame}[t]{格点算符非微扰重正化一般方法：理论结构与方法链}
  \fontsize{10}{10.8}\selectfont
\begingroup
\renewcommand{\arraystretch}{0.84}
\begin{tabularx}{0.98\linewidth}{@{}p{2.55cm}>{\raggedright\arraybackslash}X@{}}
\toprule
输入与状态 & 格点/解析输入 → 非微扰重正化、OPE 与匹配中的算符、矩阵元或场配置；具体输入见源章节。 \\
章节路线 & 引言；非微扰重正化条件；方法对一般复合算符的应用；有限算符；对数发散算符；幂次发散算符；方法在数值模拟中的实现；格点微扰论；非微扰结果；结论；附录 \\
方法摘录 & sec:casivari 上一节以标量密度与赝标密度的重正化为原型给出了非微扰方案的一般策略。本节简要描述其他情形下所需 的必要修改。我们将遵循第（交叉引用）节引入的算符分类 \\
结构式/示意 & $O^{R}(z,\mu)=Z(z,\mu,a)\,O^{\rm bare}(z,a)$\\[-0.2em]\scriptsize 该式只表达主题变量关系；论文特定推导以源文件为准。 \\
物理校验 & 重正化因子抵消截止依赖；a→0 和短距离 z→0 极限必须与匹配阶数、方案和尺度保持一致。 \\
\bottomrule
\end{tabularx}
\endgroup
  \source{方法证据：\detokenize{refer/papers/格点算符非微扰重正化一般方法_latex/chapters/section03.tex:1-7}；主题结构式为示意。}
\end{frame}

%% frame_manifest|frame_id=P19-03|paper_id=P19|claim_id=P19-结果|visual_type=result-evidence-table|width_budget=0.98linewidth|height_budget=0.86textheight|min_font_pt=10|split_allowed=no|expected_page=88|risks=title-table-source-footline
\begin{frame}[t]{格点算符非微扰重正化一般方法：结果、验证与物理含义}
  \fontsize{10}{10.8}\selectfont
\begingroup
\renewcommand{\arraystretch}{0.84}
\begin{tabularx}{0.98\linewidth}{@{}p{2.55cm}>{\raggedright\arraybackslash}X@{}}
\toprule
源文结果 & sec:concl 我们提出了一种对格点复合算符进行重正化的新非微扰方法。它适用于所有情形，在 Ward 恒等式方法 不适用的场合尤为有用。该方法使人们无需借助任何格点微扰论计算即可从格点模拟得到物理量。我们的 方案能否成功，取决于在当前（公式）值下、低动量非微扰区与离散化误差变得重要的大动量区之间是否存在 窗口。不过这一限制为所有方法所共有（Ward 恒等式方法除外——后者只……预期，随着（公式）增大，对非微扰重正化有用的动量范围将变大。我们计划把本文提出的途径应用于（公式）四费米子算符 (（交叉引用）) 以及静极限理论中重轻轴矢流的重正化 \\
验证状态 & 确证：源章节有结果/讨论摘录。 \\
库内可见证据 & 中文源约 946 行；公式命中 1，图命中 18，表/表格命中 4。 \\
物理含义 & 在本主题共同链条中，本论文把上述输入推进到可比较的中间量或观测量；具体强度、误差和外推范围不超出源文档。 \\
未验证项 & 当前工作区没有论文 PDF、原始数据和独立复现脚本；因此数值复核、图像再现和统计显著性不作额外断言。 \\
\bottomrule
\end{tabularx}
\endgroup
  \source{结果证据：\detokenize{refer/papers/格点算符非微扰重正化一般方法_latex/chapters/section07.tex:1-11}；图表/公式计数由本报告生成器对中文源的静态统计得到。}
\end{frame}

%% frame_manifest|frame_id=P19-04|paper_id=P19|claim_id=P19-局限|visual_type=risk-relation-table|width_budget=0.98linewidth|height_budget=0.86textheight|min_font_pt=10|split_allowed=no|expected_page=89|risks=title-table-source-footline
\begin{frame}[t]{格点算符非微扰重正化一般方法：局限、跨论文接口与可复查入口}
  \fontsize{10}{10.8}\selectfont
\begingroup
\renewcommand{\arraystretch}{0.84}
\begin{tabularx}{0.98\linewidth}{@{}p{2.55cm}>{\raggedright\arraybackslash}X@{}}
\toprule
源文局限 & 需要回到结果/讨论/展望章节核对适用区间；本档案不把缺失的独立数据复核补成已证实结论。 \\
跨论文接口 & 是准/赝分布从格点数据走向 MS-bar 或其他连续方案的必要中间层。 \\
极限检查 & 重正化因子抵消截止依赖；a→0 和短距离 z→0 极限必须与匹配阶数、方案和尺度保持一致。 \\
下一步核验 & 补齐 PDF 或原始数据；按源文参数复现关键图表；比较不同动量、流时间、距离、格距或方案转换后的稳定性。 \\
完整来源 & \detokenize{refer/papers/格点算符非微扰重正化一般方法_latex/main.tex}；\detokenize{refer/papers/General_Method_Nonpert_Renorm_latex/main.tex}；章节源：\detokenize{refer/papers/格点算符非微扰重正化一般方法_latex/chapters/*.tex}；英中目录：\detokenize{refer/papers/General_Method_Nonpert_Renorm_latex}. \\
\bottomrule
\end{tabularx}
\endgroup
  \source{状态：\inferred 为主题接口推断；\unverified 为当前材料边界；完整文件清单见 manifest。}
\end{frame}

%% frame_manifest|frame_id=P20-01|paper_id=P20|claim_id=P20-定位|visual_type=paper-card|width_budget=0.98linewidth|height_budget=0.86textheight|min_font_pt=10|split_allowed=no|expected_page=90|risks=title-table-source-footline
\section{准部分子分布的可重正性}
\begin{frame}[t]{准部分子分布的可重正性：研究问题与论文定位}
  \fontsize{10}{10.8}\selectfont
\begingroup
\renewcommand{\arraystretch}{0.84}
\begin{tabularx}{0.98\linewidth}{@{}p{2.55cm}>{\raggedright\arraybackslash}X@{}}
\toprule
论文编号 & P20；主题：非微扰重正化、OPE 与匹配 \\
书目信息 & Renormalizability of quasiparton distribution functions Published in Phys. Rev. D 96, 094019 (2017); DOI: 10……ty of Quasi Parton Distribution Functions''.；arXiv:1707.03107；英/中转排页数 17/15 \\
研究对象 & 近年来，准部分子分布函数在微扰 QCD 与格点 QCD 两个学界都受到了大量关注，因为它们不仅携带关于部分子分布函数的有效信息，而且可以由格点 QCD 模拟计算。然而，与部分子分布函数不同，准部分子分布函数由于不是由扭度2算符定义……的所有来源，并证明：在 QCD 微扰论的所有阶上，幂次发散以及对数发散都可以通过乘法方式重正化 \\
章节证据 & 引言；坐标空间准 PDF；一圈阶的夸克准 PDF；幂次计数；表面紫外发散；三维积分的有限性；夸克准 PDF 的发散图；重正化；总结 \\
证据状态 & \verified：摘要和章节来自现有 LaTeX；\unverified：没有论文 PDF/原始数据独立复核。 \\
\bottomrule
\end{tabularx}
\endgroup
  \source{索引：\detokenize{refer/papers/INDEX.md:20}；正文：\detokenize{refer/papers/准部分子分布的可重正性_latex/main.tex}；\detokenize{refer/papers/Renormalizability_Quasiparton_latex/main.tex}}
\end{frame}

%% frame_manifest|frame_id=P20-02|paper_id=P20|claim_id=P20-方法|visual_type=method-flow-formula|width_budget=0.98linewidth|height_budget=0.86textheight|min_font_pt=10|split_allowed=no|expected_page=91|risks=title-table-source-footline
\begin{frame}[t]{准部分子分布的可重正性：理论结构与方法链}
  \fontsize{10}{10.8}\selectfont
\begingroup
\renewcommand{\arraystretch}{0.84}
\begin{tabularx}{0.98\linewidth}{@{}p{2.55cm}>{\raggedright\arraybackslash}X@{}}
\toprule
输入与状态 & 格点/解析输入 → 非微扰重正化、OPE 与匹配中的算符、矩阵元或场配置；具体输入见源章节。 \\
章节路线 & 引言；坐标空间准 PDF；一圈阶的夸克准 PDF；幂次计数；表面紫外发散；三维积分的有限性；夸克准 PDF 的发散图；重正化；总结 \\
方法摘录 & sec:intro 部分子分布函数（PDF）（公式）定义为：在因子化标度（公式）下探测，在快速运动的强子（公式）中找到一个夸克、反夸克或胶子（分别对应（公式））、且其携带的强子动量分数介于（公式）与（公式）之间的概率分布（引文）。它们是 QCD……不与其他任何量混合，这完成了我们对坐标空间准 PDF 可重正性的证明。最后在第（交叉引用）节给出总结 \\
结构式/示意 & $O^{R}(z,\mu)=Z(z,\mu,a)\,O^{\rm bare}(z,a)$\\[-0.2em]\scriptsize 该式只表达主题变量关系；论文特定推导以源文件为准。 \\
物理校验 & 重正化因子抵消截止依赖；a→0 和短距离 z→0 极限必须与匹配阶数、方案和尺度保持一致。 \\
\bottomrule
\end{tabularx}
\endgroup
  \source{方法证据：\detokenize{refer/papers/准部分子分布的可重正性_latex/chapters/section01.tex:2-43}；主题结构式为示意。}
\end{frame}

%% frame_manifest|frame_id=P20-03|paper_id=P20|claim_id=P20-结果|visual_type=result-evidence-table|width_budget=0.98linewidth|height_budget=0.86textheight|min_font_pt=10|split_allowed=no|expected_page=92|risks=title-table-source-footline
\begin{frame}[t]{准部分子分布的可重正性：结果、验证与物理含义}
  \fontsize{10}{10.8}\selectfont
\begingroup
\renewcommand{\arraystretch}{0.84}
\begin{tabularx}{0.98\linewidth}{@{}p{2.55cm}>{\raggedright\arraybackslash}X@{}}
\toprule
源文结果 & sec:summary 我们证明了准 PDF 的紫外发散行为与 PDF 的大不相同。准 PDF 的重正化是一个简单的乘法因子，而 PDF 的重正化是卷积形式，其原因在于 PDF 的紫外发散并非完全局域，因而不同味的 PDF 在重正化下相互混合。 我们证明了坐标空间准 PDF 微扰紫外发散的时空局域性，使得坐标空间准 PDF 在 QCD 微扰论的所有阶上都可以乘法重正化，如式 (……动量空间准 PDF 在大（公式）区域是不适定的。 补充说明： 最近出现了两项关于准 PDF 重正化的独立研究（引文），虽然其方法与我们的很不一样，却得到了相同的结论 \\
验证状态 & 确证：源章节有结果/讨论摘录。 \\
库内可见证据 & 中文源约 884 行；公式命中 27，图命中 18，表/表格命中 0。 \\
物理含义 & 在本主题共同链条中，本论文把上述输入推进到可比较的中间量或观测量；具体强度、误差和外推范围不超出源文档。 \\
未验证项 & 当前工作区没有论文 PDF、原始数据和独立复现脚本；因此数值复核、图像再现和统计显著性不作额外断言。 \\
\bottomrule
\end{tabularx}
\endgroup
  \source{结果证据：\detokenize{refer/papers/准部分子分布的可重正性_latex/chapters/section06.tex:2-14}；图表/公式计数由本报告生成器对中文源的静态统计得到。}
\end{frame}

%% frame_manifest|frame_id=P20-04|paper_id=P20|claim_id=P20-局限|visual_type=risk-relation-table|width_budget=0.98linewidth|height_budget=0.86textheight|min_font_pt=10|split_allowed=no|expected_page=93|risks=title-table-source-footline
\begin{frame}[t]{准部分子分布的可重正性：局限、跨论文接口与可复查入口}
  \fontsize{10}{10.8}\selectfont
\begingroup
\renewcommand{\arraystretch}{0.84}
\begin{tabularx}{0.98\linewidth}{@{}p{2.55cm}>{\raggedright\arraybackslash}X@{}}
\toprule
源文局限 & 需要回到结果/讨论/展望章节核对适用区间；本档案不把缺失的独立数据复核补成已证实结论。 \\
跨论文接口 & 是准/赝分布从格点数据走向 MS-bar 或其他连续方案的必要中间层。 \\
极限检查 & 重正化因子抵消截止依赖；a→0 和短距离 z→0 极限必须与匹配阶数、方案和尺度保持一致。 \\
下一步核验 & 补齐 PDF 或原始数据；按源文参数复现关键图表；比较不同动量、流时间、距离、格距或方案转换后的稳定性。 \\
完整来源 & \detokenize{refer/papers/准部分子分布的可重正性_latex/main.tex}；\detokenize{refer/papers/Renormalizability_Quasiparton_latex/main.tex}；章节源：\detokenize{refer/papers/准部分子分布的可重正性_latex/chapters/*.tex}；英中目录：\detokenize{refer/papers/Renormalizability_Quasiparton_latex}. \\
\bottomrule
\end{tabularx}
\endgroup
  \source{状态：\inferred 为主题接口推断；\unverified 为当前材料边界；完整文件清单见 manifest。}
\end{frame}

%% frame_manifest|frame_id=P21-01|paper_id=P21|claim_id=P21-定位|visual_type=paper-card|width_budget=0.98linewidth|height_budget=0.86textheight|min_font_pt=10|split_allowed=no|expected_page=94|risks=title-table-source-footline
\section{威尔逊线重正化改进准分布}
\begin{frame}[t]{威尔逊线重正化改进准分布：研究问题与论文定位}
  \fontsize{10}{10.8}\selectfont
\begingroup
\renewcommand{\arraystretch}{0.84}
\begin{tabularx}{0.98\linewidth}{@{}p{2.55cm}>{\raggedright\arraybackslash}X@{}}
\toprule
论文编号 & P21；主题：非微扰重正化、OPE 与匹配 \\
书目信息 & Improved quasi parton distribution through Wilson line renormalization Nucl. Phys. B 915, 1–9 (2017); arXiv:……t]; doi:10.1016/j.nuclphysb.2016.12.011 [?].；arXiv:1609.08102；英/中转排页数 8/7 \\
研究对象 & 近期的发展表明，强子的光锥部分子分布可以在大强子动量极限下，从类空关联子——即所谓准部分子分布——中直接提取。 与通常的光锥部分子分布不同，准部分子分布含有与威尔逊线自能相联系的紫外（UV）幂次发散。 我们证明：在耦合展开的所有阶……量版本，我们给出了带格点正规化的改进非单态准夸克分布与维数正规化下相应夸克分布之间的单圈匹配核 \\
章节证据 & 引言；通过威尔逊线重正化改进准夸克分布；经由格点微扰论实现格点与连续统之间的匹配；结论；附录 \\
证据状态 & \verified：摘要和章节来自现有 LaTeX；\unverified：没有论文 PDF/原始数据独立复核。 \\
\bottomrule
\end{tabularx}
\endgroup
  \source{索引：\detokenize{refer/papers/INDEX.md:21}；正文：\detokenize{refer/papers/威尔逊线重正化改进准分布_latex/main.tex}；\detokenize{refer/papers/Wilson_Line_Renorm_Improved_QPD_latex/main.tex}}
\end{frame}

%% frame_manifest|frame_id=P21-02|paper_id=P21|claim_id=P21-方法|visual_type=method-flow-formula|width_budget=0.98linewidth|height_budget=0.86textheight|min_font_pt=10|split_allowed=no|expected_page=95|risks=title-table-source-footline
\begin{frame}[t]{威尔逊线重正化改进准分布：理论结构与方法链}
  \fontsize{10}{10.8}\selectfont
\begingroup
\renewcommand{\arraystretch}{0.84}
\begin{tabularx}{0.98\linewidth}{@{}p{2.55cm}>{\raggedright\arraybackslash}X@{}}
\toprule
输入与状态 & 格点/解析输入 → 非微扰重正化、OPE 与匹配中的算符、矩阵元或场配置；具体输入见源章节。 \\
章节路线 & 引言；通过威尔逊线重正化改进准夸克分布；经由格点微扰论实现格点与连续统之间的匹配；结论；附录 \\
方法摘录 & QCD 最重要的目标之一，是从其基本自由度——夸克与胶子——出发理解强子结构。这必然进入 QCD 的非微扰区，一般而言，难以直接从 QCD 拉氏量获得第一性原理的结果。格点 QCD 是获得此类结果的有力工具，它是一种定义在欧氏时空上的方法，已被……尔逊线自能导致的幂次发散的重正化；第 3 节给出准夸克分布的格点微扰论匹配；最后在第 4 节给出结论 \\
结构式/示意 & $O^{R}(z,\mu)=Z(z,\mu,a)\,O^{\rm bare}(z,a)$\\[-0.2em]\scriptsize 该式只表达主题变量关系；论文特定推导以源文件为准。 \\
物理校验 & 重正化因子抵消截止依赖；a→0 和短距离 z→0 极限必须与匹配阶数、方案和尺度保持一致。 \\
\bottomrule
\end{tabularx}
\endgroup
  \source{方法证据：\detokenize{refer/papers/威尔逊线重正化改进准分布_latex/chapters/section01.tex:1-12}；主题结构式为示意。}
\end{frame}

%% frame_manifest|frame_id=P21-03|paper_id=P21|claim_id=P21-结果|visual_type=result-evidence-table|width_budget=0.98linewidth|height_budget=0.86textheight|min_font_pt=10|split_allowed=no|expected_page=96|risks=title-table-source-footline
\begin{frame}[t]{威尔逊线重正化改进准分布：结果、验证与物理含义}
  \fontsize{10}{10.8}\selectfont
\begingroup
\renewcommand{\arraystretch}{0.84}
\begin{tabularx}{0.98\linewidth}{@{}p{2.55cm}>{\raggedright\arraybackslash}X@{}}
\toprule
源文结果 & 在本文中，我们证明了准 PDF 在格点正规化下出现的幂次发散可以在所有圈次上用一个质量反项消除；该反项可理解为沿威尔逊线运动的检验粒子的质量重正化。经过这样的质量重正化之后，准分布得到改进，至多只含对数发散。我们还基于一个简单的离散化规范作用量版本，利用格点微扰论给出了准夸克分布在格点与连续统之间的单圈匹配。我们的结果有助于从格点数据可靠地提取物理 PDF \\
验证状态 & 确证：源章节有结果/讨论摘录。 \\
库内可见证据 & 中文源约 568 行；公式命中 6，图命中 14，表/表格命中 0。 \\
物理含义 & 在本主题共同链条中，本论文把上述输入推进到可比较的中间量或观测量；具体强度、误差和外推范围不超出源文档。 \\
未验证项 & 当前工作区没有论文 PDF、原始数据和独立复现脚本；因此数值复核、图像再现和统计显著性不作额外断言。 \\
\bottomrule
\end{tabularx}
\endgroup
  \source{结果证据：\detokenize{refer/papers/威尔逊线重正化改进准分布_latex/chapters/section04.tex:1-4}；图表/公式计数由本报告生成器对中文源的静态统计得到。}
\end{frame}

%% frame_manifest|frame_id=P21-04|paper_id=P21|claim_id=P21-局限|visual_type=risk-relation-table|width_budget=0.98linewidth|height_budget=0.86textheight|min_font_pt=10|split_allowed=no|expected_page=97|risks=title-table-source-footline
\begin{frame}[t]{威尔逊线重正化改进准分布：局限、跨论文接口与可复查入口}
  \fontsize{10}{10.8}\selectfont
\begingroup
\renewcommand{\arraystretch}{0.84}
\begin{tabularx}{0.98\linewidth}{@{}p{2.55cm}>{\raggedright\arraybackslash}X@{}}
\toprule
源文局限 & 需要回到结果/讨论/展望章节核对适用区间；本档案不把缺失的独立数据复核补成已证实结论。 \\
跨论文接口 & 是准/赝分布从格点数据走向 MS-bar 或其他连续方案的必要中间层。 \\
极限检查 & 重正化因子抵消截止依赖；a→0 和短距离 z→0 极限必须与匹配阶数、方案和尺度保持一致。 \\
下一步核验 & 补齐 PDF 或原始数据；按源文参数复现关键图表；比较不同动量、流时间、距离、格距或方案转换后的稳定性。 \\
完整来源 & \detokenize{refer/papers/威尔逊线重正化改进准分布_latex/main.tex}；\detokenize{refer/papers/Wilson_Line_Renorm_Improved_QPD_latex/main.tex}；章节源：\detokenize{refer/papers/威尔逊线重正化改进准分布_latex/chapters/*.tex}；英中目录：\detokenize{refer/papers/Wilson_Line_Renorm_Improved_QPD_latex}. \\
\bottomrule
\end{tabularx}
\endgroup
  \source{状态：\inferred 为主题接口推断；\unverified 为当前材料边界；完整文件清单见 manifest。}
\end{frame}

%% frame_manifest|frame_id=P22-01|paper_id=P22|claim_id=P22-定位|visual_type=paper-card|width_budget=0.98linewidth|height_budget=0.86textheight|min_font_pt=10|split_allowed=no|expected_page=98|risks=title-table-source-footline
\section{大动量有效理论的重正化}
\begin{frame}[t]{大动量有效理论的重正化：研究问题与论文定位}
  \fontsize{10}{10.8}\selectfont
\begingroup
\renewcommand{\arraystretch}{0.84}
\begin{tabularx}{0.98\linewidth}{@{}p{2.55cm}>{\raggedright\arraybackslash}X@{}}
\toprule
论文编号 & P22；主题：非微扰重正化、OPE 与匹配 \\
书目信息 & Renormalization in Large Momentum Effective Theory of Parton Physics Published in Phys. Rev. Lett. 120, 1120……:1706.08962 [hep-ph]; preprint MIT-CTP/4916.；arXiv:1706.08962；英/中转排页数 7/6 \\
研究对象 & 在大动量有效场论处理部分子物理的方法中，由空间方向上的直线威尔逊线连接起来的夸克场与胶子场非局域算符的矩阵元，是在格点量子色动力学中被作为强子动量的函数加以计算的。利用重夸克有效理论的表述，我们证明这些算符在微扰论的所有阶都具有乘……是格点正规化皆如此。该结果为在蒙特卡罗模拟中通过恰当重正化的观测量抽取部分子性质提供了理论基础 \\
章节证据 & 引言；含辅助“重夸克”场的有效 QCD；维数正规化下含“重夸克”的有效 QCD 的重正化；维数正规化下非局域算符的重正化；格点正规化下非局域算符的重正化；结论 \\
证据状态 & \verified：摘要和章节来自现有 LaTeX；\unverified：没有论文 PDF/原始数据独立复核。 \\
\bottomrule
\end{tabularx}
\endgroup
  \source{索引：\detokenize{refer/papers/INDEX.md:22}；正文：\detokenize{refer/papers/大动量有效理论的重正化_latex/main.tex}；\detokenize{refer/papers/Renorm_LaMET_PRL_latex/main.tex}}
\end{frame}

%% frame_manifest|frame_id=P22-02|paper_id=P22|claim_id=P22-方法|visual_type=method-flow-formula|width_budget=0.98linewidth|height_budget=0.86textheight|min_font_pt=10|split_allowed=no|expected_page=99|risks=title-table-source-footline
\begin{frame}[t]{大动量有效理论的重正化：理论结构与方法链}
  \fontsize{10}{10.8}\selectfont
\begingroup
\renewcommand{\arraystretch}{0.84}
\begin{tabularx}{0.98\linewidth}{@{}p{2.55cm}>{\raggedright\arraybackslash}X@{}}
\toprule
输入与状态 & 格点/解析输入 → 非微扰重正化、OPE 与匹配中的算符、矩阵元或场配置；具体输入见源章节。 \\
章节路线 & 引言；含辅助“重夸克”场的有效 QCD；维数正规化下含“重夸克”的有效 QCD 的重正化；维数正规化下非局域算符的重正化；格点正规化下非局域算符的重正化；结论 \\
方法摘录 & 量子色动力学（QCD）最重要的目标之一，是从其基本自由度——夸克和胶子——出发理解强子结构。格点 QCD 是获得这类结果的一个有力工具，它已被用于较精确地研究强子的静态性质，例如质量、电荷半径等（参见文献（引文））。然而一般而言，它无法直接处理……如 LaMET 框架中定义的、用以抽取广义部分子分布、横动量依赖分布、强子分布振幅等的欧几里得观测量 \\
结构式/示意 & $O^{R}(z,\mu)=Z(z,\mu,a)\,O^{\rm bare}(z,a)$\\[-0.2em]\scriptsize 该式只表达主题变量关系；论文特定推导以源文件为准。 \\
物理校验 & 重正化因子抵消截止依赖；a→0 和短距离 z→0 极限必须与匹配阶数、方案和尺度保持一致。 \\
\bottomrule
\end{tabularx}
\endgroup
  \source{方法证据：\detokenize{refer/papers/大动量有效理论的重正化_latex/chapters/section01.tex:1-15}；主题结构式为示意。}
\end{frame}

%% frame_manifest|frame_id=P22-03|paper_id=P22|claim_id=P22-结果|visual_type=result-evidence-table|width_budget=0.98linewidth|height_budget=0.86textheight|min_font_pt=10|split_allowed=no|expected_page=100|risks=title-table-source-footline
\begin{frame}[t]{大动量有效理论的重正化：结果、验证与物理含义}
  \fontsize{10}{10.8}\selectfont
\begingroup
\renewcommand{\arraystretch}{0.84}
\begin{tabularx}{0.98\linewidth}{@{}p{2.55cm}>{\raggedright\arraybackslash}X@{}}
\toprule
源文结果 & 在本文中，我们证明了 LaMET 中计算的涉及类空威尔逊线的非局域算符，可以在微扰论所有阶进行乘法重正化。除了产生依赖（公式）的指数因子的质量抵消项之外，其余所有重正化常数都与（公式）无关，并与 HQET 中轻重夸克流的重正化相关联。这为在格点 QCD 中定义准 PDF 的连续极限提供了重要的理论基础，而后者是利用 LaMET 方法精确计算部分子物理所必需的。 注：本文完稿之际，我们获悉 DESY 团队也做了类似的考虑（引文） \\
验证状态 & 确证：源章节有结果/讨论摘录。 \\
库内可见证据 & 中文源约 393 行；公式命中 13，图命中 0，表/表格命中 0。 \\
物理含义 & 在本主题共同链条中，本论文把上述输入推进到可比较的中间量或观测量；具体强度、误差和外推范围不超出源文档。 \\
未验证项 & 当前工作区没有论文 PDF、原始数据和独立复现脚本；因此数值复核、图像再现和统计显著性不作额外断言。 \\
\bottomrule
\end{tabularx}
\endgroup
  \source{结果证据：\detokenize{refer/papers/大动量有效理论的重正化_latex/chapters/section06.tex:1-5}；图表/公式计数由本报告生成器对中文源的静态统计得到。}
\end{frame}

%% frame_manifest|frame_id=P22-04|paper_id=P22|claim_id=P22-局限|visual_type=risk-relation-table|width_budget=0.98linewidth|height_budget=0.86textheight|min_font_pt=10|split_allowed=no|expected_page=101|risks=title-table-source-footline
\begin{frame}[t]{大动量有效理论的重正化：局限、跨论文接口与可复查入口}
  \fontsize{10}{10.8}\selectfont
\begingroup
\renewcommand{\arraystretch}{0.84}
\begin{tabularx}{0.98\linewidth}{@{}p{2.55cm}>{\raggedright\arraybackslash}X@{}}
\toprule
源文局限 & 需要回到结果/讨论/展望章节核对适用区间；本档案不把缺失的独立数据复核补成已证实结论。 \\
跨论文接口 & 是准/赝分布从格点数据走向 MS-bar 或其他连续方案的必要中间层。 \\
极限检查 & 重正化因子抵消截止依赖；a→0 和短距离 z→0 极限必须与匹配阶数、方案和尺度保持一致。 \\
下一步核验 & 补齐 PDF 或原始数据；按源文参数复现关键图表；比较不同动量、流时间、距离、格距或方案转换后的稳定性。 \\
完整来源 & \detokenize{refer/papers/大动量有效理论的重正化_latex/main.tex}；\detokenize{refer/papers/Renorm_LaMET_PRL_latex/main.tex}；章节源：\detokenize{refer/papers/大动量有效理论的重正化_latex/chapters/*.tex}；英中目录：\detokenize{refer/papers/Renorm_LaMET_PRL_latex}. \\
\bottomrule
\end{tabularx}
\endgroup
  \source{状态：\inferred 为主题接口推断；\unverified 为当前材料边界；完整文件清单见 manifest。}
\end{frame}

%% frame_manifest|frame_id=P23-01|paper_id=P23|claim_id=P23-定位|visual_type=paper-card|width_budget=0.98linewidth|height_budget=0.86textheight|min_font_pt=10|split_allowed=no|expected_page=102|risks=title-table-source-footline
\section{非定域夸克双线性的非微扰重正化}
\begin{frame}[t]{非定域夸克双线性的非微扰重正化：研究问题与论文定位}
  \fontsize{10}{10.8}\selectfont
\begingroup
\renewcommand{\arraystretch}{0.84}
\begin{tabularx}{0.98\linewidth}{@{}p{2.55cm}>{\raggedright\arraybackslash}X@{}}
\toprule
论文编号 & P23；主题：非微扰重正化、OPE 与匹配 \\
书目信息 & Nonperturbative Renormalization of Nonlocal Quark Bilinears from Lattice QCD Published in Phys. Rev. Lett. 1……s on the lattice using an auxiliary field.''；arXiv:1707.07152；英/中转排页数 11/10 \\
研究对象 & 准部分子分布（quasi-PDF）为利用格点QCD对部分子分布函数（PDF）进行 从头计算提供了一条途径。准部分子分布计算中面临的问题之一是 一个非定域算符的重正化。通过引入一个辅助场，我们可以在一个扩展理论中 用一对局域算符来替……格距下的 重正化结果。重正化后的矩阵元在不同的Wilson线离散化和不同格距之间 是相互一致的 \\
章节证据 & 未验证：源文件未提供可用的散文摘录。 \\
证据状态 & \verified：摘要和章节来自现有 LaTeX；\unverified：没有论文 PDF/原始数据独立复核。 \\
\bottomrule
\end{tabularx}
\endgroup
  \source{索引：\detokenize{refer/papers/INDEX.md:23}；正文：\detokenize{refer/papers/非定域夸克双线性的非微扰重正化_latex/main.tex}；\detokenize{refer/papers/Nonpert_Renorm_Quark_Bilinears_latex/main.tex}}
\end{frame}

%% frame_manifest|frame_id=P23-02|paper_id=P23|claim_id=P23-方法|visual_type=method-flow-formula|width_budget=0.98linewidth|height_budget=0.86textheight|min_font_pt=10|split_allowed=no|expected_page=103|risks=title-table-source-footline
\begin{frame}[t]{非定域夸克双线性的非微扰重正化：理论结构与方法链}
  \fontsize{10}{10.8}\selectfont
\begingroup
\renewcommand{\arraystretch}{0.84}
\begin{tabularx}{0.98\linewidth}{@{}p{2.55cm}>{\raggedright\arraybackslash}X@{}}
\toprule
输入与状态 & 格点/解析输入 → 非微扰重正化、OPE 与匹配中的算符、矩阵元或场配置；具体输入见源章节。 \\
章节路线 & 未验证：源文件未提供可用的散文摘录。 \\
方法摘录 & 部分子分布函数（PDF）描述夸克和胶子在质子内部按纵向动量的分布。它们具有 普适性：同样的PDF出现在许多不同的散射过程中，并且是通过对撞机数据的全局 拟合在唯象上确定的。除了最低的几个Mellin矩之外，PDF一直抗拒从头 计算。格点QCD可……并行工作（引文）。我们也 注意到另一篇讨论准PDF可重正化性、但不使用辅助场方法的最新文 章（引文） \\
结构式/示意 & $O^{R}(z,\mu)=Z(z,\mu,a)\,O^{\rm bare}(z,a)$\\[-0.2em]\scriptsize 该式只表达主题变量关系；论文特定推导以源文件为准。 \\
物理校验 & 重正化因子抵消截止依赖；a→0 和短距离 z→0 极限必须与匹配阶数、方案和尺度保持一致。 \\
\bottomrule
\end{tabularx}
\endgroup
  \source{方法证据：\detokenize{refer/papers/非定域夸克双线性的非微扰重正化_latex/chapters/section01.tex:1-289}；主题结构式为示意。}
\end{frame}

%% frame_manifest|frame_id=P23-03|paper_id=P23|claim_id=P23-结果|visual_type=result-evidence-table|width_budget=0.98linewidth|height_budget=0.86textheight|min_font_pt=10|split_allowed=no|expected_page=104|risks=title-table-source-footline
\begin{frame}[t]{非定域夸克双线性的非微扰重正化：结果、验证与物理含义}
  \fontsize{10}{10.8}\selectfont
\begingroup
\renewcommand{\arraystretch}{0.84}
\begin{tabularx}{0.98\linewidth}{@{}p{2.55cm}>{\raggedright\arraybackslash}X@{}}
\toprule
源文结果 & ……在重正化矩阵元中，我们还 看到涂抹的好处：不涂抹时统计误差在大（公式）处迅速增大，（公式）之后便没有可用信号。使用涂抹后，我们能够看到矩阵元在大（公式）处回归趋于 零。在大（公式）处，五步HYP涂抹也比一步给出更精确的数据。在 图（交叉引用）中，我们比较了两个系综使用五步HYP涂抹的重正化数据。 二者符合得极好，这表明线性发散已被控制住，离散化效应不大。（公式或图表位置）最……作的初步版本之后（引文）， 我们得知存在一项基于辅助场方法的独立并行工作（引文）。我们也 注意到另一篇讨论准PDF可重正化性、但不使用辅助场方法的最新文 章（引文） \\
验证状态 & 确证：源章节有结果/讨论摘录。 \\
库内可见证据 & 中文源约 733 行；公式命中 12，图命中 13，表/表格命中 0。 \\
物理含义 & 在本主题共同链条中，本论文把上述输入推进到可比较的中间量或观测量；具体强度、误差和外推范围不超出源文档。 \\
未验证项 & 当前工作区没有论文 PDF、原始数据和独立复现脚本；因此数值复核、图像再现和统计显著性不作额外断言。 \\
\bottomrule
\end{tabularx}
\endgroup
  \source{结果证据：\detokenize{refer/papers/非定域夸克双线性的非微扰重正化_latex/chapters/section01.tex:1-289}；图表/公式计数由本报告生成器对中文源的静态统计得到。}
\end{frame}

%% frame_manifest|frame_id=P23-04|paper_id=P23|claim_id=P23-局限|visual_type=risk-relation-table|width_budget=0.98linewidth|height_budget=0.86textheight|min_font_pt=10|split_allowed=no|expected_page=105|risks=title-table-source-footline
\begin{frame}[t]{非定域夸克双线性的非微扰重正化：局限、跨论文接口与可复查入口}
  \fontsize{10}{10.8}\selectfont
\begingroup
\renewcommand{\arraystretch}{0.84}
\begin{tabularx}{0.98\linewidth}{@{}p{2.55cm}>{\raggedright\arraybackslash}X@{}}
\toprule
源文局限 & 需要回到结果/讨论/展望章节核对适用区间；本档案不把缺失的独立数据复核补成已证实结论。 \\
跨论文接口 & 是准/赝分布从格点数据走向 MS-bar 或其他连续方案的必要中间层。 \\
极限检查 & 重正化因子抵消截止依赖；a→0 和短距离 z→0 极限必须与匹配阶数、方案和尺度保持一致。 \\
下一步核验 & 补齐 PDF 或原始数据；按源文参数复现关键图表；比较不同动量、流时间、距离、格距或方案转换后的稳定性。 \\
完整来源 & \detokenize{refer/papers/非定域夸克双线性的非微扰重正化_latex/main.tex}；\detokenize{refer/papers/Nonpert_Renorm_Quark_Bilinears_latex/main.tex}；章节源：\detokenize{refer/papers/非定域夸克双线性的非微扰重正化_latex/chapters/*.tex}；英中目录：\detokenize{refer/papers/Nonpert_Renorm_Quark_Bilinears_latex}. \\
\bottomrule
\end{tabularx}
\endgroup
  \source{状态：\inferred 为主题接口推断；\unverified 为当前材料边界；完整文件清单见 manifest。}
\end{frame}

%% frame_manifest|frame_id=P24-01|paper_id=P24|claim_id=P24-定位|visual_type=paper-card|width_budget=0.98linewidth|height_budget=0.86textheight|min_font_pt=10|split_allowed=no|expected_page=106|risks=title-table-source-footline
\section{准PDF完整非微扰重正化方案}
\begin{frame}[t]{准PDF完整非微扰重正化方案：研究问题与论文定位}
  \fontsize{10}{10.8}\selectfont
\begingroup
\renewcommand{\arraystretch}{0.84}
\begin{tabularx}{0.98\linewidth}{@{}p{2.55cm}>{\raggedright\arraybackslash}X@{}}
\toprule
论文编号 & P24；主题：非微扰重正化、OPE 与匹配 \\
书目信息 & A complete non-perturbative renormalization prescription for quasi-PDFs；arXiv:1706.00265；英/中转排页数 22/19 \\
研究对象 & 在本工作中，我们首次在 RI（公式）方案下给出了非极化、螺旋度和横向性准部分子分布 （quasi-PDF）的非微扰重正化。所提出的处方同时处理重正化的各个方面：对数发散、 有限重正化，以及含 Wilson 线的费米子算符矩阵元中出……应用于一个具有（公式）动力学夸克、（公式）介子质量约 为 375 MeV 的扭转质量费米子系综 \\
章节证据 & 引言；理论设定；核子矩阵元；重正化方案；转换到（公式）方案；结果；RI（公式）方案及其到（公式）方案的转换；系统性不确定性的评估；重正化的结果；与光前 PDF 的匹配；结论与讨论 \\
证据状态 & \verified：摘要和章节来自现有 LaTeX；\unverified：没有论文 PDF/原始数据独立复核。 \\
\bottomrule
\end{tabularx}
\endgroup
  \source{索引：\detokenize{refer/papers/INDEX.md:24}；正文：\detokenize{refer/papers/准PDF完整非微扰重正化方案_latex/main.tex}；\detokenize{refer/papers/Complete_NP_Renorm_QuasiPDF_latex/main.tex}}
\end{frame}

%% frame_manifest|frame_id=P24-02|paper_id=P24|claim_id=P24-方法|visual_type=method-flow-formula|width_budget=0.98linewidth|height_budget=0.86textheight|min_font_pt=10|split_allowed=no|expected_page=107|risks=title-table-source-footline
\begin{frame}[t]{准PDF完整非微扰重正化方案：理论结构与方法链}
  \fontsize{10}{10.8}\selectfont
\begingroup
\renewcommand{\arraystretch}{0.84}
\begin{tabularx}{0.98\linewidth}{@{}p{2.55cm}>{\raggedright\arraybackslash}X@{}}
\toprule
输入与状态 & 格点/解析输入 → 非微扰重正化、OPE 与匹配中的算符、矩阵元或场配置；具体输入见源章节。 \\
章节路线 & 引言；理论设定；核子矩阵元；重正化方案；转换到（公式）方案；结果；RI（公式）方案及其到（公式）方案的转换；系统性不确定性的评估；重正化的结果；与光前 PDF 的匹配；结论与讨论 \\
方法摘录 & 本节简要介绍我们想要重正化的准 PDF 的核子矩阵元，并解释三类 PDF——非极化、螺旋度和横向性—— 的重正化处方。核子矩阵元计算的细节见文献 Refs.（引文）。 * 0.25cm \\
结构式/示意 & $O^{R}(z,\mu)=Z(z,\mu,a)\,O^{\rm bare}(z,a)$\\[-0.2em]\scriptsize 该式只表达主题变量关系；论文特定推导以源文件为准。 \\
物理校验 & 重正化因子抵消截止依赖；a→0 和短距离 z→0 极限必须与匹配阶数、方案和尺度保持一致。 \\
\bottomrule
\end{tabularx}
\endgroup
  \source{方法证据：\detokenize{refer/papers/准PDF完整非微扰重正化方案_latex/chapters/section02.tex:1-8}；主题结构式为示意。}
\end{frame}

%% frame_manifest|frame_id=P24-03|paper_id=P24|claim_id=P24-结果|visual_type=result-evidence-table|width_budget=0.98linewidth|height_budget=0.86textheight|min_font_pt=10|split_allowed=no|expected_page=108|risks=title-table-source-footline
\begin{frame}[t]{准PDF完整非微扰重正化方案：结果、验证与物理含义}
  \fontsize{10}{10.8}\selectfont
\begingroup
\renewcommand{\arraystretch}{0.84}
\begin{tabularx}{0.98\linewidth}{@{}p{2.55cm}>{\raggedright\arraybackslash}X@{}}
\toprule
源文结果 & 在本工作中，我们给出了一个具体的处方，用于非微扰地重正化准 PDF 计算所需的矩阵元。 所采用的方案是 RI（公式），随后转换到（公式）方案并演化到 2 GeV；这一步以单圈微扰完成。 我们已经论证，该重正化条件恰当地处理了矩阵元中存在的两类发散：标准的对数发散， 以及含 Wilson 线的非局域算符特有的幂发散。此外，我们给出了消除混合的重正化条件， 适用于与非极化准 PDF……通过这项工作，我们提出并讨论了准 PDF 的完整重正化程序——在此之前它一直是一个重大的不确定性来源—— 并构成了连接格点 QCD 结果与光锥 PDF 的关键里程碑 \\
验证状态 & 确证：源章节有结果/讨论摘录。 \\
库内可见证据 & 中文源约 1242 行；公式命中 8，图命中 28，表/表格命中 4。 \\
物理含义 & 在本主题共同链条中，本论文把上述输入推进到可比较的中间量或观测量；具体强度、误差和外推范围不超出源文档。 \\
未验证项 & 当前工作区没有论文 PDF、原始数据和独立复现脚本；因此数值复核、图像再现和统计显著性不作额外断言。 \\
\bottomrule
\end{tabularx}
\endgroup
  \source{结果证据：\detokenize{refer/papers/准PDF完整非微扰重正化方案_latex/chapters/section04.tex:1-31}；图表/公式计数由本报告生成器对中文源的静态统计得到。}
\end{frame}

%% frame_manifest|frame_id=P24-04|paper_id=P24|claim_id=P24-局限|visual_type=risk-relation-table|width_budget=0.98linewidth|height_budget=0.86textheight|min_font_pt=10|split_allowed=no|expected_page=109|risks=title-table-source-footline
\begin{frame}[t]{准PDF完整非微扰重正化方案：局限、跨论文接口与可复查入口}
  \fontsize{10}{10.8}\selectfont
\begingroup
\renewcommand{\arraystretch}{0.84}
\begin{tabularx}{0.98\linewidth}{@{}p{2.55cm}>{\raggedright\arraybackslash}X@{}}
\toprule
源文局限 & 需要回到结果/讨论/展望章节核对适用区间；本档案不把缺失的独立数据复核补成已证实结论。 \\
跨论文接口 & 是准/赝分布从格点数据走向 MS-bar 或其他连续方案的必要中间层。 \\
极限检查 & 重正化因子抵消截止依赖；a→0 和短距离 z→0 极限必须与匹配阶数、方案和尺度保持一致。 \\
下一步核验 & 补齐 PDF 或原始数据；按源文参数复现关键图表；比较不同动量、流时间、距离、格距或方案转换后的稳定性。 \\
完整来源 & \detokenize{refer/papers/准PDF完整非微扰重正化方案_latex/main.tex}；\detokenize{refer/papers/Complete_NP_Renorm_QuasiPDF_latex/main.tex}；章节源：\detokenize{refer/papers/准PDF完整非微扰重正化方案_latex/chapters/*.tex}；英中目录：\detokenize{refer/papers/Complete_NP_Renorm_QuasiPDF_latex}. \\
\bottomrule
\end{tabularx}
\endgroup
  \source{状态：\inferred 为主题接口推断；\unverified 为当前材料边界；完整文件清单见 manifest。}
\end{frame}

%% frame_manifest|frame_id=P25-01|paper_id=P25|claim_id=P25-定位|visual_type=paper-card|width_budget=0.98linewidth|height_budget=0.86textheight|min_font_pt=10|split_allowed=no|expected_page=110|risks=title-table-source-footline
\section{准部分子算符乘法可重正性}
\begin{frame}[t]{准部分子算符乘法可重正性：研究问题与论文定位}
  \fontsize{10}{10.8}\selectfont
\begingroup
\renewcommand{\arraystretch}{0.84}
\begin{tabularx}{0.98\linewidth}{@{}p{2.55cm}>{\raggedright\arraybackslash}X@{}}
\toprule
论文编号 & P25；主题：非微扰重正化、OPE 与匹配 \\
书目信息 & Multiplicative Renorm Quasi Operators；arXiv:1809.01836；英/中转排页数 11/9 \\
研究对象 & 从格点 QCD 计算的准部分子分布（quasi-PDF）中抽取部分子分布函数（PDF），近年来一直受到积极追求。我们把文献（引文）中关于准夸克算符乘法可重正性的证明推广到准胶子算符，并证明准胶子算符在微扰论所有阶都可以被乘法重正化……因子化，从格点 QCD 计算的准部分子算符强子矩阵元中抽取 PDF 就可以有一个坚实的理论基础 \\
章节证据 & 引言；单圈紫外发散；高阶紫外发散；胶子—规范链顶点的重正化；胶子—胶子—规范链顶点的重正化；总结 \\
证据状态 & \verified：摘要和章节来自现有 LaTeX；\unverified：没有论文 PDF/原始数据独立复核。 \\
\bottomrule
\end{tabularx}
\endgroup
  \source{索引：\detokenize{refer/papers/INDEX.md:25}；正文：\detokenize{refer/papers/准部分子算符乘法可重正性_latex/main.tex}；\detokenize{refer/papers/Multiplicative_Renorm_Quasi_Operators_latex/main.tex}}
\end{frame}

%% frame_manifest|frame_id=P25-02|paper_id=P25|claim_id=P25-方法|visual_type=method-flow-formula|width_budget=0.98linewidth|height_budget=0.86textheight|min_font_pt=10|split_allowed=no|expected_page=111|risks=title-table-source-footline
\begin{frame}[t]{准部分子算符乘法可重正性：理论结构与方法链}
  \fontsize{10}{10.8}\selectfont
\begingroup
\renewcommand{\arraystretch}{0.84}
\begin{tabularx}{0.98\linewidth}{@{}p{2.55cm}>{\raggedright\arraybackslash}X@{}}
\toprule
输入与状态 & 格点/解析输入 → 非微扰重正化、OPE 与匹配中的算符、矩阵元或场配置；具体输入见源章节。 \\
章节路线 & 引言；单圈紫外发散；高阶紫外发散；胶子—规范链顶点的重正化；胶子—胶子—规范链顶点的重正化；总结 \\
方法摘录 & sec:intro 部分子分布函数（PDF）编码了强相互作用重要的非微扰信息。基于 QCD 因子化（引文），人们已经成功地以很好的精度从高能散射数据中抽取了 PDF（引文）。然而，无论从理论还是实用的角度来看，从第一性原理的格点 QCD（LQC……我们在文献（引文）中关于准夸克算符的证明，本文针对准胶子算符的工作完成了准部分子算符乘法重正化的证明 \\
结构式/示意 & $O^{R}(z,\mu)=Z(z,\mu,a)\,O^{\rm bare}(z,a)$\\[-0.2em]\scriptsize 该式只表达主题变量关系；论文特定推导以源文件为准。 \\
物理校验 & 重正化因子抵消截止依赖；a→0 和短距离 z→0 极限必须与匹配阶数、方案和尺度保持一致。 \\
\bottomrule
\end{tabularx}
\endgroup
  \source{方法证据：\detokenize{refer/papers/准部分子算符乘法可重正性_latex/chapters/section01.tex:1-29}；主题结构式为示意。}
\end{frame}

%% frame_manifest|frame_id=P25-03|paper_id=P25|claim_id=P25-结果|visual_type=result-evidence-table|width_budget=0.98linewidth|height_budget=0.86textheight|min_font_pt=10|split_allowed=no|expected_page=112|risks=title-table-source-footline
\begin{frame}[t]{准部分子算符乘法可重正性：结果、验证与物理含义}
  \fontsize{10}{10.8}\selectfont
\begingroup
\renewcommand{\arraystretch}{0.84}
\begin{tabularx}{0.98\linewidth}{@{}p{2.55cm}>{\raggedright\arraybackslash}X@{}}
\toprule
源文结果 & 我们证明了：只要所有发散都在规范不变的方案中正规化，准胶子算符的紫外发散实际上与准夸克算符的紫外发散相似。我们证明了全部 36 个纯准胶子算符的紫外发散都定域于时空之中，并且可以被乘法重正化而不相互混合，（公式或图表位置）其中（公式）是为洛伦兹指标（公式）选出的（公式）分量个数，（公式）、（公式）、（公式）与（公式）为重正化常数。 与准夸克算符一样，由于不相互混合，准胶子算符的……的例子：它们可在 LQCD 中计算并因子化为普通 PDF；通过对第一性原理 LQCD 计算产生的这些良好“格点截面”数据进行 QCD 全局分析，即可从中抽取 PDF \\
验证状态 & 确证：源章节有结果/讨论摘录。 \\
库内可见证据 & 中文源约 631 行；公式命中 9，图命中 8，表/表格命中 0。 \\
物理含义 & 在本主题共同链条中，本论文把上述输入推进到可比较的中间量或观测量；具体强度、误差和外推范围不超出源文档。 \\
未验证项 & 当前工作区没有论文 PDF、原始数据和独立复现脚本；因此数值复核、图像再现和统计显著性不作额外断言。 \\
\bottomrule
\end{tabularx}
\endgroup
  \source{结果证据：\detokenize{refer/papers/准部分子算符乘法可重正性_latex/chapters/section06.tex:1-8}；图表/公式计数由本报告生成器对中文源的静态统计得到。}
\end{frame}

%% frame_manifest|frame_id=P25-04|paper_id=P25|claim_id=P25-局限|visual_type=risk-relation-table|width_budget=0.98linewidth|height_budget=0.86textheight|min_font_pt=10|split_allowed=no|expected_page=113|risks=title-table-source-footline
\begin{frame}[t]{准部分子算符乘法可重正性：局限、跨论文接口与可复查入口}
  \fontsize{10}{10.8}\selectfont
\begingroup
\renewcommand{\arraystretch}{0.84}
\begin{tabularx}{0.98\linewidth}{@{}p{2.55cm}>{\raggedright\arraybackslash}X@{}}
\toprule
源文局限 & 需要回到结果/讨论/展望章节核对适用区间；本档案不把缺失的独立数据复核补成已证实结论。 \\
跨论文接口 & 是准/赝分布从格点数据走向 MS-bar 或其他连续方案的必要中间层。 \\
极限检查 & 重正化因子抵消截止依赖；a→0 和短距离 z→0 极限必须与匹配阶数、方案和尺度保持一致。 \\
下一步核验 & 补齐 PDF 或原始数据；按源文参数复现关键图表；比较不同动量、流时间、距离、格距或方案转换后的稳定性。 \\
完整来源 & \detokenize{refer/papers/准部分子算符乘法可重正性_latex/main.tex}；\detokenize{refer/papers/Multiplicative_Renorm_Quasi_Operators_latex/main.tex}；章节源：\detokenize{refer/papers/准部分子算符乘法可重正性_latex/chapters/*.tex}；英中目录：\detokenize{refer/papers/Multiplicative_Renorm_Quasi_Operators_latex}. \\
\bottomrule
\end{tabularx}
\endgroup
  \source{状态：\inferred 为主题接口推断；\unverified 为当前材料边界；完整文件清单见 manifest。}
\end{frame}

%% frame_manifest|frame_id=P26-01|paper_id=P26|claim_id=P26-定位|visual_type=paper-card|width_budget=0.98linewidth|height_budget=0.86textheight|min_font_pt=10|split_allowed=no|expected_page=114|risks=title-table-source-footline
\section{大动量有效理论获取胶子分布}
\begin{frame}[t]{大动量有效理论获取胶子分布：研究问题与论文定位}
  \fontsize{10}{10.8}\selectfont
\begingroup
\renewcommand{\arraystretch}{0.84}
\begin{tabularx}{0.98\linewidth}{@{}p{2.55cm}>{\raggedright\arraybackslash}X@{}}
\toprule
论文编号 & P26；主题：非微扰重正化、OPE 与匹配 \\
书目信息 & Accessing Gluon PDF LaMET；arXiv:1808.10824；英/中转排页数 8/7 \\
研究对象 & 利用大动量有效理论（large momentum effective theory, LaMET），质子内的胶子部分子分布函数（PDFs）可以直接在欧几里得格点上计算。为实现这一目标，必须找到这样的重正化胶子准PDF：其中的幂次发……导出。我们为这些重正化的准PDF给出了一个 LaMET 因子化公式，由此可以提取胶子 PDFs \\
章节证据 & 引言；辅助伴随表示“重夸克”方法；辅助场方法中的重正化；胶子准PDF算符；胶子螺旋度分布；在格点上的实现；结论 \\
证据状态 & \verified：摘要和章节来自现有 LaTeX；\unverified：没有论文 PDF/原始数据独立复核。 \\
\bottomrule
\end{tabularx}
\endgroup
  \source{索引：\detokenize{refer/papers/INDEX.md:26}；正文：\detokenize{refer/papers/大动量有效理论获取胶子分布_latex/main.tex}；\detokenize{refer/papers/Accessing_Gluon_PDF_LaMET_latex/main.tex}}
\end{frame}

%% frame_manifest|frame_id=P26-02|paper_id=P26|claim_id=P26-方法|visual_type=method-flow-formula|width_budget=0.98linewidth|height_budget=0.86textheight|min_font_pt=10|split_allowed=no|expected_page=115|risks=title-table-source-footline
\begin{frame}[t]{大动量有效理论获取胶子分布：理论结构与方法链}
  \fontsize{10}{10.8}\selectfont
\begingroup
\renewcommand{\arraystretch}{0.84}
\begin{tabularx}{0.98\linewidth}{@{}p{2.55cm}>{\raggedright\arraybackslash}X@{}}
\toprule
输入与状态 & 格点/解析输入 → 非微扰重正化、OPE 与匹配中的算符、矩阵元或场配置；具体输入见源章节。 \\
章节路线 & 引言；辅助伴随表示“重夸克”方法；辅助场方法中的重正化；胶子准PDF算符；胶子螺旋度分布；在格点上的实现；结论 \\
方法摘录 & sec:aux 让我们从文献（引文）定义的胶子准PDF出发：（公式或图表位置）[公式] 是场强张量。除非另有说明，本节定义的量都是裸量。指标（公式）既可像胶子 PDF 中那样（引文）只对横向分量求和，也可对所有洛伦兹分量求和。由于（公式），而（……用）) 中非定域算符的重正化便归结为式 (（交叉引用）) 中两个定域复合算符（LCOs）乘积的重正化 \\
结构式/示意 & $O^{R}(z,\mu)=Z(z,\mu,a)\,O^{\rm bare}(z,a)$\\[-0.2em]\scriptsize 该式只表达主题变量关系；论文特定推导以源文件为准。 \\
物理校验 & 重正化因子抵消截止依赖；a→0 和短距离 z→0 极限必须与匹配阶数、方案和尺度保持一致。 \\
\bottomrule
\end{tabularx}
\endgroup
  \source{方法证据：\detokenize{refer/papers/大动量有效理论获取胶子分布_latex/chapters/section02.tex:1-30}；主题结构式为示意。}
\end{frame}

%% frame_manifest|frame_id=P26-03|paper_id=P26|claim_id=P26-结果|visual_type=result-evidence-table|width_budget=0.98linewidth|height_budget=0.86textheight|min_font_pt=10|split_allowed=no|expected_page=116|risks=title-table-source-footline
\begin{frame}[t]{大动量有效理论获取胶子分布：结果、验证与物理含义}
  \fontsize{10}{10.8}\selectfont
\begingroup
\renewcommand{\arraystretch}{0.84}
\begin{tabularx}{0.98\linewidth}{@{}p{2.55cm}>{\raggedright\arraybackslash}X@{}}
\toprule
源文结果 & sec:conclusion 总之，我们利用辅助伴随表示“重夸克”拉氏量，对 LaMET 中定义的胶子准PDF的重正化性质进行了系统研究。我们证明了胶子准 PDF 中的所有幂次发散都来自威尔逊线自能，并可通过一次质量重正化予以消除。我们还确定了一组适用于格点模拟的可乘法重正化胶子准PDF算符。我们的发现为从格点模拟直接抽取胶子 PDFs 提供了理论基础 \\
验证状态 & 确证：源章节有结果/讨论摘录。 \\
库内可见证据 & 中文源约 648 行；公式命中 12，图命中 0，表/表格命中 0。 \\
物理含义 & 在本主题共同链条中，本论文把上述输入推进到可比较的中间量或观测量；具体强度、误差和外推范围不超出源文档。 \\
未验证项 & 当前工作区没有论文 PDF、原始数据和独立复现脚本；因此数值复核、图像再现和统计显著性不作额外断言。 \\
\bottomrule
\end{tabularx}
\endgroup
  \source{结果证据：\detokenize{refer/papers/大动量有效理论获取胶子分布_latex/chapters/section07.tex:1-4}；图表/公式计数由本报告生成器对中文源的静态统计得到。}
\end{frame}

%% frame_manifest|frame_id=P26-04|paper_id=P26|claim_id=P26-局限|visual_type=risk-relation-table|width_budget=0.98linewidth|height_budget=0.86textheight|min_font_pt=10|split_allowed=no|expected_page=117|risks=title-table-source-footline
\begin{frame}[t]{大动量有效理论获取胶子分布：局限、跨论文接口与可复查入口}
  \fontsize{10}{10.8}\selectfont
\begingroup
\renewcommand{\arraystretch}{0.84}
\begin{tabularx}{0.98\linewidth}{@{}p{2.55cm}>{\raggedright\arraybackslash}X@{}}
\toprule
源文局限 & 需要回到结果/讨论/展望章节核对适用区间；本档案不把缺失的独立数据复核补成已证实结论。 \\
跨论文接口 & 是准/赝分布从格点数据走向 MS-bar 或其他连续方案的必要中间层。 \\
极限检查 & 重正化因子抵消截止依赖；a→0 和短距离 z→0 极限必须与匹配阶数、方案和尺度保持一致。 \\
下一步核验 & 补齐 PDF 或原始数据；按源文参数复现关键图表；比较不同动量、流时间、距离、格距或方案转换后的稳定性。 \\
完整来源 & \detokenize{refer/papers/大动量有效理论获取胶子分布_latex/main.tex}；\detokenize{refer/papers/Accessing_Gluon_PDF_LaMET_latex/main.tex}；章节源：\detokenize{refer/papers/大动量有效理论获取胶子分布_latex/chapters/*.tex}；英中目录：\detokenize{refer/papers/Accessing_Gluon_PDF_LaMET_latex}. \\
\bottomrule
\end{tabularx}
\endgroup
  \source{状态：\inferred 为主题接口推断；\unverified 为当前材料边界；完整文件清单见 manifest。}
\end{frame}

%% frame_manifest|frame_id=P27-01|paper_id=P27|claim_id=P27-定位|visual_type=paper-card|width_budget=0.98linewidth|height_budget=0.86textheight|min_font_pt=10|split_allowed=no|expected_page=118|risks=title-table-source-footline
\section{胶子伪分布的短距行为}
\begin{frame}[t]{胶子伪分布的短距行为：研究问题与论文定位}
  \fontsize{10}{10.8}\selectfont
\begingroup
\renewcommand{\arraystretch}{0.84}
\begin{tabularx}{0.98\linewidth}{@{}p{2.55cm}>{\raggedright\arraybackslash}X@{}}
\toprule
论文编号 & P27；主题：非微扰重正化、OPE 与匹配 \\
书目信息 & Gluon Pseudo-Distributions at Short Distances: Forward Case Phys. Lett. B 808 (2020) 135621; arXiv:1910.13963 [hep-lat]；arXiv:1910.13963；英/中转排页数 16/14 \\
研究对象 & 本文给出在伪分布（pseudo-PDF）方法框架内开展胶子部分子分布函数（PDF）格点计算所需的 结果。 我们对可能的双胶子关联函数进行了分类，并指出其中包含在光锥极限（公式）下决定胶子 PDF 的不变振幅的那些函数。 单圈计算在……直接联系。 我们还证明，这些结果可用于相当直接地计算准分布（quasi-PDF）的单圈匹配关系 \\
章节证据 & 引言；矩阵元；链接自能贡献与紫外发散；顶点贡献；紫外发散项；演化项；盒图；胶子自能图；DGLAP 演化结构；纯胶动力学中 DGLAP 何时为对角；约化 Ioffe 时间分布；胶子–夸克混合；匹配关系；准分布的匹配关系；总结 \\
证据状态 & \verified：摘要和章节来自现有 LaTeX；\unverified：没有论文 PDF/原始数据独立复核。 \\
\bottomrule
\end{tabularx}
\endgroup
  \source{索引：\detokenize{refer/papers/INDEX.md:27}；正文：\detokenize{refer/papers/胶子伪分布的短距行为_latex/main.tex}；\detokenize{refer/papers/Gluon_Pseudo_Short_Distances_latex/main.tex}}
\end{frame}

%% frame_manifest|frame_id=P27-02|paper_id=P27|claim_id=P27-方法|visual_type=method-flow-formula|width_budget=0.98linewidth|height_budget=0.86textheight|min_font_pt=10|split_allowed=no|expected_page=119|risks=title-table-source-footline
\begin{frame}[t]{胶子伪分布的短距行为：理论结构与方法链}
  \fontsize{10}{10.8}\selectfont
\begingroup
\renewcommand{\arraystretch}{0.84}
\begin{tabularx}{0.98\linewidth}{@{}p{2.55cm}>{\raggedright\arraybackslash}X@{}}
\toprule
输入与状态 & 格点/解析输入 → 非微扰重正化、OPE 与匹配中的算符、矩阵元或场配置；具体输入见源章节。 \\
章节路线 & 引言；矩阵元；链接自能贡献与紫外发散；顶点贡献；紫外发散项；演化项；盒图；胶子自能图；DGLAP 演化结构；纯胶动力学中 DGLAP 何时为对角；约化 Ioffe 时间分布；胶子–夸克混合；匹配关系；准分布的匹配关系；总结 \\
方法摘录 & 由两个胶子场组成的算符（四个指标均不缩并）的核子自旋平均矩阵元由（公式或图表位置）给出，其中（公式）是胶子（伴随）表示中的标准直线规范链接（公式或图表位置）对不变振幅作分解的张量结构可由两个可用四矢量（公式）、（公式）以及度规张量（公式）构成。……在单圈具有对角的 DGLAP 演化。 为回答这些问题，我们计算了 单圈 胶子交换对原双定域算符的修改 \\
结构式/示意 & $O^{R}(z,\mu)=Z(z,\mu,a)\,O^{\rm bare}(z,a)$\\[-0.2em]\scriptsize 该式只表达主题变量关系；论文特定推导以源文件为准。 \\
物理校验 & 重正化因子抵消截止依赖；a→0 和短距离 z→0 极限必须与匹配阶数、方案和尺度保持一致。 \\
\bottomrule
\end{tabularx}
\endgroup
  \source{方法证据：\detokenize{refer/papers/胶子伪分布的短距行为_latex/chapters/section02.tex:1-145}；主题结构式为示意。}
\end{frame}

%% frame_manifest|frame_id=P27-03|paper_id=P27|claim_id=P27-结果|visual_type=result-evidence-table|width_budget=0.98linewidth|height_budget=0.86textheight|min_font_pt=10|split_allowed=no|expected_page=120|risks=title-table-source-footline
\begin{frame}[t]{胶子伪分布的短距行为：结果、验证与物理含义}
  \fontsize{10}{10.8}\selectfont
\begingroup
\renewcommand{\arraystretch}{0.84}
\begin{tabularx}{0.98\linewidth}{@{}p{2.55cm}>{\raggedright\arraybackslash}X@{}}
\toprule
源文结果 & 在本文中，我们给出了为正在开展的、采用伪 PDF 方法计算胶子 PDF 的工作 奠定基础的结果。 特别地，我们对可能的双胶子关联函数作了分类。 我们指出了其中包含不变振幅（公式）的那些函数， 该振幅在光锥极限（公式）下决定胶子 PDF。 由于此极限是奇异的，需要联系（公式）与光锥 PDF（公式）的 匹配条件。 为此，我们采用 Ref.（引文）的方法， 对两个胶子场强张量的规范不……物中，我们计划给出计算的更多细节， 并给出盒图的完整结果， 特别是分布振幅与 GPD 格点计算所需的 非前向运动学。 我们还计划纳入胶子–夸克 与夸克–胶子项的计算 \\
验证状态 & 确证：源章节有结果/讨论摘录。 \\
库内可见证据 & 中文源约 1252 行；公式命中 55，图命中 11，表/表格命中 0。 \\
物理含义 & 在本主题共同链条中，本论文把上述输入推进到可比较的中间量或观测量；具体强度、误差和外推范围不超出源文档。 \\
未验证项 & 当前工作区没有论文 PDF、原始数据和独立复现脚本；因此数值复核、图像再现和统计显著性不作额外断言。 \\
\bottomrule
\end{tabularx}
\endgroup
  \source{结果证据：\detokenize{refer/papers/胶子伪分布的短距行为_latex/chapters/section08.tex:1-42}；图表/公式计数由本报告生成器对中文源的静态统计得到。}
\end{frame}

%% frame_manifest|frame_id=P27-04|paper_id=P27|claim_id=P27-局限|visual_type=risk-relation-table|width_budget=0.98linewidth|height_budget=0.86textheight|min_font_pt=10|split_allowed=no|expected_page=121|risks=title-table-source-footline
\begin{frame}[t]{胶子伪分布的短距行为：局限、跨论文接口与可复查入口}
  \fontsize{10}{10.8}\selectfont
\begingroup
\renewcommand{\arraystretch}{0.84}
\begin{tabularx}{0.98\linewidth}{@{}p{2.55cm}>{\raggedright\arraybackslash}X@{}}
\toprule
源文局限 & 需要回到结果/讨论/展望章节核对适用区间；本档案不把缺失的独立数据复核补成已证实结论。 \\
跨论文接口 & 是准/赝分布从格点数据走向 MS-bar 或其他连续方案的必要中间层。 \\
极限检查 & 重正化因子抵消截止依赖；a→0 和短距离 z→0 极限必须与匹配阶数、方案和尺度保持一致。 \\
下一步核验 & 补齐 PDF 或原始数据；按源文参数复现关键图表；比较不同动量、流时间、距离、格距或方案转换后的稳定性。 \\
完整来源 & \detokenize{refer/papers/胶子伪分布的短距行为_latex/main.tex}；\detokenize{refer/papers/Gluon_Pseudo_Short_Distances_latex/main.tex}；章节源：\detokenize{refer/papers/胶子伪分布的短距行为_latex/chapters/*.tex}；英中目录：\detokenize{refer/papers/Gluon_Pseudo_Short_Distances_latex}. \\
\bottomrule
\end{tabularx}
\endgroup
  \source{状态：\inferred 为主题接口推断；\unverified 为当前材料边界；完整文件清单见 manifest。}
\end{frame}

%% frame_manifest|frame_id=P28-01|paper_id=P28|claim_id=P28-定位|visual_type=paper-card|width_budget=0.98linewidth|height_budget=0.86textheight|min_font_pt=10|split_allowed=no|expected_page=122|risks=title-table-source-footline
\section{准部分子分布与梯度流}
\begin{frame}[t]{准部分子分布与梯度流：研究问题与论文定位}
  \fontsize{10}{10.8}\selectfont
\begingroup
\renewcommand{\arraystretch}{0.84}
\begin{tabularx}{0.98\linewidth}{@{}p{2.55cm}>{\raggedright\arraybackslash}X@{}}
\toprule
论文编号 & P28；主题：梯度流、涂抹与能动量张量 \\
书目信息 & Quasi parton distributions and the gradient flow；arXiv:1612.01584；英/中转排页数 13/11 \\
研究对象 & 我们提出一种从格点量子色动力学确定准部分子分布函数（PDFs）的新方法。 通过引入梯度流，该方法保证格点 准 PDFs 在连续极限下有限，并回避了格点上准 PDFs 重正化这一棘手且迄今尚未解决的问题。在流时间远小于核子动量所设定……F 的矩成正比。 我们利用这一关系推导出联系涂抹准 PDF 与光前 PDF 的匹配核的演化方程 \\
章节证据 & Introduction；sec:defs 分布函数；裸 PDFs；重正化的 PDFs；涂抹准 PDFs；sec:spdf 与光前分布的关系；sec:renormpdf 短距离展开；sec:match 匹配核的类 DGLAP 方程；sec:concs 结论 \\
证据状态 & \verified：摘要和章节来自现有 LaTeX；\unverified：没有论文 PDF/原始数据独立复核。 \\
\bottomrule
\end{tabularx}
\endgroup
  \source{索引：\detokenize{refer/papers/INDEX.md:28}；正文：\detokenize{refer/papers/准部分子分布与梯度流_latex/main.tex}；\detokenize{refer/papers/Quasi_PDF_Gradient_Flow_latex/main.tex}}
\end{frame}

%% frame_manifest|frame_id=P28-02|paper_id=P28|claim_id=P28-方法|visual_type=method-flow-formula|width_budget=0.98linewidth|height_budget=0.86textheight|min_font_pt=10|split_allowed=no|expected_page=123|risks=title-table-source-footline
\begin{frame}[t]{准部分子分布与梯度流：理论结构与方法链}
  \fontsize{10}{10.8}\selectfont
\begingroup
\renewcommand{\arraystretch}{0.84}
\begin{tabularx}{0.98\linewidth}{@{}p{2.55cm}>{\raggedright\arraybackslash}X@{}}
\toprule
输入与状态 & 格点/解析输入 → 梯度流、涂抹与能动量张量中的算符、矩阵元或场配置；具体输入见源章节。 \\
章节路线 & Introduction；sec:defs 分布函数；裸 PDFs；重正化的 PDFs；涂抹准 PDFs；sec:spdf 与光前分布的关系；sec:renormpdf 短距离展开；sec:match 匹配核的类 DGLAP 方程；sec:concs 结论 \\
方法摘录 & 量子色动力学（QCD）是把强子束缚态与其部分子组分——夸克和胶子——联系起来的强相互作用理论。尽管实验上无法直接探测夸克与胶子， 但强子与部分子之间的联系可以通过部分子 分布函数（PDFs）来刻画。 PDFs 描绘了强子结构中与组分夸克、胶子的……Fs 之间的关系，并在第（交叉引用）节推导匹配核的演化方程。最后我们在第（交叉引用）节给出总结与结论 \\
结构式/示意 & $\partial_t B_\mu(t,x)=D_\nu G_{\nu\mu}(t,x),\qquad t\sim L_{\rm sm}^{2}$\\[-0.2em]\scriptsize 该式只表达主题变量关系；论文特定推导以源文件为准。 \\
物理校验 & 流时间与平滑长度满足平方关系；t→0 回到原始场，有限 t 则抑制紫外模式。 \\
\bottomrule
\end{tabularx}
\endgroup
  \source{方法证据：\detokenize{refer/papers/准部分子分布与梯度流_latex/chapters/section01.tex:1-52}；主题结构式为示意。}
\end{frame}

%% frame_manifest|frame_id=P28-03|paper_id=P28|claim_id=P28-结果|visual_type=result-evidence-table|width_budget=0.98linewidth|height_budget=0.86textheight|min_font_pt=10|split_allowed=no|expected_page=124|risks=title-table-source-footline
\begin{frame}[t]{准部分子分布与梯度流：结果、验证与物理含义}
  \fontsize{10}{10.8}\selectfont
\begingroup
\renewcommand{\arraystretch}{0.84}
\begin{tabularx}{0.98\linewidth}{@{}p{2.55cm}>{\raggedright\arraybackslash}X@{}}
\toprule
源文结果 & 部分子分布函数（PDFs）借核子的组分夸克与胶子刻画核子结构。这些 PDFs 定义为光前波函数的矩阵元，无法直接在欧几里得格点 QCD 中计算。然而原则上，PDFs 的 Mellin 矩可以通过 扭度2算符的矩阵元在格点 QCD 中计算。 遗憾的是，这类计算仅限于头几个矩，因为 格点调节子的超立方对称性在不同质量量纲的扭度2算符之间诱导幂发散混合，模糊了格点上确定的矩阵元的连续……下的相应光前 PDFs。结合非微扰步标度（step-scaling）程序， 该匹配可在足够高的能量处进行， 使微扰截断误差不再失控。我们方案的非微扰实现正在进行之中 \\
验证状态 & 确证：源章节有结果/讨论摘录。 \\
库内可见证据 & 中文源约 987 行；公式命中 39，图命中 0，表/表格命中 0。 \\
物理含义 & 在本主题共同链条中，本论文把上述输入推进到可比较的中间量或观测量；具体强度、误差和外推范围不超出源文档。 \\
未验证项 & 当前工作区没有论文 PDF、原始数据和独立复现脚本；因此数值复核、图像再现和统计显著性不作额外断言。 \\
\bottomrule
\end{tabularx}
\endgroup
  \source{结果证据：\detokenize{refer/papers/准部分子分布与梯度流_latex/chapters/section05.tex:1-36}；图表/公式计数由本报告生成器对中文源的静态统计得到。}
\end{frame}

%% frame_manifest|frame_id=P28-04|paper_id=P28|claim_id=P28-局限|visual_type=risk-relation-table|width_budget=0.98linewidth|height_budget=0.86textheight|min_font_pt=10|split_allowed=no|expected_page=125|risks=title-table-source-footline
\begin{frame}[t]{准部分子分布与梯度流：局限、跨论文接口与可复查入口}
  \fontsize{10}{10.8}\selectfont
\begingroup
\renewcommand{\arraystretch}{0.84}
\begin{tabularx}{0.98\linewidth}{@{}p{2.55cm}>{\raggedright\arraybackslash}X@{}}
\toprule
源文局限 & 需要回到结果/讨论/展望章节核对适用区间；本档案不把缺失的独立数据复核补成已证实结论。 \\
跨论文接口 & 把连续极限、微扰有限性和能动量张量重构连接到可计算的格点观测量。 \\
极限检查 & 流时间与平滑长度满足平方关系；t→0 回到原始场，有限 t 则抑制紫外模式。 \\
下一步核验 & 补齐 PDF 或原始数据；按源文参数复现关键图表；比较不同动量、流时间、距离、格距或方案转换后的稳定性。 \\
完整来源 & \detokenize{refer/papers/准部分子分布与梯度流_latex/main.tex}；\detokenize{refer/papers/Quasi_PDF_Gradient_Flow_latex/main.tex}；章节源：\detokenize{refer/papers/准部分子分布与梯度流_latex/chapters/*.tex}；英中目录：\detokenize{refer/papers/Quasi_PDF_Gradient_Flow_latex}. \\
\bottomrule
\end{tabularx}
\endgroup
  \source{状态：\inferred 为主题接口推断；\unverified 为当前材料边界；完整文件清单见 manifest。}
\end{frame}

%% frame_manifest|frame_id=P29-01|paper_id=P29|claim_id=P29-定位|visual_type=paper-card|width_budget=0.98linewidth|height_budget=0.86textheight|min_font_pt=10|split_allowed=no|expected_page=126|risks=title-table-source-footline
\section{准光前关联函数的自重正化}
\begin{frame}[t]{准光前关联函数的自重正化：研究问题与论文定位}
  \fontsize{10}{10.8}\selectfont
\begingroup
\renewcommand{\arraystretch}{0.84}
\begin{tabularx}{0.98\linewidth}{@{}p{2.55cm}>{\raggedright\arraybackslash}X@{}}
\toprule
论文编号 & P29；主题：非微扰重正化、OPE 与匹配 \\
书目信息 & 1ex] Lattice Parton Collaboration (LPC)；arXiv:2103.02965；英/中转排页数 48/44 \\
研究对象 & 在大动量有效理论的应用中，由于 Wilson 线自能量中的线性发散，欧氏关联函数在格点正规化下的重正化是一个挑战。基于准 PDF 算符在格点间距（公式）= 0.03 fm（公式）0.12 fm、价夸克分别为 clover 与 ov……用的 RI/MOM 与比值重正化方案中存在大的非微扰效应，这倾向于支持近期提出的混合重正化流程 \\
章节证据 & 引言；微扰论中的重正化；格点矩阵元与模拟设置；线性发散的简单检验；格点矩阵元的自重正化；方法策略；重整子不确定性；调节参数 d 以匹配连续方案；自重正化流程；不同作用量与矩阵元结果的比较；其……DF 算符的真空矩阵元；Landau 规范固定的 Wilson 链接；总结 \\
证据状态 & \verified：摘要和章节来自现有 LaTeX；\unverified：没有论文 PDF/原始数据独立复核。 \\
\bottomrule
\end{tabularx}
\endgroup
  \source{索引：\detokenize{refer/papers/INDEX.md:29}；正文：\detokenize{refer/papers/准光前关联函数的自重正化_latex/main.tex}；\detokenize{refer/papers/Self_Renorm_Quasi_LightFront_latex/main.tex}}
\end{frame}

%% frame_manifest|frame_id=P29-02|paper_id=P29|claim_id=P29-方法|visual_type=method-flow-formula|width_budget=0.98linewidth|height_budget=0.86textheight|min_font_pt=10|split_allowed=no|expected_page=127|risks=title-table-source-footline
\begin{frame}[t]{准光前关联函数的自重正化：理论结构与方法链}
  \fontsize{10}{10.8}\selectfont
\begingroup
\renewcommand{\arraystretch}{0.84}
\begin{tabularx}{0.98\linewidth}{@{}p{2.55cm}>{\raggedright\arraybackslash}X@{}}
\toprule
输入与状态 & 格点/解析输入 → 非微扰重正化、OPE 与匹配中的算符、矩阵元或场配置；具体输入见源章节。 \\
章节路线 & 引言；微扰论中的重正化；格点矩阵元与模拟设置；线性发散的简单检验；格点矩阵元的自重正化；方法策略；重整子不确定性；调节参数 d 以匹配连续方案；自重正化流程；不同作用量与矩阵元结果的比较；其……DF 算符的真空矩阵元；Landau 规范固定的 Wilson 链接；总结 \\
方法摘录 & sec:setup 格点矩阵元标准的非微扰重正化方法是：计算该算符的一些辅助矩阵元（它们也可以在微扰论中计算），并用它们作为重正化因子来趋近连续极限。本文聚焦于研究可能对准光前关联的重正化有用的辅助矩阵元。我们主要考虑以下几种情形下准 PDF ……（公式）指数衰减，因而在（公式）时无法获得任何有意义的连续极限。必须进行重正化才能恢复良好的连续极限 \\
结构式/示意 & $O^{R}(z,\mu)=Z(z,\mu,a)\,O^{\rm bare}(z,a)$\\[-0.2em]\scriptsize 该式只表达主题变量关系；论文特定推导以源文件为准。 \\
物理校验 & 重正化因子抵消截止依赖；a→0 和短距离 z→0 极限必须与匹配阶数、方案和尺度保持一致。 \\
\bottomrule
\end{tabularx}
\endgroup
  \source{方法证据：\detokenize{refer/papers/准光前关联函数的自重正化_latex/chapters/section03.tex:1-113}；主题结构式为示意。}
\end{frame}

%% frame_manifest|frame_id=P29-03|paper_id=P29|claim_id=P29-结果|visual_type=result-evidence-table|width_budget=0.98linewidth|height_budget=0.86textheight|min_font_pt=10|split_allowed=no|expected_page=128|risks=title-table-source-footline
\begin{frame}[t]{准光前关联函数的自重正化：结果、验证与物理含义}
  \fontsize{10}{10.8}\selectfont
\begingroup
\renewcommand{\arraystretch}{0.84}
\begin{tabularx}{0.98\linewidth}{@{}p{2.55cm}>{\raggedright\arraybackslash}X@{}}
\toprule
源文结果 & sec:summary 在格点上计算的准光前关联的正确重正化，一直是把 LaMET 应用于部分子物理研究的主要挑战。这类关联含有与 Wilson 线相伴的线性发散，必须以高精度消除它们才能提取想要的物理信息。本文提出一种自重正化方法，从一个可在短距离处匹配到连续方案的准光前矩阵元中，同时消除线性发散、对数发散以及离散化误差。作为一个典型的范例，我们证明该方法对（公式）、核子和离……何算符矩阵元， 也可以推广到 LaMET 应用中含有 Wilson 线的任意矩阵元（引文）。当然，在这些情形中还需要仔细考虑矩阵元的物理解释以及与格点计算相关的细节 \\
验证状态 & 确证：源章节有结果/讨论摘录。 \\
库内可见证据 & 中文源约 1881 行；公式命中 15，图命中 83，表/表格命中 10。 \\
物理含义 & 在本主题共同链条中，本论文把上述输入推进到可比较的中间量或观测量；具体强度、误差和外推范围不超出源文档。 \\
未验证项 & 当前工作区没有论文 PDF、原始数据和独立复现脚本；因此数值复核、图像再现和统计显著性不作额外断言。 \\
\bottomrule
\end{tabularx}
\endgroup
  \source{结果证据：\detokenize{refer/papers/准光前关联函数的自重正化_latex/chapters/section09.tex:1-81}；图表/公式计数由本报告生成器对中文源的静态统计得到。}
\end{frame}

%% frame_manifest|frame_id=P29-04|paper_id=P29|claim_id=P29-局限|visual_type=risk-relation-table|width_budget=0.98linewidth|height_budget=0.86textheight|min_font_pt=10|split_allowed=no|expected_page=129|risks=title-table-source-footline
\begin{frame}[t]{准光前关联函数的自重正化：局限、跨论文接口与可复查入口}
  \fontsize{10}{10.8}\selectfont
\begingroup
\renewcommand{\arraystretch}{0.84}
\begin{tabularx}{0.98\linewidth}{@{}p{2.55cm}>{\raggedright\arraybackslash}X@{}}
\toprule
源文局限 & 需要回到结果/讨论/展望章节核对适用区间；本档案不把缺失的独立数据复核补成已证实结论。 \\
跨论文接口 & 是准/赝分布从格点数据走向 MS-bar 或其他连续方案的必要中间层。 \\
极限检查 & 重正化因子抵消截止依赖；a→0 和短距离 z→0 极限必须与匹配阶数、方案和尺度保持一致。 \\
下一步核验 & 补齐 PDF 或原始数据；按源文参数复现关键图表；比较不同动量、流时间、距离、格距或方案转换后的稳定性。 \\
完整来源 & \detokenize{refer/papers/准光前关联函数的自重正化_latex/main.tex}；\detokenize{refer/papers/Self_Renorm_Quasi_LightFront_latex/main.tex}；章节源：\detokenize{refer/papers/准光前关联函数的自重正化_latex/chapters/*.tex}；英中目录：\detokenize{refer/papers/Self_Renorm_Quasi_LightFront_latex}. \\
\bottomrule
\end{tabularx}
\endgroup
  \source{状态：\inferred 为主题接口推断；\unverified 为当前材料边界；完整文件清单见 manifest。}
\end{frame}

%% frame_manifest|frame_id=P30-01|paper_id=P30|claim_id=P30-定位|visual_type=paper-card|width_budget=0.98linewidth|height_budget=0.86textheight|min_font_pt=10|split_allowed=no|expected_page=130|risks=title-table-source-footline
\section{准光前关联的混合重正化方案}
\begin{frame}[t]{准光前关联的混合重正化方案：研究问题与论文定位}
  \fontsize{10}{10.8}\selectfont
\begingroup
\renewcommand{\arraystretch}{0.84}
\begin{tabularx}{0.98\linewidth}{@{}p{2.55cm}>{\raggedright\arraybackslash}X@{}}
\toprule
论文编号 & P30；主题：非微扰重正化、OPE 与匹配 \\
书目信息 & A Hybrid Renormalization Scheme for Quasi Light-Front Correlations in Large-Momentum Effective Theory Nucl. ……8.03886; doi:10.1016/j.nuclphysb.2021.115311；arXiv:2008.03886；英/中转排页数 28/24 \\
研究对象 & 在大动量有效理论（LaMET）中，部分子物理的计算从格点 QCD 中计算动量为（公式）的强子内坐标空间（公式）方向关联函数（公式）出发。这类关联函数同时 含有对格距（公式）的线性发散与对数发散，因而必须加以恰当的重正化。我们引入一……们建议分别 利用坐标空间关联在中等及较大（公式）下的普遍性质，把格点数据外推到渐近（公式）区域 \\
章节证据 & 引言；无穷大动量态中的部分子量子与大动量展开；混合重正化方案；短距离（公式）的重正化；大距离（公式）的重正化；渐近距离处的数据分析策略；中等大（公式）下的指数外推；超大（公式）下的代数外推；大动量展开对短距离展开；结论 \\
证据状态 & \verified：摘要和章节来自现有 LaTeX；\unverified：没有论文 PDF/原始数据独立复核。 \\
\bottomrule
\end{tabularx}
\endgroup
  \source{索引：\detokenize{refer/papers/INDEX.md:30}；正文：\detokenize{refer/papers/准光前关联的混合重正化方案_latex/main.tex}；\detokenize{refer/papers/Hybrid_Renorm_Quasi_LightFront_latex/main.tex}}
\end{frame}

%% frame_manifest|frame_id=P30-02|paper_id=P30|claim_id=P30-方法|visual_type=method-flow-formula|width_budget=0.98linewidth|height_budget=0.86textheight|min_font_pt=10|split_allowed=no|expected_page=131|risks=title-table-source-footline
\begin{frame}[t]{准光前关联的混合重正化方案：理论结构与方法链}
  \fontsize{10}{10.8}\selectfont
\begingroup
\renewcommand{\arraystretch}{0.84}
\begin{tabularx}{0.98\linewidth}{@{}p{2.55cm}>{\raggedright\arraybackslash}X@{}}
\toprule
输入与状态 & 格点/解析输入 → 非微扰重正化、OPE 与匹配中的算符、矩阵元或场配置；具体输入见源章节。 \\
章节路线 & 引言；无穷大动量态中的部分子量子与大动量展开；混合重正化方案；短距离（公式）的重正化；大距离（公式）的重正化；渐近距离处的数据分析策略；中等大（公式）下的指数外推；超大（公式）下的代数外推；大动量展开对短距离展开；结论 \\
方法摘录 & sec:intro 部分子物理对于理解强子高能碰撞的动力学以及研究强子的内部结构都十分重要（引文）。最熟悉的例子是夸克与胶子的部分子分布函数（PDF）：它们一方面为对撞机上的高能产生过程提供束流信息（引文），另一方面描述了对撞强子的束缚态物理。……如何利用相关关联函数在大光前（LF）距离处的渐近行为，消除截断傅里叶变换在动量分布中引起的非物理振荡 \\
结构式/示意 & $O^{R}(z,\mu)=Z(z,\mu,a)\,O^{\rm bare}(z,a)$\\[-0.2em]\scriptsize 该式只表达主题变量关系；论文特定推导以源文件为准。 \\
物理校验 & 重正化因子抵消截止依赖；a→0 和短距离 z→0 极限必须与匹配阶数、方案和尺度保持一致。 \\
\bottomrule
\end{tabularx}
\endgroup
  \source{方法证据：\detokenize{refer/papers/准光前关联的混合重正化方案_latex/chapters/section01.tex:1-16}；主题结构式为示意。}
\end{frame}

%% frame_manifest|frame_id=P30-03|paper_id=P30|claim_id=P30-结果|visual_type=result-evidence-table|width_budget=0.98linewidth|height_budget=0.86textheight|min_font_pt=10|split_allowed=no|expected_page=132|risks=title-table-source-footline
\begin{frame}[t]{准光前关联的混合重正化方案：结果、验证与物理含义}
  \fontsize{10}{10.8}\selectfont
\begingroup
\renewcommand{\arraystretch}{0.84}
\begin{tabularx}{0.98\linewidth}{@{}p{2.55cm}>{\raggedright\arraybackslash}X@{}}
\toprule
源文结果 & sec:conclusion 总之，我们讨论了格点上准 LF 关联的重正化与匹配中的若干进一步细节。我们提出了分别处理短距离与长距离关联的混合重正化方案：短距离关联可以通过把同一关联子夹在不同外部态之间作比值来重正化，而长距离关联则用 Wilson 线质量重正化来处理，并以一个连续性条件与短距离区域匹配。这样，我们在重正化阶段避免了在大距离引入额外的非微扰效应。我们还提出如何利……数据时的大动量展开与 CSF 方法作了比较：前者是提取 PDF（公式）依赖的系统展开，后者则不是。 我们在此的提议有望大幅改进当前 LaMET 格点应用中的计算策略 \\
验证状态 & 确证：源章节有结果/讨论摘录。 \\
库内可见证据 & 中文源约 1462 行；公式命中 47，图命中 8，表/表格命中 0。 \\
物理含义 & 在本主题共同链条中，本论文把上述输入推进到可比较的中间量或观测量；具体强度、误差和外推范围不超出源文档。 \\
未验证项 & 当前工作区没有论文 PDF、原始数据和独立复现脚本；因此数值复核、图像再现和统计显著性不作额外断言。 \\
\bottomrule
\end{tabularx}
\endgroup
  \source{结果证据：\detokenize{refer/papers/准光前关联的混合重正化方案_latex/chapters/section06.tex:1-6}；图表/公式计数由本报告生成器对中文源的静态统计得到。}
\end{frame}

%% frame_manifest|frame_id=P30-04|paper_id=P30|claim_id=P30-局限|visual_type=risk-relation-table|width_budget=0.98linewidth|height_budget=0.86textheight|min_font_pt=10|split_allowed=no|expected_page=133|risks=title-table-source-footline
\begin{frame}[t]{准光前关联的混合重正化方案：局限、跨论文接口与可复查入口}
  \fontsize{10}{10.8}\selectfont
\begingroup
\renewcommand{\arraystretch}{0.84}
\begin{tabularx}{0.98\linewidth}{@{}p{2.55cm}>{\raggedright\arraybackslash}X@{}}
\toprule
源文局限 & 需要回到结果/讨论/展望章节核对适用区间；本档案不把缺失的独立数据复核补成已证实结论。 \\
跨论文接口 & 是准/赝分布从格点数据走向 MS-bar 或其他连续方案的必要中间层。 \\
极限检查 & 重正化因子抵消截止依赖；a→0 和短距离 z→0 极限必须与匹配阶数、方案和尺度保持一致。 \\
下一步核验 & 补齐 PDF 或原始数据；按源文参数复现关键图表；比较不同动量、流时间、距离、格距或方案转换后的稳定性。 \\
完整来源 & \detokenize{refer/papers/准光前关联的混合重正化方案_latex/main.tex}；\detokenize{refer/papers/Hybrid_Renorm_Quasi_LightFront_latex/main.tex}；章节源：\detokenize{refer/papers/准光前关联的混合重正化方案_latex/chapters/*.tex}；英中目录：\detokenize{refer/papers/Hybrid_Renorm_Quasi_LightFront_latex}. \\
\bottomrule
\end{tabularx}
\endgroup
  \source{状态：\inferred 为主题接口推断；\unverified 为当前材料边界；完整文件清单见 manifest。}
\end{frame}

%% frame_manifest|frame_id=P31-01|paper_id=P31|claim_id=P31-定位|visual_type=paper-card|width_budget=0.98linewidth|height_budget=0.86textheight|min_font_pt=10|split_allowed=no|expected_page=134|risks=title-table-source-footline
\section{CT18全局分析}
\begin{frame}[t]{CT18全局分析：研究问题与论文定位}
  \fontsize{10}{10.8}\selectfont
\begingroup
\renewcommand{\arraystretch}{0.84}
\begin{tabularx}{0.98\linewidth}{@{}p{2.55cm}>{\raggedright\arraybackslash}X@{}}
\toprule
论文编号 & P31；主题：胶子、强子部分子分布与 TMD \\
书目信息 & New CTEQ global analysis of quantum chromodynamics with high-precision data from the LHC；arXiv:1912.10053；英/中转排页数 146/143 \\
研究对象 & 我们给出 CTEQ-TEA 合作组的新部分子分布函数（PDF）。它们是在 CT14 全球 QCD 分析既有数据集的基础上，加入 HERA I+II 合并深度非弹性散射数据集以及种类繁多的大型强子对撞机（LHC）高精度数据获得的。 ……的 CT18Z。文中还给出了 LHC 标准烛光截面（如（公式）融合希格斯玻色子截面）的理论计算 \\
章节证据 & sec:Introduction 引言；sec:Datasets CT18 全球 QCD 分析概览；sec:summary 概要总结；CT18 拟合的实验数据集；替代 PDF 拟合：CT18……sian 剖析；补充材料；PDF 组之间的更多比较；更多直方图与数据的比较 \\
证据状态 & \verified：摘要和章节来自现有 LaTeX；\unverified：没有论文 PDF/原始数据独立复核。 \\
\bottomrule
\end{tabularx}
\endgroup
  \source{索引：\detokenize{refer/papers/INDEX.md:31}；正文：\detokenize{refer/papers/CT18全局分析_latex/main.tex}；\detokenize{refer/papers/CT18_Global_Analysis_latex/main.tex}}
\end{frame}

%% frame_manifest|frame_id=P31-02|paper_id=P31|claim_id=P31-方法|visual_type=method-flow-formula|width_budget=0.98linewidth|height_budget=0.86textheight|min_font_pt=10|split_allowed=no|expected_page=135|risks=title-table-source-footline
\begin{frame}[t]{CT18全局分析：理论结构与方法链}
  \fontsize{10}{10.8}\selectfont
\begingroup
\renewcommand{\arraystretch}{0.84}
\begin{tabularx}{0.98\linewidth}{@{}p{2.55cm}>{\raggedright\arraybackslash}X@{}}
\toprule
输入与状态 & 格点/解析输入 → 胶子、强子部分子分布与 TMD中的算符、矩阵元或场配置；具体输入见源章节。 \\
章节路线 & sec:Introduction 引言；sec:Datasets CT18 全球 QCD 分析概览；sec:summary 概要总结；CT18 拟合的实验数据集；替代 PDF 拟合：CT18……sian 剖析；补充材料；PDF 组之间的更多比较；更多直方图与数据的比较 \\
方法摘录 & 为确定适合全球拟合的候选实验数据集，我们开发了两个快速初步分析程序。PDFSense 程序（引文）由南方卫理公会大学（SMU）开发，用于在做拟合之前定量预测哪些数据集会对全球 PDF 拟合产生影响。ePump 程序（引文）由密歇根州立大学（MS……算 NNLO 截面。APPLgrid 表与其他组的同类公开表格做了交叉验证，并在速度与精度上作了优化 \\
结构式/示意 & $\mathcal O_{\rm Euclid}\xrightarrow[\rm matching\,/\,evolution]{}f_{g/h}(x,\mu)\;{\rm or}\;F_{\rm TMD}(x,b_T,\mu,\zeta)$\\[-0.2em]\scriptsize 该式只表达主题变量关系；论文特定推导以源文件为准。 \\
物理校验 & 纵向分数、横向距离和重正化尺度是不同变量；不能把有限动量或有限 b\_T 的结果直接当作光锥分布。 \\
\bottomrule
\end{tabularx}
\endgroup
  \source{方法证据：\detokenize{refer/papers/CT18全局分析_latex/chapters/section02.tex:74-81}；主题结构式为示意。}
\end{frame}

%% frame_manifest|frame_id=P31-03|paper_id=P31|claim_id=P31-结果|visual_type=result-evidence-table|width_budget=0.98linewidth|height_budget=0.86textheight|min_font_pt=10|split_allowed=no|expected_page=136|risks=title-table-source-footline
\begin{frame}[t]{CT18全局分析：结果、验证与物理含义}
  \fontsize{10}{10.8}\selectfont
\begingroup
\renewcommand{\arraystretch}{0.84}
\begin{tabularx}{0.98\linewidth}{@{}p{2.55cm}>{\raggedright\arraybackslash}X@{}}
\toprule
源文结果 & 本文给出了 CT18 族部分子分布函数（PDF），包括 CT18Z、CT18A 与 CT18X 替代拟合。CT18 是 CTEQ-TEA 组全球分析给出的下一代 NNLO（以及 NLO） 质子 PDF。它承接 CT14 与 NNLO 分布发布之后的下一次更新，后者正是 CT14 发表后因 HERA I 与 II 精密合并数据发布而催生。 CT18 是名义上的 CTEQ-TEA ……APDF5 格式的 CT18 网格，以及 LHAPDF5 库 CTEQ-TEA 模块的更新版——编译时必须一并包含，以支持调用 CT18 附带的全部本征矢组（引文） \\
验证状态 & 确证：源章节有结果/讨论摘录。 \\
库内可见证据 & 中文源约 5705 行；公式命中 60，图命中 322，表/表格命中 35。 \\
物理含义 & 在本主题共同链条中，本论文把上述输入推进到可比较的中间量或观测量；具体强度、误差和外推范围不超出源文档。 \\
未验证项 & 当前工作区没有论文 PDF、原始数据和独立复现脚本；因此数值复核、图像再现和统计显著性不作额外断言。 \\
\bottomrule
\end{tabularx}
\endgroup
  \source{结果证据：\detokenize{refer/papers/CT18全局分析_latex/chapters/section07.tex:1-114}；图表/公式计数由本报告生成器对中文源的静态统计得到。}
\end{frame}

%% frame_manifest|frame_id=P31-04|paper_id=P31|claim_id=P31-局限|visual_type=risk-relation-table|width_budget=0.98linewidth|height_budget=0.86textheight|min_font_pt=10|split_allowed=no|expected_page=137|risks=title-table-source-footline
\begin{frame}[t]{CT18全局分析：局限、跨论文接口与可复查入口}
  \fontsize{10}{10.8}\selectfont
\begingroup
\renewcommand{\arraystretch}{0.84}
\begin{tabularx}{0.98\linewidth}{@{}p{2.55cm}>{\raggedright\arraybackslash}X@{}}
\toprule
源文局限 & 需要回到结果/讨论/展望章节核对适用区间；本档案不把缺失的独立数据复核补成已证实结论。 \\
跨论文接口 & 检验 LaMET/赝 PDF/准 TMD 方法是否能够覆盖真实强子结构与实验拟合所需的观测量。 索引备注：全量图已按 arXiv 源补齐（234 图，原参数还原） \\
极限检查 & 纵向分数、横向距离和重正化尺度是不同变量；不能把有限动量或有限 b\_T 的结果直接当作光锥分布。 \\
下一步核验 & 补齐 PDF 或原始数据；按源文参数复现关键图表；比较不同动量、流时间、距离、格距或方案转换后的稳定性。 \\
完整来源 & \detokenize{refer/papers/CT18全局分析_latex/main.tex}；\detokenize{refer/papers/CT18_Global_Analysis_latex/main.tex}；章节源：\detokenize{refer/papers/CT18全局分析_latex/chapters/*.tex}；英中目录：\detokenize{refer/papers/CT18_Global_Analysis_latex}. \\
\bottomrule
\end{tabularx}
\endgroup
  \source{状态：\inferred 为主题接口推断；\unverified 为当前材料边界；完整文件清单见 manifest。}
\end{frame}

%% frame_manifest|frame_id=P32-01|paper_id=P32|claim_id=P32-定位|visual_type=paper-card|width_budget=0.98linewidth|height_budget=0.86textheight|min_font_pt=10|split_allowed=no|expected_page=138|risks=title-table-source-footline
\section{NNPDF17高精度数据部分子分布}
\begin{frame}[t]{NNPDF17高精度数据部分子分布：研究问题与论文定位}
  \fontsize{10}{10.8}\selectfont
\begingroup
\renewcommand{\arraystretch}{0.84}
\begin{tabularx}{0.98\linewidth}{@{}p{2.55cm}>{\raggedright\arraybackslash}X@{}}
\toprule
论文编号 & P32；主题：胶子、强子部分子分布与 TMD \\
书目信息 & NNPDF17 High Precision Data；arXiv:1706.00428；英/中转排页数 104/100 \\
研究对象 & 0.25cm] （Parton distributions from high-precision collider data） 0.5cm Eur. Phys. J. C 77 (2017) 663（公式）arXiv:1706.……F3.1 PDF 首次同时以 Hessian 集合和经优化的压缩副本数蒙特卡罗集合两种形式发布 \\
章节证据 & 引言；实验与理论输入；实验数据：总体概览；深度非弹性结构函数；固定靶 Drell-Yan 产生；单举喷注；强子对撞机上的 Drell-Yan 产生；（公式）玻色子的横向动量；（公式）产生的微……展望；NNPDF3.1 缩减集合的验证；交付清单；展望；程序开发与基准测试 \\
证据状态 & \verified：摘要和章节来自现有 LaTeX；\unverified：没有论文 PDF/原始数据独立复核。 \\
\bottomrule
\end{tabularx}
\endgroup
  \source{索引：\detokenize{refer/papers/INDEX.md:32}；正文：\detokenize{refer/papers/NNPDF17高精度数据部分子分布_latex/main.tex}；\detokenize{refer/papers/NNPDF17_High_Precision_Data_latex/main.tex}}
\end{frame}

%% frame_manifest|frame_id=P32-02|paper_id=P32|claim_id=P32-方法|visual_type=method-flow-formula|width_budget=0.98linewidth|height_budget=0.86textheight|min_font_pt=10|split_allowed=no|expected_page=139|risks=title-table-source-footline
\begin{frame}[t]{NNPDF17高精度数据部分子分布：理论结构与方法链}
  \fontsize{10}{10.8}\selectfont
\begingroup
\renewcommand{\arraystretch}{0.84}
\begin{tabularx}{0.98\linewidth}{@{}p{2.55cm}>{\raggedright\arraybackslash}X@{}}
\toprule
输入与状态 & 格点/解析输入 → 胶子、强子部分子分布与 TMD中的算符、矩阵元或场配置；具体输入见源章节。 \\
章节路线 & 引言；实验与理论输入；实验数据：总体概览；深度非弹性结构函数；固定靶 Drell-Yan 产生；单举喷注；强子对撞机上的 Drell-Yan 产生；（公式）玻色子的横向动量；（公式）产生的微……展望；NNPDF3.1 缩减集合的验证；交付清单；展望；程序开发与基准测试 \\
方法摘录 & sec:expdata NNPDF3.1 PDF 集合纳入了大量新的实验数据。我们既用 NNPDF 3.0 中已有观测量的改进测定扩充了数据集（如（公式）与（公式）快度分布），也加入了两个新过程：顶夸克微分分布，以及首次进入全局 PDF 确定的…… 可用的以及即将到来的 13 TeV LHC Run II 数据将成为未来 NNPDF 版本的一部分 \\
结构式/示意 & $\mathcal O_{\rm Euclid}\xrightarrow[\rm matching\,/\,evolution]{}f_{g/h}(x,\mu)\;{\rm or}\;F_{\rm TMD}(x,b_T,\mu,\zeta)$\\[-0.2em]\scriptsize 该式只表达主题变量关系；论文特定推导以源文件为准。 \\
物理校验 & 纵向分数、横向距离和重正化尺度是不同变量；不能把有限动量或有限 b\_T 的结果直接当作光锥分布。 \\
\bottomrule
\end{tabularx}
\endgroup
  \source{方法证据：\detokenize{refer/papers/NNPDF17高精度数据部分子分布_latex/chapters/section02.tex:1-12}；主题结构式为示意。}
\end{frame}

%% frame_manifest|frame_id=P32-03|paper_id=P32|claim_id=P32-结果|visual_type=result-evidence-table|width_budget=0.98linewidth|height_budget=0.86textheight|min_font_pt=10|split_allowed=no|expected_page=140|risks=title-table-source-footline
\begin{frame}[t]{NNPDF17高精度数据部分子分布：结果、验证与物理含义}
  \fontsize{10}{10.8}\selectfont
\begingroup
\renewcommand{\arraystretch}{0.84}
\begin{tabularx}{0.98\linewidth}{@{}p{2.55cm}>{\raggedright\arraybackslash}X@{}}
\toprule
源文结果 & sec:outlook 这里给出的 NNPDF3.1 PDF 确定是对 NNPDF3.0 的更新，却在方法学与数据集两方面都包含实质创新， 产生了精度与准确性都显著提高的 PDF 集合。得益于此，可以推进多项衍生项目： 或旨在更新以往基于旧版 NNPDF 发布的结果，或因为更高的准确性使以前不可能或不感兴趣的项目变得可行。 这些衍生项目包括： itemize 强耦合常数（公式）……度很可能很快变得不可忽略，在某些运动学区域甚至会占据主导。 所有这些改进都将成为未来一次重大 PDF 发布的组成部分。 center 5cm .1pt center \\
验证状态 & 确证：源章节有结果/讨论摘录。 \\
库内可见证据 & 中文源约 5225 行；公式命中 3，图命中 339，表/表格命中 43。 \\
物理含义 & 在本主题共同链条中，本论文把上述输入推进到可比较的中间量或观测量；具体强度、误差和外推范围不超出源文档。 \\
未验证项 & 当前工作区没有论文 PDF、原始数据和独立复现脚本；因此数值复核、图像再现和统计显著性不作额外断言。 \\
\bottomrule
\end{tabularx}
\endgroup
  \source{结果证据：\detokenize{refer/papers/NNPDF17高精度数据部分子分布_latex/chapters/section06.tex:292-345}；图表/公式计数由本报告生成器对中文源的静态统计得到。}
\end{frame}

%% frame_manifest|frame_id=P32-04|paper_id=P32|claim_id=P32-局限|visual_type=risk-relation-table|width_budget=0.98linewidth|height_budget=0.86textheight|min_font_pt=10|split_allowed=no|expected_page=141|risks=title-table-source-footline
\begin{frame}[t]{NNPDF17高精度数据部分子分布：局限、跨论文接口与可复查入口}
  \fontsize{10}{10.8}\selectfont
\begingroup
\renewcommand{\arraystretch}{0.84}
\begin{tabularx}{0.98\linewidth}{@{}p{2.55cm}>{\raggedright\arraybackslash}X@{}}
\toprule
源文局限 & 需要回到结果/讨论/展望章节核对适用区间；本档案不把缺失的独立数据复核补成已证实结论。 \\
跨论文接口 & 检验 LaMET/赝 PDF/准 TMD 方法是否能够覆盖真实强子结构与实验拟合所需的观测量。 索引备注：中文版5类大表已原生还原（xeCJK 兼容修复） \\
极限检查 & 纵向分数、横向距离和重正化尺度是不同变量；不能把有限动量或有限 b\_T 的结果直接当作光锥分布。 \\
下一步核验 & 补齐 PDF 或原始数据；按源文参数复现关键图表；比较不同动量、流时间、距离、格距或方案转换后的稳定性。 \\
完整来源 & \detokenize{refer/papers/NNPDF17高精度数据部分子分布_latex/main.tex}；\detokenize{refer/papers/NNPDF17_High_Precision_Data_latex/main.tex}；章节源：\detokenize{refer/papers/NNPDF17高精度数据部分子分布_latex/chapters/*.tex}；英中目录：\detokenize{refer/papers/NNPDF17_High_Precision_Data_latex}. \\
\bottomrule
\end{tabularx}
\endgroup
  \source{状态：\inferred 为主题接口推断；\unverified 为当前材料边界；完整文件清单见 manifest。}
\end{frame}

%% frame_manifest|frame_id=P33-01|paper_id=P33|claim_id=P33-定位|visual_type=paper-card|width_budget=0.98linewidth|height_budget=0.86textheight|min_font_pt=10|split_allowed=no|expected_page=142|risks=title-table-source-footline
\section{基于流的格点MCMC生成模型}
\begin{frame}[t]{基于流的格点MCMC生成模型：研究问题与论文定位}
  \fontsize{10}{10.8}\selectfont
\begingroup
\renewcommand{\arraystretch}{0.84}
\begin{tabularx}{0.98\linewidth}{@{}p{2.55cm}>{\raggedright\arraybackslash}X@{}}
\toprule
论文编号 & P33；主题：机器学习采样与随机量化 \\
书目信息 & Flow Generative MCMC Lattice；arXiv:1904.12072；英/中转排页数 20/17 \\
研究对象 & 本文提出一种在格点场论中进行蒙特卡洛采样的马尔可夫链更新方案，它使用一个机器学习的基于流的生成模型。该生成模型可被优化（训练），以从近似于所研究理论的格点作用量所决定的期望玻尔兹曼分布中产生样本。对模型的训练能系统地改善马尔可夫链……可训练。在二维（公式）理论中，将该算法与 HMC 及局域 Metropolis 抽样进行了比较 \\
章节证据 & 引言；基于流的马尔可夫链蒙特卡洛算法；结合生成模型的 Metropolis-Hastings 方法；利用归一化流采样；生成式 Metropolis-Hastings 的自相关时间；临界慢化；……误差标度；临界慢化；训练成本；总结；自相关的接受率估计量；平均绝对误差损失 \\
证据状态 & \verified：摘要和章节来自现有 LaTeX；\unverified：没有论文 PDF/原始数据独立复核。 \\
\bottomrule
\end{tabularx}
\endgroup
  \source{索引：\detokenize{refer/papers/INDEX.md:33}；正文：\detokenize{refer/papers/基于流的格点MCMC生成模型_latex/main.tex}；\detokenize{refer/papers/Flow_Generative_MCMC_Lattice_latex/main.tex}}
\end{frame}

%% frame_manifest|frame_id=P33-02|paper_id=P33|claim_id=P33-方法|visual_type=method-flow-formula|width_budget=0.98linewidth|height_budget=0.86textheight|min_font_pt=10|split_allowed=no|expected_page=143|risks=title-table-source-footline
\begin{frame}[t]{基于流的格点MCMC生成模型：理论结构与方法链}
  \fontsize{10}{10.8}\selectfont
\begingroup
\renewcommand{\arraystretch}{0.84}
\begin{tabularx}{0.98\linewidth}{@{}p{2.55cm}>{\raggedright\arraybackslash}X@{}}
\toprule
输入与状态 & 格点/解析输入 → 机器学习采样与随机量化中的算符、矩阵元或场配置；具体输入见源章节。 \\
章节路线 & 引言；基于流的马尔可夫链蒙特卡洛算法；结合生成模型的 Metropolis-Hastings 方法；利用归一化流采样；生成式 Metropolis-Hastings 的自相关时间；临界慢化；……误差标度；临界慢化；训练成本；总结；自相关的接受率估计量；平均绝对误差损失 \\
方法摘录 & subsec:generative-metropolis 给定一个可以从已知概率分布（公式）抽样的模型，可以通过独立 Metropolis 采样器——Metropolis-Hastings 方法的一个特例（引文）——为期望概率分布（公式）构造一……以概率 1 被接受，此时马尔可夫链过程等价于对期望分布的直接抽样。这一点将在第（交叉引用）节精确阐述 \\
结构式/示意 & $z\sim p(z),\qquad U=f_\theta(z),\qquad p_\theta(U)\;\text{与目标系综比较}$\\[-0.2em]\scriptsize 该式只表达主题变量关系；论文特定推导以源文件为准。 \\
物理校验 & 规范不变性/等变性是结构约束；采样质量须由可观测量、接受率、自相关或分布距离验证，不能只看训练损失。 \\
\bottomrule
\end{tabularx}
\endgroup
  \source{方法证据：\detokenize{refer/papers/基于流的格点MCMC生成模型_latex/chapters/section02.tex:61-85}；主题结构式为示意。}
\end{frame}

%% frame_manifest|frame_id=P33-03|paper_id=P33|claim_id=P33-结果|visual_type=result-evidence-table|width_budget=0.98linewidth|height_budget=0.86textheight|min_font_pt=10|split_allowed=no|expected_page=144|risks=title-table-source-footline
\begin{frame}[t]{基于流的格点MCMC生成模型：结果、验证与物理含义}
  \fontsize{10}{10.8}\selectfont
\begingroup
\renewcommand{\arraystretch}{0.84}
\begin{tabularx}{0.98\linewidth}{@{}p{2.55cm}>{\raggedright\arraybackslash}X@{}}
\toprule
源文结果 & sec:summary 本工作定义了一种从期望概率分布中抽取格点场构型样本的基于流的 MCMC 算法： enumerate 训练一个 real NVP 流模型，使其近似产生期望分布； 由训练好的模型提出样本； 从任意初始构型出发，用 Metropolis-Hastings 算法对每个提议做接受或拒绝以推进马尔可夫链。 enumerate 本文证明了该方案定义一条遍历且平衡的马尔……外，为某个作用量训练的模型既可以重新训练，也可以作为针对参数值相近的另一作用量的流模型的先验。这将使参数调节更加高效，并能在现有模型之间插值或从其外推，生成更多系综 \\
验证状态 & 确证：源章节有结果/讨论摘录。 \\
库内可见证据 & 中文源约 1088 行；公式命中 32，图命中 18，表/表格命中 2。 \\
物理含义 & 在本主题共同链条中，本论文把上述输入推进到可比较的中间量或观测量；具体强度、误差和外推范围不超出源文档。 \\
未验证项 & 当前工作区没有论文 PDF、原始数据和独立复现脚本；因此数值复核、图像再现和统计显著性不作额外断言。 \\
\bottomrule
\end{tabularx}
\endgroup
  \source{结果证据：\detokenize{refer/papers/基于流的格点MCMC生成模型_latex/chapters/section04.tex:1-17}；图表/公式计数由本报告生成器对中文源的静态统计得到。}
\end{frame}

%% frame_manifest|frame_id=P33-04|paper_id=P33|claim_id=P33-局限|visual_type=risk-relation-table|width_budget=0.98linewidth|height_budget=0.86textheight|min_font_pt=10|split_allowed=no|expected_page=145|risks=title-table-source-footline
\begin{frame}[t]{基于流的格点MCMC生成模型：局限、跨论文接口与可复查入口}
  \fontsize{10}{10.8}\selectfont
\begingroup
\renewcommand{\arraystretch}{0.84}
\begin{tabularx}{0.98\linewidth}{@{}p{2.55cm}>{\raggedright\arraybackslash}X@{}}
\toprule
源文局限 & 需要回到结果/讨论/展望章节核对适用区间；本档案不把缺失的独立数据复核补成已证实结论。 \\
跨论文接口 & 为大规模格点系综生成提供算法补充，但必须回到物理可观测量和系综正确性检验。 \\
极限检查 & 规范不变性/等变性是结构约束；采样质量须由可观测量、接受率、自相关或分布距离验证，不能只看训练损失。 \\
下一步核验 & 补齐 PDF 或原始数据；按源文参数复现关键图表；比较不同动量、流时间、距离、格距或方案转换后的稳定性。 \\
完整来源 & \detokenize{refer/papers/基于流的格点MCMC生成模型_latex/main.tex}；\detokenize{refer/papers/Flow_Generative_MCMC_Lattice_latex/main.tex}；章节源：\detokenize{refer/papers/基于流的格点MCMC生成模型_latex/chapters/*.tex}；英中目录：\detokenize{refer/papers/Flow_Generative_MCMC_Lattice_latex}. \\
\bottomrule
\end{tabularx}
\endgroup
  \source{状态：\inferred 为主题接口推断；\unverified 为当前材料边界；完整文件清单见 manifest。}
\end{frame}

%% frame_manifest|frame_id=P34-01|paper_id=P34|claim_id=P34-定位|visual_type=paper-card|width_budget=0.98linewidth|height_budget=0.86textheight|min_font_pt=10|split_allowed=no|expected_page=146|risks=title-table-source-footline
\section{规范等变的流采样}
\begin{frame}[t]{规范等变的流采样：研究问题与论文定位}
  \fontsize{10}{10.8}\selectfont
\begingroup
\renewcommand{\arraystretch}{0.84}
\begin{tabularx}{0.98\linewidth}{@{}p{2.55cm}>{\raggedright\arraybackslash}X@{}}
\toprule
论文编号 & P34；主题：机器学习采样与随机量化 \\
书目信息 & Equivariant flow-based sampling for lattice gauge theory；arXiv:2003.06413；英/中转排页数 12/11 \\
研究对象 & 我们定义了一类机器学习的基于流的格点规范理论抽样算法，它们在构造上就是规范不变的。我们演示了该框架在两维时空 U(1) 规范理论中的应用，并发现在参数空间的临界点附近，与杂交蒙特卡罗（Hybrid Monte Carlo）和热浴（Heat Bath）等更传统的抽样方法相比，该方法在抽取拓扑量样本时的效率高出若干个数量级 \\
章节证据 & 基于流的采样；规范不变的流；规范等变耦合层；在 U(1) 规范理论中的应用；总结 \\
证据状态 & \verified：摘要和章节来自现有 LaTeX；\unverified：没有论文 PDF/原始数据独立复核。 \\
\bottomrule
\end{tabularx}
\endgroup
  \source{索引：\detokenize{refer/papers/INDEX.md:34}；正文：\detokenize{refer/papers/规范等变的流采样_latex/main.tex}；\detokenize{refer/papers/Equivariant_Flow_Sampling_latex/main.tex}}
\end{frame}

%% frame_manifest|frame_id=P34-02|paper_id=P34|claim_id=P34-方法|visual_type=method-flow-formula|width_budget=0.98linewidth|height_budget=0.86textheight|min_font_pt=10|split_allowed=no|expected_page=147|risks=title-table-source-footline
\begin{frame}[t]{规范等变的流采样：理论结构与方法链}
  \fontsize{10}{10.8}\selectfont
\begingroup
\renewcommand{\arraystretch}{0.84}
\begin{tabularx}{0.98\linewidth}{@{}p{2.55cm}>{\raggedright\arraybackslash}X@{}}
\toprule
输入与状态 & 格点/解析输入 → 机器学习采样与随机量化中的算符、矩阵元或场配置；具体输入见源章节。 \\
章节路线 & 基于流的采样；规范不变的流；规范等变耦合层；在 U(1) 规范理论中的应用；总结 \\
方法摘录 & 定义在两维时空中的、规范群为（公式）的规范理论是（公式）电动力学的淬火极限，即 Schwinger 模型（引文）。完整的 Schwinger 模型在解析上易于处理，同时复现了量子色动力学的许多特征（禁闭、轴反常、拓扑以及手征对称性破缺）。即使在……weep 混合起来，从而兼得标准马尔可夫链步骤对局域可观测量之利和基于流算法对延展及拓扑可观测量之利 \\
结构式/示意 & $z\sim p(z),\qquad U=f_\theta(z),\qquad p_\theta(U)\;\text{与目标系综比较}$\\[-0.2em]\scriptsize 该式只表达主题变量关系；论文特定推导以源文件为准。 \\
物理校验 & 规范不变性/等变性是结构约束；采样质量须由可观测量、接受率、自相关或分布距离验证，不能只看训练损失。 \\
\bottomrule
\end{tabularx}
\endgroup
  \source{方法证据：\detokenize{refer/papers/规范等变的流采样_latex/chapters/section05.tex:1-55}；主题结构式为示意。}
\end{frame}

%% frame_manifest|frame_id=P34-03|paper_id=P34|claim_id=P34-结果|visual_type=result-evidence-table|width_budget=0.98linewidth|height_budget=0.86textheight|min_font_pt=10|split_allowed=no|expected_page=148|risks=title-table-source-footline
\begin{frame}[t]{规范等变的流采样：结果、验证与物理含义}
  \fontsize{10}{10.8}\selectfont
\begingroup
\renewcommand{\arraystretch}{0.84}
\begin{tabularx}{0.98\linewidth}{@{}p{2.55cm}>{\raggedright\arraybackslash}X@{}}
\toprule
源文结果 & 当格点规范理论中的抽样发生临界慢化时，在理论的临界极限附近精确估计人们感兴趣的物理量就会受到阻碍。基于流的模型使我们能够直接从目标分布的一个近似中抽样，由此导出的物理可观测量估计值在无限统计极限下是 exact 的。本文引入了构造上满足精确规范不变性的基于流的模型，并发现把该方法应用于二维阿贝尔规范理论时，对拓扑量的估计比 HMC、热浴等现有算法更为高效。 这里提出的方法一般地……）定义了李群上的满射（虽然不是双射）映射，二者都使用了神经网络参数化。未来的工作将探索在这些方法的推广之上构造核，从而为 QCD 这样的非阿贝尔理论产生规范不变的流 \\
验证状态 & 确证：源章节有结果/讨论摘录。 \\
库内可见证据 & 中文源约 969 行；公式命中 8，图命中 8，表/表格命中 0。 \\
物理含义 & 在本主题共同链条中，本论文把上述输入推进到可比较的中间量或观测量；具体强度、误差和外推范围不超出源文档。 \\
未验证项 & 当前工作区没有论文 PDF、原始数据和独立复现脚本；因此数值复核、图像再现和统计显著性不作额外断言。 \\
\bottomrule
\end{tabularx}
\endgroup
  \source{结果证据：\detokenize{refer/papers/规范等变的流采样_latex/chapters/section06.tex:1-5}；图表/公式计数由本报告生成器对中文源的静态统计得到。}
\end{frame}

%% frame_manifest|frame_id=P34-04|paper_id=P34|claim_id=P34-局限|visual_type=risk-relation-table|width_budget=0.98linewidth|height_budget=0.86textheight|min_font_pt=10|split_allowed=no|expected_page=149|risks=title-table-source-footline
\begin{frame}[t]{规范等变的流采样：局限、跨论文接口与可复查入口}
  \fontsize{10}{10.8}\selectfont
\begingroup
\renewcommand{\arraystretch}{0.84}
\begin{tabularx}{0.98\linewidth}{@{}p{2.55cm}>{\raggedright\arraybackslash}X@{}}
\toprule
源文局限 & 需要回到结果/讨论/展望章节核对适用区间；本档案不把缺失的独立数据复核补成已证实结论。 \\
跨论文接口 & 为大规模格点系综生成提供算法补充，但必须回到物理可观测量和系综正确性检验。 \\
极限检查 & 规范不变性/等变性是结构约束；采样质量须由可观测量、接受率、自相关或分布距离验证，不能只看训练损失。 \\
下一步核验 & 补齐 PDF 或原始数据；按源文参数复现关键图表；比较不同动量、流时间、距离、格距或方案转换后的稳定性。 \\
完整来源 & \detokenize{refer/papers/规范等变的流采样_latex/main.tex}；\detokenize{refer/papers/Equivariant_Flow_Sampling_latex/main.tex}；章节源：\detokenize{refer/papers/规范等变的流采样_latex/chapters/*.tex}；英中目录：\detokenize{refer/papers/Equivariant_Flow_Sampling_latex}. \\
\bottomrule
\end{tabularx}
\endgroup
  \source{状态：\inferred 为主题接口推断；\unverified 为当前材料边界；完整文件清单见 manifest。}
\end{frame}

%% frame_manifest|frame_id=P35-01|paper_id=P35|claim_id=P35-定位|visual_type=paper-card|width_budget=0.98linewidth|height_budget=0.86textheight|min_font_pt=10|split_allowed=no|expected_page=150|risks=title-table-source-footline
\section{含费米子场的格点能动量张量}
\begin{frame}[t]{含费米子场的格点能动量张量：研究问题与论文定位}
  \fontsize{10}{10.8}\selectfont
\begingroup
\renewcommand{\arraystretch}{0.84}
\begin{tabularx}{0.98\linewidth}{@{}p{2.55cm}>{\raggedright\arraybackslash}X@{}}
\toprule
论文编号 & P35；主题：梯度流、涂抹与能动量张量 \\
书目信息 & Lattice energy–momentum tensor from the Yang–Mills gradient flow—inclusion of fermion fields；arXiv:1403.4772；英/中转排页数 30/31 \\
研究对象 & 经所谓 Yang–Mills 梯度流变形之后的场的局域乘积成为重正化复合算符。这一事实已被用来在纯 Yang–Mills 理论的格点表述中构造正确归一且守恒的能动量张量。本文将该构造进一步推广到含 费米子的矢量型规范理论 \\
章节证据 & 引言与主要结果；维数正规化下的能动量张量；Yang–Mills 梯度流；流方程与微扰展开；加圈费米子场；由流动场构造的能动量张量；小流时间展开与重正化群论证；（公式）至单圈阶；一致性检验：迹……积分公式；（公式）的胶子贡献；我们计算方案的合理性论证；修改后的能动量张量 \\
证据状态 & \verified：摘要和章节来自现有 LaTeX；\unverified：没有论文 PDF/原始数据独立复核。 \\
\bottomrule
\end{tabularx}
\endgroup
  \source{索引：\detokenize{refer/papers/INDEX.md:35}；正文：\detokenize{refer/papers/含费米子场的格点能动量张量_latex/main.tex}；\detokenize{refer/papers/Lattice_EMT_Fermion_Fields_latex/main.tex}}
\end{frame}

%% frame_manifest|frame_id=P35-02|paper_id=P35|claim_id=P35-方法|visual_type=method-flow-formula|width_budget=0.98linewidth|height_budget=0.86textheight|min_font_pt=10|split_allowed=no|expected_page=151|risks=title-table-source-footline
\begin{frame}[t]{含费米子场的格点能动量张量：理论结构与方法链}
  \fontsize{10}{10.8}\selectfont
\begingroup
\renewcommand{\arraystretch}{0.84}
\begin{tabularx}{0.98\linewidth}{@{}p{2.55cm}>{\raggedright\arraybackslash}X@{}}
\toprule
输入与状态 & 格点/解析输入 → 梯度流、涂抹与能动量张量中的算符、矩阵元或场配置；具体输入见源章节。 \\
章节路线 & 引言与主要结果；维数正规化下的能动量张量；Yang–Mills 梯度流；流方程与微扰展开；加圈费米子场；由流动场构造的能动量张量；小流时间展开与重正化群论证；（公式）至单圈阶；一致性检验：迹……积分公式；（公式）的胶子贡献；我们计算方案的合理性论证；修改后的能动量张量 \\
方法摘录 & sec:1 能量与动量是物理学中的基本概念。然而在格点场论（量子场论发展最成熟的非微扰表述）中， 与之对应的诺特流，即能动量 张量（引文），的构造却 并非直截了当（引文），因为平移不变性被格点结构 显式破缺。在近期的一篇论文（引文）中，人们基于……数定义（公式）。例如，对（公式）的基本 表示（公式）（此时（公式）），常规归一化为（公式或图表位置） \\
结构式/示意 & $\partial_t B_\mu(t,x)=D_\nu G_{\nu\mu}(t,x),\qquad t\sim L_{\rm sm}^{2}$\\[-0.2em]\scriptsize 该式只表达主题变量关系；论文特定推导以源文件为准。 \\
物理校验 & 流时间与平滑长度满足平方关系；t→0 回到原始场，有限 t 则抑制紫外模式。 \\
\bottomrule
\end{tabularx}
\endgroup
  \source{方法证据：\detokenize{refer/papers/含费米子场的格点能动量张量_latex/chapters/section01.tex:1-93}；主题结构式为示意。}
\end{frame}

%% frame_manifest|frame_id=P35-03|paper_id=P35|claim_id=P35-结果|visual_type=result-evidence-table|width_budget=0.98linewidth|height_budget=0.86textheight|min_font_pt=10|split_allowed=no|expected_page=152|risks=title-table-source-footline
\begin{frame}[t]{含费米子场的格点能动量张量：结果、验证与物理含义}
  \fontsize{10}{10.8}\selectfont
\begingroup
\renewcommand{\arraystretch}{0.84}
\begin{tabularx}{0.98\linewidth}{@{}p{2.55cm}>{\raggedright\arraybackslash}X@{}}
\toprule
源文结果 & sec:6 本文基于 Yang–Mills 梯度流，构造了公式（交叉引用），它为在含费米子的格点规范理论 中计算包含能动量张量的关联函数提供了一种可能的方法。这是文献（引文）针对纯 Yang–Mills 理论的构造的自然推广。尽管其在格点蒙特卡罗模拟中应用的可行性尚待仔细 考察，但 quenched QCD 热力学方面的经验（引文）强烈表明：即便采用目前可得 到的格点参数，也存……）中（公式）的外推。我们期待本表述的各种应用。其一是具有红外不动点的 多味规范理论（此类理论正受到近期的活跃研究；最近一次格点会议上的报 告（引文）可作为近期综述） \\
验证状态 & 确证：源章节有结果/讨论摘录。 \\
库内可见证据 & 中文源约 2667 行；公式命中 121，图命中 64，表/表格命中 10。 \\
物理含义 & 在本主题共同链条中，本论文把上述输入推进到可比较的中间量或观测量；具体强度、误差和外推范围不超出源文档。 \\
未验证项 & 当前工作区没有论文 PDF、原始数据和独立复现脚本；因此数值复核、图像再现和统计显著性不作额外断言。 \\
\bottomrule
\end{tabularx}
\endgroup
  \source{结果证据：\detokenize{refer/papers/含费米子场的格点能动量张量_latex/chapters/section06.tex:1-11}；图表/公式计数由本报告生成器对中文源的静态统计得到。}
\end{frame}

%% frame_manifest|frame_id=P35-04|paper_id=P35|claim_id=P35-局限|visual_type=risk-relation-table|width_budget=0.98linewidth|height_budget=0.86textheight|min_font_pt=10|split_allowed=no|expected_page=153|risks=title-table-source-footline
\begin{frame}[t]{含费米子场的格点能动量张量：局限、跨论文接口与可复查入口}
  \fontsize{10}{10.8}\selectfont
\begingroup
\renewcommand{\arraystretch}{0.84}
\begin{tabularx}{0.98\linewidth}{@{}p{2.55cm}>{\raggedright\arraybackslash}X@{}}
\toprule
源文局限 & 需要回到结果/讨论/展望章节核对适用区间；本档案不把缺失的独立数据复核补成已证实结论。 \\
跨论文接口 & 把连续极限、微扰有限性和能动量张量重构连接到可计算的格点观测量。 \\
极限检查 & 流时间与平滑长度满足平方关系；t→0 回到原始场，有限 t 则抑制紫外模式。 \\
下一步核验 & 补齐 PDF 或原始数据；按源文参数复现关键图表；比较不同动量、流时间、距离、格距或方案转换后的稳定性。 \\
完整来源 & \detokenize{refer/papers/含费米子场的格点能动量张量_latex/main.tex}；\detokenize{refer/papers/Lattice_EMT_Fermion_Fields_latex/main.tex}；章节源：\detokenize{refer/papers/含费米子场的格点能动量张量_latex/chapters/*.tex}；英中目录：\detokenize{refer/papers/Lattice_EMT_Fermion_Fields_latex}. \\
\bottomrule
\end{tabularx}
\endgroup
  \source{状态：\inferred 为主题接口推断；\unverified 为当前材料边界；完整文件清单见 manifest。}
\end{frame}

%% frame_manifest|frame_id=P36-01|paper_id=P36|claim_id=P36-定位|visual_type=paper-card|width_budget=0.98linewidth|height_budget=0.86textheight|min_font_pt=10|split_allowed=no|expected_page=154|risks=title-table-source-footline
\section{梯度流形式的双圈能动量张量}
\begin{frame}[t]{梯度流形式的双圈能动量张量：研究问题与论文定位}
  \fontsize{10}{10.8}\selectfont
\begingroup
\renewcommand{\arraystretch}{0.84}
\begin{tabularx}{0.98\linewidth}{@{}p{2.55cm}>{\raggedright\arraybackslash}X@{}}
\toprule
论文编号 & P36；主题：梯度流、涂抹与能动量张量 \\
书目信息 & The two-loop energy-momentum tensor within the gradient-flow formalism；arXiv:1808.09837；英/中转排页数 22/19 \\
研究对象 & QCD 能量-动量张量的梯度流表述被推广到 NNLO 微扰论。也就是说，在常规 能量-动量张量的相应表达式中乘在流动算符上的 Wilson 系数被计算到了这一阶。 结果的获得得益于近代正规微扰论工具的应用：把出现的双圈积分——其中还包含 流时间积分——约化为一小组可以解析计算的主积分 \\
章节证据 & 引言；微扰论中的 QCD 梯度流；能量-动量张量；Wilson 系数的计算；投影方法；计算方法；到 NNLO QCD 的系数函数；迹反常；算符重正化；结论；费曼规则 \\
证据状态 & \verified：摘要和章节来自现有 LaTeX；\unverified：没有论文 PDF/原始数据独立复核。 \\
\bottomrule
\end{tabularx}
\endgroup
  \source{索引：\detokenize{refer/papers/INDEX.md:36}；正文：\detokenize{refer/papers/梯度流形式的双圈能动量张量_latex/main.tex}；\detokenize{refer/papers/Two_Loop_EMT_Gradient_Flow_latex/main.tex}}
\end{frame}

%% frame_manifest|frame_id=P36-02|paper_id=P36|claim_id=P36-方法|visual_type=method-flow-formula|width_budget=0.98linewidth|height_budget=0.86textheight|min_font_pt=10|split_allowed=no|expected_page=155|risks=title-table-source-footline
\begin{frame}[t]{梯度流形式的双圈能动量张量：理论结构与方法链}
  \fontsize{10}{10.8}\selectfont
\begingroup
\renewcommand{\arraystretch}{0.84}
\begin{tabularx}{0.98\linewidth}{@{}p{2.55cm}>{\raggedright\arraybackslash}X@{}}
\toprule
输入与状态 & 格点/解析输入 → 梯度流、涂抹与能动量张量中的算符、矩阵元或场配置；具体输入见源章节。 \\
章节路线 & 引言；微扰论中的 QCD 梯度流；能量-动量张量；Wilson 系数的计算；投影方法；计算方法；到 NNLO QCD 的系数函数；迹反常；算符重正化；结论；费曼规则 \\
方法摘录 & 为计算系数（公式），我们使用所谓的``投影方法''（引文）： 构造外态（公式）和微分算符（公式），使得（公式或图表位置）这里为方便起见略去了洛伦兹指标；并且我们把矩阵元定义为只包含关于（w.r.t.） QCD 粒子为单粒子不可约（1 PI ）的…… eq:Omega 中出现的不同味道求和后，（公式）与（公式）之间的关系容易建立为（公式或图表位置） \\
结构式/示意 & $\partial_t B_\mu(t,x)=D_\nu G_{\nu\mu}(t,x),\qquad t\sim L_{\rm sm}^{2}$\\[-0.2em]\scriptsize 该式只表达主题变量关系；论文特定推导以源文件为准。 \\
物理校验 & 流时间与平滑长度满足平方关系；t→0 回到原始场，有限 t 则抑制紫外模式。 \\
\bottomrule
\end{tabularx}
\endgroup
  \source{方法证据：\detokenize{refer/papers/梯度流形式的双圈能动量张量_latex/chapters/section04.tex:3-170}；主题结构式为示意。}
\end{frame}

%% frame_manifest|frame_id=P36-03|paper_id=P36|claim_id=P36-结果|visual_type=result-evidence-table|width_budget=0.98linewidth|height_budget=0.86textheight|min_font_pt=10|split_allowed=no|expected_page=156|risks=title-table-source-footline
\begin{frame}[t]{梯度流形式的双圈能动量张量：结果、验证与物理含义}
  \fontsize{10}{10.8}\selectfont
\begingroup
\renewcommand{\arraystretch}{0.84}
\begin{tabularx}{0.98\linewidth}{@{}p{2.55cm}>{\raggedright\arraybackslash}X@{}}
\toprule
源文结果 & sec:conclusions 我们给出了梯度流定义的能动量张量的普适 Wilson 系数到 QCD。 修正使三个数值上占主导的系数（公式）、（公式）、（公式）在中心标度（公式）GeV（（公式）GeV）处改变约 10\%（1–2\%），其中（公式）通过 eq:mu0 与流时间（公式）相联系。相对于把重正化标度绕其中心值变化 因子 2 所导出的 结果，我们观察到理论不确定性的减小。第…… 总之，我们希望这些结果有助于改进格点上 的研究。它们是梯度流形式内 高阶微扰计算系统化框架（引文）的第一个成果，该框架在这一理论 框架的其他应用中也应证明是有用的 \\
验证状态 & 确证：源章节有结果/讨论摘录。 \\
库内可见证据 & 中文源约 1437 行；公式命中 64，图命中 16，表/表格命中 2。 \\
物理含义 & 在本主题共同链条中，本论文把上述输入推进到可比较的中间量或观测量；具体强度、误差和外推范围不超出源文档。 \\
未验证项 & 当前工作区没有论文 PDF、原始数据和独立复现脚本；因此数值复核、图像再现和统计显著性不作额外断言。 \\
\bottomrule
\end{tabularx}
\endgroup
  \source{结果证据：\detokenize{refer/papers/梯度流形式的双圈能动量张量_latex/chapters/section08.tex:1-16}；图表/公式计数由本报告生成器对中文源的静态统计得到。}
\end{frame}

%% frame_manifest|frame_id=P36-04|paper_id=P36|claim_id=P36-局限|visual_type=risk-relation-table|width_budget=0.98linewidth|height_budget=0.86textheight|min_font_pt=10|split_allowed=no|expected_page=157|risks=title-table-source-footline
\begin{frame}[t]{梯度流形式的双圈能动量张量：局限、跨论文接口与可复查入口}
  \fontsize{10}{10.8}\selectfont
\begingroup
\renewcommand{\arraystretch}{0.84}
\begin{tabularx}{0.98\linewidth}{@{}p{2.55cm}>{\raggedright\arraybackslash}X@{}}
\toprule
源文局限 & 需要回到结果/讨论/展望章节核对适用区间；本档案不把缺失的独立数据复核补成已证实结论。 \\
跨论文接口 & 把连续极限、微扰有限性和能动量张量重构连接到可计算的格点观测量。 \\
极限检查 & 流时间与平滑长度满足平方关系；t→0 回到原始场，有限 t 则抑制紫外模式。 \\
下一步核验 & 补齐 PDF 或原始数据；按源文参数复现关键图表；比较不同动量、流时间、距离、格距或方案转换后的稳定性。 \\
完整来源 & \detokenize{refer/papers/梯度流形式的双圈能动量张量_latex/main.tex}；\detokenize{refer/papers/Two_Loop_EMT_Gradient_Flow_latex/main.tex}；章节源：\detokenize{refer/papers/梯度流形式的双圈能动量张量_latex/chapters/*.tex}；英中目录：\detokenize{refer/papers/Two_Loop_EMT_Gradient_Flow_latex}. \\
\bottomrule
\end{tabularx}
\endgroup
  \source{状态：\inferred 为主题接口推断；\unverified 为当前材料边界；完整文件清单见 manifest。}
\end{frame}

%% frame_manifest|frame_id=P37-01|paper_id=P37|claim_id=P37-定位|visual_type=paper-card|width_budget=0.98linewidth|height_budget=0.86textheight|min_font_pt=10|split_allowed=no|expected_page=158|risks=title-table-source-footline
\section{梯度流高阶计算的技术与结果}
\begin{frame}[t]{梯度流高阶计算的技术与结果：研究问题与论文定位}
  \fontsize{10}{10.8}\selectfont
\begingroup
\renewcommand{\arraystretch}{0.84}
\begin{tabularx}{0.98\linewidth}{@{}p{2.55cm}>{\raggedright\arraybackslash}X@{}}
\toprule
论文编号 & P37；主题：梯度流、涂抹与能动量张量 \\
书目信息 & Results and techniques for higher order calculations within the gradient-flow formalism JHEP 06 (2019) 121. ……e except where noted in CONVERSION\_GUIDE.md.；arXiv:1905.00882；英/中转排页数 30/30 \\
研究对象 & 我们详细描述了在 QCD 梯度流形式中对可观测量实施系统微扰计算的实现方法。 其中包括本文所考虑的五维场论以及各复合算符的全部相关费曼规则的汇总。 我们使用标准微扰计算中的工具，在微扰论的更高阶得到有限流时间（公式）处的格林函数。……ged) 夸克场到 方案的转换因子。我们还重新计算了早先给出的三圈胶子凝聚结果， 提高了其精度 \\
章节证据 & 引言；微扰论中的 梯度流；梯度流的定义；流方程的微扰解；费曼规则；自动化实现；费曼图表达式的生成；流时间圈积分的 IBP 约化；流时间圈积分的数值计算；可观测量；三圈胶子凝聚的结果；三圈夸克凝聚的结果；梯度流耦合与质量；梯度流耦合；梯度流质量；总结；费曼规则；解析结果 \\
证据状态 & \verified：摘要和章节来自现有 LaTeX；\unverified：没有论文 PDF/原始数据独立复核。 \\
\bottomrule
\end{tabularx}
\endgroup
  \source{索引：\detokenize{refer/papers/INDEX.md:37}；正文：\detokenize{refer/papers/梯度流高阶计算的技术与结果_latex/main.tex}；\detokenize{refer/papers/Higher_Order_Gradient_Flow_Techniques_latex/main.tex}}
\end{frame}

%% frame_manifest|frame_id=P37-02|paper_id=P37|claim_id=P37-方法|visual_type=method-flow-formula|width_budget=0.98linewidth|height_budget=0.86textheight|min_font_pt=10|split_allowed=no|expected_page=159|risks=title-table-source-footline
\begin{frame}[t]{梯度流高阶计算的技术与结果：理论结构与方法链}
  \fontsize{10}{10.8}\selectfont
\begingroup
\renewcommand{\arraystretch}{0.84}
\begin{tabularx}{0.98\linewidth}{@{}p{2.55cm}>{\raggedright\arraybackslash}X@{}}
\toprule
输入与状态 & 格点/解析输入 → 梯度流、涂抹与能动量张量中的算符、矩阵元或场配置；具体输入见源章节。 \\
章节路线 & 引言；微扰论中的 梯度流；梯度流的定义；流方程的微扰解；费曼规则；自动化实现；费曼图表达式的生成；流时间圈积分的 IBP 约化；流时间圈积分的数值计算；可观测量；三圈胶子凝聚的结果；三圈夸克凝聚的结果；梯度流耦合与质量；梯度流耦合；梯度流质量；总结；费曼规则；解析结果 \\
方法摘录 & 下面我们在（公式）维欧氏时空中工作，（公式）。 把常规 （regular） 我们用``流动的''（flowed）与``常规的''（regular） 来区分在（公式）与在（公式）定义的量。 的胶子场与夸克场（公式）、（公式）延拓为（公式）维场（公……合圈。后面会看清楚：仅由流动场构 成的闭合圈为零，因此可以在拉格朗日量层面就把（公式）与（公式）略去 \\
结构式/示意 & $\partial_t B_\mu(t,x)=D_\nu G_{\nu\mu}(t,x),\qquad t\sim L_{\rm sm}^{2}$\\[-0.2em]\scriptsize 该式只表达主题变量关系；论文特定推导以源文件为准。 \\
物理校验 & 流时间与平滑长度满足平方关系；t→0 回到原始场，有限 t 则抑制紫外模式。 \\
\bottomrule
\end{tabularx}
\endgroup
  \source{方法证据：\detokenize{refer/papers/梯度流高阶计算的技术与结果_latex/chapters/section02.tex:7-157}；主题结构式为示意。}
\end{frame}

%% frame_manifest|frame_id=P37-03|paper_id=P37|claim_id=P37-结果|visual_type=result-evidence-table|width_budget=0.98linewidth|height_budget=0.86textheight|min_font_pt=10|split_allowed=no|expected_page=160|risks=title-table-source-footline
\begin{frame}[t]{梯度流高阶计算的技术与结果：结果、验证与物理含义}
  \fontsize{10}{10.8}\selectfont
\begingroup
\renewcommand{\arraystretch}{0.84}
\begin{tabularx}{0.98\linewidth}{@{}p{2.55cm}>{\raggedright\arraybackslash}X@{}}
\toprule
源文结果 & sec:conclusions 本文给出了在微扰梯度流形式中计算关联函数的一个框架，它使用了标准微扰 QCD 计算的工具与技术。该框架基于 Luscher:2011bx 的五维场论 表述，以费曼图方法处理梯度流形式。我们实现了梯度流形式中相应的 费曼规 则，以及与本文所考虑的量相关的复合算符的费曼规则，并在很大程度上自动化了代 数表达式与圈积分的约化和计算。我们框架的一个重要要……已在这一 方向迈出第一步。 或许还有夸克质量的精密第一性原理测定。 本文聚焦真空矩阵元的计算；原则上，我们的框架也应适用于更普遍的问题，例如可 能涉及更多标度的情形 \\
验证状态 & 确证：源章节有结果/讨论摘录。 \\
库内可见证据 & 中文源约 2436 行；公式命中 102，图命中 55，表/表格命中 4。 \\
物理含义 & 在本主题共同链条中，本论文把上述输入推进到可比较的中间量或观测量；具体强度、误差和外推范围不超出源文档。 \\
未验证项 & 当前工作区没有论文 PDF、原始数据和独立复现脚本；因此数值复核、图像再现和统计显著性不作额外断言。 \\
\bottomrule
\end{tabularx}
\endgroup
  \source{结果证据：\detokenize{refer/papers/梯度流高阶计算的技术与结果_latex/chapters/section06.tex:1-25}；图表/公式计数由本报告生成器对中文源的静态统计得到。}
\end{frame}

%% frame_manifest|frame_id=P37-04|paper_id=P37|claim_id=P37-局限|visual_type=risk-relation-table|width_budget=0.98linewidth|height_budget=0.86textheight|min_font_pt=10|split_allowed=no|expected_page=161|risks=title-table-source-footline
\begin{frame}[t]{梯度流高阶计算的技术与结果：局限、跨论文接口与可复查入口}
  \fontsize{10}{10.8}\selectfont
\begingroup
\renewcommand{\arraystretch}{0.84}
\begin{tabularx}{0.98\linewidth}{@{}p{2.55cm}>{\raggedright\arraybackslash}X@{}}
\toprule
源文局限 & 需要回到结果/讨论/展望章节核对适用区间；本档案不把缺失的独立数据复核补成已证实结论。 \\
跨论文接口 & 把连续极限、微扰有限性和能动量张量重构连接到可计算的格点观测量。 \\
极限检查 & 流时间与平滑长度满足平方关系；t→0 回到原始场，有限 t 则抑制紫外模式。 \\
下一步核验 & 补齐 PDF 或原始数据；按源文参数复现关键图表；比较不同动量、流时间、距离、格距或方案转换后的稳定性。 \\
完整来源 & \detokenize{refer/papers/梯度流高阶计算的技术与结果_latex/main.tex}；\detokenize{refer/papers/Higher_Order_Gradient_Flow_Techniques_latex/main.tex}；章节源：\detokenize{refer/papers/梯度流高阶计算的技术与结果_latex/chapters/*.tex}；英中目录：\detokenize{refer/papers/Higher_Order_Gradient_Flow_Techniques_latex}. \\
\bottomrule
\end{tabularx}
\endgroup
  \source{状态：\inferred 为主题接口推断；\unverified 为当前材料边界；完整文件清单见 manifest。}
\end{frame}

%% frame_manifest|frame_id=P38-01|paper_id=P38|claim_id=P38-定位|visual_type=paper-card|width_budget=0.98linewidth|height_budget=0.86textheight|min_font_pt=10|split_allowed=no|expected_page=162|risks=title-table-source-footline
\section{局域涂抹算符乘积展开}
\begin{frame}[t]{局域涂抹算符乘积展开：研究问题与论文定位}
  \fontsize{10}{10.8}\selectfont
\begingroup
\renewcommand{\arraystretch}{0.84}
\begin{tabularx}{0.98\linewidth}{@{}p{2.55cm}>{\raggedright\arraybackslash}X@{}}
\toprule
论文编号 & P38；主题：梯度流、涂抹与能动量张量 \\
书目信息 & Locally smeared operator product expansions in scalar field theory Phys. Rev. D 91, 074513 (2015). DOI: 10.1……evD.91.074513; arXiv:1501.05348v2 [hep-lat].；arXiv:1501.05348；英/中转排页数 19/19 \\
研究对象 & 我们提出一种新的局域涂抹算符乘积展开，将非定域算符按一组涂抹算符构成的基底分解。 这种涂抹算符乘积展开在形式上把由格点场论数值计算确定的非微扰矩阵元与连续理论中 非定域算符的矩阵元联系起来。只要在连续极限下保持涂抹尺度固定，这些非……系数的联系变得复杂， 需要构造一套合适的理论框架。我们以实标量场论中的例子演示了该方法的可行性 \\
章节证据 & 引言；sec:ope 算符乘积展开；sec:phi4 标量场理论的梯度流；sec:lattsmear 格点上的混合；sec:wc sOPE 的威尔逊系数；非连通贡献；连通贡献；sec:rgeqns 重正化群方程；sOPE 威尔逊系数的重正化群方程；结论 \\
证据状态 & \verified：摘要和章节来自现有 LaTeX；\unverified：没有论文 PDF/原始数据独立复核。 \\
\bottomrule
\end{tabularx}
\endgroup
  \source{索引：\detokenize{refer/papers/INDEX.md:38}；正文：\detokenize{refer/papers/局域涂抹算符乘积展开_latex/main.tex}；\detokenize{refer/papers/Locally_Smeared_OPE_latex/main.tex}}
\end{frame}

%% frame_manifest|frame_id=P38-02|paper_id=P38|claim_id=P38-方法|visual_type=method-flow-formula|width_budget=0.98linewidth|height_budget=0.86textheight|min_font_pt=10|split_allowed=no|expected_page=163|risks=title-table-source-footline
\begin{frame}[t]{局域涂抹算符乘积展开：理论结构与方法链}
  \fontsize{10}{10.8}\selectfont
\begingroup
\renewcommand{\arraystretch}{0.84}
\begin{tabularx}{0.98\linewidth}{@{}p{2.55cm}>{\raggedright\arraybackslash}X@{}}
\toprule
输入与状态 & 格点/解析输入 → 梯度流、涂抹与能动量张量中的算符、矩阵元或场配置；具体输入见源章节。 \\
章节路线 & 引言；sec:ope 算符乘积展开；sec:phi4 标量场理论的梯度流；sec:lattsmear 格点上的混合；sec:wc sOPE 的威尔逊系数；非连通贡献；连通贡献；sec:rgeqns 重正化群方程；sOPE 威尔逊系数的重正化群方程；结论 \\
方法摘录 & 我们在四维欧氏标量场论中工作，该理论具有四次自相互作用与裸质量（公式），由作用量（公式或图表位置）定义。 为研究 sOPE，我们引入标量梯度流，它由流时间演化方程（公式或图表位置）定义，其中（公式）是非零流时间（公式）处的标量场，（公式）是欧氏……因此任何相互作用都必须发生在零流时间处。四点顶点相应的费曼规则就是标准的 （欧氏）四点顶点，（公式） \\
结构式/示意 & $\partial_t B_\mu(t,x)=D_\nu G_{\nu\mu}(t,x),\qquad t\sim L_{\rm sm}^{2}$\\[-0.2em]\scriptsize 该式只表达主题变量关系；论文特定推导以源文件为准。 \\
物理校验 & 流时间与平滑长度满足平方关系；t→0 回到原始场，有限 t 则抑制紫外模式。 \\
\bottomrule
\end{tabularx}
\endgroup
  \source{方法证据：\detokenize{refer/papers/局域涂抹算符乘积展开_latex/chapters/section03.tex:1-43}；主题结构式为示意。}
\end{frame}

%% frame_manifest|frame_id=P38-03|paper_id=P38|claim_id=P38-结果|visual_type=result-evidence-table|width_budget=0.98linewidth|height_budget=0.86textheight|min_font_pt=10|split_allowed=no|expected_page=164|risks=title-table-source-footline
\begin{frame}[t]{局域涂抹算符乘积展开：结果、验证与物理含义}
  \fontsize{10}{10.8}\selectfont
\begingroup
\renewcommand{\arraystretch}{0.84}
\begin{tabularx}{0.98\linewidth}{@{}p{2.55cm}>{\raggedright\arraybackslash}X@{}}
\toprule
源文结果 & 我们提出了一种新方法——涂抹算符乘积展开（sOPE）——用以从数值非微扰计算中提取矩阵元， 而不受幂次发散混合之苦。涂抹算符乘积展开是一个普遍的框架，适用于任何存在遭受幂次发散混合的 非微扰矩阵元的渐近自由理论。在 QCD 之内，最直接的应用是深度非弹性散射；其他应用还包括（公式）衰变（引文）与（公式）介子混合（引文）的非微扰确定。在 QCD 之外，应用还包括海森堡模型、 自旋……发散混合。 所得矩阵元是两个尺度的函数：重正化尺度与涂抹长度。sOPE 系统地把这些矩阵元与可在微扰论中计算的 涂抹威尔逊系数联系起来，从而给出非定域算符的完整确定 \\
验证状态 & 确证：源章节有结果/讨论摘录。 \\
库内可见证据 & 中文源约 1577 行；公式命中 58，图命中 16，表/表格命中 0。 \\
物理含义 & 在本主题共同链条中，本论文把上述输入推进到可比较的中间量或观测量；具体强度、误差和外推范围不超出源文档。 \\
未验证项 & 当前工作区没有论文 PDF、原始数据和独立复现脚本；因此数值复核、图像再现和统计显著性不作额外断言。 \\
\bottomrule
\end{tabularx}
\endgroup
  \source{结果证据：\detokenize{refer/papers/局域涂抹算符乘积展开_latex/chapters/section07.tex:1-15}；图表/公式计数由本报告生成器对中文源的静态统计得到。}
\end{frame}

%% frame_manifest|frame_id=P38-04|paper_id=P38|claim_id=P38-局限|visual_type=risk-relation-table|width_budget=0.98linewidth|height_budget=0.86textheight|min_font_pt=10|split_allowed=no|expected_page=165|risks=title-table-source-footline
\begin{frame}[t]{局域涂抹算符乘积展开：局限、跨论文接口与可复查入口}
  \fontsize{10}{10.8}\selectfont
\begingroup
\renewcommand{\arraystretch}{0.84}
\begin{tabularx}{0.98\linewidth}{@{}p{2.55cm}>{\raggedright\arraybackslash}X@{}}
\toprule
源文局限 & 需要回到结果/讨论/展望章节核对适用区间；本档案不把缺失的独立数据复核补成已证实结论。 \\
跨论文接口 & 把连续极限、微扰有限性和能动量张量重构连接到可计算的格点观测量。 \\
极限检查 & 流时间与平滑长度满足平方关系；t→0 回到原始场，有限 t 则抑制紫外模式。 \\
下一步核验 & 补齐 PDF 或原始数据；按源文参数复现关键图表；比较不同动量、流时间、距离、格距或方案转换后的稳定性。 \\
完整来源 & \detokenize{refer/papers/局域涂抹算符乘积展开_latex/main.tex}；\detokenize{refer/papers/Locally_Smeared_OPE_latex/main.tex}；章节源：\detokenize{refer/papers/局域涂抹算符乘积展开_latex/chapters/*.tex}；英中目录：\detokenize{refer/papers/Locally_Smeared_OPE_latex}. \\
\bottomrule
\end{tabularx}
\endgroup
  \source{状态：\inferred 为主题接口推断；\unverified 为当前材料边界；完整文件清单见 manifest。}
\end{frame}

%% frame_manifest|frame_id=P39-01|paper_id=P39|claim_id=P39-定位|visual_type=paper-card|width_budget=0.98linewidth|height_budget=0.86textheight|min_font_pt=10|split_allowed=no|expected_page=166|risks=title-table-source-footline
\section{微扰论中的涂抹准分布}
\begin{frame}[t]{微扰论中的涂抹准分布：研究问题与论文定位}
  \fontsize{10}{10.8}\selectfont
\begingroup
\renewcommand{\arraystretch}{0.84}
\begin{tabularx}{0.98\linewidth}{@{}p{2.55cm}>{\raggedright\arraybackslash}X@{}}
\toprule
论文编号 & P39；主题：胶子、强子部分子分布与 TMD \\
书目信息 & Smeared quasidistributions in perturbation theory INT-PUB-17-043. arXiv:1710.04607v3 [hep-lat]; Phys. Rev. D……8); Erratum Phys. Rev. D 110, 059902 (2024).；arXiv:1710.04607；英/中转排页数 20/17 \\
研究对象 & 准分布与赝分布为从 QCD 的第一性原理计算确定部分子分布函数提供了一条新途径。本文在微扰论单圈水平计算味非单态非极化准分布，利用梯度流移除紫外发散。我证明，正如预期，梯度流在单圈水平不改变准分布的红外结构，并利用所得结果把涂抹矩…… QCD 计算得到的数值结果与从实验数据全局拟合分析中提取的光前部分子分布函数联系起来所必需的 \\
章节证据 & 引言；sec:defns 光前分布与准分布；光前 PDF；涂抹准分布；涂抹准分布与 PDF 的联系；sec:pertth 微扰论中的矩阵元；顶点型图；波函数型图；单圈重正化参数；单圈涂抹矩阵……s 数值结果；sec:summary 总结；维数正规化下的矩阵元；数值积分 \\
证据状态 & \verified：摘要和章节来自现有 LaTeX；\unverified：没有论文 PDF/原始数据独立复核。 \\
\bottomrule
\end{tabularx}
\endgroup
  \source{索引：\detokenize{refer/papers/INDEX.md:39}；正文：\detokenize{refer/papers/微扰论中的涂抹准分布_latex/main.tex}；\detokenize{refer/papers/Smeared_Quasidists_Perturbation_latex/main.tex}}
\end{frame}

%% frame_manifest|frame_id=P39-02|paper_id=P39|claim_id=P39-方法|visual_type=method-flow-formula|width_budget=0.98linewidth|height_budget=0.86textheight|min_font_pt=10|split_allowed=no|expected_page=167|risks=title-table-source-footline
\begin{frame}[t]{微扰论中的涂抹准分布：理论结构与方法链}
  \fontsize{10}{10.8}\selectfont
\begingroup
\renewcommand{\arraystretch}{0.84}
\begin{tabularx}{0.98\linewidth}{@{}p{2.55cm}>{\raggedright\arraybackslash}X@{}}
\toprule
输入与状态 & 格点/解析输入 → 胶子、强子部分子分布与 TMD中的算符、矩阵元或场配置；具体输入见源章节。 \\
章节路线 & 引言；sec:defns 光前分布与准分布；光前 PDF；涂抹准分布；涂抹准分布与 PDF 的联系；sec:pertth 微扰论中的矩阵元；顶点型图；波函数型图；单圈重正化参数；单圈涂抹矩阵……s 数值结果；sec:summary 总结；维数正规化下的矩阵元；数值积分 \\
方法摘录 & 我在图（交叉引用）中展示涂抹非单态矩阵元的领头阶（即树图级）贡献。在图（交叉引用）与图（交叉引用）中，我采用如下约定：实线表示费米子传播子；双线是费米子流核；空心方块是任意流时（公式）处的流顶点；灰圈表示固定流时（公式）的费米子场；灰方块表示固……对外部动量、夸克质量与流时的函数关系数值求值所有九个依赖流时的图，以展示单圈微扰论下完整矩阵元的行为 \\
结构式/示意 & $\mathcal O_{\rm Euclid}\xrightarrow[\rm matching\,/\,evolution]{}f_{g/h}(x,\mu)\;{\rm or}\;F_{\rm TMD}(x,b_T,\mu,\zeta)$\\[-0.2em]\scriptsize 该式只表达主题变量关系；论文特定推导以源文件为准。 \\
物理校验 & 纵向分数、横向距离和重正化尺度是不同变量；不能把有限动量或有限 b\_T 的结果直接当作光锥分布。 \\
\bottomrule
\end{tabularx}
\endgroup
  \source{方法证据：\detokenize{refer/papers/微扰论中的涂抹准分布_latex/chapters/section03.tex:1-27}；主题结构式为示意。}
\end{frame}

%% frame_manifest|frame_id=P39-03|paper_id=P39|claim_id=P39-结果|visual_type=result-evidence-table|width_budget=0.98linewidth|height_budget=0.86textheight|min_font_pt=10|split_allowed=no|expected_page=168|risks=title-table-source-footline
\begin{frame}[t]{微扰论中的涂抹准分布：结果、验证与物理含义}
  \fontsize{10}{10.8}\selectfont
\begingroup
\renewcommand{\arraystretch}{0.84}
\begin{tabularx}{0.98\linewidth}{@{}p{2.55cm}>{\raggedright\arraybackslash}X@{}}
\toprule
源文结果 & 准分布与赝分布为核子结构打开了一扇新的窗口，并首次提供了直接从格点 QCD 提取部分子分布函数的机会。非微扰计算的初步结果令人鼓舞，但在格点学界能够宣称稳健地确定光前 PDF 之前，仍有挑战需要克服。格点 QCD 的一个特别困难，是与 Wilson 线长度相联系的幂次发散的存在，在取连续极限之前必须将其非微扰地移除。在文献（引文）中，我们引入涂抹准分布作为规避 Wilson 线……阵元。把这一研究扩展到两圈很可能需要自动化的方法，因为梯度流诱导的额外顶点使图的数目随微扰阶迅速增长。与涂抹准分布和赝分布相联系的系统不确定度的非微扰研究正在进行中 \\
验证状态 & 确证：源章节有结果/讨论摘录。 \\
库内可见证据 & 中文源约 1161 行；公式命中 46，图命中 16，表/表格命中 0。 \\
物理含义 & 在本主题共同链条中，本论文把上述输入推进到可比较的中间量或观测量；具体强度、误差和外推范围不超出源文档。 \\
未验证项 & 当前工作区没有论文 PDF、原始数据和独立复现脚本；因此数值复核、图像再现和统计显著性不作额外断言。 \\
\bottomrule
\end{tabularx}
\endgroup
  \source{结果证据：\detokenize{refer/papers/微扰论中的涂抹准分布_latex/chapters/section05.tex:1-8}；图表/公式计数由本报告生成器对中文源的静态统计得到。}
\end{frame}

%% frame_manifest|frame_id=P39-04|paper_id=P39|claim_id=P39-局限|visual_type=risk-relation-table|width_budget=0.98linewidth|height_budget=0.86textheight|min_font_pt=10|split_allowed=no|expected_page=169|risks=title-table-source-footline
\begin{frame}[t]{微扰论中的涂抹准分布：局限、跨论文接口与可复查入口}
  \fontsize{10}{10.8}\selectfont
\begingroup
\renewcommand{\arraystretch}{0.84}
\begin{tabularx}{0.98\linewidth}{@{}p{2.55cm}>{\raggedright\arraybackslash}X@{}}
\toprule
源文局限 & 需要回到结果/讨论/展望章节核对适用区间；本档案不把缺失的独立数据复核补成已证实结论。 \\
跨论文接口 & 检验 LaMET/赝 PDF/准 TMD 方法是否能够覆盖真实强子结构与实验拟合所需的观测量。 \\
极限检查 & 纵向分数、横向距离和重正化尺度是不同变量；不能把有限动量或有限 b\_T 的结果直接当作光锥分布。 \\
下一步核验 & 补齐 PDF 或原始数据；按源文参数复现关键图表；比较不同动量、流时间、距离、格距或方案转换后的稳定性。 \\
完整来源 & \detokenize{refer/papers/微扰论中的涂抹准分布_latex/main.tex}；\detokenize{refer/papers/Smeared_Quasidists_Perturbation_latex/main.tex}；章节源：\detokenize{refer/papers/微扰论中的涂抹准分布_latex/chapters/*.tex}；英中目录：\detokenize{refer/papers/Smeared_Quasidists_Perturbation_latex}. \\
\bottomrule
\end{tabularx}
\endgroup
  \source{状态：\inferred 为主题接口推断；\unverified 为当前材料边界；完整文件清单见 manifest。}
\end{frame}

%% frame_manifest|frame_id=P40-01|paper_id=P40|claim_id=P40-定位|visual_type=paper-card|width_budget=0.98linewidth|height_budget=0.86textheight|min_font_pt=10|split_allowed=no|expected_page=170|risks=title-table-source-footline
\section{梯度流中的非光锥威尔逊线算符}
\begin{frame}[t]{梯度流中的非光锥威尔逊线算符：研究问题与论文定位}
  \fontsize{10}{10.8}\selectfont
\begingroup
\renewcommand{\arraystretch}{0.84}
\begin{tabularx}{0.98\linewidth}{@{}p{2.55cm}>{\raggedright\arraybackslash}X@{}}
\toprule
论文编号 & P40；主题：梯度流、涂抹与能动量张量 \\
书目信息 & Off LightCone Wilson Line Flow；arXiv:2312.05032；英/中转排页数 27/27 \\
研究对象 & ： 非光锥威尔逊线算符由通过类时或类空威尔逊线相连的局域算符构成，威尔逊线保证了算符的规范不变性。非光锥威尔逊线算符在多种语境中有广泛应用。例如，类空威尔逊线算符在确定准分布函数（quasi-PDF）方面起着关键作用，而类时威尔逊……D 中夸克偶素衰变和产生相联系的胶子关联函数，以及以色电场与色磁场插入威尔逊圈表述的自旋依赖势 \\
章节证据 & 引言；一维辅助场形式与非光锥威尔逊线算符的重正化；一维辅助场形式；非光锥威尔逊线算符的重正化；梯度流与小流时展开；梯度流形式；小流时展开；计算与结果；（公式）\&（公式）；（公式）；（公式）；……准PDF算符；威尔逊线自能图（公式）；顶点图（公式）；图（公式）；最终结果 \\
证据状态 & \verified：摘要和章节来自现有 LaTeX；\unverified：没有论文 PDF/原始数据独立复核。 \\
\bottomrule
\end{tabularx}
\endgroup
  \source{索引：\detokenize{refer/papers/INDEX.md:40}；正文：\detokenize{refer/papers/梯度流中的非光锥威尔逊线算符_latex/main.tex}；\detokenize{refer/papers/Off_LightCone_Wilson_Line_Flow_latex/main.tex}}
\end{frame}

%% frame_manifest|frame_id=P40-02|paper_id=P40|claim_id=P40-方法|visual_type=method-flow-formula|width_budget=0.98linewidth|height_budget=0.86textheight|min_font_pt=10|split_allowed=no|expected_page=171|risks=title-table-source-footline
\begin{frame}[t]{梯度流中的非光锥威尔逊线算符：理论结构与方法链}
  \fontsize{10}{10.8}\selectfont
\begingroup
\renewcommand{\arraystretch}{0.84}
\begin{tabularx}{0.98\linewidth}{@{}p{2.55cm}>{\raggedright\arraybackslash}X@{}}
\toprule
输入与状态 & 格点/解析输入 → 梯度流、涂抹与能动量张量中的算符、矩阵元或场配置；具体输入见源章节。 \\
章节路线 & 引言；一维辅助场形式与非光锥威尔逊线算符的重正化；一维辅助场形式；非光锥威尔逊线算符的重正化；梯度流与小流时展开；梯度流形式；小流时展开；计算与结果；（公式）\&（公式）；（公式）；（公式）；……准PDF算符；威尔逊线自能图（公式）；顶点图（公式）；图（公式）；最终结果 \\
方法摘录 & 路径有序相位因子（公式），即所谓威尔逊线 在本工作中，我们只关注非光锥威尔逊线与威尔逊圈；只要不产生歧义，就把它们简称为威尔逊线。 ，是规范理论中的一个基本对象，它在把局域场连接起来以构造规范不变算符方面起着关键作用。对威尔逊线和威尔逊圈重正化……详细匹配计算。我们将这些结果与文献（引文）的已有结果以及第（交叉引用）节在小流时极限下的结果进行比较 \\
结构式/示意 & $\partial_t B_\mu(t,x)=D_\nu G_{\nu\mu}(t,x),\qquad t\sim L_{\rm sm}^{2}$\\[-0.2em]\scriptsize 该式只表达主题变量关系；论文特定推导以源文件为准。 \\
物理校验 & 流时间与平滑长度满足平方关系；t→0 回到原始场，有限 t 则抑制紫外模式。 \\
\bottomrule
\end{tabularx}
\endgroup
  \source{方法证据：\detokenize{refer/papers/梯度流中的非光锥威尔逊线算符_latex/chapters/section01.tex:1-26}；主题结构式为示意。}
\end{frame}

%% frame_manifest|frame_id=P40-03|paper_id=P40|claim_id=P40-结果|visual_type=result-evidence-table|width_budget=0.98linewidth|height_budget=0.86textheight|min_font_pt=10|split_allowed=no|expected_page=172|risks=title-table-source-footline
\begin{frame}[t]{梯度流中的非光锥威尔逊线算符：结果、验证与物理含义}
  \fontsize{10}{10.8}\selectfont
\begingroup
\renewcommand{\arraystretch}{0.84}
\begin{tabularx}{0.98\linewidth}{@{}p{2.55cm}>{\raggedright\arraybackslash}X@{}}
\toprule
源文结果 & sec:summary 我们发展了在小流时极限下从梯度流方案匹配到（公式）方案、适用于威尔逊线算符的系统方法。该方法可推广到多种应用，包括准PDF、pNRQCD 框架下夸克偶素衰变与产生中出现的胶子关联函数，以及以色磁场插入威尔逊圈表述的自旋依赖势。 在一维辅助场形式中，本工作所讨论的威尔逊线算符从梯度流方案到（公式）方案的匹配，在小流时极限下可以约化为对相应局域流算符的匹配。……从梯度流方案匹配到（公式）方案的系统方法， 把当前研究推广到两圈阶是直截了当的，这将更有助于用非光锥威尔逊线算符构造的矩阵元的格点梯度流计算。我们把它留作未来的研究 \\
验证状态 & 确证：源章节有结果/讨论摘录。 \\
库内可见证据 & 中文源约 1392 行；公式命中 99，图命中 14，表/表格命中 0。 \\
物理含义 & 在本主题共同链条中，本论文把上述输入推进到可比较的中间量或观测量；具体强度、误差和外推范围不超出源文档。 \\
未验证项 & 当前工作区没有论文 PDF、原始数据和独立复现脚本；因此数值复核、图像再现和统计显著性不作额外断言。 \\
\bottomrule
\end{tabularx}
\endgroup
  \source{结果证据：\detokenize{refer/papers/梯度流中的非光锥威尔逊线算符_latex/chapters/section05.tex:1-25}；图表/公式计数由本报告生成器对中文源的静态统计得到。}
\end{frame}

%% frame_manifest|frame_id=P40-04|paper_id=P40|claim_id=P40-局限|visual_type=risk-relation-table|width_budget=0.98linewidth|height_budget=0.86textheight|min_font_pt=10|split_allowed=no|expected_page=173|risks=title-table-source-footline
\begin{frame}[t]{梯度流中的非光锥威尔逊线算符：局限、跨论文接口与可复查入口}
  \fontsize{10}{10.8}\selectfont
\begingroup
\renewcommand{\arraystretch}{0.84}
\begin{tabularx}{0.98\linewidth}{@{}p{2.55cm}>{\raggedright\arraybackslash}X@{}}
\toprule
源文局限 & 需要回到结果/讨论/展望章节核对适用区间；本档案不把缺失的独立数据复核补成已证实结论。 \\
跨论文接口 & 把连续极限、微扰有限性和能动量张量重构连接到可计算的格点观测量。 \\
极限检查 & 流时间与平滑长度满足平方关系；t→0 回到原始场，有限 t 则抑制紫外模式。 \\
下一步核验 & 补齐 PDF 或原始数据；按源文参数复现关键图表；比较不同动量、流时间、距离、格距或方案转换后的稳定性。 \\
完整来源 & \detokenize{refer/papers/梯度流中的非光锥威尔逊线算符_latex/main.tex}；\detokenize{refer/papers/Off_LightCone_Wilson_Line_Flow_latex/main.tex}；章节源：\detokenize{refer/papers/梯度流中的非光锥威尔逊线算符_latex/chapters/*.tex}；英中目录：\detokenize{refer/papers/Off_LightCone_Wilson_Line_Flow_latex}. \\
\bottomrule
\end{tabularx}
\endgroup
  \source{状态：\inferred 为主题接口推断；\unverified 为当前材料边界；完整文件清单见 manifest。}
\end{frame}

%% frame_manifest|frame_id=P41-01|paper_id=P41|claim_id=P41-定位|visual_type=paper-card|width_budget=0.98linewidth|height_budget=0.86textheight|min_font_pt=10|split_allowed=no|expected_page=174|risks=title-table-source-footline
\section{首个核子胶子部分子分布}
\begin{frame}[t]{首个核子胶子部分子分布：研究问题与论文定位}
  \fontsize{10}{10.8}\selectfont
\begingroup
\renewcommand{\arraystretch}{0.84}
\begin{tabularx}{0.98\linewidth}{@{}p{2.55cm}>{\raggedright\arraybackslash}X@{}}
\toprule
论文编号 & P41；主题：非微扰重正化、OPE 与匹配 \\
书目信息 & First Nucleon Gluon PDF from Large Momentum Effective Theory；arXiv:2505.13321；英/中转排页数 16/15 \\
研究对象 & 我们报告了利用大动量有效理论（Large-Momentum Effective Theory, LaMET）得到的首个核子胶子部分子分布函数（PDF）。 我们聚焦于前人工作（引文）中已被证明在用 LaMET 计算胶子 PDF 时具……量与规范链涂抹带来的物理效应，本计算提供了极佳的原理性验证，并与若干选定的全局拟合结果合理相符 \\
章节证据 & 引言；LaMET 与混合重正化的表述；数值结果与讨论；结论与展望；外推拟合范围效应；大-（公式）外推 \\
证据状态 & \verified：摘要和章节来自现有 LaTeX；\unverified：没有论文 PDF/原始数据独立复核。 \\
\bottomrule
\end{tabularx}
\endgroup
  \source{索引：\detokenize{refer/papers/INDEX.md:41}；正文：\detokenize{refer/papers/首个核子胶子部分子分布_latex/main.tex}；\detokenize{refer/papers/First_Nucleon_Gluon_PDF_LaMET_latex/main.tex}}
\end{frame}

%% frame_manifest|frame_id=P41-02|paper_id=P41|claim_id=P41-方法|visual_type=method-flow-formula|width_budget=0.98linewidth|height_budget=0.86textheight|min_font_pt=10|split_allowed=no|expected_page=175|risks=title-table-source-footline
\begin{frame}[t]{首个核子胶子部分子分布：理论结构与方法链}
  \fontsize{10}{10.8}\selectfont
\begingroup
\renewcommand{\arraystretch}{0.84}
\begin{tabularx}{0.98\linewidth}{@{}p{2.55cm}>{\raggedright\arraybackslash}X@{}}
\toprule
输入与状态 & 格点/解析输入 → 非微扰重正化、OPE 与匹配中的算符、矩阵元或场配置；具体输入见源章节。 \\
章节路线 & 引言；LaMET 与混合重正化的表述；数值结果与讨论；结论与展望；外推拟合范围效应；大-（公式）外推 \\
方法摘录 & sec:intro 部分子分布函数（PDF）在高能粒子物理中扮演着基础性角色，它编码了质子与中子内部夸克和胶子的动量分布。 这些函数是在诸如大型强子对撞机（LHC）等强子对撞机上作出精确预言的重要输入，直接影响从标准模型（SM）到超越标准模型（……子 PDF，并与唯象拟合进行比较。 我们在 Sec.（交叉引用）中总结并展望胶子 PDF 的未来前景 \\
结构式/示意 & $O^{R}(z,\mu)=Z(z,\mu,a)\,O^{\rm bare}(z,a)$\\[-0.2em]\scriptsize 该式只表达主题变量关系；论文特定推导以源文件为准。 \\
物理校验 & 重正化因子抵消截止依赖；a→0 和短距离 z→0 极限必须与匹配阶数、方案和尺度保持一致。 \\
\bottomrule
\end{tabularx}
\endgroup
  \source{方法证据：\detokenize{refer/papers/首个核子胶子部分子分布_latex/chapters/section01.tex:1-36}；主题结构式为示意。}
\end{frame}

%% frame_manifest|frame_id=P41-03|paper_id=P41|claim_id=P41-结果|visual_type=result-evidence-table|width_budget=0.98linewidth|height_budget=0.86textheight|min_font_pt=10|split_allowed=no|expected_page=176|risks=title-table-source-footline
\begin{frame}[t]{首个核子胶子部分子分布：结果、验证与物理含义}
  \fontsize{10}{10.8}\selectfont
\begingroup
\renewcommand{\arraystretch}{0.84}
\begin{tabularx}{0.98\linewidth}{@{}p{2.55cm}>{\raggedright\arraybackslash}X@{}}
\toprule
源文结果 & sec:conclusion 我们展示了通过 LaMET 并结合混合比值重正化计算得到的首个核子胶子 PDF，采用两个重于物理值的$\pi$介子质量（（公式）与（公式）MeV），以及两种规范链涂抹技术：流时间（公式）的 Wilson 流与 5 步超立方涂抹。 基于迈向 LaMET 胶子 PDF 的先前工作（引文），我们聚焦于单一胶子算符（公式）（定义于文献（引文）式 6），该算符已被用……。 这将使我们能更好地理解离散化与涂抹效应，并通过取物理连续极限改善胶子 PDF 的系统误差。 随着误差变小，还需要理解标度变化对 LaMET 胶子 PDF 的影响 \\
验证状态 & 确证：源章节有结果/讨论摘录。 \\
库内可见证据 & 中文源约 694 行；公式命中 11，图命中 26，表/表格命中 0。 \\
物理含义 & 在本主题共同链条中，本论文把上述输入推进到可比较的中间量或观测量；具体强度、误差和外推范围不超出源文档。 \\
未验证项 & 当前工作区没有论文 PDF、原始数据和独立复现脚本；因此数值复核、图像再现和统计显著性不作额外断言。 \\
\bottomrule
\end{tabularx}
\endgroup
  \source{结果证据：\detokenize{refer/papers/首个核子胶子部分子分布_latex/chapters/section04.tex:1-19}；图表/公式计数由本报告生成器对中文源的静态统计得到。}
\end{frame}

%% frame_manifest|frame_id=P41-04|paper_id=P41|claim_id=P41-局限|visual_type=risk-relation-table|width_budget=0.98linewidth|height_budget=0.86textheight|min_font_pt=10|split_allowed=no|expected_page=177|risks=title-table-source-footline
\begin{frame}[t]{首个核子胶子部分子分布：局限、跨论文接口与可复查入口}
  \fontsize{10}{10.8}\selectfont
\begingroup
\renewcommand{\arraystretch}{0.84}
\begin{tabularx}{0.98\linewidth}{@{}p{2.55cm}>{\raggedright\arraybackslash}X@{}}
\toprule
源文局限 & 需要回到结果/讨论/展望章节核对适用区间；本档案不把缺失的独立数据复核补成已证实结论。 \\
跨论文接口 & 是准/赝分布从格点数据走向 MS-bar 或其他连续方案的必要中间层。 索引备注：Good-Yao-Lin；与库内PDF同文已核 \\
极限检查 & 重正化因子抵消截止依赖；a→0 和短距离 z→0 极限必须与匹配阶数、方案和尺度保持一致。 \\
下一步核验 & 补齐 PDF 或原始数据；按源文参数复现关键图表；比较不同动量、流时间、距离、格距或方案转换后的稳定性。 \\
完整来源 & \detokenize{refer/papers/首个核子胶子部分子分布_latex/main.tex}；\detokenize{refer/papers/First_Nucleon_Gluon_PDF_LaMET_latex/main.tex}；章节源：\detokenize{refer/papers/首个核子胶子部分子分布_latex/chapters/*.tex}；英中目录：\detokenize{refer/papers/First_Nucleon_Gluon_PDF_LaMET_latex}. \\
\bottomrule
\end{tabularx}
\endgroup
  \source{状态：\inferred 为主题接口推断；\unverified 为当前材料边界；完整文件清单见 manifest。}
\end{frame}

%% frame_manifest|frame_id=P42-01|paper_id=P42|claim_id=P42-定位|visual_type=paper-card|width_budget=0.98linewidth|height_budget=0.86textheight|min_font_pt=10|split_allowed=no|expected_page=178|risks=title-table-source-footline
\section{$\pi$介子胶子部分子分布}
\begin{frame}[t]{$\pi$介子胶子部分子分布：研究问题与论文定位}
  \fontsize{10}{10.8}\selectfont
\begingroup
\renewcommand{\arraystretch}{0.84}
\begin{tabularx}{0.98\linewidth}{@{}p{2.55cm}>{\raggedright\arraybackslash}X@{}}
\toprule
论文编号 & P42；主题：胶子、强子部分子分布与 TMD \\
书目信息 & Gluon PDF Pion Lattice；arXiv:2104.06372；英/中转排页数 18/17 \\
研究对象 & 我们利用赝PDF方法首次由格点QCD确定了（公式）介子胶子分布的（公式）依赖。我们使用 MILC 合作组产生的、具有2+1+1味高度改进 staggered 夸克（HISQ）的格点系综，涵盖两种晶格间距（公式）与0.15 fm、三……。我们将最轻（公式）介子质量系综的结果与 JAM 及 xFitter 全局拟合的测定进行了比较 \\
章节证据 & 引言；用赝PDF方法由格点计算得到胶子PDF；结果与讨论；总结 \\
证据状态 & \verified：摘要和章节来自现有 LaTeX；\unverified：没有论文 PDF/原始数据独立复核。 \\
\bottomrule
\end{tabularx}
\endgroup
  \source{索引：\detokenize{refer/papers/INDEX.md:42}；正文：\detokenize{refer/papers/}$\pi$\detokenize{介子胶子部分子分布_latex/main.tex}；\detokenize{refer/papers/Gluon_PDF_Pion_Lattice_latex/main.tex}}
\end{frame}

%% frame_manifest|frame_id=P42-02|paper_id=P42|claim_id=P42-方法|visual_type=method-flow-formula|width_budget=0.98linewidth|height_budget=0.86textheight|min_font_pt=10|split_allowed=no|expected_page=179|risks=title-table-source-footline
\begin{frame}[t]{$\pi$介子胶子部分子分布：理论结构与方法链}
  \fontsize{10}{10.8}\selectfont
\begingroup
\renewcommand{\arraystretch}{0.84}
\begin{tabularx}{0.98\linewidth}{@{}p{2.55cm}>{\raggedright\arraybackslash}X@{}}
\toprule
输入与状态 & 格点/解析输入 → 胶子、强子部分子分布与 TMD中的算符、矩阵元或场配置；具体输入见源章节。 \\
章节路线 & 引言；用赝PDF方法由格点计算得到胶子PDF；结果与讨论；总结 \\
方法摘录 & sec:cal-details 本工作中，我们使用文献（引文）定义的非极化胶子算符，（公式或图表位置）其中算符（公式），（公式）为 Wilson 线长度，场强（公式）定义为（公式或图表位置）其中（公式）是晶格间距，（公式）是强耦合常数， pla……，我们从 two-sim 拟合得到的基态RITD是稳定的，并用它们提取胶子PDF。（公式或图表位置） \\
结构式/示意 & $\mathcal O_{\rm Euclid}\xrightarrow[\rm matching\,/\,evolution]{}f_{g/h}(x,\mu)\;{\rm or}\;F_{\rm TMD}(x,b_T,\mu,\zeta)$\\[-0.2em]\scriptsize 该式只表达主题变量关系；论文特定推导以源文件为准。 \\
物理校验 & 纵向分数、横向距离和重正化尺度是不同变量；不能把有限动量或有限 b\_T 的结果直接当作光锥分布。 \\
\bottomrule
\end{tabularx}
\endgroup
  \source{方法证据：\detokenize{refer/papers/}$\pi$\detokenize{介子胶子部分子分布_latex/chapters/section02.tex:1-219}；主题结构式为示意。}
\end{frame}

%% frame_manifest|frame_id=P42-03|paper_id=P42|claim_id=P42-结果|visual_type=result-evidence-table|width_budget=0.98linewidth|height_budget=0.86textheight|min_font_pt=10|split_allowed=no|expected_page=180|risks=title-table-source-footline
\begin{frame}[t]{$\pi$介子胶子部分子分布：结果、验证与物理含义}
  \fontsize{10}{10.8}\selectfont
\begingroup
\renewcommand{\arraystretch}{0.84}
\begin{tabularx}{0.98\linewidth}{@{}p{2.55cm}>{\raggedright\arraybackslash}X@{}}
\toprule
源文结果 & sec:summary 本工作中，我们用赝PDF方法首次从格点QCD计算了（公式）介子胶子PDF，并研究了其（公式）介子质量与晶格间距依赖。 我们在具有（公式）味高度改进 staggered 夸克（HISQ）的系综上使用 clover 价费米子，涵盖两种晶格间距（（公式）与0.15 fm）和三种（公式）介子质量（220、310与690 MeV）。 这些系综使我们得以考察（公式）……）。 未来精度更高、系统误差控制更好的格点QCD（公式）介子胶子PDF研究，与预期实验（引文）的结果一道纳入全局拟合分析时，将为（公式）介子内的胶子成分提供最佳测定 \\
验证状态 & 确证：源章节有结果/讨论摘录。 \\
库内可见证据 & 中文源约 1890 行；公式命中 13，图命中 17，表/表格命中 2。 \\
物理含义 & 在本主题共同链条中，本论文把上述输入推进到可比较的中间量或观测量；具体强度、误差和外推范围不超出源文档。 \\
未验证项 & 当前工作区没有论文 PDF、原始数据和独立复现脚本；因此数值复核、图像再现和统计显著性不作额外断言。 \\
\bottomrule
\end{tabularx}
\endgroup
  \source{结果证据：\detokenize{refer/papers/}$\pi$\detokenize{介子胶子部分子分布_latex/chapters/section04.tex:1-20}；图表/公式计数由本报告生成器对中文源的静态统计得到。}
\end{frame}

%% frame_manifest|frame_id=P42-04|paper_id=P42|claim_id=P42-局限|visual_type=risk-relation-table|width_budget=0.98linewidth|height_budget=0.86textheight|min_font_pt=10|split_allowed=no|expected_page=181|risks=title-table-source-footline
\begin{frame}[t]{$\pi$介子胶子部分子分布：局限、跨论文接口与可复查入口}
  \fontsize{10}{10.8}\selectfont
\begingroup
\renewcommand{\arraystretch}{0.84}
\begin{tabularx}{0.98\linewidth}{@{}p{2.55cm}>{\raggedright\arraybackslash}X@{}}
\toprule
源文局限 & 需要回到结果/讨论/展望章节核对适用区间；本档案不把缺失的独立数据复核补成已证实结论。 \\
跨论文接口 & 检验 LaMET/赝 PDF/准 TMD 方法是否能够覆盖真实强子结构与实验拟合所需的观测量。 \\
极限检查 & 纵向分数、横向距离和重正化尺度是不同变量；不能把有限动量或有限 b\_T 的结果直接当作光锥分布。 \\
下一步核验 & 补齐 PDF 或原始数据；按源文参数复现关键图表；比较不同动量、流时间、距离、格距或方案转换后的稳定性。 \\
完整来源 & \detokenize{refer/papers/}$\pi$\detokenize{介子胶子部分子分布_latex/main.tex}；\detokenize{refer/papers/Gluon_PDF_Pion_Lattice_latex/main.tex}；章节源：\detokenize{refer/papers/}$\pi$\detokenize{介子胶子部分子分布_latex/chapters/*.tex}；英中目录：\detokenize{refer/papers/Gluon_PDF_Pion_Lattice_latex}. \\
\bottomrule
\end{tabularx}
\endgroup
  \source{状态：\inferred 为主题接口推断；\unverified 为当前材料边界；完整文件清单见 manifest。}
\end{frame}

%% frame_manifest|frame_id=P43-01|paper_id=P43|claim_id=P43-定位|visual_type=paper-card|width_budget=0.98linewidth|height_budget=0.86textheight|min_font_pt=10|split_allowed=no|expected_page=182|risks=title-table-source-footline
\section{K介子胶子分布初探}
\begin{frame}[t]{K介子胶子分布初探：研究问题与论文定位}
  \fontsize{10}{10.8}\selectfont
\begingroup
\renewcommand{\arraystretch}{0.84}
\begin{tabularx}{0.98\linewidth}{@{}p{2.55cm}>{\raggedright\arraybackslash}X@{}}
\toprule
论文编号 & P43；主题：胶子、强子部分子分布与 TMD \\
书目信息 & First Glimpse into the Kaon Gluon Parton Distribution Using Lattice QCD Preprint MSUHEP-21-033; arXiv:2112.03124 [hep-lat]; DOI: 10.1103/PhysRevD.106.094510.；arXiv:2112.03124；英/中转排页数 14/14 \\
研究对象 & 在本工作中，我们给出了由格点量子色动力学得到的 K 介子胶子部分子分布的首批结果。 我们在（公式）介子质量约 310 MeV、两个格距（0.15 与 0.12 fm）下开展格点计算，采用由 MILC 合作组生成的（公式）味 HIS……）方案、（公式）GeV 处提取了 K 介子的胶子分布函数。 在较小格距处的结果与唯象测定相一致 \\
章节证据 & 引言；格点矩阵元；由赝 PDF 方法提取 K 介子胶子 PDF；总结 \\
证据状态 & \verified：摘要和章节来自现有 LaTeX；\unverified：没有论文 PDF/原始数据独立复核。 \\
\bottomrule
\end{tabularx}
\endgroup
  \source{索引：\detokenize{refer/papers/INDEX.md:43}；正文：\detokenize{refer/papers/K介子胶子分布初探_latex/main.tex}；\detokenize{refer/papers/First_Glimpse_Kaon_Gluon_PDF_latex/main.tex}}
\end{frame}

%% frame_manifest|frame_id=P43-02|paper_id=P43|claim_id=P43-方法|visual_type=method-flow-formula|width_budget=0.98linewidth|height_budget=0.86textheight|min_font_pt=10|split_allowed=no|expected_page=183|risks=title-table-source-footline
\begin{frame}[t]{K介子胶子分布初探：理论结构与方法链}
  \fontsize{10}{10.8}\selectfont
\begingroup
\renewcommand{\arraystretch}{0.84}
\begin{tabularx}{0.98\linewidth}{@{}p{2.55cm}>{\raggedright\arraybackslash}X@{}}
\toprule
输入与状态 & 格点/解析输入 → 胶子、强子部分子分布与 TMD中的算符、矩阵元或场配置；具体输入见源章节。 \\
章节路线 & 引言；格点矩阵元；由赝 PDF 方法提取 K 介子胶子 PDF；总结 \\
方法摘录 & sec:lattice-details 对于 K 介子基态（公式），我们在若干助推动量（公式）与 Wilson 线位移长度（公式）下计算其格点胶子矩阵元，所用非极化胶子算符最早由文献（引文）引入：（公式或图表位置）其中（公式），Wilson 链……公式）的选取较为敏感。 分析中所采用的（公式）确实能可靠地提取基态。（公式或图表位置）[公式/图表] \\
结构式/示意 & $\mathcal O_{\rm Euclid}\xrightarrow[\rm matching\,/\,evolution]{}f_{g/h}(x,\mu)\;{\rm or}\;F_{\rm TMD}(x,b_T,\mu,\zeta)$\\[-0.2em]\scriptsize 该式只表达主题变量关系；论文特定推导以源文件为准。 \\
物理校验 & 纵向分数、横向距离和重正化尺度是不同变量；不能把有限动量或有限 b\_T 的结果直接当作光锥分布。 \\
\bottomrule
\end{tabularx}
\endgroup
  \source{方法证据：\detokenize{refer/papers/K介子胶子分布初探_latex/chapters/section02.tex:1-147}；主题结构式为示意。}
\end{frame}

%% frame_manifest|frame_id=P43-03|paper_id=P43|claim_id=P43-结果|visual_type=result-evidence-table|width_budget=0.98linewidth|height_budget=0.86textheight|min_font_pt=10|split_allowed=no|expected_page=184|risks=title-table-source-footline
\begin{frame}[t]{K介子胶子分布初探：结果、验证与物理含义}
  \fontsize{10}{10.8}\selectfont
\begingroup
\renewcommand{\arraystretch}{0.84}
\begin{tabularx}{0.98\linewidth}{@{}p{2.55cm}>{\raggedright\arraybackslash}X@{}}
\toprule
源文结果 & sec:summary 在本工作中，我们利用赝 PDF 方法首次对 K 介子胶子部分子分布进行了格点研究。 我们以 clover 费米子为价作用，使用 MILC 合作组生成的（公式）介子质量 310 MeV、两个格距（0.15 与 0.12 fm）的 2+1+1 味 HISQ 规范组态。 我们采用动量涂抹与高统计测量（最多达 O(324,000) 次），使 K 介子的助推动量达……限体积效应与非物理（公式）介子质量效应，未包含在这项开创性研究中。 未来研究应致力于改进这些系统误差，并为即将开展的实验工作提供对 K 介子胶子 PDF 更好的测定 \\
验证状态 & 确证：源章节有结果/讨论摘录。 \\
库内可见证据 & 中文源约 1251 行；公式命中 11，图命中 23，表/表格命中 2。 \\
物理含义 & 在本主题共同链条中，本论文把上述输入推进到可比较的中间量或观测量；具体强度、误差和外推范围不超出源文档。 \\
未验证项 & 当前工作区没有论文 PDF、原始数据和独立复现脚本；因此数值复核、图像再现和统计显著性不作额外断言。 \\
\bottomrule
\end{tabularx}
\endgroup
  \source{结果证据：\detokenize{refer/papers/K介子胶子分布初探_latex/chapters/section04.tex:1-14}；图表/公式计数由本报告生成器对中文源的静态统计得到。}
\end{frame}

%% frame_manifest|frame_id=P43-04|paper_id=P43|claim_id=P43-局限|visual_type=risk-relation-table|width_budget=0.98linewidth|height_budget=0.86textheight|min_font_pt=10|split_allowed=no|expected_page=185|risks=title-table-source-footline
\begin{frame}[t]{K介子胶子分布初探：局限、跨论文接口与可复查入口}
  \fontsize{10}{10.8}\selectfont
\begingroup
\renewcommand{\arraystretch}{0.84}
\begin{tabularx}{0.98\linewidth}{@{}p{2.55cm}>{\raggedright\arraybackslash}X@{}}
\toprule
源文局限 & 需要回到结果/讨论/展望章节核对适用区间；本档案不把缺失的独立数据复核补成已证实结论。 \\
跨论文接口 & 检验 LaMET/赝 PDF/准 TMD 方法是否能够覆盖真实强子结构与实验拟合所需的观测量。 \\
极限检查 & 纵向分数、横向距离和重正化尺度是不同变量；不能把有限动量或有限 b\_T 的结果直接当作光锥分布。 \\
下一步核验 & 补齐 PDF 或原始数据；按源文参数复现关键图表；比较不同动量、流时间、距离、格距或方案转换后的稳定性。 \\
完整来源 & \detokenize{refer/papers/K介子胶子分布初探_latex/main.tex}；\detokenize{refer/papers/First_Glimpse_Kaon_Gluon_PDF_latex/main.tex}；章节源：\detokenize{refer/papers/K介子胶子分布初探_latex/chapters/*.tex}；英中目录：\detokenize{refer/papers/First_Glimpse_Kaon_Gluon_PDF_latex}. \\
\bottomrule
\end{tabularx}
\endgroup
  \source{状态：\inferred 为主题接口推断；\unverified 为当前材料边界；完整文件清单见 manifest。}
\end{frame}

%% frame_manifest|frame_id=P44-01|paper_id=P44|claim_id=P44-定位|visual_type=paper-card|width_budget=0.98linewidth|height_budget=0.86textheight|min_font_pt=10|split_allowed=no|expected_page=186|risks=title-table-source-footline
\section{从准TMD波函数确定CS核}
\begin{frame}[t]{从准TMD波函数确定CS核：研究问题与论文定位}
  \fontsize{10}{10.8}\selectfont
\begingroup
\renewcommand{\arraystretch}{0.84}
\begin{tabularx}{0.98\linewidth}{@{}p{2.55cm}>{\raggedright\arraybackslash}X@{}}
\toprule
论文编号 & P44；主题：胶子、强子部分子分布与 TMD \\
书目信息 & Nonperturbative Determination of Collins-Soper Kernel from Quasi Transverse-Momentum Dependent Wave Functions；arXiv:2204.00200；英/中转排页数 26/24 \\
研究对象 & 在大动量有效理论框架下并以一圈匹配精度，我们对支配横动量依赖（TMD）分布快度演化的 Collins-Soper 核进行了格点计算。我们首先在 MILC 的价 clover 夸克置于动力学 HISQ 海上、晶格间距（公式）fm 的……定度，以及更高幂次修正的影响。我们的结果原则上允许以一圈精度确定软函数及其他横动量依赖的物理量 \\
章节证据 & 引言；理论框架；Collins-Soper 核与快度演化；准 TMD 波函数；准 TMDWFs 的因子化；由准 TMDWFs 确定 Collins-Soper 核；数值模拟与结果；格点设置；……（公式）的依赖；准 TMDWFs 对规范链长度（公式）的依赖；幂次修正效应 \\
证据状态 & \verified：摘要和章节来自现有 LaTeX；\unverified：没有论文 PDF/原始数据独立复核。 \\
\bottomrule
\end{tabularx}
\endgroup
  \source{索引：\detokenize{refer/papers/INDEX.md:44}；正文：\detokenize{refer/papers/从准TMD波函数确定CS核_latex/main.tex}；\detokenize{refer/papers/CS_Kernel_Quasi_TMDWF_latex/main.tex}}
\end{frame}

%% frame_manifest|frame_id=P44-02|paper_id=P44|claim_id=P44-方法|visual_type=method-flow-formula|width_budget=0.98linewidth|height_budget=0.86textheight|min_font_pt=10|split_allowed=no|expected_page=187|risks=title-table-source-footline
\begin{frame}[t]{从准TMD波函数确定CS核：理论结构与方法链}
  \fontsize{10}{10.8}\selectfont
\begingroup
\renewcommand{\arraystretch}{0.84}
\begin{tabularx}{0.98\linewidth}{@{}p{2.55cm}>{\raggedright\arraybackslash}X@{}}
\toprule
输入与状态 & 格点/解析输入 → 胶子、强子部分子分布与 TMD中的算符、矩阵元或场配置；具体输入见源章节。 \\
章节路线 & 引言；理论框架；Collins-Soper 核与快度演化；准 TMD 波函数；准 TMDWFs 的因子化；由准 TMDWFs 确定 Collins-Soper 核；数值模拟与结果；格点设置；……（公式）的依赖；准 TMDWFs 对规范链长度（公式）的依赖；幂次修正效应 \\
方法摘录 & sec:framework 在本节中，我们回顾本次计算所需的必要理论背景。我们给出 CS 核与快度演化的定义，并引入准 TMD 波函数。随后我们讨论准 TMDWFs 的因子化 及其与 CS 核的联系 \\
结构式/示意 & $\mathcal O_{\rm Euclid}\xrightarrow[\rm matching\,/\,evolution]{}f_{g/h}(x,\mu)\;{\rm or}\;F_{\rm TMD}(x,b_T,\mu,\zeta)$\\[-0.2em]\scriptsize 该式只表达主题变量关系；论文特定推导以源文件为准。 \\
物理校验 & 纵向分数、横向距离和重正化尺度是不同变量；不能把有限动量或有限 b\_T 的结果直接当作光锥分布。 \\
\bottomrule
\end{tabularx}
\endgroup
  \source{方法证据：\detokenize{refer/papers/从准TMD波函数确定CS核_latex/chapters/section02.tex:1-8}；主题结构式为示意。}
\end{frame}

%% frame_manifest|frame_id=P44-03|paper_id=P44|claim_id=P44-结果|visual_type=result-evidence-table|width_budget=0.98linewidth|height_budget=0.86textheight|min_font_pt=10|split_allowed=no|expected_page=188|risks=title-table-source-footline
\begin{frame}[t]{从准TMD波函数确定CS核：结果、验证与物理含义}
  \fontsize{10}{10.8}\selectfont
\begingroup
\renewcommand{\arraystretch}{0.84}
\begin{tabularx}{0.98\linewidth}{@{}p{2.55cm}>{\raggedright\arraybackslash}X@{}}
\toprule
源文结果 & sec:summary 在本工作中，我们在大动量有效理论框架下于 MILC 格点组态上计算了 CS 核。与我们先前的研究（引文）相比，本研究采用了一圈匹配 kernel，并使用了多个强子动量来提取 CS 核。我们发现，在小（公式）区域，我们的结果与微扰 QCD 一致。 在大（公式）区域，我们的结果在不确定度范围内似乎与文献中其他格点计算相一致。 在今后的研究中，我们需要 使用具……海夸克一致的价夸克质量以减小非幺正性效应。这类效应之一 可能是介子 波函数的虚部——目前它与微扰计算的 结果看起来不一致。显然，所有这些探索都将 耗费更多的计算资源 \\
验证状态 & 确证：源章节有结果/讨论摘录。 \\
库内可见证据 & 中文源约 913 行；公式命中 33，图命中 52，表/表格命中 0。 \\
物理含义 & 在本主题共同链条中，本论文把上述输入推进到可比较的中间量或观测量；具体强度、误差和外推范围不超出源文档。 \\
未验证项 & 当前工作区没有论文 PDF、原始数据和独立复现脚本；因此数值复核、图像再现和统计显著性不作额外断言。 \\
\bottomrule
\end{tabularx}
\endgroup
  \source{结果证据：\detokenize{refer/papers/从准TMD波函数确定CS核_latex/chapters/section04.tex:1-16}；图表/公式计数由本报告生成器对中文源的静态统计得到。}
\end{frame}

%% frame_manifest|frame_id=P44-04|paper_id=P44|claim_id=P44-局限|visual_type=risk-relation-table|width_budget=0.98linewidth|height_budget=0.86textheight|min_font_pt=10|split_allowed=no|expected_page=189|risks=title-table-source-footline
\begin{frame}[t]{从准TMD波函数确定CS核：局限、跨论文接口与可复查入口}
  \fontsize{10}{10.8}\selectfont
\begingroup
\renewcommand{\arraystretch}{0.84}
\begin{tabularx}{0.98\linewidth}{@{}p{2.55cm}>{\raggedright\arraybackslash}X@{}}
\toprule
源文局限 & 需要回到结果/讨论/展望章节核对适用区间；本档案不把缺失的独立数据复核补成已证实结论。 \\
跨论文接口 & 检验 LaMET/赝 PDF/准 TMD 方法是否能够覆盖真实强子结构与实验拟合所需的观测量。 \\
极限检查 & 纵向分数、横向距离和重正化尺度是不同变量；不能把有限动量或有限 b\_T 的结果直接当作光锥分布。 \\
下一步核验 & 补齐 PDF 或原始数据；按源文参数复现关键图表；比较不同动量、流时间、距离、格距或方案转换后的稳定性。 \\
完整来源 & \detokenize{refer/papers/从准TMD波函数确定CS核_latex/main.tex}；\detokenize{refer/papers/CS_Kernel_Quasi_TMDWF_latex/main.tex}；章节源：\detokenize{refer/papers/从准TMD波函数确定CS核_latex/chapters/*.tex}；英中目录：\detokenize{refer/papers/CS_Kernel_Quasi_TMDWF_latex}. \\
\bottomrule
\end{tabularx}
\endgroup
  \source{状态：\inferred 为主题接口推断；\unverified 为当前材料边界；完整文件清单见 manifest。}
\end{frame}

%% frame_manifest|frame_id=P45-01|paper_id=P45|claim_id=P45-定位|visual_type=paper-card|width_budget=0.98linewidth|height_budget=0.86textheight|min_font_pt=10|split_allowed=no|expected_page=190|risks=title-table-source-footline
\section{LaMET计算TMD软函数}
\begin{frame}[t]{LaMET计算TMD软函数：研究问题与论文定位}
  \fontsize{10}{10.8}\selectfont
\begingroup
\renewcommand{\arraystretch}{0.84}
\begin{tabularx}{0.98\linewidth}{@{}p{2.55cm}>{\raggedright\arraybackslash}X@{}}
\toprule
论文编号 & P45；主题：胶子、强子部分子分布与 TMD \\
书目信息 & Lattice-QCD Calculations of TMD Soft Function Through Large-Momentum Effective Theory；arXiv:2005.14572；英/中转排页数 16/14 \\
研究对象 & 横向动量依赖（TMD）软函数是 Drell-Yan 及其他具有较小横向动量的过程中 QCD 因子化的关键要素。我们在一个（公式）fm 的 2+1 味 CLS 动力学系综上，对该函数在中等偏大的快度区域进行了格点 QCD 研究。我们……MD 波函数的大动量演化，研究了软函数中与快度相关的部分——Collins-Soper 演化核 \\
章节证据 & 引言；理论框架；模拟设置；数值结果；总结与展望；补充材料；模拟检验；TMDWF 的（公式）依赖；TMDWF 算符混合的估计；内禀软函数的列表结果；形状因子的两态拟合；提取 Collins-Soper 核时可能出现的虚部 \\
证据状态 & \verified：摘要和章节来自现有 LaTeX；\unverified：没有论文 PDF/原始数据独立复核。 \\
\bottomrule
\end{tabularx}
\endgroup
  \source{索引：\detokenize{refer/papers/INDEX.md:45}；正文：\detokenize{refer/papers/LaMET计算TMD软函数_latex/main.tex}；\detokenize{refer/papers/TMD_Soft_Function_LaMET_latex/main.tex}}
\end{frame}

%% frame_manifest|frame_id=P45-02|paper_id=P45|claim_id=P45-方法|visual_type=method-flow-formula|width_budget=0.98linewidth|height_budget=0.86textheight|min_font_pt=10|split_allowed=no|expected_page=191|risks=title-table-source-footline
\begin{frame}[t]{LaMET计算TMD软函数：理论结构与方法链}
  \fontsize{10}{10.8}\selectfont
\begingroup
\renewcommand{\arraystretch}{0.84}
\begin{tabularx}{0.98\linewidth}{@{}p{2.55cm}>{\raggedright\arraybackslash}X@{}}
\toprule
输入与状态 & 格点/解析输入 → 胶子、强子部分子分布与 TMD中的算符、矩阵元或场配置；具体输入见源章节。 \\
章节路线 & 引言；理论框架；模拟设置；数值结果；总结与展望；补充材料；模拟检验；TMDWF 的（公式）依赖；TMDWF 算符混合的估计；内禀软函数的列表结果；形状因子的两态拟合；提取 Collins-Soper 核时可能出现的虚部 \\
方法摘录 & 内禀软函数（（公式））可以从一个非单态轻赝标介子大动量形状因子的 QCD 因子化中得到，该介子由（公式）组成，跃迁流由两个夸克双线性算符构成，其横向间隔固定为（公式），（公式或图表位置）这里（公式）是不同味的轻夸克场，（公式）。 为了提取软因子……；而式 (（交叉引用）) 与 (（交叉引用）) 是这项开创性工作所用的领头阶近似。（公式或图表位置） \\
结构式/示意 & $\mathcal O_{\rm Euclid}\xrightarrow[\rm matching\,/\,evolution]{}f_{g/h}(x,\mu)\;{\rm or}\;F_{\rm TMD}(x,b_T,\mu,\zeta)$\\[-0.2em]\scriptsize 该式只表达主题变量关系；论文特定推导以源文件为准。 \\
物理校验 & 纵向分数、横向距离和重正化尺度是不同变量；不能把有限动量或有限 b\_T 的结果直接当作光锥分布。 \\
\bottomrule
\end{tabularx}
\endgroup
  \source{方法证据：\detokenize{refer/papers/LaMET计算TMD软函数_latex/chapters/section02.tex:1-82}；主题结构式为示意。}
\end{frame}

%% frame_manifest|frame_id=P45-03|paper_id=P45|claim_id=P45-结果|visual_type=result-evidence-table|width_budget=0.98linewidth|height_budget=0.86textheight|min_font_pt=10|split_allowed=no|expected_page=192|risks=title-table-source-footline
\begin{frame}[t]{LaMET计算TMD软函数：结果、验证与物理含义}
  \fontsize{10}{10.8}\selectfont
\begingroup
\renewcommand{\arraystretch}{0.84}
\begin{tabularx}{0.98\linewidth}{@{}p{2.55cm}>{\raggedright\arraybackslash}X@{}}
\toprule
源文结果 & 本工作通过模拟四夸克非局域算符的轻介子形状因子与准 TMD 波函数，给出了内禀软函数的一项探索性格点计算。结果显示出温和的强子动量依赖，这使得未来的精密研究可以借助微扰匹配消除大动量依赖（引文）。 作为可靠性检验，由准 TMDWF 结果得到的 CS 核与已有计算的一致性表明： 包括部分退禁闭（partially quenching）效应在内的 系统不确定度、仅用领头阶微扰匹配以及 LaMET 展开中缺失的（公式）幂次修正 可能 都是次领头的。我们的计算为在第一性原理上预言小横向动量下 Drell-Yan、希格斯产生等物理截面前进了重要一步 \\
验证状态 & 确证：源章节有结果/讨论摘录。 \\
库内可见证据 & 中文源约 782 行；公式命中 21，图命中 38，表/表格命中 4。 \\
物理含义 & 在本主题共同链条中，本论文把上述输入推进到可比较的中间量或观测量；具体强度、误差和外推范围不超出源文档。 \\
未验证项 & 当前工作区没有论文 PDF、原始数据和独立复现脚本；因此数值复核、图像再现和统计显著性不作额外断言。 \\
\bottomrule
\end{tabularx}
\endgroup
  \source{结果证据：\detokenize{refer/papers/LaMET计算TMD软函数_latex/chapters/section05.tex:1-5}；图表/公式计数由本报告生成器对中文源的静态统计得到。}
\end{frame}

%% frame_manifest|frame_id=P45-04|paper_id=P45|claim_id=P45-局限|visual_type=risk-relation-table|width_budget=0.98linewidth|height_budget=0.86textheight|min_font_pt=10|split_allowed=no|expected_page=193|risks=title-table-source-footline
\begin{frame}[t]{LaMET计算TMD软函数：局限、跨论文接口与可复查入口}
  \fontsize{10}{10.8}\selectfont
\begingroup
\renewcommand{\arraystretch}{0.84}
\begin{tabularx}{0.98\linewidth}{@{}p{2.55cm}>{\raggedright\arraybackslash}X@{}}
\toprule
源文局限 & 需要回到结果/讨论/展望章节核对适用区间；本档案不把缺失的独立数据复核补成已证实结论。 \\
跨论文接口 & 检验 LaMET/赝 PDF/准 TMD 方法是否能够覆盖真实强子结构与实验拟合所需的观测量。 \\
极限检查 & 纵向分数、横向距离和重正化尺度是不同变量；不能把有限动量或有限 b\_T 的结果直接当作光锥分布。 \\
下一步核验 & 补齐 PDF 或原始数据；按源文参数复现关键图表；比较不同动量、流时间、距离、格距或方案转换后的稳定性。 \\
完整来源 & \detokenize{refer/papers/LaMET计算TMD软函数_latex/main.tex}；\detokenize{refer/papers/TMD_Soft_Function_LaMET_latex/main.tex}；章节源：\detokenize{refer/papers/LaMET计算TMD软函数_latex/chapters/*.tex}；英中目录：\detokenize{refer/papers/TMD_Soft_Function_LaMET_latex}. \\
\bottomrule
\end{tabularx}
\endgroup
  \source{状态：\inferred 为主题接口推断；\unverified 为当前材料边界；完整文件清单见 manifest。}
\end{frame}

%% frame_manifest|frame_id=P46-01|paper_id=P46|claim_id=P46-定位|visual_type=paper-card|width_budget=0.98linewidth|height_budget=0.86textheight|min_font_pt=10|split_allowed=no|expected_page=194|risks=title-table-source-footline
\section{扩散模型即随机量化}
\begin{frame}[t]{扩散模型即随机量化：研究问题与论文定位}
  \fontsize{10}{10.8}\selectfont
\begingroup
\renewcommand{\arraystretch}{0.84}
\begin{tabularx}{0.98\linewidth}{@{}p{2.55cm}>{\raggedright\arraybackslash}X@{}}
\toprule
论文编号 & P46；主题：机器学习采样与随机量化 \\
书目信息 & Diffusion Models as Stochastic Quantization in Lattice Field Theory；arXiv:2309.17082；英/中转排页数 27/23 \\
研究对象 & 在本工作中，我们建立了生成式扩散模型（diffusion models，DMs）与随机量化（stochastic quantization，SQ）之间的直接联系。该扩散模型通过近似由朗之万方程支配的随机过程的逆过程来实现：从先验分……展，特别是在生成大系综代价高昂的情形。 0.5em 关键词： 随机量化，朗之万动力学，扩散模型 \\
章节证据 & 引言；随机量化；Fokker–Planck 方程与观测量；朗之万模拟；扩散模型；去噪模型；作为随机量化的扩散模型；构建有效作用量；玩具模型中的有效作用量；用于格点（公式）标量场的扩散模型；组……t 结构；Skilling–Hutchinson 估计器；统计检验；接受率 \\
证据状态 & \verified：摘要和章节来自现有 LaTeX；\unverified：没有论文 PDF/原始数据独立复核。 \\
\bottomrule
\end{tabularx}
\endgroup
  \source{索引：\detokenize{refer/papers/INDEX.md:46}；正文：\detokenize{refer/papers/扩散模型即随机量化_latex/main.tex}；\detokenize{refer/papers/Diffusion_Stochastic_Quantization_latex/main.tex}}
\end{frame}

%% frame_manifest|frame_id=P46-02|paper_id=P46|claim_id=P46-方法|visual_type=method-flow-formula|width_budget=0.98linewidth|height_budget=0.86textheight|min_font_pt=10|split_allowed=no|expected_page=195|risks=title-table-source-footline
\begin{frame}[t]{扩散模型即随机量化：理论结构与方法链}
  \fontsize{10}{10.8}\selectfont
\begingroup
\renewcommand{\arraystretch}{0.84}
\begin{tabularx}{0.98\linewidth}{@{}p{2.55cm}>{\raggedright\arraybackslash}X@{}}
\toprule
输入与状态 & 格点/解析输入 → 机器学习采样与随机量化中的算符、矩阵元或场配置；具体输入见源章节。 \\
章节路线 & 引言；随机量化；Fokker–Planck 方程与观测量；朗之万模拟；扩散模型；去噪模型；作为随机量化的扩散模型；构建有效作用量；玩具模型中的有效作用量；用于格点（公式）标量场的扩散模型；组……t 结构；Skilling–Hutchinson 估计器；统计检验；接受率 \\
方法摘录 & subsec:sim 在大多数情形下，Fokker–Planck 方程无法解析求解。作为替代，人们通过对朗之万方程 式（交叉引用）离散化来数值求解该随机过程。 使用 Itô 微积分，采用时间步长（公式）的最简离散化朗之万方程为（公式或图表位置）……类理论中的应用通常称为复朗之万动力学（引文），至今仍是活跃的研究领域（引文）， 但本文不作进一步探讨 \\
结构式/示意 & $z\sim p(z),\qquad U=f_\theta(z),\qquad p_\theta(U)\;\text{与目标系综比较}$\\[-0.2em]\scriptsize 该式只表达主题变量关系；论文特定推导以源文件为准。 \\
物理校验 & 规范不变性/等变性是结构约束；采样质量须由可观测量、接受率、自相关或分布距离验证，不能只看训练损失。 \\
\bottomrule
\end{tabularx}
\endgroup
  \source{方法证据：\detokenize{refer/papers/扩散模型即随机量化_latex/chapters/section02.tex:86-112}；主题结构式为示意。}
\end{frame}

%% frame_manifest|frame_id=P46-03|paper_id=P46|claim_id=P46-结果|visual_type=result-evidence-table|width_budget=0.98linewidth|height_budget=0.86textheight|min_font_pt=10|split_allowed=no|expected_page=196|risks=title-table-source-footline
\begin{frame}[t]{扩散模型即随机量化：结果、验证与物理含义}
  \fontsize{10}{10.8}\selectfont
\begingroup
\renewcommand{\arraystretch}{0.84}
\begin{tabularx}{0.98\linewidth}{@{}p{2.55cm}>{\raggedright\arraybackslash}X@{}}
\toprule
源文结果 & sec:sum 在本工作中，我们提出了一种利用 ML 社区中流行的扩散模型生成量子场组态的新方法。我们着重指出了它与随机量化（SQ）的联系——后者基于沿虚构时间方向的随机过程来量子化场论。在 SQ 中，随机过程的漂移项是已知的，因为它由希望抽样的概率分布导出，而该分布由所考虑的量子理论决定。在 DMs 中，漂移项未知，而是在前向随机过程中从数据学得。估计出漂移项之后，训练好的 ……）动力学。众所周知，CL 在随机过程的长轨迹上可能变得不可靠。把 CL 与 DMs 结合，可以考虑相当短的轨迹，随后用所得组态训练 DM 以生成更多组态来提高统计量 \\
验证状态 & 确证：源章节有结果/讨论摘录。 \\
库内可见证据 & 中文源约 1224 行；公式命中 52，图命中 29，表/表格命中 2。 \\
物理含义 & 在本主题共同链条中，本论文把上述输入推进到可比较的中间量或观测量；具体强度、误差和外推范围不超出源文档。 \\
未验证项 & 当前工作区没有论文 PDF、原始数据和独立复现脚本；因此数值复核、图像再现和统计显著性不作额外断言。 \\
\bottomrule
\end{tabularx}
\endgroup
  \source{结果证据：\detokenize{refer/papers/扩散模型即随机量化_latex/chapters/section05.tex:1-22}；图表/公式计数由本报告生成器对中文源的静态统计得到。}
\end{frame}

%% frame_manifest|frame_id=P46-04|paper_id=P46|claim_id=P46-局限|visual_type=risk-relation-table|width_budget=0.98linewidth|height_budget=0.86textheight|min_font_pt=10|split_allowed=no|expected_page=197|risks=title-table-source-footline
\begin{frame}[t]{扩散模型即随机量化：局限、跨论文接口与可复查入口}
  \fontsize{10}{10.8}\selectfont
\begingroup
\renewcommand{\arraystretch}{0.84}
\begin{tabularx}{0.98\linewidth}{@{}p{2.55cm}>{\raggedright\arraybackslash}X@{}}
\toprule
源文局限 & 需要回到结果/讨论/展望章节核对适用区间；本档案不把缺失的独立数据复核补成已证实结论。 \\
跨论文接口 & 为大规模格点系综生成提供算法补充，但必须回到物理可观测量和系综正确性检验。 \\
极限检查 & 规范不变性/等变性是结构约束；采样质量须由可观测量、接受率、自相关或分布距离验证，不能只看训练损失。 \\
下一步核验 & 补齐 PDF 或原始数据；按源文参数复现关键图表；比较不同动量、流时间、距离、格距或方案转换后的稳定性。 \\
完整来源 & \detokenize{refer/papers/扩散模型即随机量化_latex/main.tex}；\detokenize{refer/papers/Diffusion_Stochastic_Quantization_latex/main.tex}；章节源：\detokenize{refer/papers/扩散模型即随机量化_latex/chapters/*.tex}；英中目录：\detokenize{refer/papers/Diffusion_Stochastic_Quantization_latex}. \\
\bottomrule
\end{tabularx}
\endgroup
  \source{状态：\inferred 为主题接口推断；\unverified 为当前材料边界；完整文件清单见 manifest。}
\end{frame}

%% frame_manifest|frame_id=P47-01|paper_id=P47|claim_id=P47-定位|visual_type=paper-card|width_budget=0.98linewidth|height_budget=0.86textheight|min_font_pt=10|split_allowed=no|expected_page=198|risks=title-table-source-footline
\section{任意阶部分子分布矩}
\begin{frame}[t]{任意阶部分子分布矩：研究问题与论文定位}
  \fontsize{10}{10.8}\selectfont
\begingroup
\renewcommand{\arraystretch}{0.84}
\begin{tabularx}{0.98\linewidth}{@{}p{2.55cm}>{\raggedright\arraybackslash}X@{}}
\toprule
论文编号 & P47；主题：梯度流、涂抹与能动量张量 \\
书目信息 & Moments of parton distribution functions of any order from lattice QCD；arXiv:2311.18704；英/中转排页数 15/13 \\
研究对象 & 我们描述一种在格点量子色动力学（QCD）中确定任意阶部分子分布函数矩的方法。 该方法基于费米子场与规范场的梯度流。带流的 twist-2 算符矩阵元是乘法重正化的， 与物理矩阵元的匹配可以借助连续对称性和欧几里得四维转动的不可约表……配。因此可以只使用时间方向指标的 twist-2 算符， 从而显著改善强子矩阵元计算中的信噪比 \\
章节证据 & 引言；twist-2 算符；梯度流；O(（公式）) 改进；匹配系数；应用；总结；补充材料；有限体积效应；离散化误差；微扰匹配 \\
证据状态 & \verified：摘要和章节来自现有 LaTeX；\unverified：没有论文 PDF/原始数据独立复核。 \\
\bottomrule
\end{tabularx}
\endgroup
  \source{索引：\detokenize{refer/papers/INDEX.md:47}；正文：\detokenize{refer/papers/任意阶部分子分布矩_latex/main.tex}；\detokenize{refer/papers/PDF_Moments_Any_Order_latex/main.tex}}
\end{frame}

%% frame_manifest|frame_id=P47-02|paper_id=P47|claim_id=P47-方法|visual_type=method-flow-formula|width_budget=0.98linewidth|height_budget=0.86textheight|min_font_pt=10|split_allowed=no|expected_page=199|risks=title-table-source-footline
\begin{frame}[t]{任意阶部分子分布矩：理论结构与方法链}
  \fontsize{10}{10.8}\selectfont
\begingroup
\renewcommand{\arraystretch}{0.84}
\begin{tabularx}{0.98\linewidth}{@{}p{2.55cm}>{\raggedright\arraybackslash}X@{}}
\toprule
输入与状态 & 格点/解析输入 → 梯度流、涂抹与能动量张量中的算符、矩阵元或场配置；具体输入见源章节。 \\
章节路线 & 引言；twist-2 算符；梯度流；O(（公式）) 改进；匹配系数；应用；总结；补充材料；有限体积效应；离散化误差；微扰匹配 \\
方法摘录 & 量子色动力学（QCD）描述夸克—胶子在很宽能量范围内的相互作用，并在低能区提出 非微扰挑战。强子结构研究依赖于部分子分布函数（PDF），而 PDF 的确定又依赖于 精确的 QCD 计算，影响着 LHC 以及未来项目（EIC、HL-LHC、LHe……互补的途径。 我们描述的方法同时解决了过去阻碍从格点 QCD 计算任意阶矩的 理论与数值两方面的挑战 \\
结构式/示意 & $\partial_t B_\mu(t,x)=D_\nu G_{\nu\mu}(t,x),\qquad t\sim L_{\rm sm}^{2}$\\[-0.2em]\scriptsize 该式只表达主题变量关系；论文特定推导以源文件为准。 \\
物理校验 & 流时间与平滑长度满足平方关系；t→0 回到原始场，有限 t 则抑制紫外模式。 \\
\bottomrule
\end{tabularx}
\endgroup
  \source{方法证据：\detokenize{refer/papers/任意阶部分子分布矩_latex/chapters/section01.tex:1-53}；主题结构式为示意。}
\end{frame}

%% frame_manifest|frame_id=P47-03|paper_id=P47|claim_id=P47-结果|visual_type=result-evidence-table|width_budget=0.98linewidth|height_budget=0.86textheight|min_font_pt=10|split_allowed=no|expected_page=200|risks=title-table-source-footline
\begin{frame}[t]{任意阶部分子分布矩：结果、验证与物理含义}
  \fontsize{10}{10.8}\selectfont
\begingroup
\renewcommand{\arraystretch}{0.84}
\begin{tabularx}{0.98\linewidth}{@{}p{2.55cm}>{\raggedright\arraybackslash}X@{}}
\toprule
源文结果 & 我们描述了一种方法，它同时解决了从格点 QCD 计算部分子分布函数矩时 所面临的理论与数值挑战。 格点调节把欧几里得转动对称性破缺到超立方群， 因此格点上 twist-2 算符的重正化是确定部分子分布函数矩（公式）的一大障碍。 虽然当（公式）时相对格距的幂次发散不可避免， 但对（公式），通过选取超立方群合适的不可约表示仍可消去幂次发散； 代价是强子矩阵元的计算必须使用空间外动量……若干系统不确定度制约。 补充材料（引文）中讨论了这些系统误差， 但只有透彻的数值研究才能最终回答其适用范围， 以及这一结果是否为确定任意阶部分子分布函数矩开辟了道路 \\
验证状态 & 确证：源章节有结果/讨论摘录。 \\
库内可见证据 & 中文源约 1377 行；公式命中 4，图命中 12，表/表格命中 2。 \\
物理含义 & 在本主题共同链条中，本论文把上述输入推进到可比较的中间量或观测量；具体强度、误差和外推范围不超出源文档。 \\
未验证项 & 当前工作区没有论文 PDF、原始数据和独立复现脚本；因此数值复核、图像再现和统计显著性不作额外断言。 \\
\bottomrule
\end{tabularx}
\endgroup
  \source{结果证据：\detokenize{refer/papers/任意阶部分子分布矩_latex/chapters/section06.tex:1-28}；图表/公式计数由本报告生成器对中文源的静态统计得到。}
\end{frame}

%% frame_manifest|frame_id=P47-04|paper_id=P47|claim_id=P47-局限|visual_type=risk-relation-table|width_budget=0.98linewidth|height_budget=0.86textheight|min_font_pt=10|split_allowed=no|expected_page=201|risks=title-table-source-footline
\begin{frame}[t]{任意阶部分子分布矩：局限、跨论文接口与可复查入口}
  \fontsize{10}{10.8}\selectfont
\begingroup
\renewcommand{\arraystretch}{0.84}
\begin{tabularx}{0.98\linewidth}{@{}p{2.55cm}>{\raggedright\arraybackslash}X@{}}
\toprule
源文局限 & 需要回到结果/讨论/展望章节核对适用区间；本档案不把缺失的独立数据复核补成已证实结论。 \\
跨论文接口 & 把连续极限、微扰有限性和能动量张量重构连接到可计算的格点观测量。 \\
极限检查 & 流时间与平滑长度满足平方关系；t→0 回到原始场，有限 t 则抑制紫外模式。 \\
下一步核验 & 补齐 PDF 或原始数据；按源文参数复现关键图表；比较不同动量、流时间、距离、格距或方案转换后的稳定性。 \\
完整来源 & \detokenize{refer/papers/任意阶部分子分布矩_latex/main.tex}；\detokenize{refer/papers/PDF_Moments_Any_Order_latex/main.tex}；章节源：\detokenize{refer/papers/任意阶部分子分布矩_latex/chapters/*.tex}；英中目录：\detokenize{refer/papers/PDF_Moments_Any_Order_latex}. \\
\bottomrule
\end{tabularx}
\endgroup
  \source{状态：\inferred 为主题接口推断；\unverified 为当前材料边界；完整文件清单见 manifest。}
\end{frame}

%% frame_manifest|frame_id=P48-01|paper_id=P48|claim_id=P48-定位|visual_type=paper-card|width_budget=0.98linewidth|height_budget=0.86textheight|min_font_pt=10|split_allowed=no|expected_page=202|risks=title-table-source-footline
\section{准分布的幂修正与renormalon}
\begin{frame}[t]{准分布的幂修正与renormalon：研究问题与论文定位}
  \fontsize{10}{10.8}\selectfont
\begingroup
\renewcommand{\arraystretch}{0.84}
\begin{tabularx}{0.98\linewidth}{@{}p{2.55cm}>{\raggedright\arraybackslash}X@{}}
\toprule
论文编号 & P48；主题：非微扰重正化、OPE 与匹配 \\
书目信息 & Power corrections and renormalons in parton quasi-distributions；arXiv:1810.00048；英/中转排页数 24/24 \\
研究对象 & QCD 中短距离量的微扰展开是阶乘发散的，而这一缺陷可以转化为研究非微扰修正的有用工具。在本工作中，我们用这一方法研究格点部分子分布函数计算中出现的部分子准分布（quasi-distribution）与赝分布（pseudo-dis…… Bjorken（公式）变量的函数依赖关系的预言。我们还证明这些修正会受到归一化步骤的强烈影响 \\
章节证据 & 引言；一般框架；部分子准分布；光线 OPE 与靶质量修正；Borel 变换与 renormalon；大-（公式）系数函数；Borel 变换的奇点；幂修正；Ioffe 时间准分布；部分子准分布……ion in the bubble-chain approximation \\
证据状态 & \verified：摘要和章节来自现有 LaTeX；\unverified：没有论文 PDF/原始数据独立复核。 \\
\bottomrule
\end{tabularx}
\endgroup
  \source{索引：\detokenize{refer/papers/INDEX.md:48}；正文：\detokenize{refer/papers/准分布的幂修正与renormalon_latex/main.tex}；\detokenize{refer/papers/Power_Corrections_Renormalons_Quasi_latex/main.tex}}
\end{frame}

%% frame_manifest|frame_id=P48-02|paper_id=P48|claim_id=P48-方法|visual_type=method-flow-formula|width_budget=0.98linewidth|height_budget=0.86textheight|min_font_pt=10|split_allowed=no|expected_page=203|risks=title-table-source-footline
\begin{frame}[t]{准分布的幂修正与renormalon：理论结构与方法链}
  \fontsize{10}{10.8}\selectfont
\begingroup
\renewcommand{\arraystretch}{0.84}
\begin{tabularx}{0.98\linewidth}{@{}p{2.55cm}>{\raggedright\arraybackslash}X@{}}
\toprule
输入与状态 & 格点/解析输入 → 非微扰重正化、OPE 与匹配中的算符、矩阵元或场配置；具体输入见源章节。 \\
章节路线 & 引言；一般框架；部分子准分布；光线 OPE 与靶质量修正；Borel 变换与 renormalon；大-（公式）系数函数；Borel 变换的奇点；幂修正；Ioffe 时间准分布；部分子准分布……ion in the bubble-chain approximation \\
方法摘录 & sect\_introduction1 QCD 的格点计算已经展示了与大量实验测量互补的能力，在某些情况下其精度甚至超过实验测量。如今，部分子分布函数（PDF）的格点计算已提上日程。除了允许计算 PDF 前几个 Mellin 矩的标准途径之外，人……集相关的解析表达式，并用价夸克 PDF 的现实参数化给出数值研究结果。最后的第 5 节用于总结与结论 \\
结构式/示意 & $O^{R}(z,\mu)=Z(z,\mu,a)\,O^{\rm bare}(z,a)$\\[-0.2em]\scriptsize 该式只表达主题变量关系；论文特定推导以源文件为准。 \\
物理校验 & 重正化因子抵消截止依赖；a→0 和短距离 z→0 极限必须与匹配阶数、方案和尺度保持一致。 \\
\bottomrule
\end{tabularx}
\endgroup
  \source{方法证据：\detokenize{refer/papers/准分布的幂修正与renormalon_latex/chapters/section01.tex:1-66}；主题结构式为示意。}
\end{frame}

%% frame_manifest|frame_id=P48-03|paper_id=P48|claim_id=P48-结果|visual_type=result-evidence-table|width_budget=0.98linewidth|height_budget=0.86textheight|min_font_pt=10|split_allowed=no|expected_page=204|risks=title-table-source-footline
\begin{frame}[t]{准分布的幂修正与renormalon：结果、验证与物理含义}
  \fontsize{10}{10.8}\selectfont
\begingroup
\renewcommand{\arraystretch}{0.84}
\begin{tabularx}{0.98\linewidth}{@{}p{2.55cm}>{\raggedright\arraybackslash}X@{}}
\toprule
源文结果 & 我们基于对相应系数函数中阶乘发散（renormalon）的研究，在气泡链近似下分析了 qPDF 与 pPDF 的幂次压低（高扭度）贡献。阶乘渐近行为意味着级数之和仅在幂次精度内有定义，因此 QCD 微扰论必须补充非微扰的幂次压低贡献才能给出无歧义的预言。我们的结果应被视作高扭度修正的一个``最小模型''：它包含了理论自洽性所必需的效应，但可能遗漏其他非微扰修正，例如受对称性保护……给出了系数函数的全阶结果，即（公式）的项。 这些结果可用于估计未计算高阶效应及采用 BLM 型步骤设定尺度的影响。 相应分析超出了本文的范围，将在未来的出版物中给出 \\
验证状态 & 确证：源章节有结果/讨论摘录。 \\
库内可见证据 & 中文源约 1598 行；公式命中 86，图命中 22，表/表格命中 0。 \\
物理含义 & 在本主题共同链条中，本论文把上述输入推进到可比较的中间量或观测量；具体强度、误差和外推范围不超出源文档。 \\
未验证项 & 当前工作区没有论文 PDF、原始数据和独立复现脚本；因此数值复核、图像再现和统计显著性不作额外断言。 \\
\bottomrule
\end{tabularx}
\endgroup
  \source{结果证据：\detokenize{refer/papers/准分布的幂修正与renormalon_latex/chapters/section05.tex:1-23}；图表/公式计数由本报告生成器对中文源的静态统计得到。}
\end{frame}

%% frame_manifest|frame_id=P48-04|paper_id=P48|claim_id=P48-局限|visual_type=risk-relation-table|width_budget=0.98linewidth|height_budget=0.86textheight|min_font_pt=10|split_allowed=no|expected_page=205|risks=title-table-source-footline
\begin{frame}[t]{准分布的幂修正与renormalon：局限、跨论文接口与可复查入口}
  \fontsize{10}{10.8}\selectfont
\begingroup
\renewcommand{\arraystretch}{0.84}
\begin{tabularx}{0.98\linewidth}{@{}p{2.55cm}>{\raggedright\arraybackslash}X@{}}
\toprule
源文局限 & 需要回到结果/讨论/展望章节核对适用区间；本档案不把缺失的独立数据复核补成已证实结论。 \\
跨论文接口 & 是准/赝分布从格点数据走向 MS-bar 或其他连续方案的必要中间层。 \\
极限检查 & 重正化因子抵消截止依赖；a→0 和短距离 z→0 极限必须与匹配阶数、方案和尺度保持一致。 \\
下一步核验 & 补齐 PDF 或原始数据；按源文参数复现关键图表；比较不同动量、流时间、距离、格距或方案转换后的稳定性。 \\
完整来源 & \detokenize{refer/papers/准分布的幂修正与renormalon_latex/main.tex}；\detokenize{refer/papers/Power_Corrections_Renormalons_Quasi_latex/main.tex}；章节源：\detokenize{refer/papers/准分布的幂修正与renormalon_latex/chapters/*.tex}；英中目录：\detokenize{refer/papers/Power_Corrections_Renormalons_Quasi_latex}. \\
\bottomrule
\end{tabularx}
\endgroup
  \source{状态：\inferred 为主题接口推断；\unverified 为当前材料边界；完整文件清单见 manifest。}
\end{frame}

%% frame_manifest|frame_id=P49-01|paper_id=P49|claim_id=P49-定位|visual_type=paper-card|width_budget=0.98linewidth|height_budget=0.86textheight|min_font_pt=10|split_allowed=no|expected_page=206|risks=title-table-source-footline
\section{伪部分子分布的单圈演化}
\begin{frame}[t]{伪部分子分布的单圈演化：研究问题与论文定位}
  \fontsize{10}{10.8}\selectfont
\begingroup
\renewcommand{\arraystretch}{0.84}
\begin{tabularx}{0.98\linewidth}{@{}p{2.55cm}>{\raggedright\arraybackslash}X@{}}
\toprule
论文编号 & P49；主题：LaMET、准分布与赝分布 \\
书目信息 & One-loop evolution of parton pseudo-distribution functions on the lattice arXiv:1801.02427v3 [hep-ph], 3 Jul……019 (2018); DOI: 10.1103/PhysRevD.98.014019.；arXiv:1801.02427；英/中转排页数 17/14 \\
研究对象 & 我们把关于约化 Ioffe 时间伪分布（公式）的单圈修正的近期计算纳入进来，以扩展 Orginos 等人对格点数据所作的领头对数分析。我们观察到，单圈修正包含一个大项，它反映了如下事实：最重要的图形所涉及的有效距离远小于名义距离（……展开是可控的。如此提取出的部分子分布在（公式）区域与全局拟合相当接近，但在（公式）区域偏离它们 \\
章节证据 & 引言；Ioffe 时间分布与伪 PDF；单圈修正的结构；禁闭与红外截断；重标度关系；单圈修正；格点数据中的演化；总体特征；构建（公式）ITD；虚部；总结与结论 \\
证据状态 & \verified：摘要和章节来自现有 LaTeX；\unverified：没有论文 PDF/原始数据独立复核。 \\
\bottomrule
\end{tabularx}
\endgroup
  \source{索引：\detokenize{refer/papers/INDEX.md:49}；正文：\detokenize{refer/papers/伪部分子分布的单圈演化_latex/main.tex}；\detokenize{refer/papers/One_Loop_Evolution_Pseudo_PDF_latex/main.tex}}
\end{frame}

%% frame_manifest|frame_id=P49-02|paper_id=P49|claim_id=P49-方法|visual_type=method-flow-formula|width_budget=0.98linewidth|height_budget=0.86textheight|min_font_pt=10|split_allowed=no|expected_page=207|risks=title-table-source-footline
\begin{frame}[t]{伪部分子分布的单圈演化：理论结构与方法链}
  \fontsize{10}{10.8}\selectfont
\begingroup
\renewcommand{\arraystretch}{0.84}
\begin{tabularx}{0.98\linewidth}{@{}p{2.55cm}>{\raggedright\arraybackslash}X@{}}
\toprule
输入与状态 & 格点/解析输入 → LaMET、准分布与赝分布中的算符、矩阵元或场配置；具体输入见源章节。 \\
章节路线 & 引言；Ioffe 时间分布与伪 PDF；单圈修正的结构；禁闭与红外截断；重标度关系；单圈修正；格点数据中的演化；总体特征；构建（公式）ITD；虚部；总结与结论 \\
方法摘录 & 费曼的部分子分布函数（PDF） [引文]（公式）是量子色动力学（QCD）中描述硬包容过程的关键构件。作为累积强子结构非微扰信息的对象，PDF 是格点研究的天然课题。然而，PDF 的直接定义涉及光锥（公式）上双定域算符的矩阵元——这样的间隔在欧氏…… ITD（公式）的数据转换为在（公式）方案中定义的标准部分子密度（公式）。全文总结与结论见第 V 节 \\
结构式/示意 & $\widetilde q(x,P_z,\mu)=\int dy\,C(x,y,\mu/P_z)q(y,\mu)+\mathcal O(\Lambda_{\rm QCD}^{2}/P_z^{2})$\\[-0.2em]\scriptsize 该式只表达主题变量关系；论文特定推导以源文件为准。 \\
物理校验 & 大动量极限 P\_z→∞ 压低幂修正；有限动量、有限间距和核子激发态是必须单独控制的误差源。 \\
\bottomrule
\end{tabularx}
\endgroup
  \source{方法证据：\detokenize{refer/papers/伪部分子分布的单圈演化_latex/chapters/section01.tex:1-26}；主题结构式为示意。}
\end{frame}

%% frame_manifest|frame_id=P49-03|paper_id=P49|claim_id=P49-结果|visual_type=result-evidence-table|width_budget=0.98linewidth|height_budget=0.86textheight|min_font_pt=10|split_allowed=no|expected_page=208|risks=title-table-source-footline
\begin{frame}[t]{伪部分子分布的单圈演化：结果、验证与物理含义}
  \fontsize{10}{10.8}\selectfont
\begingroup
\renewcommand{\arraystretch}{0.84}
\begin{tabularx}{0.98\linewidth}{@{}p{2.55cm}>{\raggedright\arraybackslash}X@{}}
\toprule
源文结果 & 本文扩展了文献（引文）对部分子伪分布及约化伪 ITD 格点数据所作的领头对数分析。为此，我们纳入了约化伪 ITD 单圈层次的近期结果（引文）（另见（引文））。 我们发现该修正包含一个大项，导致与 LLA 相比的重要数值变化。大修正之所以出现，是因为最重要图形所涉及的有效距离远小于名义距离（公式）。这改变了重标度关系（公式）中的系数（公式）（在我们的特定 ITD 情形从（公式）变……物理（公式）介子质量的模拟会在小（公式）区域给出与全局拟合更好的符合。这一期望得到近期在物理（公式）介子质量下用准 PDF 格点模拟提取（公式）的工作（引文）的支持 \\
验证状态 & 确证：源章节有结果/讨论摘录。 \\
库内可见证据 & 中文源约 841 行；公式命中 27，图命中 22，表/表格命中 0。 \\
物理含义 & 在本主题共同链条中，本论文把上述输入推进到可比较的中间量或观测量；具体强度、误差和外推范围不超出源文档。 \\
未验证项 & 当前工作区没有论文 PDF、原始数据和独立复现脚本；因此数值复核、图像再现和统计显著性不作额外断言。 \\
\bottomrule
\end{tabularx}
\endgroup
  \source{结果证据：\detokenize{refer/papers/伪部分子分布的单圈演化_latex/chapters/section05.tex:1-12}；图表/公式计数由本报告生成器对中文源的静态统计得到。}
\end{frame}

%% frame_manifest|frame_id=P49-04|paper_id=P49|claim_id=P49-局限|visual_type=risk-relation-table|width_budget=0.98linewidth|height_budget=0.86textheight|min_font_pt=10|split_allowed=no|expected_page=209|risks=title-table-source-footline
\begin{frame}[t]{伪部分子分布的单圈演化：局限、跨论文接口与可复查入口}
  \fontsize{10}{10.8}\selectfont
\begingroup
\renewcommand{\arraystretch}{0.84}
\begin{tabularx}{0.98\linewidth}{@{}p{2.55cm}>{\raggedright\arraybackslash}X@{}}
\toprule
源文局限 & 需要回到结果/讨论/展望章节核对适用区间；本档案不把缺失的独立数据复核补成已证实结论。 \\
跨论文接口 & 承接格点矩阵元，向重正化、匹配、演化和最终 PDF/TMD 现象学输出。 \\
极限检查 & 大动量极限 P\_z→∞ 压低幂修正；有限动量、有限间距和核子激发态是必须单独控制的误差源。 \\
下一步核验 & 补齐 PDF 或原始数据；按源文参数复现关键图表；比较不同动量、流时间、距离、格距或方案转换后的稳定性。 \\
完整来源 & \detokenize{refer/papers/伪部分子分布的单圈演化_latex/main.tex}；\detokenize{refer/papers/One_Loop_Evolution_Pseudo_PDF_latex/main.tex}；章节源：\detokenize{refer/papers/伪部分子分布的单圈演化_latex/chapters/*.tex}；英中目录：\detokenize{refer/papers/One_Loop_Evolution_Pseudo_PDF_latex}. \\
\bottomrule
\end{tabularx}
\endgroup
  \source{状态：\inferred 为主题接口推断；\unverified 为当前材料边界；完整文件清单见 manifest。}
\end{frame}

%% frame_manifest|frame_id=P50-01|paper_id=P50|claim_id=P50-定位|visual_type=paper-card|width_budget=0.98linewidth|height_budget=0.86textheight|min_font_pt=10|split_allowed=no|expected_page=210|risks=title-table-source-footline
\section{同位旋矢量分布的重正化匹配系统分析}
\begin{frame}[t]{同位旋矢量分布的重正化匹配系统分析：研究问题与论文定位}
  \fontsize{10}{10.8}\selectfont
\begingroup
\renewcommand{\arraystretch}{0.84}
\begin{tabularx}{0.98\linewidth}{@{}p{2.55cm}>{\raggedright\arraybackslash}X@{}}
\toprule
论文编号 & P50；主题：非微扰重正化、OPE 与匹配 \\
书目信息 & 1.2ex] Lattice Parton Collaboration (LPC)；arXiv:1807.06566；英/中转排页数 24/22 \\
研究对象 & 我们利用大动量有效理论（large-momentum effective theory）框架，对非极化同位旋矢量夸克部分子分布函数（PDF）进行了一项细致的格点 QCD 研究。我们选取由空间关联函数定义的准 PDF，它与同量纲的其……细分析了来自重正化和微扰匹配的系统误差。我们得到的光锥 PDF 结果与最新的唯象分析符合得较好 \\
章节证据 & 引言；非微扰重正化与匹配；格点上的 RI/MOM 重正化；准 PDF 与 PDF 的一圈匹配；PDF 的格点计算；格点矩阵元素；系统 不确定度；PDF 的最终结果；总结；一般协变规范下带（公式）的一圈准 PDF；朗道规范下（公式）的一圈准 PDF \\
证据状态 & \verified：摘要和章节来自现有 LaTeX；\unverified：没有论文 PDF/原始数据独立复核。 \\
\bottomrule
\end{tabularx}
\endgroup
  \source{索引：\detokenize{refer/papers/INDEX.md:50}；正文：\detokenize{refer/papers/同位旋矢量分布的重正化匹配系统分析_latex/main.tex}；\detokenize{refer/papers/Isovector_Systematic_Renorm_Matching_latex/main.tex}}
\end{frame}

%% frame_manifest|frame_id=P50-02|paper_id=P50|claim_id=P50-方法|visual_type=method-flow-formula|width_budget=0.98linewidth|height_budget=0.86textheight|min_font_pt=10|split_allowed=no|expected_page=211|risks=title-table-source-footline
\begin{frame}[t]{同位旋矢量分布的重正化匹配系统分析：理论结构与方法链}
  \fontsize{10}{10.8}\selectfont
\begingroup
\renewcommand{\arraystretch}{0.84}
\begin{tabularx}{0.98\linewidth}{@{}p{2.55cm}>{\raggedright\arraybackslash}X@{}}
\toprule
输入与状态 & 格点/解析输入 → 非微扰重正化、OPE 与匹配中的算符、矩阵元或场配置；具体输入见源章节。 \\
章节路线 & 引言；非微扰重正化与匹配；格点上的 RI/MOM 重正化；准 PDF 与 PDF 的一圈匹配；PDF 的格点计算；格点矩阵元素；系统 不确定度；PDF 的最终结果；总结；一般协变规范下带（公式）的一圈准 PDF；朗道规范下（公式）的一圈准 PDF \\
方法摘录 & subsec:ME 在本小节中，我们给出在 CLS（公式）、（公式）味 clover 费米子系综 H102 上的格点 QCD 计算结果：格距（公式）fm，pion 质量（公式）356 MeV，盒子尺寸（公式）fm（（公式））（引文）。价 clo……变化的曲线。不同（公式）的结果在统计不确定度内彼此一致。这表明来自更高扭度效应的幂次修正可能并不可观 \\
结构式/示意 & $O^{R}(z,\mu)=Z(z,\mu,a)\,O^{\rm bare}(z,a)$\\[-0.2em]\scriptsize 该式只表达主题变量关系；论文特定推导以源文件为准。 \\
物理校验 & 重正化因子抵消截止依赖；a→0 和短距离 z→0 极限必须与匹配阶数、方案和尺度保持一致。 \\
\bottomrule
\end{tabularx}
\endgroup
  \source{方法证据：\detokenize{refer/papers/同位旋矢量分布的重正化匹配系统分析_latex/chapters/section03.tex:3-112}；主题结构式为示意。}
\end{frame}

%% frame_manifest|frame_id=P50-03|paper_id=P50|claim_id=P50-结果|visual_type=result-evidence-table|width_budget=0.98linewidth|height_budget=0.86textheight|min_font_pt=10|split_allowed=no|expected_page=212|risks=title-table-source-footline
\begin{frame}[t]{同位旋矢量分布的重正化匹配系统分析：结果、验证与物理含义}
  \fontsize{10}{10.8}\selectfont
\begingroup
\renewcommand{\arraystretch}{0.84}
\begin{tabularx}{0.98\linewidth}{@{}p{2.55cm}>{\raggedright\arraybackslash}X@{}}
\toprule
源文结果 & sec:summary 本文研究了以（公式）定义的准 PDF，它在（公式）阶无算符混合。我们利用（公式）MeV 的格点数据演示了匹配流程，并表明激发态污染得到了良好控制。我们计算了一圈匹配系数，并详细讨论了系统误差来源以及投影的选取。 我们发现来自 FT 截断方法和（公式）依赖的系统不确定度相当可观。但在（公式）时，来自（公式）、匹配反演以及投影选取的不确定度相对较小。与此同时……子动量外推进行估计。 我们关于光锥 PDF 的最终结果在大（公式）区与全局分析一致，这是一个令人鼓舞的信号，表明 LaMET 或许能使我们在未来精确地触及部分子物理 \\
验证状态 & 确证：源章节有结果/讨论摘录。 \\
库内可见证据 & 中文源约 1560 行；公式命中 49，图命中 31，表/表格命中 2。 \\
物理含义 & 在本主题共同链条中，本论文把上述输入推进到可比较的中间量或观测量；具体强度、误差和外推范围不超出源文档。 \\
未验证项 & 当前工作区没有论文 PDF、原始数据和独立复现脚本；因此数值复核、图像再现和统计显著性不作额外断言。 \\
\bottomrule
\end{tabularx}
\endgroup
  \source{结果证据：\detokenize{refer/papers/同位旋矢量分布的重正化匹配系统分析_latex/chapters/section05.tex:1-12}；图表/公式计数由本报告生成器对中文源的静态统计得到。}
\end{frame}

%% frame_manifest|frame_id=P50-04|paper_id=P50|claim_id=P50-局限|visual_type=risk-relation-table|width_budget=0.98linewidth|height_budget=0.86textheight|min_font_pt=10|split_allowed=no|expected_page=213|risks=title-table-source-footline
\begin{frame}[t]{同位旋矢量分布的重正化匹配系统分析：局限、跨论文接口与可复查入口}
  \fontsize{10}{10.8}\selectfont
\begingroup
\renewcommand{\arraystretch}{0.84}
\begin{tabularx}{0.98\linewidth}{@{}p{2.55cm}>{\raggedright\arraybackslash}X@{}}
\toprule
源文局限 & 需要回到结果/讨论/展望章节核对适用区间；本档案不把缺失的独立数据复核补成已证实结论。 \\
跨论文接口 & 是准/赝分布从格点数据走向 MS-bar 或其他连续方案的必要中间层。 \\
极限检查 & 重正化因子抵消截止依赖；a→0 和短距离 z→0 极限必须与匹配阶数、方案和尺度保持一致。 \\
下一步核验 & 补齐 PDF 或原始数据；按源文参数复现关键图表；比较不同动量、流时间、距离、格距或方案转换后的稳定性。 \\
完整来源 & \detokenize{refer/papers/同位旋矢量分布的重正化匹配系统分析_latex/main.tex}；\detokenize{refer/papers/Isovector_Systematic_Renorm_Matching_latex/main.tex}；章节源：\detokenize{refer/papers/同位旋矢量分布的重正化匹配系统分析_latex/chapters/*.tex}；英中目录：\detokenize{refer/papers/Isovector_Systematic_Renorm_Matching_latex}. \\
\bottomrule
\end{tabularx}
\endgroup
  \source{状态：\inferred 为主题接口推断；\unverified 为当前材料边界；完整文件清单见 manifest。}
\end{frame}

%% frame_manifest|frame_id=C-01|paper_id=GLOBAL|claim_id=C1|visual_type=synthesis-flow|width_budget=0.98linewidth|height_budget=0.86textheight|min_font_pt=10|split_allowed=no|expected_page=214|risks=title-table-source-footline
\begin{frame}[t]{交叉结论：不同方法最终都在控制“可计算性”和“光锥物理”的间隙}
  \small
  \begingroup
\renewcommand{\arraystretch}{0.84}
\begin{tabularx}{0.98\linewidth}{@{}p{2.55cm}>{\raggedright\arraybackslash}X@{}}
\toprule
格点底座 & 规范不变离散化与强子关联函数使欧氏时空中的非微扰量可计算。 \\
紫外控制 & 涂抹、梯度流和非微扰重正化抑制或拆分短距离发散，但引入流时间、方案和匹配尺度。 \\
运动学桥梁 & LaMET、准/赝 PDF 和 TMD 因子化用大动量或短距离展开连接等时算符与光锥分布。 \\
物理输出 & 只有在动量、距离、格距、体积和匹配阶数等系统性检验闭合后，才可提升为强子结构结论。 \\
\bottomrule
\end{tabularx}
\endgroup
  \source{综合判断：50 篇论文的章节和结论；跨主题关系属于\inferred，不替代单篇论文证据。}
\end{frame}

%% frame_manifest|frame_id=C-02|paper_id=GLOBAL|claim_id=C2|visual_type=error-matrix|width_budget=0.98linewidth|height_budget=0.86textheight|min_font_pt=10|split_allowed=no|expected_page=215|risks=title-table-source-footline
\begin{frame}[t]{交叉结论：系统误差必须按物理极限分层，而不能用单一误差条代替}
  \small
  \begingroup
\renewcommand{\arraystretch}{0.84}
\begin{tabularx}{0.98\linewidth}{@{}p{2.55cm}>{\raggedright\arraybackslash}X@{}}
\toprule
离散化 & 改变格距或使用改进作用量，检验 a→0；静态源文件未提供统一跨论文极限数据。 \\
有限体积 & 检查 \(m_\pi L\)、长距离关联和强子动量；各论文参数不同，不能直接合并百分比。 \\
激发态 & 多态拟合、源汇时间和矩阵元平台共同决定基态隔离；缺少原始相关函数时标记未验证。 \\
匹配/重正化 & 方案、尺度、截断阶数和算符混合影响连续分布；必须与对应论文的公式和附录共同复核。 \\
统计/采样 & 自相关、系综覆盖和相关误差影响有效样本数；机器学习采样必须回到物理可观测量。 \\
\bottomrule
\end{tabularx}
\endgroup
  \source{综合依据：各篇的设置、结果、误差和讨论章节；数值量级未从不同源文档强行合并。}
\end{frame}

%% frame_manifest|frame_id=C-03|paper_id=GLOBAL|claim_id=C3|visual_type=topic-relation|width_budget=0.98linewidth|height_budget=0.86textheight|min_font_pt=10|split_allowed=no|expected_page=216|risks=title-table-source-footline
\begin{frame}[t]{主题接口：从 Wilson 线到 TMD 的路径由重正化和匹配连接}
  \small
  \begingroup
\renewcommand{\arraystretch}{0.84}
\begin{tabularx}{0.98\linewidth}{@{}p{2.55cm}>{\raggedright\arraybackslash}X@{}}
\toprule
基础 → 流 & 规范不变格点和 Wilson 线定义输入；梯度流/涂抹改变短距离行为并提供平滑尺度。 \\
流 → 重正化 & 流时间、OPE 和能动量张量构造有限或可匹配的复合算符；方案转换仍需明确。 \\
重正化→ & 线性发散、算符混合和 RI/MS-bar 转换决定准分布能否与连续 PDF 对接。 \\
LaMET→ & 核子、$\pi$/K 胶子分布和准 TMD 计算把形式体系推进到现象学所需的对象。 \\
采样 → 全链条 & 流模型、扩散和规范等变采样影响系综成本，但不能绕过可观测量验证。 \\
\bottomrule
\end{tabularx}
\endgroup
  \source{综合来源：各主题档案页；箭头表示方法接口，不表示所有论文都有直接引文关系。}
\end{frame}

%% frame_manifest|frame_id=C-04|paper_id=GLOBAL|claim_id=C4|visual_type=coverage-table|width_budget=0.98linewidth|height_budget=0.86textheight|min_font_pt=10|split_allowed=no|expected_page=217|risks=title-table-source-footline
\begin{frame}[t]{证据覆盖结论：本报告能完整回指源文本，但不能替代原始数据复现}
  \small
  \begingroup
\renewcommand{\arraystretch}{0.84}
\begin{tabularx}{0.98\linewidth}{@{}p{2.55cm}>{\raggedright\arraybackslash}X@{}}
\toprule
已覆盖 & 50 篇索引论文的英中目录、摘要/章节结构、方法与结果摘录、索引页数和源路径。 \\
静态统计 & 中文源的行数、公式环境、图形命令和表格命中数；这些是文件结构统计，不是物理误差分析。 \\
未覆盖 & 论文 PDF 版面逐页核验、原始数据、外部数据库书目信息、独立编译每篇源目录和数值再现。 \\
解释边界 & 主题结构式和跨论文接口是带标签的综合推断；单篇论文的细节以对应源文件为准。 \\
\bottomrule
\end{tabularx}
\endgroup
  \source{直接证据：\detokenize{refer/papers/INDEX.md} 与各英中 LaTeX 目录；限制来自当前工作区材料。}
\end{frame}

%% frame_manifest|frame_id=C-05|paper_id=GLOBAL|claim_id=C5|visual_type=roadmap-table|width_budget=0.98linewidth|height_budget=0.86textheight|min_font_pt=10|split_allowed=no|expected_page=218|risks=title-table-source-footline
\begin{frame}[t]{后续核验清单：每项都有可验收产物，而不是泛泛的“继续优化”}
  \small
  \begingroup
\renewcommand{\arraystretch}{0.84}
\begin{tabularx}{0.98\linewidth}{@{}p{2.55cm}>{\raggedright\arraybackslash}X@{}}
\toprule
材料补齐 & 为 50 篇论文补充可访问 PDF/官方源包或原始数据；验收：书目、图表和源文本逐项对账。 \\
编译复核 & 逐目录重新执行双遍 XeLaTeX；验收：零错误、页数记录完整、英中编号对应。 \\
数值复现 & 按每篇设置复现代表性图表；验收：参数、单位、误差和基线均可追溯。 \\
系统误差 & 分离格距、体积、动量、流时间/距离、激发态和匹配截断；验收：稳定性表格闭合。 \\
物理输出 & 仅在极限和误差证据闭合后更新 PDF/TMD/采样结论；验收：结论—证据—限制一一对应。 \\
\bottomrule
\end{tabularx}
\endgroup
  \source{计划依据：各篇结果/讨论中的未解决问题；具体排期需由研究者另行决定。}
\end{frame}

%% frame_manifest|frame_id=C-06|paper_id=GLOBAL|claim_id=C6|visual_type=source-catalog|width_budget=0.98linewidth|height_budget=0.86textheight|min_font_pt=10|split_allowed=no|expected_page=219|risks=title-table-source-footline
\begin{frame}[t]{来源目录：50 篇论文均有唯一编号和英中目录入口}
  \scriptsize
  \begin{tabularx}{0.98\linewidth}{@{}lXl@{}}
    \toprule
    编号 & 中文目录 & arXiv/索引来源 \\
    \midrule
    P01 & 夸克禁闭 & books/复用(Wilson74) \\
    P02 & SU3链接变量解析涂抹 & arXiv:hep-lat/0311018 \\
    P03 & 夸克场产生算符构造 & arXiv:0905.2160 \\
    P04 & 高动量强子新夸克涂抹 & arXiv:1602.05525 \\
    P05 & 平凡化映射威尔逊流与HMC & arXiv:0907.5491 \\
    P06 & 连续威尔逊圈的无穷N相变 & arXiv:hep-th/0601210 \\
    P07 & 威尔逊流的性质与应用 & arXiv:1006.4518 \\
    P08 & 梯度流的微扰分析 & arXiv:1101.0963 \\
    P09 & 手征对称性与杨米尔斯梯度流 & arXiv:1302.5246 \\
    P10 & 杨米尔斯梯度流提取能动量张量 & arXiv:1304.0533 \\
    P11 & 欧氏格点上的部分子物理 & arXiv:1305.1539 \\
    P12 & 大动量有效理论与部分子物理 & arXiv:1404.6680 \\
    P13 & 准部分子分布单圈匹配 & arXiv:1310.7471 \\
    P14 & 大动量有效理论再论 & arXiv:1706.07416 \\
    P15 & 准分布动量分布与伪分布 & arXiv:1705.01488 \\
    P16 & 部分子伪分布的格点探索 & arXiv:1706.05373 \\
    P17 & 欧氏与光锥分布因子化定理 & arXiv:1801.03917 \\
    \bottomrule
  \end{tabularx}
  \vspace{0.2em}
  \scriptsize 本页列 P01–P17；后续页面继续列出其余论文。
  \source{索引：\detokenize{refer/papers/INDEX.md}；机器可读完整清单：\detokenize{docs/report_refer_papers_all_20260828.manifest.tsv}。}
\end{frame}

%% frame_manifest|frame_id=C-07|paper_id=GLOBAL|claim_id=C7|visual_type=source-catalog|width_budget=0.98linewidth|height_budget=0.86textheight|min_font_pt=10|split_allowed=no|expected_page=220|risks=title-table-source-footline
\begin{frame}[t]{来源目录（续）：中段论文与可追溯状态}
  \scriptsize
  \begin{tabularx}{0.98\linewidth}{@{}lXl@{}}
    \toprule
    编号 & 中文目录 & arXiv/索引来源 \\
    \midrule
    P18 & 大动量有效理论综述 & arXiv:2004.03543 \\
    P19 & 格点算符非微扰重正化一般方法 & arXiv:hep-lat/9411010 \\
    P20 & 准部分子分布的可重正性 & arXiv:1707.03107 \\
    P21 & 威尔逊线重正化改进准分布 & arXiv:1609.08102 \\
    P22 & 大动量有效理论的重正化 & arXiv:1706.08962 \\
    P23 & 非定域夸克双线性的非微扰重正化 & arXiv:1707.07152 \\
    P24 & 准PDF完整非微扰重正化方案 & arXiv:1706.00265 \\
    P25 & 准部分子算符乘法可重正性 & arXiv:1809.01836 \\
    P26 & 大动量有效理论获取胶子分布 & arXiv:1808.10824 \\
    P27 & 胶子伪分布的短距行为 & arXiv:1910.13963 \\
    P28 & 准部分子分布与梯度流 & arXiv:1612.01584 \\
    P29 & 准光前关联函数的自重正化 & arXiv:2103.02965 \\
    P30 & 准光前关联的混合重正化方案 & arXiv:2008.03886 \\
    P31 & CT18全局分析 & arXiv:1912.10053 \\
    P32 & NNPDF17高精度数据部分子分布 & arXiv:1706.00428 \\
    P33 & 基于流的格点MCMC生成模型 & arXiv:1904.12072 \\
    P34 & 规范等变的流采样 & arXiv:2003.06413 \\
    \bottomrule
  \end{tabularx}
  \vspace{0.2em}
  \scriptsize 本页列 P18–P34；最后一页列出 P35–P50。
  \source{索引：\detokenize{refer/papers/INDEX.md}；完整源文件字段见 manifest。}
\end{frame}

%% frame_manifest|frame_id=C-08|paper_id=GLOBAL|claim_id=C8|visual_type=source-catalog|width_budget=0.98linewidth|height_budget=0.86textheight|min_font_pt=10|split_allowed=no|expected_page=221|risks=title-table-source-footline
\begin{frame}[t]{来源目录（续）：后段论文与可追溯状态}
  \scriptsize
  \begin{tabularx}{0.98\linewidth}{@{}lXl@{}}
    \toprule
    编号 & 中文目录 & arXiv/索引来源 \\
    \midrule
    P35 & 含费米子场的格点能动量张量 & arXiv:1403.4772 \\
    P36 & 梯度流形式的双圈能动量张量 & arXiv:1808.09837 \\
    P37 & 梯度流高阶计算的技术与结果 & arXiv:1905.00882 \\
    P38 & 局域涂抹算符乘积展开 & arXiv:1501.05348 \\
    P39 & 微扰论中的涂抹准分布 & arXiv:1710.04607 \\
    P40 & 梯度流中的非光锥威尔逊线算符 & arXiv:2312.05032 \\
    P41 & 首个核子胶子部分子分布 & arXiv:2505.13321 \\
    P42 & $\pi$介子胶子部分子分布 & arXiv:2104.06372 \\
    P43 & K介子胶子分布初探 & arXiv:2112.03124 \\
    P44 & 从准TMD波函数确定CS核 & arXiv:2204.00200 \\
    P45 & LaMET计算TMD软函数 & arXiv:2005.14572 \\
    P46 & 扩散模型即随机量化 & arXiv:2309.17082 \\
    P47 & 任意阶部分子分布矩 & arXiv:2311.18704 \\
    P48 & 准分布的幂修正与renormalon & arXiv:1810.00048 \\
    P49 & 伪部分子分布的单圈演化 & arXiv:1801.02427 \\
    P50 & 同位旋矢量分布的重正化匹配系统分析 & arXiv:1807.06566 \\
    \bottomrule
  \end{tabularx}
  \vspace{0.2em}
  \scriptsize 本页列 P35–P50；每篇档案第 4 页给出中文/英文 main.tex 与章节入口。
  \source{索引：\detokenize{refer/papers/INDEX.md}；机器可读完整清单：\detokenize{docs/report_refer_papers_all_20260828.manifest.tsv}。}
\end{frame}

\end{document}
